\makeatletter
\def\input@path{{template/}}
\makeatother
\documentclass{csnotes}

% ==============================================================================
% DOCUMENT SPECIFIC CUSTOMIZATIONS
% ==============================================================================

% --------------------------- START CUSTOMIZATIONS -----------------------------

% You can add any custom packages or commands specific to this document here

\NewDocumentCommand{\Ass}{m}{\operatorname{Ass}\left(#1\right)}

\NewDocumentCommand{\Sub}{m}{\operatorname{Sub}\left(#1\right)}


\NewDocumentCommand{\AP}{o}{
  \IfNoValueTF{#1}
      {\operatorname{AP}}
      {\operatorname{AP}\left(#1\right)}
}

\NewDocumentCommand{\AF}{o}{
  \IfNoValueTF{#1}
      {\operatorname{AF}}
      {\operatorname{AF}_{#1}}
}

\NewDocumentCommand{\FOL}{o}{
  \IfNoValueTF{#1}
      {\operatorname{FOL}}
      {\operatorname{FOL}_{#1}}
}

\NewDocumentCommand{\Sat}{o}{
  \IfNoValueTF{#1}
      {\operatorname{SAT}}
      {\operatorname{Sat}\left(#1\right)}
}

\NewDocumentCommand{\Unsat}{o}{
  \IfNoValueTF{#1}
      {\operatorname{UNSAT}}
      {\operatorname{Unsat}\left(#1\right)}
}

\NewDocumentCommand{\Val}{m o}{
  \IfNoValueTF{#2}
  {\operatorname{Val}^{#1}}
  {\operatorname{Val}^{#1}\left(#2\right)}
}

\NewDocumentCommand{\Valid}{o}{
  \IfNoValueTF{#1}
      {\operatorname{VALID}}
      {\operatorname{Valid}\left(#1\right)}
}

\NewDocumentCommand{\Taut}{o}{
  \IfNoValueTF{#1}
      {\operatorname{TAUT}}
      {\operatorname{Taut}\left(#1\right)}
}

\NewDocumentCommand{\den}{m o o}{
  \IfNoValueTF{#2}
    {\IfNoValueTF{#3}
      {\llbracket#1\rrbracket}
      {\llbracket#1\rrbracket_{#3}}}
    {\IfNoValueTF{#3}
      {\llbracket#1\rrbracket^{#2}}
      {\llbracket#1\rrbracket^{#2}_{#3}}}
}

\newcommand{\WFF}{\operatorname{WFF}}
\newcommand{\eqdef}{\stackrel{\triangle}{=}}
\newcommand{\PL}{\operatorname{PL}}
\newcommand{\PNF}{\operatorname{PNF}}
\newcommand{\NNF}{\operatorname{NNF}}
\newcommand{\DNF}{\operatorname{DNF}}
\newcommand{\CNF}{\operatorname{CNF}}
\newcommand{\Axi}{\operatorname{Axi}}
\newcommand{\Rules}{\mathcal{R}}
\newcommand{\Rule}{\operatorname{R}}
\newcommand{\Var}{\operatorname{Var}}
\newcommand{\RAA}{\operatorname{RAA}}
\newcommand{\Th}[1]{\operatorname{Th}_{#1}}
\newcommand{\subs}[3]{{{\left(#1\right)}^{#2}_{#3}}}
\newcommand{\FuncSyms}{\mathcal{F}}
\newcommand{\PredSyms}{\mathcal{P}}
\newcommand{\ConstSyms}{\mathcal{C}}
\newcommand{\VarSyms}{\mathcal{X}}
\newcommand{\Terms}{\mathcal{T}}
\newcommand{\Interp}{\mathfrak{I}}
\newcommand{\arty}{\operatorname{\textit{arity}}}


% ---------------------------- END CUSTOMIZATIONS ------------------------------


% ==============================================================================
% DOCUMENT INFORMATION
% ==============================================================================
\title{Logic for Computer Science}
\author{Riccardo Elena, Alfredo Volpe, Riccardo Regina}
\date{\today}

% Additional information for the title page
\university{Università degli Studi Federico II di Napoli}
\department{Dipartimento d'Ingegneria Elettrica e Tecnologie dell'Informazione}
\course{Corso di Laurea Magistrale in Computer Science}
\subject{Logic for Computer Science}
\academicyear{Anno Accademico 2025/2026}

% ==============================================================================
% BIBLIOGRAPHY
% ==============================================================================
% Specify the .bib file with your bibliographic references
% Uncomment the following line when you have created the bibliography.bib file
% \addbibresource{bibliography.bib}

% ==============================================================================
% INIZIO DOCUMENTO
% ==============================================================================

\begin{document}

\pagenumbering{gobble}

\maketitlepage{}
\whitepage{}

\begin{abstract}
  Questo file contiene appunti e materiali di studio relativi al corso di Logic for Computer Science tenuto dal Prof.~M.~Benerecetti alla Università degli Studi Federico II di Napoli nell'A.A.~25/26.
  
  Questo corso tratterà di logica formale, in particolare di logica proposizionale e logica del primo ordine, con un focus particolare sulle applicazioni in informatica e scienze computazionali.
  %todo: update repo
  \textbf{Questi appunti sono destinati a supportare gli studenti nel loro percorso di apprendimento e a fornire una risorsa di riferimento per gli argomenti discussi durante il corso, ma non sostituiscono in alcun modo il materiale ufficiale e le lezioni del docente.}
  Questo documento è prodotto da studenti, per studenti, e potrebbe dunque contenere errori o imprecisioni. Feedback e correzioni sono ben accetti e incentivati, e possono essere inviati tramite pull request o issue nella \href{https://github.com/RiccardoElena/wvaevwv}{repository GitHub del documento}.

  A tal proposito, nel documento saranno presenti alcuni contenuti non provenienti direttamente dalle lezioni, ma aggiunti o modificati dagli autori per migliorare la comprensibilità o l'approfondimento degli argomenti trattati.
  Ci impegniamo a verificare col docente ogni modifica sostanziale al materiale originale, ma in caso tale verifica non fosse ancora avvenuta, verranno inseriti degli indicatori per segnalare le eventuali divergenze dal materiale ufficiale.
  In particolare:
  \begin{itemize}
    \item I contenuti aggiunti o modificati in maniera sostanziale dagli autori saranno segnalati con il simbolo \warningsymbol{}
    \item I contenuti che invece hanno subito esclusivamente riorganizzazione strutturale o di formattazione verranno segnalati col simbolo \infosymbol{}
    \item Eventuali contenuti aggiuntivi, non presentati a lezione ma verificati col docente, saranno segnalati con il simbolo \extrasymbol{}
  \end{itemize}
  
  % Firma autori / link repo (posizionata come firma in basso a destra)
  \begin{flushright}
    \vspace{2cm}
    \textbf{Gli Autori:}\\
    \href{https://github.com/RiccardoElena}{Riccardo Elena}\\
  \end{flushright}
\end{abstract}
\whitepage{}

\newgeometry{top=2cm}
\tableofcontents
\restoregeometry{}

\whitepage{}

\pagenumbering{arabic}
\setcounter{page}{1}

\chapter{Introduzione}

\section{Ragionamenti}

Il ragionamento è l'attività fondamentale attraverso cui l'essere umano processa informazioni per giungere a nuove conclusioni.
Nella vita quotidiana ragioniamo continuamente: dalle decisioni più semplici alle inferenze più complesse, utilizziamo schemi di ragionamento che ci permettono di derivare nuova conoscenza a partire da ciò che già sappiamo.
La logica formale studia proprio questi schemi, astraendo dalla particolare natura del contenuto per concentrarsi sulla forma e sulla validità delle inferenze.
Un esempio classico è il \textit{modus ponens}, una regola inferenziale che permette di concludere $B$ a partire dalle premesse $A \rightarrow B$ e $A$: se sappiamo che \qi{se piove, la strada è bagnata} e osserviamo che \qi{piove}, possiamo validamente concludere che \qi{la strada è bagnata}.
In questo caso, la validità del ragionamento non dipende dal contenuto specifico delle proposizioni (pioggia, strada bagnata), ma dalla loro struttura logica: ogni volta che abbiamo un'implicazione e il suo antecedente, possiamo inferire il conseguente.

Non tutti i ragionamenti, tuttavia, seguono la stessa logica, né hanno lo stesso grado di certezza nelle loro conclusioni.
Possiamo distinguere tre forme principali di ragionamento, ciascuna con caratteristiche e applicazioni diverse, che riflettono diversi modi in cui gli esseri umani costruiscono e giustificano la conoscenza.

\begin{definition}{Ragionamento Deduttivo}
  Il \keyword{ragionamento deduttivo} parte da premesse generali per derivare conclusioni necessarie: se le premesse sono vere, la conclusione non può che essere vera.
\end{definition}

È il ragionamento tipico della matematica e della logica formale, dove la validità è garantita dalla sola struttura dell'argomento.
In un certo senso, il ragionamento deduttivo non aggiunge nuova informazione, ma piuttosto esplicita ciò che è già implicito nelle premesse.

\begin{definition}{Ragionamento Abduttivo}
  Il \keyword{ragionamento abduttivo} procede dall'effetto alla causa, cercando la spiegazione più plausibile per un fenomeno osservato.
\end{definition}

È il ragionamento tipico della diagnosi medica o dell'investigazione scientifica, dove si osservano dei dati e si cerca di inferire la causa più probabile, pur sapendo che potrebbero esserci altre spiegazioni possibili.

\begin{definition}{Ragionamento Induttivo}
  Il \keyword{ragionamento induttivo} procede dal particolare al generale: osservando un certo numero di casi, si induce una regola o una legge generale.
\end{definition}
Qui la conclusione è probabile ma non certa, e la forza dell'inferenza dipende dal numero e dalla qualità delle osservazioni.

Ognuno di questi tipi di ragionamento ha un ruolo fondamentale nella nostra comprensione del mondo e nella nostra capacità di agire efficacemente e per ognuno di essi esistono formulazioni diverse del concetto di validità e di conseguenza logica.
In questo corso ci concentreremo principalmente sul ragionamento deduttivo e sui sistemi formali che ne catturano la struttura. Studieremo come rappresentare il linguaggio naturale attraverso linguaggi formali, quali la \keyword{logica proposizionale} e la \keyword{logica dei predicati}, come definire rigorosamente la nozione di conseguenza logica, e come costruire sistemi di deduzione che ci permettano di derivare meccanicamente conclusioni valide.

\section{Validità e Conseguenza Logica}

Abbiamo visto che esistono diverse forme di ragionamento, ma come possiamo stabilire quando un ragionamento è corretto?
Nel contesto della logica deduttiva, la nozione fondamentale è quella di \textit{validità}: un ragionamento valido è tale che, se le premesse sono vere, la conclusione non può essere falsa.
Questa proprietà dipende unicamente dalla struttura logica dell'argomento, non dal contenuto specifico delle proposizioni coinvolte.

\begin{definition}{Ragionamento valido (informalmente)}
  Un ragionamento è detto \keyword{valido} se, in ogni situazione in cui le premesse sono tutte vere, anche la conclusione è vera
\end{definition}

La validità è strettamente connessa a un altro concetto fondamentale: quello di consistenza (o coerenza) di un insieme di enunciati.
Intuitivamente, un insieme di affermazioni è consistente se è possibile che siano tutte vere contemporaneamente, mentre è inconsistente se contiene una contraddizione che rende impossibile la loro verità simultanea.

\begin{definition}{Inconsistenza (informalmente)}
  Un insieme di enunciati è detto \keyword{inconsistente} se non esiste nessuna \qi{situazione} immaginabile in cui gli enunciati sono tutti veri
\end{definition}

È importante notare che l'inconsistenza è una proprietà strutturale, non empirica: dipende dalle relazioni logiche tra gli enunciati, non dal fatto che siano effettivamente veri o falsi nel mondo reale.

\begin{note}{}
  L'inconsistenza non riguarda la verità o falsità degli enunciati, ma la possibilità di rendere tutti gli enunciati veri contemporaneamente.
  Anche un insieme di enunciati tutti falsi può non essere inconsistente, se esiste una situazione in cui tutti gli enunciati sono veri.
\end{note}

Il concetto duale all'inconsistenza è naturalmente quello di consistenza, che caratterizza gli insiemi di enunciati privi di contraddizioni logiche.

\begin{definition}{Consistenza (informalmente)}
  Un insieme di enunciati è detto \keyword{consistente} se esiste almeno una \qi{situazione} immaginabile in cui gli enunciati sono tutti veri
\end{definition}

\begin{note}{}
  Essere consistenti significa non essere inconsistente.
\end{note}

Per comprendere meglio questi concetti, consideriamo un esempio concreto che illustra come un insieme di enunciati apparentemente innocui possa in realtà essere inconsistente.

\begin{example}{}
  Consideriamo l'insieme di enunciati:
  \begin{enumerate}
    \item O Carlo non dorme, o Carlo sogna\label{ex_item:inconsistency_1}
    \item Carlo dorme\label{ex_item:inconsistency_2}
    \item Se Carlo sogna allora Carlo è felice\label{ex_item:inconsistency_3}
    \item Carlo non è felice\label{ex_item:inconsistency_4}
  \end{enumerate}

  Questi enunciati sono chiaramente inconsistenti. Infatti, dato che Carlo dorme (\ref{ex_item:inconsistency_2}), deve necessariamente sognare (\ref{ex_item:inconsistency_1}). 
  Ma se Carlo sogna, allora deve essere felice (\ref{ex_item:inconsistency_3}), il che contraddice l'enunciato (\ref{ex_item:inconsistency_4}) secondo cui Carlo non è felice.
  Abbiamo quindi una contraddizione che rende l'insieme inconsistente.
\end{example}

I concetti di validità e consistenza che abbiamo introdotto ci permettono ora di definire in modo preciso la nozione centrale della logica: quella di \textit{conseguenza logica}.


Un enunciato è conseguenza logica di un insieme di premesse se segue necessariamente da esse, ovvero se non è possibile che le premesse siano tutte vere e la conclusione falsa.

Questa definizione può essere formulata elegantemente in termini di inconsistenza: un enunciato segue logicamente da un insieme di premesse se assumere la sua negazione, insieme alle premesse, porta a una contraddizione.

\begin{definition}{Conseguenza Logica}
  Dato un insieme \(\Gamma\) di enunciati e un enunciato \(\phi\), diremo che \(\phi\) è \keyword{conseguenza logica} di \(\Gamma\) se l'insieme \(\Gamma \cup \{\neg \phi\}\) è inconsistente. % NDR: Excluded middle assumed
\end{definition}

Si noti che, se $\phi$ è conseguenza logica di $\Gamma$, allora, secondo la definizione: se tutti gli enunciati in $\Gamma$ sono veri, allora $\neg \phi$ è necessariamente falso, ovvero $\phi$ è vero; se $\neg \phi$ è vero, allora almeno uno degli enunciati di $\Gamma$ dovrà essere falso. Il primo caso cattura l'idea della dimostrazione diretta, mentre il secondo caso cattura l'idea della dimostrazione per assurdo.
In particolare, vediamo come il ragionamento per assurdo si applica ad un caso concreto.

\begin{example}{}
  Dati gli enunciati \(\Gamma\):
  \begin{enumerate}
    \item Se Carlo è italiano e Giovani non è spagnolo, allora Rachele è tedesca\label{ex_item:logical_consequence_1}
    \item Se Rachele è tedesca allora Lucia è francese oppure Giovanni è spagnolo\label{ex_item:logical_consequence_2}
    \item Se Lucia non è francese allora Carlo è italiano\label{ex_item:logical_consequence_3}
    \item Giovanni non è spagnolo\label{ex_item:logical_consequence_4}
  \end{enumerate}
  L'enunciato \(\phi\):
  \begin{center}
    Lucia è Francese
  \end{center}
  è conseguenza logica di \(\Gamma\)?
  
  Per verificarlo, dobbiamo valutare se l'insieme \(\Gamma \cup \{\neg \phi\}\) è inconsistente.
  In altre parole, procediamo con una dimostrazione per assurdo.
  
  Assumiamo che Lucia non sia francese (\(\neg \phi\)).

  Allora, per l'enunciato (\ref{ex_item:logical_consequence_3}), Carlo è italiano. 

  Dato che Giovanni non è spagnolo (\ref{ex_item:logical_consequence_4}) e Carlo è italiano, dall'enunciato (\ref{ex_item:logical_consequence_1}) concludiamo che Rachele è tedesca.

  Ma se Rachele è tedesca, per l'enunciato (\ref{ex_item:logical_consequence_2}), Lucia è francese oppure Giovanni è spagnolo.
  Poiché Giovanni non è spagnolo (\ref{ex_item:logical_consequence_4}), deve necessariamente essere vero che Lucia è francese.
  Questo contraddice la nostra ipotesi iniziale che Lucia non sia francese.
  Abbiamo quindi ottenuto una contraddizione, il che dimostra che \(\Gamma \cup \{\neg \phi\}\) è inconsistente, e quindi \(\phi\) è conseguenza logica di \(\Gamma\).
\end{example}

Andiamo dunque a definire formalmente la nozione di \keyword{logica}, per andare a formalizzare i concetti visti finora in modo rigoroso e sistematico.

\begin{definition}{Logica}
  Una \keyword{logica} è un sistema formale caratterizzato da:
  \begin{itemize}
    \item Un insieme di simboli e regole sintattiche che definiscono un \keyword{linguaggio} formale, che ne delinea la sintassi e la struttura delle espressioni;
    \item Una \keyword{semantica}, che assegna significati alle espressioni del linguaggio;
    \item Un \keyword{calcolo}, che definisce le regole di inferenza.
  \end{itemize}

  Il \keyword{calcolo} funge in qualche modo da ponte tra la sintassi e la semantica, essendo un insieme di regole di manipolazione simbolica (trasformazioni sintattiche) che preservano certe proprietà semantiche, come la verità o la validità.
\end{definition}

Esempi di logiche sono:
\begin{itemize}
  \item Logica Proposizionale
  \item Logiche Modali (Tra cui la logica temporale)
  \item Logica del Primo Ordine
\end{itemize}

Durante il corso, ci concentreremo sulle logiche dette \textit{classiche}, che si basano su principi come il \textit{terzo escluso} e il \textit{non contraddittorio}, e che sono caratterizzate da una semantica bivalente (ogni proposizione è vera o falsa).
Tra queste in particolare, studieremo la logica proposizionale e la logica del primo ordine, che sono le fondamenta su cui si basano molte altre logiche più complesse.
\chapter{Logica Proposizionale}

Abbiamo visto, nel capitolo precedente, che una logica è caratterizzata da tre componenti fondamentali: un linguaggio formale, una semantica e un calcolo.
In questo capitolo, studieremo in dettaglio la \textit{logica proposizionale}, la più semplice tra le logiche classiche, che costituisce il punto di partenza per comprendere sistemi logici più complessi.
Procederemo sistematicamente, partendo dalla definizione del linguaggio (sintassi), per poi introdurre la semantica e infine studiare i sistemi di calcolo per la derivazione di conseguenze logiche.

\section{Sintassi della Logica Proposizionale}

Il primo passo nella costruzione di una logica formale è definire il suo linguaggio, anche detto alfabeto. Il linguaggio è l'insieme dei simboli da cui attingere per la creazione di espressioni.
In generale, un linguaggio logico è composto da due tipi di simboli: i \textit{simboli non logici}, su cui la logica non fa alcuna assunzione e il cui significato varia a seconda della \qi{situazione}, e i \textit{simboli logici}, i cui significati sono fissati dalla specifica logica.

\begin{definition}{Alfabeto della logica proposizionale}
  L'alfabeto della logica proposizionale è costituito da:
  \begin{itemize}
    \item Un insieme di \textit{simboli non logici}, detti \keyword{proposizioni atomiche} e denotati con \(\AP = \set{A,B,C,\ldots}\).
    \item Un insieme di \textit{simboli logici}, tra cui i connettivi logici \(\set{\land, \lor, \rightarrow, \neg,\leftrightarrow}\) e le parentesi \(\set{(,)}\) per disambiguare la struttura delle formule.
  \end{itemize}
\end{definition}

\begin{note}{}
L'insieme dei simboli non logici è potenzialmente infinito.
\end{note}

Un'\keyword{espressione} è una sequenza finita di simboli del linguaggio.
Tuttavia, non tutte le espressioni sono significative dal punto di vista logico: vogliamo considerare solo quelle che rispettano le regole grammaticali del linguaggio.

\begin{definition}{Formula ben formata (\(\WFF\))}
  Una formula ben formata (\(\WFF\)) è un'espressione che può essere costruita usando una grammatica definita sui simboli del linguaggio.
\end{definition}

Nel caso specifico della logica proposizionale, l'insieme \(\WFF\) viene indicato con \(PL\), che sta per \textit{Propositional Logic}.
In generale, l'insieme delle formule ben formate di una logica viene indicato con le iniziali di tale logica. 

La definizione precisa di quali espressioni siano formule ben formate nella logica proposizionale è data ricorsivamente attraverso le seguenti regole di formazione.

\begin{definition}{\(\WFF\) Proposizionale}
  Nella logica proposizionale, una formula è ben formata se è costruita secondo le seguenti regole:
\begin{itemize}
  \item \(\AP \subseteq \WFF\)
  \item Se \(\phi \in \WFF\), allora \(\neg \phi \in \WFF\)
  \item Se \(\phi,\psi \in \WFF\), allora \((\phi \odot \psi)\in \WFF\), dove \(\odot \in \set{\land, \lor, \rightarrow, \leftrightarrow}\) 
  \item Nessuna altra espressione è una formula ben formata
\end{itemize}
  
\end{definition}

Come si può notare, la definizione è una grammatica libera da contesto, e, di conseguenza, ogni formula ben formata può essere rappresentata come un albero sintattico, in cui i nodi interni corrispondono a connettivi logici e le foglie a proposizioni atomiche.
Tale albero sintattico cattura la struttura gerarchica della formula, mostrando come essa sia stata costruita a partire dalle sue componenti più semplici.

\begin{example}{Albero sintattico}
  Consideriamo la formula \(\phi = (A \rightarrow (B \land \neg C))\). 
  
  Secondo la definizione ricorsiva, questa formula è ben formata: a partire dalle variabili proposizionali \(A\), \(B\), \(C\) (che, per la prima regola, sono \(\WFF\)), costruiamo \(\neg C\) (seconda regola), poi \((B \land \neg C)\) (terza regola), e infine \((A \rightarrow (B \land \neg C))\) (terza regola).
  
  L'albero sintattico di \(\phi\) rappresenta questa costruzione gerarchica:
  
  \begin{center}
    \begin{forest}
      [\(\rightarrow\)
        [\(A\)]
        [\(\land\)
          [\(B\)]
          [\(\neg\)
            [\(C\)]
          ]
        ]
      ]
    \end{forest}
  \end{center}
  
  In questo albero:
  \begin{itemize}
    \item Le foglie sono le variabili proposizionali \(A\), \(B\), \(C\)
    \item I nodi interni sono i connettivi logici \(\rightarrow\), \(\land\), \(\neg\)
    \item La struttura dell'albero riflette l'ordine di applicazione delle regole di costruzione delle \(\WFF\)
  \end{itemize}
\end{example}

Tale costruzione ci permette di utilizzare un potente strumento di ragionamento: l'\textit{induzione strutturale}.
L'induzione strutturale consente di dimostrare proprietà delle formule procedendo per casi: il caso base verifica che la proprietà in esame sia valida per le formule atomiche (variabili proposizionali), mentre i passi induttivi verificano la proprietà su formule costruite a partire da sottoformule più semplici.

A differenza dell'induzione matematica semplice sui numeri naturali, dove per dimostrare \(P(n)\) assumiamo solo \(P(n-1)\), l'induzione strutturale è basata sul principio di \textit{induzione completa} (o forte), secondo cui, per dimostrare una proprietà \(P(n)\) assumiamo che valga \(P(m), \forall m < n\). 
Nel caso dell'induzione strutturale, per dimostrare una proprietà di \(\phi\), assumiamo che tale proprietà valga per ogni sottoformula propria di \(\phi\).
L'induzione strutturale su una formula è equivalente all'induzione completa sull'altezza del relativo albero sintattico. Poiché l'altezza di un albero sintattico fornisce un ordine ben fondato\footnote{L'altezza di un albero indica il numero massimo di nodi lungo un percorso dalla radice a una foglia. L'altezza è un intero non negativo: il minimo è \(0\), ovvero l'altezza di una foglia. Essendo l'insieme degli interi non negativi, assieme alla relazione d'ordine \(\leq\), ben fondato, anche l'insieme delle altezze, con la stessa relazione d'ordine, è ben fondato.}, anche l'induzione strutturale è ben fondata.

La definizione ricorsiva di formula ben formata impone che i sottoalberi di un albero sintattico siano a loro volta alberi sintattici.
Tali sottoalberi rappresentano le \keyword{sottoformule} di una formula, il cui insieme è definito come segue.

\begin{definition}{Sottoformule}
  Dato \(\phi \in WFF\), l'\keyword{insieme delle sottoformule} di \(\phi\), denotato con \(\Sub{\phi}\), è definito ricorsivamente come segue:
  \[\Sub{\phi} = \begin{cases}
    \{\phi\} & ,\phi \in \AP\\
    \Sub{\phi'} \cup \{\neg \phi'\} & ,\phi = \neg \phi'\\
    \Sub{\phi'} \cup \Sub{\psi'} \cup \{(\phi' \odot \psi')\} & ,\phi = (\phi' \odot \psi')
  \end{cases}
  \]
\end{definition}

Possiamo ora osservare una prima applicazione dell'induzione strutturale, dimostrando una proprietà fondamentale delle formule ben formate: il numero di sottoformule di una formula \(\phi\) è al più pari alla lunghezza della formula stessa.

\begin{note}{}
  Tale proprietà evidenzia l'importanza delle regole di formazione delle WFF, differenziandole da generiche successioni di simboli.
  Infatti, data una stringa qualsiasi, il numero di sottostringhe è quadratico rispetto alla lunghezza della stringa\footnote{Si prenda ad esempio la formula \(ABC\), che non è ben formata. In essa, troviamo una sola sottostringa di lunghezza \(3\) (\(ABC\)), due sottostringhe di lunghezza \(2\) (\(AB,BC\)), tre sottostringhe di lunghezza \(1\) (\(A,B,C\)). Si vede quindi che, per una stringa qualsiasi di lunghezza \(n\), il numero di sue sottoformule è \(n(n-1)/2=\mathcal{O}(n^{2})\).}, mentre, presa una WFF, il numero di sottoformule è solamente lineare rispetto alla lunghezza della formula.
\end{note}

\begin{proposition}[namedlabel={prop:subformula-upperbound}{Limite superiore al numero di sottoformule}]{Limite superiore al numero di sottoformule}
  Sia \(\phi \in WFF\). Siano inoltre \(\abs{\phi}\) la lunghezza della formula \(\phi\) in numero di simboli e \(\abs{\Sub{\phi}}\) la cardinalità dell'insieme delle sottoformule di \(\phi\). Allora
  \[\abs{\Sub{\phi}}\leq \abs{\phi}\]
\end{proposition}

\begin{proof}
  Per dimostrare questa proposizione, procediamo per induzione strutturale sulla formula \(\phi\).
  \begin{description}
    \item[\(\phi \in \AP\)]: In questo caso, \(\phi\) è una formula atomica, quindi \(\Sub{\phi} = \{\phi\}\) e \(\abs{\Sub{\phi}} = 1\). Dato che \(\abs{\phi} = 1\) (ogni simbolo ha lunghezza 1), abbiamo
          \[\abs{\Sub{\phi}} = 1 \leq 1 = \abs{\phi}\]
    \item[\(\phi = \neg \phi'\)]: In questo caso, per induzione strutturale, assumiamo che \(\abs{\Sub{\phi'}} \leq \abs{\phi'}\). Dato che \(\Sub{\phi} = \Sub{\phi'} \cup \{\neg \phi'\}\), abbiamo
          \[\abs{\Sub{\phi}} = \abs{\Sub{\phi'}} + 1\]
          Inoltre, \(\abs{\phi} = \abs{\phi'} + 1\) (perché aggiungiamo il simbolo di negazione). Quindi,
          \[\abs{\Sub{\phi}} = \abs{\Sub{\phi'}} + 1 \leq \abs{\phi'} + 1 = \abs{\phi}\]
    \item[\(\phi = (\phi' \odot \psi')\)]: Infine, sempre per induzione strutturale, assumiamo che \(\abs{\Sub{\phi'}} \leq \abs{\phi'}\) e \(\abs{\Sub{\psi'}} \leq \abs{\psi'}\). Dato che \(\Sub{\phi} = \Sub{\phi'} \cup \Sub{\psi'} \cup \{(\phi' \odot \psi')\}\), abbiamo
          \[\abs{\Sub{\phi}} \leq \abs{\Sub{\phi'}} + \abs{\Sub{\psi'}} + 1\]
          In questo caso non è detto valga l'uguaglianza, perché è possibile che \(\Sub{\phi'}\) e \(\Sub{\psi'}\) non siano disgiunti.
          Inoltre, \(\abs{\phi} = \abs{\phi'} + \abs{\psi'} + 1\) (perché aggiungiamo il connettivo). Quindi partendo da:
          \[\abs{\Sub{\phi'}} \leq \abs{\phi'}\]
          sommiamo \(\abs{\Sub{\psi'}}\) ad ambo i membri dell'disequazione:
          \[\abs{\Sub{\phi'}} +\abs{\Sub{\psi'}} \leq \abs{\phi'} + \abs{\Sub{\psi'}} \leq \abs{\phi'} + \abs{\psi'}\]
          e dunque
          \[\abs{\Sub{\phi}} \leq \abs{\Sub{\phi'}} + \abs{\Sub{\psi'}} + 1 \leq \abs{\phi'} + \abs{\psi'} + 1 = \abs{\phi}\]
  \end{description}
\end{proof}

\section{Semantica Proposizionale tramite Modelli}

Definita la sintassi della logica proposizionale, diamo una definizione formale di semantica. La \keyword{semantica} della logica proposizionale formalizza il concetto intuitivo di \q{situazione}, che determina quali proposizioni elementari sono vere e quali false.
Un modo per definire la semantica è attraverso il concetto di \keyword{modello}.
Un modello, in generale, è una struttura matematica che interpreta i simboli del linguaggio e stabilisce le condizioni di verità per le formule.
Nel caso della logica proposizionale, una possibile implementazione di modello è mediante insiemi, per cui un modello è un qualsiasi sottoinsieme di \(\AP\).

\begin{definition}{Modello per la logica proposizionale}
  In logica proposizionale, un modello \(m\) è un qualsiasi sottoinsieme di \(\AP\). In altre parole, \(m \in 2^{\AP}\).\footnote{\(2^{\AP}\) non è altro che una notazione alternativa per l'insieme di tutti i sottoinsiemi di \(\AP\), chiamato anche insieme potenza indicato anche con \(\mathcal{P}(\AP)\).}
\end{definition}

Essendo la logica proposizionale molto semplice, per darle significato è sufficiente una struttura altrettanto semplice: un insieme di proposizioni atomiche che consideriamo vere.
Infatti, un modello non è altro che un insieme di proposizioni elementari tale per cui una proposizione è vera se e soltanto se appartiene al modello.
Affinché un modello sia ben definito su formule più complesse, è necessaria una regola ricorsiva per il soddisfacimento.

\begin{note}{}
Essendo \(\AP\) potenzialmente infinito, anche i modelli possono essere infiniti. Da un punto di vista matematico non è particolarmente importante, ma modelli infiniti portano a problemi computazionali di rappresentazione e manipolazione.
\end{note}

\begin{definition}{Soddisfacimento di una formula da parte di un modello}
  Dato un modello \(m \in 2^{\AP}\) e una formula \(\phi \in \PL\), definiamo la \keyword{relazione di soddisfazione}, \(\models \subseteq 2^{\AP} \times \PL\), come segue:

  \(m\models \phi \) se e solo se vale uno dei seguenti casi:
  \begin{enumerate}
    \item \(\phi \in m\) se \(\phi \in \AP\).
    \item \(m \not\models \phi'\) se \(\phi = \neg \phi'\).
    \item \(m \models \phi_1\) e \(m \models \phi_2\) se \(\phi = \phi_1 \land \phi_2\).
    \item \(m \models \phi_1\) o \(m \models \phi_2\) se \(\phi = \phi_1 \lor \phi_2\).
    \item Se \(m \models \phi_1\) allora \(m \models \phi_2\) se \(\phi = \phi_1 \rightarrow \phi_2\).
    \item \(m \models \phi_1\) se e solo se \(m \models \phi_2\) se \(\phi = \phi_1 \leftrightarrow \phi_2\).
  \end{enumerate}
\end{definition}

Tale definizione rappresenta una procedura ricorsiva per determinare se una formula è soddisfatta da un modello, partendo dalle proposizioni atomiche e risalendo l'albero sintattico della formula.

Il problema decisionale di verificare se \(m \models \phi\) è detto \keyword{model checking}, e, a patto di risolvere il problema di rappresentazione di modelli potenzialmente infiniti, è facilmente decidibile\footnote{Un problema è decidibile quando per esso esiste un algoritmo, la cui terminazione è garantita, che lo risolve. In particolare, un problema è facilmente decidibile quando esiste un tale algoritmo che abbia complessità polinomiale.}, poiché equivale a un'esplorazione in post-ordine dell'albero sintattico di \(\phi\), che ha un numero di nodi pari al numero di sottoformule di \(\phi\), che è al più pari alla lunghezza di \(\phi\).

Avendo finalmente definito la semantica possiamo andare a formalizzare le definizioni informali date nel primo capitolo.
\begin{definition}{Soddisfacibilità}
  Una formula \(\phi\) è \keyword{soddisfacibile} (\(\Sat[\phi]\) o \(\phi \in \Sat\)) se e soltanto se esiste un modello \(m\in 2^{\AP}\) tale che \(m \models \phi\).
\end{definition}

Il problema decisionale di verificare \(\Sat[\phi]\) è detto \keyword{soddisfacibilità} ed è un problema centrale in logica e informatica teorica, poiché molti problemi di interesse, come pianificazione delle azioni, colorazione dei grafi, etc, possono essere ridotti a esso. 
Ciò significa che, se riuscissimo a risolvere il problema di soddisfacibilità in modo efficiente, potremmo risolvere anche molti altri problemi in modo efficiente.
Rispetto al model checking, il problema di soddisfacibilità è spesso più complesso, poiché richiede l'esplorazione dello spazio dei modelli, che cresce esponenzialmente con il numero di proposizioni atomiche coinvolte nella formula\footnote{Se una formula è composta da una sola proposizione elementare \(A\), allora sono necessari due modelli per esplorare tutte le situazioni possibili: uno in cui \(A\) è vera, uno in cui \(A\) è falsa. Se \(\phi\) è composta da due formule elementari \(A\) e \(B\), allora, per esplorare tutte le situazioni possibili, ovvero tutte le combinazioni di verità e falsità delle di \(A\) e \(B\), sono necessari quattro modelli. Se le formule sono tre, servono otto modelli e così via. In generale, per descrivere tutti casi possibili per una formula con \(n\) proposizioni elementari, sono necessari \(2^{n}\) modelli.}. 
Più semplicemente, il problema di soddisfacibilità è più complesso perché costituito da tanti problemi di model checking per quanti sono i modelli possibili.

\begin{definition}{Validità}
  Una formula \(\phi\) è \keyword{valida} (\(\Valid[\phi]\) o \(\phi \in \Valid\)) se e soltanto se per ogni modello \(m\in 2^{\AP}\) vale \(m \models \phi\).
\end{definition}

La validità e la soddisfacibilità sono concetti duali, nel senso che \(\Sat[\phi]\iff \neg \Valid[\neg \phi]\) e \(\Valid[\phi] \iff \neg \Sat[\neg \phi]\).\footnote{I simboli \(\iff\) e \(\implies\) sono simboli di \textbf{metalinguaggio}. Tali simboli non appartengono al linguaggio oggetto del discorso (in questo caso la logica proposizionale), ma sono simboli sono utilizzati (con l'interpretazione usuale) per parlare di proprietà del linguaggio stesso.}
Questo è diretta conseguenza del fatto che, data una qualsiasi proprietà \(P\), \(\exists x:P(x) \iff \forall x, \neg P(\neg x)\). Ciò verrà dimostrato formalmente quando tratteremo la logica del primo ordine.

\begin{definition}{Estensione della soddisfacibilità}
  Dato \(\Gamma \subseteq \PL\) e \(m \in 2^{\AP}\), diremo che \(m\) soddisfa \(\Gamma\) (\(m \models \Gamma\)) se e soltanto se \(m \models \phi\) per ogni \(\phi \in \Gamma\).
\end{definition}
Ciò consente una scrittura più compatta per intendere che un modello soddisfa un insieme di formule.

\begin{definition}{Conseguenza logica}
  Dato \(\Gamma \subseteq \PL\), una formula \(\phi\) è conseguenza logica di \(\Gamma\) (\(\Gamma \models \phi\)) se e soltanto se per ogni \(m \in 2^{\AP}\) che soddisfa \(\Gamma\) allora \(m \models \phi\).

  Alternativamente, \(\Gamma \models \phi\) se e soltanto se non esiste \(m \in 2^{\AP}\) tale che \(m \models \Gamma\) e \(m \not\models \phi\).
\end{definition}

\begin{note}{}
  Nonostante si usi la stessa notazione \(\models\) per indicare sia la soddisfacibilità che la conseguenza logica, è importante sottolineare che si tratta di due concetti distinti, anche se tra loro molto legati.
  La soddisfacibilità \(m \models \phi\) è una relazione contenuta in \(2^{\AP}\times \PL\), mentre la conseguenza logica \(\Gamma \models \phi\) è una relazione contenuta in \(2^{\PL}\times \PL\).
\end{note}

Dalla definizione di conseguenza logica segue che se \(\emptyset \models \phi\), denotato alternativamente con \(\models \phi\), allora \(\phi\in\Valid\).
Se consideriamo \(\Gamma\) come l'insieme dei vincoli a cui il modello è sottoposto, allora è chiaro che, se \(\Gamma\) è vuoto, e quindi non vi sono vincoli, allora ogni modello lo soddisfa.
Quindi, la relazione di conseguenza logica dall'insieme vuoto si riduce alla soddisfacibilità di \(\phi\) da parte di ogni modello, ovvero la validità di \(\phi\).
In altre parole, potendo risolvere il problema decisionale di conseguenza logica, potremmo risolvere quello di validità.

Inoltre, se \(\Gamma\) è un insieme finito di formule, si verifica anche il viceversa, ovvero è possibile risolvere il problema decisionale di validità per risolvere quello di conseguenza logica. Infatti, in tal caso, è possibile esprimere \(\Gamma \models \phi\), ovvero il concetto di conseguenza logica di \(\phi\) rispetto a \(\Gamma\), in termini di validità della formula
\[\left(\bigwedge_{\gamma \in \Gamma} \gamma\right) \rightarrow \phi.\]

Ciò, oltre a evidenziare la stretta connessione tra conseguenza logica e validità (e quindi soddisfacibilità, per dualità), sottolinea come la relazione \(\models\) sia correlata alla semantica del simbolo di implicazione \(\rightarrow\), cosa facilmente evincibile dall'analogia tra definizione di conseguenza logica e di semantica dell'implicazione.

\section{Semantica Proposizionale tramite Assegnamenti}

Definito formalmente un primo concetto di semantica, andiamo a risolvere il problema della sua rappresentabilità definendo un altro tipo di semantica, che dimostreremo essere equivalente a quella precedente.

La semantica mediante modelli è semplice ed elegante. 
Tuttavia, data una formula soddisfacibile, esistono infiniti modelli che la soddisfano: un modello potrebbe contenere anche proposizioni elementari non contenute in tale formula.
L'infinitezza dei modelli è un problema non trascurabile per la computazione.
Per risolverlo, possiamo semplicemente limitare la semantica alle sole proposizioni atomiche che compaiono nella formula d'interesse.

\begin{definition}{Proposizioni atomiche di una formula}
  Data \(\phi \in \PL\), possiamo definire \(\AP[\phi] = \Sub{\phi} \cap \AP\) l'insieme delle proposizioni atomiche contenute in \(\phi\).
\end{definition}

Essendo \(\abs{\Sub{\phi}}\leq \abs{\phi}\), si ha \(\abs{\AP[\phi]}\leq \abs{\phi}\). Essendo \(\abs{\phi}\) finito, anche \(\abs{\AP[\phi]}\) lo sarà.

\begin{definition}{Assegnamento}
  Dato \(\phi \in \PL\), un assegnamento \(\alpha\) per \(\phi\) è una funzione \(\alpha: \Delta_\phi \rightarrow \{0,1\}\), tale che \(\Delta_\phi\) è finito e \(\AP[\phi] \subseteq \Delta_\phi \subseteq \AP\).\footnote{Il dominio \(\Delta_{\phi}\) di un assegnamento è un qualsiasi insieme di proposizioni che contenga le proposizioni elementari di \(\phi\). Un assegnamento per cui \(\Delta_{\phi}=\AP[\phi]\) è un assegnamento minimale, ovvero un assegnamento che copre solo le proposizioni strettamente necessarie.}

  Chiameremo l'insieme di tali assegnamenti \(\Ass{\phi}\).\footnote{Analogamente alla notazione \(2^S\) per indicare l'insieme \(\mathcal{P}(S)\), una notazione molto in uso per indicare un insieme di funzioni \(f: D \to C\) è \(C^D\) essendo che la cardinalità di tale insieme (per insiemi finiti) è proprio \(\abs{C}^{\abs{D}}\). Di conseguenza, una notazione alternativa per \(\Ass{\phi}\) è \(\set{0,1}^{\Delta_\phi}\)}
\end{definition}

Come visto nella definizione, data una formula \(\phi\), l'insieme degli assegnamenti compatibili è finito.
In particolare, essi ne sono \(\abs{\set{0,1}}^{\abs{\AP[\phi]}}=2^{\abs{\AP[\phi]}}\).
Tali funzioni di assegnamento ci permettono di assegnare il valore di verità a ciascuna proposizione atomica considerata. 
Tuttavia, per calcolare il valore di verità di una formula complessa, definiamo ora la sua \keyword{valutazione} in modo ricorsivo, a partire dall'assegnamento, un po' come avevamo fatto per la soddisfacibilità nella semantica tramite modelli.

\begin{definition}{Valutazione di una formula}
  Siano \(\phi \in \PL\) e \(\alpha \in \Ass{\phi}\). Definiamo dunque la \keyword{valutazione basata su \(\alpha\)} come una funzione \(\Val{\alpha}: \PL \to \set{0,1}\) come\footnote{Per essere precisi, \(\Val{}\) ha una dipendenza implicita dall'assegnamento, e quindi sarebbe più corretto descriverla come una funzione \(\Val{}: \Ass{\phi} \to (\PL \to \set{0,1})\). Nel pratico però, l'assegnamento sarà fissato a priori.}:
  \[\Val{\alpha}[\phi] = \begin{cases}
  \alpha(\phi) &, \phi \in \AP[\phi] \\
  1 - \Val{\alpha}(\phi') &, \phi = \neg \phi' \\
  \min\set{\Val{\alpha}[\phi_1], \Val{\alpha}[\phi_2]} &, \phi = \phi_1 \land \phi_2 \\
  \max\set{\Val{\alpha}[\phi_1], \Val{\alpha}[\phi_2]} &, \phi = \phi_1 \lor \phi_2 \\
  1 &, \phi = \phi_1 \rightarrow \phi_2 \text{ e } \Val{\alpha}[\phi_1] \leq \Val{\alpha}[\phi_2] \\
  0 &, \phi = \phi_1 \rightarrow \phi_2 \text{ e } \Val{\alpha}[\phi_1] > \Val{\alpha}[\phi_2]
  \end{cases}\]
\end{definition}

A questo punto, verificare se una certa formula assume valore \(0\) o \(1\) rispetto a un assegnamento compatibile diventa un problema facilmente decidibile. 
Infatti, la valutazione di una formula \(\phi\) ha complessità \(\BigTheta{\abs{\phi}}\).

Per poter dire di aver risolto definitivamente il problema dell'infinità dei modelli, bisogna dimostrare che questa nuova semantica, basata sugli assegnamenti, sia equivalente a quella basata su modelli. 
In altre parole, dobbiamo mostrare che a ogni modello corrisponde un assegnamento equivalente, e viceversa, e che la relazione di soddisfacimento \(m \models \phi\) è equivalente alla valutazione \(\Val{\alpha}[\phi] = 1\).
Per fare ciò, andremo a mostrare che una formula è tautologica se e solo se è valida. 
Facendo ciò, otterremo anche un modo per correlare \(m \models \phi\) e \(\Val{\alpha}[\phi] = 1\).
Prima, però, definiamo il concetto di \keyword{tautologia}.

\begin{definition}{Tautologia}
  Una formula \(\phi\) è una \keyword{tautologia} (\(\Taut[\phi]\)) quando, per ogni assegnamento compatibile \(\alpha\), si ha \(\Val{\alpha}[\phi]=1\).
\end{definition}

\begin{theorem}[label=thm:validity-equals-taut]{Equivalenza tra validità e tautologicità}
  Data \(\phi \in \PL\), \(\models \phi \) se e solo se \(\Taut[\phi]\).
\end{theorem}

In precedenza, abbiamo osservato che \(\models \phi \iff \forall m, m \models \phi\).
Per definizione di tautologia, si ha che
\(\Taut[\phi]\iff\forall \alpha ,\Val{\alpha}[\phi] = 1\).
In altre parole, per dimostrare il teorema, basta mostrare che \(\forall m, m \models \phi \iff \forall \alpha, \Val{\alpha}[\phi] =1\).

\begin{lemma}[listtitle={Estrapolato dalla dimostrazione del \Cref{thm:validity-equals-taut}}]{\infosymbol{} Estrapolato dalla dimostrazione del \Crefalt[yellow]{thm:validity-equals-taut}}
  Dato un modello \(m \in 2^{AP}\), esiste un \(\alpha \in Ass(\phi)\) tale per cui \(m \models \phi \iff \Val{\alpha}[\phi]=1\).
\end{lemma}

\begin{proof}
  Dato \(m\), possiamo definire un assegnamento \(\alpha_{m}\) tale che
  \[\alpha_{m}(P) = \begin{cases}
    1, & P \in m\\
    0, & \text{ altrimenti}
  \end{cases}\]
  Possiamo mostrare quindi, per induzione strutturale, che \(m\models \phi \iff \Val{\alpha_m}[\phi] = 1\).
  \begin{description}
    \item[\(\phi \in \AP\):]
      \[m \models \phi \iff \phi \in m \eqdef {\iff}\alpha_m(\phi) = 1 \iff \Val{\alpha_m}[\phi] = 1\]
    \item[\(\phi = \neg \phi'\):]
      \[m \models \phi \iff m \not\models \phi' \stackrel{ind. hp.}{\iff} \Val{\alpha_m}[\phi'] = 0\]
      Ma allora dalla definizione di valutazione, \(\Val{\alpha_m}[\phi] = 1-\Val{\alpha_m}[\phi'] = 1-0 = 1\).
    \item[\(\phi = \phi_1 \land \phi_2\):]
      \[m \models \phi \iff m \models \phi_1 \text{ e } m \models \phi_2 \overset{ind. hp.}{\iff} \Val{\alpha_m}[\phi_1] = 1 \text{ e } \Val{\alpha_m}[\phi_2] = 1\]
      Ma allora dalla definizione di valutazione, \(\Val{\alpha_m}[\phi] = \min\set{\Val{\alpha_m}[\phi_1], \Val{\alpha_m}[\phi_2]} = \min\set{1, 1} = 1\).
    \item[\(\phi = \phi_1 \lor \phi_2\):]
      \[m \models \phi \iff m \models \phi_1 \text{ o } m \models \phi_2 \overset{ind. hp.}{\iff} \Val{\alpha_m}[\phi_1] = 1 \text{ o } \Val{\alpha_m}[\phi_2] = 1\]
      Ma allora dalla definizione di valutazione, \(\Val{\alpha_m}[\phi] = \max\set{\Val{\alpha_m}[\phi_1], \Val{\alpha_m}[\phi_2]} = \max\set{1, 0} = 1\).
    \item[\(\phi = \phi_1 \rightarrow \phi_2\):] 
      \[m \models \phi \iff \text{ se } m \models \phi_1 \text{ allora } m \models \phi_2\]
      Per ipotesi induttiva, se \(m \not\models \phi_1\) allora \(\Val{\alpha_m}[\phi_1] = 0\), e quindi a prescindere dal valore di \(\Val{\alpha_m}[\phi_2]\) si ha \(\Val{\alpha_m}[\phi] = 1\), essendo che \(0 \leq \Val{\alpha_m}[\phi_2]\).
      Se invece \(m \models \phi_1\) allora \(\Val{\alpha_m}[\phi_1] = 1\), ma sempre per ipotesi induttiva \(m \models \phi_2\) e quindi \(\Val{\alpha_m}[\phi_2] = 1\), e dunque \(\Val{\alpha_m}[\phi] = 1\), essendo che \(1 \leq 1\).
  \end{description}
\end{proof}

\begin{note}{}
  Tale dimostrazione e quella della \Namedcref{prop:subformula-upperbound} sebbene entrambe basate sull'induzione strutturale, sono molto diverse tra loro. Nella dimostrazione del \Namedonlyhref{prop:subformula-upperbound} infatti è stato possibile \qi{accorpare} i casi di operatori binari, mentre nella dimostrazione del lemma è stato necessario trattare separatamente ogni connettivo, a causa della diversa semantica di ciascuno di essi.

  Questo perché la proposizione riguardava una proprietà \textbf{sintattica} delle formule, per la quale ogni simbolo non ha altra caratterizzazione che la sua presenza nella grammatica, mentre il lemma riguarda una proprietà \textbf{semantica} delle formule, per la quale ogni simbolo ha una caratterizzazione specifica che deve essere rispettata.

  Tale differenza tra proprietà sintattiche e semantiche è fondamentale per comprendere la natura dei risultati che andremo a dimostrare, oltre a essere un fondamentale aiuto per prevedere che tipo di dimostrazione sarà necessario per dimostrare un certo risultato.
\end{note}

\begin{lemma}[listtitle={Estrapolato dalla dimostrazione del \Cref{thm:validity-equals-taut}}]{\infosymbol{} Estrapolato dalla dimostrazione del \Crefalt[yellow]{thm:validity-equals-taut}}
  Dato un assegnamento \(\alpha \in Ass(\phi)\), esiste un \(m \in 2^{AP}\) tale per cui \(m \models \phi \iff \Val{\alpha}[\phi]=1\).
\end{lemma}

\begin{proof}
  Dato un assegnamento \(\alpha\), è possibile definire un modello \(m_{\alpha}\) t.c.
  \[m_\alpha = \set{P \in \AP[\phi]}[\alpha(P) = 1]\]
  Si dimostra per induzione strutturale, in maniera del tutto analoga alla dimostrazione del precedente lemma, che \(\Val{\alpha}[\phi] = 1 \iff m_\alpha \models \phi\).
\end{proof}

Procediamo adesso, per contrapposizione, alla dimostrazione del teorema, ovvero a dimostrare \(\not\models \phi \iff \neg \Taut(\phi)\).
\begin{proof}[Dimostrazione (\ref{thm:validity-equals-taut})]
  \begin{description}
    \item[\qi{\(\implies\)}]
      Sia \(m\) tale che \(m \models \neg \phi\).
      Costruendo \(\alpha_{m}\) secondo il lemma, si ha:
      \[m\models \neg \phi \iff \Val{\alpha_{m}}[\neg\phi]=1 \iff \Val{\alpha_{m}}[\phi]=0\]
      Allora, abbiamo trovato un assegnamento \(\alpha_{m}\) tale per cui la valutazione di \(\phi\) è \(0\). Quindi, \(\neg\Taut[\phi]\).
    \item[\qi{\(\impliedby\)}]
      Sia \(\alpha\) un assegnamento per \(\phi\) tale che \(\Val{\alpha}[\phi]=0\).
      Costruiamo \(m_{\alpha}\) secondo il lemma.
      Allora, si ha:
      \[\Val{\alpha}[\neg\phi]=1 \iff m \not\models \phi\]
      Allora, abbiamo trovato un modello \(m_{\alpha}\) che non soddisfa \(\phi\). Quindi, \(\not\models \phi\).
  \end{description}
\end{proof}

Questo teorema ci dice che risolvere \(\Taut\) equivale a risolvere \(\Valid\) e che, quindi, essendo \(\Taut\) computabile, in quanto una ricerca in un insieme finito, anche \(\Valid\) lo è.
Il costo computazionale di \(\Valid\) è però esponenziale, in quanto \(\abs{\Ass{\phi}} = 2^{\abs{\phi}}\) e dunque \(\Valid\) ha complessità \(\BigO{\abs{\phi}2^{\abs{\phi}}}\leq 2^{\BigTheta{\abs{\phi}}}\).

\begin{example}[namedlabel={ex:notable_tautologies}{Tautologie notevoli}]{Tautologie notevoli}
  Vediamo alcune tautologie notevoli che ci saranno successivamente utili per la dimostrazione di altri risultati.
  \begin{enumerate}
    \item \(P \lor \neg P\) (legge del terzo escluso).
    \item \(\neg (P \land \neg P)\) (principio di non contraddizione \(\bot\))
    \item \(P \lor Q \leftrightarrow \neg(\neg P \land \neg Q)\) (legge di De Morgan)
    \item \(P\land Q \leftrightarrow \neg(\neg P \lor \neg Q)\) (legge di De Morgan)
    \item \(\neg \neg P \iff P\) (eliminazione della doppia negazione)
    \item \(P \land(Q \lor R) \leftrightarrow (P \land Q) \lor (P \land R)\) (distributività)
    \item \((P \lor(Q\land R)) \leftrightarrow (P \lor Q) \land (P \lor R)\) (distributività)
    \item \((P \rightarrow Q) \leftrightarrow (\neg P \lor Q)\) (implicazione materiale)
  \end{enumerate}
\end{example}

\section{Sostituzioni nella Logica Proposizionale}

Per avvicinarci gradualmente verso la formalizzazione del concetto di \q{ragionamento}, o deduzione, anticipato all'inizio del corso, introduciamo una prima forma di manipolazione delle formule: la sostituzione, che permette di costruire nuove formule a partire da formule già esistenti, preservando alcune proprietà semantiche.

\begin{definition}{Sostituzione}
  Definiamo una \keyword{sostituzione} come una funzione \(\sigma: \AP \to \PL\).
\end{definition}

La sostituzione permette di sostituire ogni proposizione atomica con una formula arbitraria, e, dunque, di costruire nuove formule a partire da formule già esistenti. Procediamo, adesso, con la definizione formale.

\begin{definition}{Applicazione di una sostituzione a una formula}
  Dato \(\phi \in \PL\) e una sostituzione \(\sigma\), definiamo l'applicazione di \(\sigma\) a \(\phi\), denotata con \(\phi\sigma\), come segue:
  \begin{itemize}
   \item se \(\phi \in \AP\) allora \(\phi\sigma = \sigma(\phi)\)
   \item se \(\phi = \neg \phi'\) allora \(\phi\sigma = (\neg \phi ')\sigma = \neg (\phi'\sigma)\)
   \item se \(\phi = \phi_1 \odot \phi_2\) allora \(\phi\sigma = (\phi_1 \odot \phi_2)\sigma = (\phi_1\sigma) \odot (\phi_2\sigma)\)
  \end{itemize}
\end{definition}

\begin{note}{}
  È importante sottolineare che la sostituzione è da intendersi come applicata simultaneamente a tutte le proposizioni atomiche. Se una proposizione atomica \(P\) compare sia come simbolo da sostituire che come simbolo di una formula sostituita, la formula sostituita a \(P\) non viene ulteriormente sostituita anche internamente alle formule sostituite contenenti \(P\).
\end{note}

\begin{example}{Sostituzioni e loro applicazione}  
  Consideriamo \(\sigma = \set{A \leftarrow (A\lor B), B \leftarrow \neg A, C \leftarrow C}\) e \(\phi = (A \land B) \lor C\).
  
  Allora, \(\phi\sigma\) sarà:
  \[((A\lor B) \land \neg A) \lor C\]

  e non:
  \[((A\lor \neg A)\land \neg A)\lor C\]

  Nel secondo caso infatti si è applicata la sostituzione per \(B\) anche all'interno della formula sostituita a \(A\). Tale comportamento è quello che si sarebbe ottenuto considerando un ulteriore sostituzione \(\tau = \set{B\leftarrow \neg A}\) e applicando \((\phi\sigma)\tau\).
\end{example}

Per mostrare come la sostituzione sia una prima forma di deduzione logica molto elementare, e che quindi preservi alcune proprietà semantiche, andiamo a enunciare e successivamente dimostrare il \Namedonly{thm:substitution}.

\begin{theorem}[namedlabel={thm:substitution}{Teorema di Sostituzione}]{Teorema di Sostituzione}
  Sia \(\phi\in \Taut\) e \(\sigma\) un'arbitraria sostituzione. Allora \(\phi\sigma\in \Taut\).
\end{theorem}

Per dimostrare questo teorema, è necessario dare carattere semantico all'operazione di sostituzione, altrimenti esclusivamente sintattica, e mostrare che questa operazione preserva la proprietà di essere una tautologia.

\begin{definition}{Semantica della sostituzione}
  Data \(\phi \in \PL\), una sostituzione \(\sigma\) e un assegnamento \(\alpha \in \Ass{\phi\sigma}\), definiamo l'assegnamento \(\alpha^\sigma \in \Ass{\phi}\) come:
  \[\alpha^\sigma(P)=\Val{\alpha}[\sigma(P)]\]
\end{definition}

In altre parole, \(\alpha^{\sigma}\) è un assegnamento che, a ciascuna proposizione elementare \(P\in Sub(\phi)\), assegna la valutazione della formula \(\sigma(P)\) ottenuta dopo la sostituzione \(\sigma\).

\begin{note}{}
  Si noti che \(\alpha^\sigma\) è un assegnamento per \(\phi\), mentre \(\alpha\) è un assegnamento per \(\phi\sigma\).
\end{note}

Vediamo dunque un risultato intermedio che ci sarà utile per dimostrare il \Namedonlyhref{thm:substitution}.
\begin{lemma}[namedlabel={lem:substitution}{Lemma di Sostituzione}]{Lemma di Sostituzione}
  \(\forall \phi \in \PL, \sigma \in \PL^{\AP}, \alpha \in \Ass{\phi\sigma}\) vale che
  \[\Val{\alpha}[\phi\sigma]=\Val{\alpha^\sigma}[\phi]\]
\end{lemma}

\begin{proof}
  Dimostriamo per induzione strutturale.

  \begin{description}
    \item[\(\phi \in \AP\):] in questo caso, che rappresenta il caso base, \(\phi\sigma = \sigma(\phi)\). Ma quindi
      \[\Val{\alpha}[\phi\sigma] = \Val{\alpha}[\sigma(\phi)] \eqdef \alpha^\sigma(\phi) = \Val{\alpha^\sigma}[\phi]\]
    \item[\(\phi = \neg \psi\):] in questo primo caso induttivo, dalla definizione di sostituzione abbiamo che \(\phi\sigma = \neg (\psi\sigma)\). Allora:
      \[\Val{\alpha}[\phi\sigma] = \Val{\alpha}[\neg (\psi\sigma)] = 1 - \Val{\alpha}[\psi\sigma]\]
      Ma per ipotesi induttiva \(\Val{\alpha}[\psi\sigma] = \Val{\alpha^\sigma}[\psi]\), quindi
      \[\Val{\alpha}[\phi\sigma] = 1 - \Val{\alpha^\sigma}[\psi] = \Val{\alpha^\sigma}[\neg \psi] = \Val{\alpha^\sigma}[\phi]\]
    \item[\(\phi = \psi_1 \odot \psi_2\):] Per tutti gli operatori binari, che rappresentano i restanti casi induttivi, la struttura di dimostrazione è analoga, e si basa sempre sull'ipotesi induttiva. Ad esempio, se \(\phi = \psi_1 \land \psi_2\) allora
      \[\Val{\alpha}[\phi\sigma] = \Val{\alpha}[(\psi_1\sigma) \land (\psi_2\sigma)] = \min\set{\Val{\alpha}[\psi_1\sigma], \Val{\alpha}[\psi_2\sigma]}\]
      Ma per ipotesi induttiva \(\Val{\alpha}[\psi_i\sigma] = \Val{\alpha^\sigma}[\psi_i]\) per \(i\in\set{1,2}\), quindi
      \[\Val{\alpha}[\phi\sigma] = \min\set{\Val{\alpha^\sigma}[\psi_1], \Val{\alpha^\sigma}[\psi_2]} = \Val{\alpha^\sigma}[\psi_1 \land \psi_2] = \Val{\alpha^\sigma}[\phi]\]
  \end{description}
\end{proof}

\begin{note}{}
  Si noti che questa dimostrazione non contraddice la nota precedente riguardo alla differenza tra proprietà sintattiche e semantiche, in quanto nonostante per brevità siano state omesse le dimostrazioni per i diversi operatori binari, esse sono comunque diverse tra loro, in quanto ogni connettivo ha una semantica diversa che deve essere rispettata. Per mostrare infatti la dimostrazione di tale caso è stato necessario istanziare \(\odot\) con un operatore specifico, e mostrare che la semantica di tale operatore è preservata dalla sostituzione.
\end{note}

Forti del lemma, possiamo finalmente procedere con la dimostrazione del \Namedonlyhref{thm:substitution}, che ora risulta quasi immediata.
Vogliamo dimostrare che \(\forall \sigma \in \PL^{\AP}\), \(\phi \in \Taut \implies \phi\sigma \in \Taut\).
Procederemo per contrapposizione, ovvero assumeremo che \(\phi\sigma \not\in \Taut\) e dimostreremo che \(\phi \not\in \Taut\).

\begin{proof}[Dimostrazione del (\ref{thm:substitution})]
  Se \(\phi\sigma \not\in \Taut\), allora \(\exists \alpha \in \Ass{\phi\sigma} \st \Val{\alpha}[\phi\sigma] = 0\).

  Per il \Namedonlyhref{lem:substitution}, \(\Val{\alpha}[\phi\sigma] = \Val{\alpha^\sigma}[\phi]\), quindi \(\Val{\alpha^\sigma}[\phi] = 0\), ovvero esiste un assegnamento \(\alpha^\sigma\) tale che \(\phi\) è falsa, e quindi \(\phi \not\in \Taut\).
\end{proof}


\begin{observation}{}
  È importante notare che, nonostante la sostituzione preservi validità e tautologicità, essa non preserva la soddisfacibilità o, in generale, la verità.
  Formalmente:
  \[\phi \in \Sat \centernot\implies \phi\sigma \in \Sat\]

  Questo è facile da vedere considerando \(\phi = P \in \AP\). Si osserva \(\phi \in \Sat\) per qualsiasi assegnamento \(\alpha\) tale per cui \(\alpha(P)=1\).
  Scegliendo però \(\sigma : P \to P\land \neg P\), che è chiaramente insoddisfacibile, e fissato un modello \(m\), si ha che
  \[m\models\phi \centernot\implies m\models\phi\sigma\]
  Altrimenti, si otterrebbe una contraddizione.

  Invece, la sostituzione conserva l'insoddisfacibilità. Infatti,
  \[\phi \in \Unsat \iff \neg \phi \in \Taut \implies (\neg\phi)\sigma \in \Taut\iff \]
  \[\iff \neg(\phi\sigma) \in \Taut \implies \phi\sigma \in \Unsat\]
\end{observation}

Il teorema di sostituzione è un primo esempio di come sia possibile manipolare le formule in modo da preservare alcune proprietà semantiche. 
In particolare, tale teorema ci dice che è possibile creare infinite formule valide, poiché è sufficiente partire da una formula valida, applicare una qualsiasi sostituzione e ottenere una nuova formula valida.

\subsection{Deduzione Semantica}

Possiamo notare che \(A \to (B\to C)\equiv (A\land B)\to C\).
Infatti, chiamando \(\phi_1 = A \to (B\to C)\) e la seconda \(\phi_2 = (A\land B)\to C\) possiamo andare ad analizzare quali assegnamenti rendono vere o false ciascuna di queste formule.

La prima formula \(\phi_1\) ha come connettivo principale l'implicazione, e dunque è falsa se e soltanto se \(A\) è vera e \(B\to C\) è falsa. Ma \(B\to C\) è un altra implicazione, dunque è falsa se e soltanto se \(B\) è vera e \(C\) è falsa. Di conseguenza, \(\phi_1\) è falsa se e soltanto se \(A\) è vera, \(B\) è vera e \(C\) è falsa.

D'altro canto, anche la seconda formula \(\phi_2\) ha come connettivo principale l'implicazione, e dunque è falsa se e soltanto se \(A\land B\) è vera e \(C\) è falsa. Ma \(A\land B\) è vera se e soltanto se \(A\) è vera e \(B\) è vera. Di conseguenza, \(\phi_2\) è falsa se e soltanto se \(A\) è vera, \(B\) è vera e \(C\) è falsa, ovvero le due formule sono false, e quindi vere, negli stessi assegnamenti, e dunque sono logicamente equivalenti.

Possiamo dunque trasformare questa equivalenza in una tautologia, scrivendo
\[A \to (B\to C) \leftrightarrow (A\land B)\to C\]

Questa tautologia è alla base del concetto di deduzione semantica, ovvero la possibilità di dedurre una formula \(\phi\) sapendo che esiste una formula \(\psi\) vera tale che \(\psi \to \phi\) è vera.

\begin{theorem}[namedlabel={thm:semantic_deduction}{Teorema di Deduzione Semantica}]{Teorema di Deduzione Semantica}
  Sia \(\Gamma \subseteq \PL\) e \(\phi,\psi \in \PL\). Allora
  \[\Gamma,\phi \models \psi \iff \Gamma \models \phi\to\psi\]
\end{theorem}

Il teorema permette di \qi{estrarre} una premessa, spostandola nella conseguenza logica da dimostrare.
Ciò consente spesso di semplificare le dimostrazioni. 
Infatti, per dimostrare \(\Gamma,\phi \models \psi\), è necessario, in primo luogo, trovare un modello per cui valga \(\Gamma, \phi\).
Abbiamo visto che la ricerca di un modello ha complessità esponenziale rispetto al numero di proposizioni elementari coinvolte. A ciò, si aggiungono eventuali dipendenze tra queste. 
Pertanto, cercare un modello che soddisfi \(\Gamma\) è sicuramente più facile che cercare un modello che soddisfi \(\Gamma, \phi\).
Inoltre, la deduzione logica permette di soddisfare le proposizioni in maniera sequenziale.
Ad esempio, invece di dire \qi{Se piove e fa freddo, allora prendo l'ombrello e la sciarpa}, possiamo \qi{Se piove, allora se fa freddo prendo l'ombrello e la sciarpa}, il che ci consente di affrontare un problema alla volta.

\begin{proof}
  Dalla definizione di conseguenza logica, \(\Gamma,\phi \models \psi\) significa che
  \[\forall m \in 2^{\AP}, m\models \Gamma \cup \set{\phi} \implies m\models \psi\]

  L'argomento sinistro dell'implicazione è però equivalente a dire che \(m\models \Gamma\) e \(m\models \phi\). Quindi
  \[\forall m \in 2^{\AP}, m\models \Gamma \text{ e } m\models \phi \implies m\models \psi\]

  Dando quindi semantica agli operatori di metalivello nel modo usuale (ovvero identico a quelli della logica proposizionale), possiamo riscrivere la precedente formula con variabili proposizionali \(A\), \(B\) e \(C\) al posto di \(m\models \Gamma\), \(m\models \phi\) e \(m\models \psi\) rispettivamente, ottenendo
  \[\forall m \in 2^{\AP}, A\text{ e }B \implies C\]
  che, per la tautologia vista in precedenza, è equivalente a
  \[\forall m \in 2^{\AP}, A \implies (B \implies C)\]
  Applicando la sostituzione \(\sigma=\set{ A \to m\models \Gamma, B \to m\models \phi, C \to m\models \psi }\), si ottiene 
  \[\forall m \in 2^{\AP}, m\models \Gamma \implies (m\models \phi \implies m\models \psi)\]

  Per semantica dell'implicazione si ha che \((m\models \phi \implies m\models \psi) \iff (m\models \phi \to \psi)\), dunque
  \[\forall m \in 2^{\AP}, m\models \Gamma \implies m\models \phi\to\psi\]
\end{proof}

\section{Forme normali}

Il linguaggio che abbiamo utilizzato fino ad adesso contiene molti simboli, ma come abbiamo già osservato tramite alcune tautologie è possibile esprimere alcuni operatori in funzione di altri.

Un esempio potrebbero essere le leggi di De Morgan
\[\neg(\phi\land\psi) \equiv \neg\phi \lor \neg\psi\]
\[\neg(\phi\lor\psi) \equiv \neg\phi \land \neg\psi\]

Interpretando queste equivalenze come regole di trasformazione, questo ci permette di spingere le negazioni fino alle proposizioni atomiche mantenendo l'equivalenza logica.

In questo modo, possiamo definire una prima forma normale, chiamata \(\PNF\) (Positive Normal Form) o \(\NNF\) (Negation Normal Form), che contiene tutte le formule di \(\PL\) in cui le negazioni occorrono solo davanti a proposizioni atomiche.

Per definire formalmente questa forma normale, introduciamo il concetto di letterale.

\begin{definition}{Letterale}
  Una formula \(\phi\) è un \keyword{letterale} se \(\phi \in AP\) o \(\phi = \neg P \), con \(P \in AP\).
\end{definition}

\begin{definition}{Forma normale positiva}
  Definiamo la forma normale positiva (\(\PNF\) o \(\NNF\)) come l'insieme di tutte le formule costruite come in \(\PL\), senza usare negazioni e utilizzando l'insieme dei letterali \(\mathcal{L}\) come \(\AP\).
\end{definition}

Si ha ovviamente che \(\PNF \subset \PL\), ma dalle leggi di De Morgan è ovvio che
\[\forall \phi \in \PL, \exists \phi' \in \PNF: \phi \iff \phi'\]

In generale, il concetto di forma normale è un sottoinsieme di formule, costruibile in maniera induttiva, che mantiene la completezza semantica.
Ciò significa che per ogni formula esiste una formula equivalente in forma normale.

Le due forme normali più importanti sono la \keyword{forma normale congiuntiva} (\(\CNF\)) e la \keyword{forma normale disgiuntiva} (\(\DNF\)).

\begin{definition}{Disjunctive Normal Form}
  Definiamo la \keyword{forma normale disgiuntiva} (\(\DNF\)) come l'insieme di tutte le formule con la seguente struttura:
  \[\phi = \bigvee_{i=1}^k\left(\bigwedge_{j=1}^{r_i} l_{ij}\right)\]
  dove \(l_{ij} \in \mathcal{L}\) per ogni \(i\) e \(j\).
\end{definition}

\begin{example}{}
  Un esempio di formula in DNF è la seguente:
  \[(P\land \neg Q) \lor (\neg P \land \neg Q \land R) \lor \ldots\]

  Ognuna delle parentesi (quindi ogni congiunzione) è detta \keyword{clausola congiuntiva}.
\end{example}

\begin{definition}{Conjunctive Normal Form}
  Definiamo la \keyword{forma normale congiuntiva} (\(\CNF\)) come l'insieme di tutte le formule con la seguente struttura:
  \[\phi = \bigwedge_{i=1}^k\left(\bigvee_{j=1}^{r_i} l_{ij}\right)\]
  dove \(l_{ij} \in \mathcal{L}\) per ogni \(i\) e \(j\).
\end{definition}

\begin{example}{}
  Un esempio di formula in CNF è la seguente:
  \[(P\lor \neg Q) \land (\neg P \lor \neg Q \lor R) \land \ldots\]

  Ognuna delle parentesi (quindi ogni disgiunzione) è detta \keyword{clausola disgiuntiva}.
\end{example}

Essendoci dualità tra congiunzione e disgiunzione, è facile vedere che le due forme normali sono duali.

Ovviamente \(\CNF \cup \DNF \subseteq \PNF\). È possibile mostrare facilmente come, per ogni \(\phi \in \Sat\), sia possibile costruire \(\psi \in \DNF\) tale che \(\phi \equiv \psi\) partendo dalla tabella di verità di \(\phi\).
Infatti, basterà aggiungere una clausola congiuntiva per ogni assegnamento che rende vera \(\phi\).
In ciascuna clausola congiuntiva, un \(P\in \AP[\phi]\) corrisponderà al letterale \(P\) se \(P\) è vero nell'assegnamento, al letterale \(\neg P\) altrimenti.
Il seguente esempio chiarirà sicuramente il procedimento descritto.

\begin{example}{}
  Consideriamo la formula \(\phi \) con la seguente tabella di verità:

  \begin{center}
    \begin{tabular}{ccc|c}
      \(A\) & \(B\) & \(C\) & \(\phi\)\\
      \hline
      0 & 0 & 0 & 1\\
      0 & 0 & 1 & 0\\
      0 & 1 & 0 & 1\\
      0 & 1 & 1 & 0\\
      1 & 0 & 0 & 1\\
      1 & 0 & 1 & 0\\
      1 & 1 & 0 & 1\\
      1 & 1 & 1 & 0\\
    \end{tabular}
  \end{center}

  Dalla tabella di verità, possiamo costruire la formula in DNF equivalente a \(\phi\) come segue:
  \[(\neg A \land \neg B \land \neg C) \lor (\neg A \land B \land \neg C) \lor (A \land \neg B \land \neg C) \lor (A \land B \land \neg C)\]
  
\end{example}

\begin{note}{}
  Il procedimento mostrato non permette di costruire formule non soddisfacibili. Se non riuscissimo a esprimere le formule insoddisfacibili in \(\DNF\), allora la \(\DNF\) non manterrebbe, come abbiamo detto, la completezza semantica.
  Per risolvere il \qi{buco} lasciato dal procedimento descritto, è sufficiente associare, a ogni formula insoddisfacibile, la formula \(P\land \neg P\), che è \(\DNF\).
  Ciò è possibile dal momento in cui tutte le formule insoddisfacibili sono equivalenti a una contraddizione.
\end{note}

Per le \(\CNF\) il ragionamento è simile. 
In questo caso, però, poiché dobbiamo esprimere la formula come una congiunzione di clausole disgiuntive, non possiamo semplicemente elencare degli assegnamenti come in precedenza, poiché gli assegnamenti sono essenzialmente delle congiunzioni.
Dunque, al posto di elencare gli assegnamenti che rendono vera la formula, elenchiamo in congiunzione la negazione di ogni assegnamento che la rende falsa. 
In questo modo, ci troviamo con una formula della forma:
\[\phi' = \neg(l_1^1 \land \ldots \land \ldots \land l_1^{r_1})\land\neg(l_2^{1}\land \ldots \land \ldots\land l_2^{r_2})\land\ldots\land \neg(l_k^1\land\ldots\land l_k^{r_k})\]
A questo punto, applicando le leggi di De Morgan, possiamo ottenere una formula della forma:
\[\phi'' = (\neg l_1^1 \lor \ldots \lor \neg l_1^{r_1})\land(\neg l_2^1 \lor \ldots \lor \neg l_2^{r_2})\land\ldots\land (\neg l_k^1 \lor \ldots \lor \neg l_k^{r_k})\]
che è esattamente una formula in \(\CNF\).
Si noti che, se, prima dell'applicazione della legge di De Morgan, un letterale era della forma \(\neg l\), allora l'applicazione di De Morgan porta il letterale a essere \(\neg \neg l\), che, per doppia negazione, equivale a \(l\).

\begin{example}{}
  Riprendendo la tabella di verità dell'esempio precedente, possiamo costruire la formula in \(\CNF\) equivalente a \(\phi\) come segue:
  \[\neg(\neg A \land \neg B \land C) \land \neg(\neg A \land B \land C) \land \neg(A \land \neg B \land C) \land \neg(A \land B \land C)\]
  \[\equiv (A \lor B \lor \neg C) \land (A \lor \neg B \lor \neg C) \land (\neg A \lor B \lor \neg C) \land (\neg A \lor \neg B \lor \neg C)\]
\end{example}

\begin{note}{}
  In modo duale alle \(\DNF\), il procedimento descritto non consente di costruire una formula \(\CNF\) equivalente a formule valide.
  Poiché tutte le formule valide sono equivalenti a una tautologia, è possibile associare a esse la formula \(P \lor \neg P\), che è \(\CNF\).
\end{note}

Tuttavia, queste costruzioni sono estremamente inefficienti e portano a formule di dimensione esponenziale rispetto alla formula di partenza. 
Ciò avviene perché la dimensione della tabella di verità di una formula \(\phi\) è \(2^{\abs{\AP[\phi]}}\), dovendo considerare tutti gli assegnamenti possibili, e dunque la formula in \(\DNF\) o \(\CNF\) costruita a partire da tale tabella può avere un numero di clausole esponenziale. 
In particolare, se \(\phi\) è già in \(\DNF\) o in \(\CNF\), allora passare alla forma normale duale usando le leggi di De Morgan e la distributività porta a una formula di dimensione esponenziale.

\begin{example}{}
  Sia \(\phi \in\DNF\) tale che \(\phi\) sia:
  \[(P_1 \land P_2)\lor (P_3 \land P_4)\]

  Andiamo quindi a trasformarla in \(\CNF\) usando la proprietà distributiva.
  \[(P_1 \land P_2)\lor (P_3 \land P_4)\equiv\]
  \[\equiv ((P_1 \land P_2)\lor P_3)\land((P_1 \lor P_4)\lor P_4)\equiv\]
  \[\equiv (P_1 \lor P_3) \land (P_2 \lor P_3) \land (P_1 \lor P_4) \land (P_2 \lor P_4)\]

  In questo caso si è passati da una formula con \(2\) clausole a una formula con \(4\) clausole, che non è un evidenza sufficiente per la crescita esponenziale.

  Consideriamo però una formula:
  \[\phi = (l_1^1 \land l_2^1) \lor (l_1^2 \land l_2^2) \lor\ldots\lor (l_1^n \land l_2^n)\]
  che è in \(\DNF\) e contiene \(n\) clausole congiuntive binarie.

  Analizzando la distributività però, possiamo notare come venga generata una clausola disgiuntiva per ogni possibile combinazione di letterali.

  In generale, quindi, ogni clausola disgiuntiva prodotta con la distributività corrisponde a una scelta di un letterale per ogni clausola congiuntiva, ed è dunque equivalente a una funzione \(f:\set{1,2,\ldots,n} \to \{1,2\}\). Dunque, il numero di clausole disgiuntive generate è pari al numero di tali funzioni, che è \(\abs{\set{1,2}}^n = 2^n\).
\end{example}

Il motivo per cui abbiamo introdotto queste forme è che, se una formula è in \(\DNF\), la sua soddisfacibilità è facilmente verificabile. 
Dualmente, se una formula è in \(\CNF\), la sua validità è facilmente verificabile.
Infatti, sia
\[\phi = C_1 \lor C_2 \lor \ldots \lor C_k\]
con \(C_i = (l_1^i \land \ldots \land l_{r_i}^i)\), per ogni \(i \in \set{1,\ldots,k}\).
Allora \(\phi\) è soddisfacibile se e solo se almeno una clausola congiuntiva \(C_i\) è soddisfacibile.
Inoltre, una clausola congiuntiva \(C_i\) è insoddisfacibile se e solo se contiene due letterali complementari, cioè esistono \(l_j^i\) e \(l_h^i\) tali che \(l_j^i = \neg l_h^i\).
Infatti, in presenza di una coppia complementare compare una contraddizione del tipo \(P \land \neg P\); viceversa, se non compaiono letterali complementari, è sempre possibile costruire un assegnamento che rende vera la clausola.

Dualmente, sia
\[\phi = D_1 \land D_2 \land \ldots \land D_k\]
con \(D_i = (l_1^i \lor \ldots \lor l_{r_i}^i)\), per ogni \(i \in \set{1,\ldots,k}\).
Allora \(\phi\) è valida se e solo se tutte le clausole disgiuntive \(D_i\) sono valide.
Inoltre, una clausola disgiuntiva \(D_i\) è valida se e solo se contiene due letterali complementari, cioè esistono \(l_j^i\) e \(l_h^i\) tali che \(l_j^i = \neg l_h^i\).

\chapter{Calcolo deduttivo proposizionale}

Nel capitolo precedente abbiamo definito una logica come una tripla \((\mathcal{L}, \models, \vdash)\), dove \(\mathcal{L}\) è un linguaggio, \(\models\) è la relazione di soddisfacibilità che ne esprime la semantica, e \(\vdash\) è un calcolo, che ci permette di dedurre formule a partire da altre, o verificare se una formula è valida o meno.

Per quanto concerne la logica proposizionale, abbiamo definito il linguaggio \(\PL\) e fornito ben due semantiche, ma nonostante abbiamo visto una forma di calcolo tramite le tabelle di verità, questo non è sufficiente.
Infatti, questo tipo di calcolo non \qi{scala} a logiche più complesse ed espressive, poiché assume il potersi ridurre a un numero finito di casi, cosa che già per il prim'ordine non è possibile.

Vedremo in questo capitolo due calcoli, un primo detto \keyword{calcolo di deduzione naturale} più affine al ragionamento umano, e un secondo detto \keyword{calcolo di risoluzione} più affine a un'implementazione automatica, e che è alla base di molti sistemi di dimostrazione automatica.

Entrambi questi calcoli sono di tipo puramente sintattico, e non fanno riferimento alla semantica, ma le regole di trasformazione che li caratterizzano sono formulate in modo da preservare le proprietà semantiche, e quindi ci permettono di dedurre conseguenze logiche, ovvero di ottenere la stessa relazione di deducibilità che abbiamo definito tramite la semantica, ma in maniera puramente sintattica.

\begin{definition}{Calcolo inferenziale}
  Definiamo un \keyword{calcolo inferenziale} come una coppia \((\Axi, \Rules)\) dove \(\Axi \subseteq \PL\) è un insieme di formule detto \keyword{assiomi}, e \(\Rules\) è un insieme di regole di trasformazione dette \keyword{regole di inferenza}.

  Le formule in \(\Axi\) sono proposizioni considerate vere a priori, e le regole in \(\Rules\) sono regole che ci permettono a partire da un numero finito di formule, dette \keyword{premesse}, di dedurre una nuova formula, detta \keyword{conclusione}.
\end{definition}

Ovviamente, come anticipato non tutte le regole di trasformazioni sono regole d'inferenza valide, poiché vorremmo che queste regole preservassero delle proprietà semantiche.

A prescindere da ciò però, a partire da un calcolo inferenziale \((\Axi, \Rules)\) possiamo definire una relazione di deducibilità.

\begin{definition}[namedlabel={def:deducibility}{Deducibilità}]{Deducibilità}
  Dati un calcolo inferenziale \((\Axi, \Rules)\) e un insieme di formule \(\Gamma\), diciamo che \(\Gamma\) deduce una formula \(\phi\) tramite \(\Rules\), e scriviamo \(\Gamma \vdash_{\Rules} \phi\), se è possibile dedurre come conclusione \(\phi\) a partire da un numero finito di formule in \(\Gamma\) e \(\Axi\), applicando un numero finito di regole in \(\Rules\).

  Per indicare l'uso di una specifica regola di inferenza \(\Rule \in \Rules\), scriveremo \(\Gamma \vdash_{\Rule} \phi\) per indicare che \(\Gamma\) deduce \(\phi\) tramite una singola applicazione di \(\Rule\).
  In caso non ci sia ambiguità sul calcolo inferenziale, scriveremo semplicemente \(\Gamma \vdash \phi\).
\end{definition}

Intuitivamente, l'obiettivo che vorremmo ottenere con un calcolo inferenziale, e quindi con la relazione di deducibilità, sarebbe quello di catturare, in maniera puramente sintattica, la conseguenza logica, ovvero vorremmo che:
\[\Gamma \models \phi \iff \Gamma \vdash \phi\]

Per quanto riguarda gli assiomi, sono un insieme minimale di tautologie dalle quali tramite il teorema di sostituzione varranno in diverse istanziazioni.\footnote{In un certo senso la regola di sostituzione fa parte di \(R\).}

Il primo calcolo che vedremo però non ha assiomi, ma solo regole di inferenza, e si chiama \keyword{calcolo di deduzione naturale}.

\section{Calcolo di deduzione naturale}

Avendo già anticipato l'assenza di assiomi, il calcolo di deduzione naturale esclusivamente dalle regole di inferenza, che sono formulate in modo da catturare il significato semantico degli operatori logici, e quindi ci permettono di dedurre conseguenze logiche.\footnote{
  La versione del calcolo di deduzione naturale presentato qui non è la sua versione minimale, e conterrà delle regole che non sono strettamente necessarie, ma che rendono più semplice la deduzione di formule, e quindi più affine al ragionamento umano.}

In particolare, per ogni operatore logico avremo una regola detta di \keyword{introduzione}, che ci permetterà di dedurre una formula che ha quell'operatore come connettivo principale, e una regola detta di \keyword{eliminazione}, che ci permetterà di dedurre formule che non hanno quell'operatore come connettivo principale, a partire da formule che invece lo hanno.

Nel nostro caso, considereremo i seguenti operatori logici: \(\land\), \(\lor\), \(\to\), \(\neg\) e \(\bot\), dove \(\bot\) è un operatore \(0\)-ario che rappresenta la contraddizione.

\begin{note}{}
  La notazione presentata è da intendersi come segue: sopra la linea ci sono le premesse, sotto la linea c'è la conclusione, e a destra della regola c'è il nome della regola, che è puramente indicativo, e serve solo a identificare la regola, ma non ha alcun significato formale.

  Inoltre, quando ci sono delle assunzioni che è possibile scaricare, queste sono indicate tra parentesi quadre, mentre i puntini di sospensione indicano l'applicazione di un numero finito di regole da quelle assunzioni fino alla conclusione.
\end{note}

\paragraph{Regole per \(\land\):}

Le regole di introduzione ed eliminazione per \(\land\) sono le più intuitive, e sono le seguenti:

\begin{equation*}
  \infer[(\land I)]{A \land B}{A & B} \quad\quad\quad \infer[(\land E^l)]{A}{A \land B} \quad\quad\quad \infer[(\land E^r)]{B}{A \land B}
\end{equation*}

\paragraph{Regole per \(\to\):}

Per quanto riguarda \(\to\), la regola di eliminazione è abbastanza intuitiva oltre che a essere molto conosciuta col nome di \textit{Modus Ponens}, mentre la regola di introduzione è un po' più complessa.

\begin{equation*}
  \infer[(\to E)]{B}{A \to B & A} \quad\quad\quad
  \infer[(\to I)]{A \to B}{
    \begin{array}{c}
    [A] \\
    \vdots \\
    B
  \end{array}
    } 
\end{equation*}

Quello che ci dice è che, se assumendo \(A\) riusciamo a dedurre \(B\), allora possiamo dedurre \(A \to B\) senza più assumere \(A\) \keyword{scaricando l'assunzione}. Tale regola cattura il procedimento usuale di dimostrazione di un teorema che prevede un'implicazione, ovvero volendo dimostrare \(A \to B\), allora si assume \(A\) e si tenta di dedurre \(B\). In caso affermativo, si potrà concludere \(A \to B\).

\paragraph{Regole per \(\lor\):}

Le regole di introduzione per \(\lor\) sono abbastanza intuitive, mentre la regola di eliminazione è un po' più complessa.

\begin{equation*}
  \infer[(\lor I^l)]{A \lor B}{A} \quad\quad\quad \infer[(\lor I^r)]{A \lor B}{B} \quad\quad\quad
  \infer[(\lor E)]{C}{
    \begin{array}[b]{c}
    A \lor B
    \end{array} &
    \begin{array}[b]{c}
    [A] \\
    \vdots \\
    C
  \end{array} &
    \begin{array}[b]{c}
    [B] \\
    \vdots \\
    C
  \end{array}
  }
\end{equation*}

Quest'ultima regola ci dice che, se abbiamo \(A \lor B\), e se assumendo \(A\) e \(B\) \textbf{separatamente} riusciamo a dedurre \(C\) in entrambi i casi, allora possiamo dedurre \(C\) senza più assumere \(A\) e \(B\), ovvero scaricando entrambe le assunzioni nelle rispettive deduzioni.
Questa regola cattura il concetto di dimostrazione per casi, ovvero se si vuole dimostrare \(C\) a partire da \(A \lor B\), sarà sufficiente dedurre \(C\) sia nel caso in cui sia vero \(A\), che nel caso in cui sia vero \(B\). In caso affermativo, si potrà concludere \(C\).
 
\paragraph{Regole per \(\neg\):}

Le regole di introduzione ed eliminazione per \(\neg\) sono le seguenti:

\begin{equation*}
  \infer[(\neg I)]{\neg A}{
    \begin{array}{c}
    [A] \\
    \vdots \\
    \bot
  \end{array}
  } \quad\quad\quad
  \infer[(\neg E)]{\bot}{\neg A & A}
\end{equation*}

La prima regola ci dice che, se assumendo \(A\) riusciamo a dedurre \(\bot\), allora possiamo dedurre \(\neg A\) senza più assumere \(A\) scaricando l'assunzione. Tale regola cattura il procedimento usuale di dimostrazione di un teorema che prevede una negazione, ovvero volendo dimostrare \(\neg A\), allora si assume \(A\) e si tenta di dedurre \(\bot\). In caso affermativo, si potrà concludere \(\neg A\).

\paragraph{Regola di riduzione all'assurdo:}

Infine, vediamo due regole che potremmo definire di eliminazione di \(\bot\), di cui una generalizzazione dell'altra.

\begin{equation*}
  \infer[(\bot E)]{A}{\bot} \quad\quad\quad
  \infer[(\RAA)]{A}{
    \begin{array}{c}
    [\neg A] \\
    \vdots \\
    \bot
  \end{array}
  }
\end{equation*}

La seconda regola è esattamente la \textit{Reductio ad Absurdum}, e ci dice che, se assumendo \(\neg A\) riusciamo a dedurre \(\bot\), allora possiamo dedurre \(A\) senza più assumere \(\neg A\) scaricando l'assunzione. Tale regola cattura il procedimento usuale di dimostrazione per assurdo, in cui negando la tesi si riesce a concludere una contraddizione. L'altra regola invece, che è una generalizzazione della seconda, cattura il concetto di implicazione da premesse false.
Essenzialmente quello che esprime è che, a partire da premesse contraddittorie, è possibile dedurre qualsiasi formula, e quindi in particolare \(A\). Un analogia per avere più chiara una giustificazione intuitiva di questa regola è considerare la semantica usuale dell'implicazione, nella quale, con antecedente falso, non importa il valore di verità del conseguente, ovvero a partire da premesse false è vero tutto e il contrario di tutto.

Il processo di deduzione naturale può quindi essere visto come una costruzione incrementale: si parte da un insieme di assunzioni iniziali e, mediante passi atomici costituiti da applicazioni di regole di inferenza, si ottiene progressivamente la formula obiettivo.
Durante tale processo, alcune assunzioni vengono scaricate tramite le regole che lo consentono (ad esempio \((\to I)\), \((\neg I)\), \((\lor E)\), \((\RAA)\)), mentre le altre restano attive fino alla fine della derivazione.
Queste ultime costituiscono le assunzioni residue, e sono precisamente quelle che compaiono a sinistra del simbolo di deducibilità.
In altri termini, una derivazione conclude \(\Gamma \vdash \phi\) quando \(\phi\) è ottenuta a partire da assunzioni non scaricate esattamente pari a \(\Gamma\).

Vediamo quindi una formalizzazione più rigorosa di questo processo, e in particolare di come si costruisce un albero di deduzione.

\begin{definition}{Albero di deduzione ben formato}
  Un \keyword{albero di deduzione} è un albero finito in cui nodi sono etichettati con formule di \(\PL\).
  
  Un tale albero è detto \keyword{ben formato} se ogni nodo interno è ottenuto dai suoi figli mediante un'istanza di una regola del calcolo naturale.
  
  Le etichette delle foglie sono \keyword{occorrenze di assunzioni}.
  Se un nodo interno è ottenuto da una regola che scarica un'assunzione, allora tutte le occorrenze di quell'assunzione nel sottoalbero radicato in quel nodo sono dette \keyword{scaricate}.
\end{definition}

\begin{definition}{Deduzione nel calcolo naturale}
  Dato \(\Gamma\subseteq \PL\) e \(\phi\in\PL\), il \keyword{sequente} \(\Gamma \vdash \phi\) è detto \keyword{deducibile} se esiste un albero di deduzione ben formato con radice etichettata da \(\phi\) e foglie etichettate da formule in \(\Gamma\) o assunzioni scaricate.
\end{definition}

Andiamo ora a vedere un primo esempio di deduzione naturale, per poi procedere con esempi più complessi che ci aiuteranno a costruire l'intuizione necessaria per affrontare deduzioni più complesse, guidando la scelta delle regole da applicare.

\begin{example}[namedlabel={ex:natural_deduction}{Deduzione Naturale}]{Deduzione naturale}
  Supponiamo di voler dimostrare che \(\vdash A \to \neg \neg A\).\footnote{Indichiamo con \(\vdash A\) per indicare che \(A\) è dimostrabile senza ipotesi, ovvero che \(A\) è un teorema.}

  Chiaramente \(A \equiv \neg \neg A\), quindi ci aspettiamo di poter dimostrare questa formula.

  Iniziamo assumendo \(A\) e \(\neg A\), quindi:
  \begin{prooftree}\label{prtree:natural_deduction_example}
    \AxiomC{\({[A]}_2\)}
    \AxiomC{\({[\neg A]}_1\)}
    \BinaryInfC{\(\bot\)}
    \RuleLabel{\(\neg I_1\)}
    \UnaryInfC{\(\neg\neg A\)}
    \RuleLabel{\(\to I_2\)}
    \UnaryInfC{\(A \to \neg\neg A\)}
  \end{prooftree}

  Andiamo a etichettare le inferenze scaricate con i numeri delle regole di inferenza che le scaricano, in modo da poterle identificare facilmente.

  Possiamo vedere la rappresentazione come albero della deduzione come segue:
  \begin{center}
    \begin{tikzpicture}
      \node (root) {\(A \to \neg\neg A\)};
      \node (n1) [above =of root] {\(\neg\neg A\)};
      \node (n2) [above =of n1] {\(\bot\)};
      \node (n3) [above left=of n2] {\([A]\)};
      \node (n4) [above right=of n2] {\([\neg A]\)};
      \draw (n1) -- (root);
      \draw (n2) -- (n1);
      \draw (n3) -- (n2);
      \draw (n4) -- (n2);
    \end{tikzpicture}
  \end{center}
\end{example}

\subsection{Applicazione del calcolo naturale proposizionale}

In questa sezione andremo ad analizzare nel dettaglio alcuni esempi di deduzioni naturali che esploreranno tutte le regole di inferenza che abbiamo introdotto, in modo da avere un'idea più chiara di come funziona questo calcolo, e di come scegliere le regole da applicare per dedurre formule più complesse.

\subsubsection*{Dimostrazione di \(\vdash A \leftrightarrow \neg \neg A\).}


Sebbene nel calcolo non abbiamo considerato regole di deduzione per \(\leftrightarrow\), possiamo considerare sempre la formulazione equivalente mediante \(\land\) e \(\to\), ovvero \((A \to \neg \neg A) \land (\neg \neg A \to A)\).

Possiamo facilmente notare che per dimostrare una formula con \(\land\) come connettivo principale, è ragionevole pensare di usare \(\land I\), che ci permette di dedurre una formula con \(\land\) come connettivo principale a partire da due formule che hanno come connettivo principale i due operandi di \(\land\). Quindi come primo passo, possiamo pensare di usare \(\land I\) per ottenere \((A \to \neg \neg A) \land (\neg \neg A \to A)\) a partire da \(A \to \neg \neg A\) e \(\neg \neg A \to A\).

\begin{prooftree}
  \AxiomC{\(A \to \neg \neg A\)}
  \AxiomC{\(\neg \neg A \to A\)}
  \RuleLabel{\(\land I\)}
  \BinaryInfC{\((A \to \neg \neg A) \land (\neg \neg A \to A)\)}
\end{prooftree}

Ovviamente, \(A \to \neg \neg A\) è già stato dimostrato in \Cref{prtree:natural_deduction_example}, quindi possiamo banalmente sostituire quella foglia con l'albero di deduzione che abbiamo già costruito, e quindi ci resta da dimostrare \(\neg \neg A \to A\).

\begin{prooftree}
  \AxiomC{\Namedcref{ex:natural_deduction}}
  \UnaryInfC{\(A \to \neg \neg A\)}
  \AxiomC{\(\neg \neg A \to A\)}
  \RuleLabel{\(\land I\)}
  \BinaryInfC{\((A \to \neg \neg A) \land (\neg \neg A \to A)\)}
\end{prooftree}

Dobbiamo quindi dimostrare \(\neg \neg A \to A\). Per fare questo, possiamo andare a retroso, e pensare di usare la regola \(\to I\), che ci permette di dedurre un'implicazione a partire dal suo conseguente. Inoltre tale regola è particolarmente importante perché ci \qi{regala} un'assunzione scaricabile, ovvero ci permette di usare liberamente come assunzione il suo antecedente nel sottoalbero radicato in essa, avendo la possibilità di scaricarla alla fine della deduzione. Le assunzioni scaricabili possono essere da guida per la scelta dei passi successivi, poiché se abbiamo un'assunzione scaricabile di una formula \(A\), allora è ragionevole pensare di usare tale formula come premessa per dedurre altre formule, e quindi di usare regole di inferenza che hanno \(A\) come premessa.

\begin{prooftree}
  \AxiomC{\Namedcref{ex:natural_deduction}}
  \UnaryInfC{\(A \to \neg \neg A\)}
  \AxiomC{\(A\)}
  \RuleLabel{\(\to I_1\)}[\({[\neg\neg A]}_1\)]
  \UnaryInfC{\(\neg \neg A \to A\)}
  \RuleLabel{\(\land I\)}
  \BinaryInfC{\((A \to \neg \neg A) \land (\neg \neg A \to A)\)}
\end{prooftree}

A questo punto dobbiamo dimostrare \(A\), sapendo di poter usare liberamente \(\neg \neg A\) come assunzione. Osservando le regole che abbiamo a disposizione, nessuna ci permette direttamente di dedurre \(A\) da \(\neg \neg A\), ma \(\RAA\) ci permette di dedurre \(A\) da \(\bot\) scaricando l'assunzione di \(\neg A\).

\begin{prooftree}
  \AxiomC{\Namedcref{ex:natural_deduction}}
  \UnaryInfC{\(A \to \neg \neg A\)}
  \AxiomC{\(\bot\)}
  \RuleLabel{\(\RAA_2\)}[\({[\neg A]}_2\)]
  \UnaryInfC{\(A\)}
  \RuleLabel{\(\to I_1\)}[\({[\neg\neg A]}_1\)]
  \UnaryInfC{\(\neg \neg A \to A\)}
  \RuleLabel{\(\land I\)}
  \BinaryInfC{\((A \to \neg \neg A) \land (\neg \neg A \to A)\)}
\end{prooftree}

A questo punto il gioco è fatto, poiché ci resta da dimostrare \(\bot\), ma abbiamo come assunzioni scaricabili sia \(\neg A\) che \(A\), e quindi possiamo trivialmente usare \(\neg E\).

\begin{prooftree}
  \AxiomC{\Namedcref{ex:natural_deduction}}
  \UnaryInfC{\(A \to \neg \neg A\)}
  \AxiomC{\({[\neg\neg A]}_1\)}
  \AxiomC{\({[\neg A]}_2\)}
  \RuleLabel{\(\neg E\)}
  \BinaryInfC{\(\bot\)}
  \RuleLabel{\(\RAA_2\)}
  \UnaryInfC{\(A\)}
  \RuleLabel{\(\to I_1\)}
  \UnaryInfC{\(\neg \neg A \to A\)}
  \RuleLabel{\(\land I\)}
  \BinaryInfC{\((A \to \neg \neg A) \land (\neg \neg A \to A)\)}
\end{prooftree}

\subsubsection*{Dimostrazione di \((P\to Q)\to((Q\to R) \to (P\to R))\).}

Anche qui, il connettivo principale è \(\to\), che vogliamo dimostrare partendo da un insieme vuoto di premesse, quindi dovremo scaricare ogni assunzione che facciamo.
Possiamo ragionevolmente iniziare con \(\to I\) per ottenere la nostra formula a partire da \((Q\to R) \to (P\to R)\) e scaricando l'assunzione di \(P\to Q\).

\begin{prooftree}
  \AxiomC{\((Q\to R) \to (P\to R)\)}
  \RuleLabel{\(\to I_1\)}[\({[P\to Q]}_1\)]
  \UnaryInfC{\((P\to Q)\to((Q\to R) \to (P\to R))\)}
\end{prooftree}

Ci troviamo adesso a dover dimostrare un altra implicazione, e quindi possiamo pensare di usare nuovamente \(\to I\) per ottenere \((Q\to R) \to (P\to R)\) a partire da \(P\to R\) scaricando l'assunzione di \(Q\to R\).

\begin{prooftree}
  \AxiomC{\(P\to R\)}
  \RuleLabel{\(\to I_2\)}[\({[Q\to R]}_2\)]
  \UnaryInfC{\((Q\to R) \to (P\to R)\)}
  \RuleLabel{\(\to I_1\)}[\({[P\to Q]}_1\)]
  \UnaryInfC{\((P\to Q)\to((Q\to R) \to (P\to R))\)}
\end{prooftree}

Ancora una volta, ci troviamo a dover dimostrare un'altra implicazione, e quindi possiamo pensare di usare nuovamente \(\to I\) per ottenere \(P\to R\) a partire da \(R\) scaricando l'assunzione di \(P\).

\begin{prooftree}
  \AxiomC{\(R\)}
  \RuleLabel{\(\to I_3\)}[\({[P]}_3\)]
  \UnaryInfC{\(P\to R\)}
  \RuleLabel{\(\to I_2\)}[\({[Q\to R]}_2\)]
  \UnaryInfC{\((Q\to R) \to (P\to R)\)}
  \RuleLabel{\(\to I_1\)}[\({[P\to Q]}_1\)]
  \UnaryInfC{\((P\to Q)\to((Q\to R) \to (P\to R))\)}
\end{prooftree}

Adesso, facendoci guidare dalle assunzioni scaricabili, possiamo notare che usando \(\to E\) possiamo dedurre \(R\) a partire da \(Q\to R\) assumendo \(Q\).

\begin{prooftree}
  \AxiomC{\(Q\)}
  \AxiomC{\({[Q\to R]}_2\)}
  \RuleLabel{\(\to E\)}
  \BinaryInfC{\(R\)}
  \RuleLabel{\(\to I_3\)}[\({[P]}_3\)]
  \UnaryInfC{\(P\to R\)}
  \RuleLabel{\(\to I_2\)}
  \UnaryInfC{\((Q\to R) \to (P\to R)\)}
  \RuleLabel{\(\to I_1\)}[\({[P\to Q]}_1\)]
  \UnaryInfC{\((P\to Q)\to((Q\to R) \to (P\to R))\)}
\end{prooftree}

Se terminassimo ora la dimostrazione, avremmo una ottenuto una deduzione per il sequente \(Q\vdash (P\to Q)\to((Q\to R) \to (P\to R))\), poiché l'unica assunzione non scaricata è \(Q\). Usando però le assunzioni scaricabili, possiamo dedurre \(Q\) usando ancora una volta \(\to E\) a partire da \(P\to Q\) e \(P\).

\begin{prooftree}
  \AxiomC{\({[P\to Q]}_1\)}
  \AxiomC{\({[P]}_3\)}
  \RuleLabel{\(\to E\)}
  \BinaryInfC{\(Q\)}
  \AxiomC{\({[Q\to R]}_2\)}
  \RuleLabel{\(\to E\)}
  \BinaryInfC{\(R\)}
  \RuleLabel{\(\to I_3\)}
  \UnaryInfC{\(P\to R\)}
  \RuleLabel{\(\to I_2\)}
  \UnaryInfC{\((Q\to R) \to (P\to R)\)}
  \RuleLabel{\(\to I_1\)}
  \UnaryInfC{\((P\to Q)\to((Q\to R) \to (P\to R))\)}
\end{prooftree}

Solo ora, possiamo concludere la dimostrazione, poiché ogni assunzione è stata scaricata.

\subsubsection*{Dimostrazione di \(\neg(P\land\neg P)\).}

Il connettivo principale ora è \(\neg\), e quindi la strada più promettente è quella di usare \(\neg I\) per ottenere \(\neg(P\land\neg P)\) a partire da \(\bot\) scaricando l'assunzione di \(P\land\neg P\).

\begin{prooftree}
  \AxiomC{\(\bot\)}
  \RuleLabel{\(\neg I\)}[\({[P\land\neg P]}_1\)]
  \UnaryInfC{\(\neg(P\land\neg P)\)}
\end{prooftree}

Dobbiamo quindi dimostrare \(\bot\) avendo a disposizione \(P\land\neg P\) come assunzione scaricabile. Essendo \(\neg E\) l'unica regola che ci permette di dedurre \(\bot\), è ragionevole pensare di usarla. In particolare, per poterla applicare è necessario avere una formula \(P\) e la sua negazione.

\begin{prooftree}
  \AxiomC{\(P\)}
  \AxiomC{\(\neg P\)}
  \RuleLabel{\(\neg E\)}
  \BinaryInfC{\(\bot\)}
  \RuleLabel{\(\neg I\)}[\({[P\land\neg P]}_1\)]
  \UnaryInfC{\(\neg(P\land\neg P)\)}
\end{prooftree}

Possiamo notare però che entrambe queste assunzioni sono deducibili dalle regole di eliminazione per \(\land\) a partire da \(P\land\neg P\), che sappiamo essere un'assunzione scaricabile, e quindi possiamo usarla come premessa per dedurre sia \(P\) che \(\neg P\).

\begin{prooftree}
  \AxiomC{\({[P\land\neg P]}_1\)}
  \RuleLabel{\(\land E^l\)}
  \UnaryInfC{\(P\)}
  \AxiomC{\({[P\land\neg P]}_1\)}
  \RuleLabel{\(\land E^r\)}
  \UnaryInfC{\(\neg P\)}
  \RuleLabel{\(\neg E\)}
  \BinaryInfC{\(\bot\)}
  \RuleLabel{\(\neg I\)}[\({[P\land\neg P]}_1\)]
  \UnaryInfC{\(\neg(P\land\neg P)\)}
\end{prooftree}

Questo esempio ci mostra come un assunzione scaricabile possa essere usata più volte come premessa per dedurre formule diverse, a patto che la regola che la scarica sia radice di un sottoalbero che contiene tutte le deduzioni che la usano come premessa. In questo caso, essendo che \(\neg I\) è la radice di tutto l'albero, possiamo usare \(P\land\neg P\) liberamente, ma se tale regola fosse stata un altro nodo interno, avremmo potuto scaricarla solo nei rami che discendevano da quel nodo, e non in altri rami.

\begin{note}{}
  Le deduzioni che abbiamo visto fino a ora sono tutte deduzioni di tautologie note. Un altro modo per farci guidare nella scelta delle regole da applicare è quello di pensare alla validità semantica della formula che vogliamo dimostrare o a eventuali equivalenze, per capire da quale assunzione \q{estrarre} l'informazione che ci manca per la dimostrazione.

  Formalmente questo processo non è corretto, poiché assume che il nostro calcolo inferenziale catturi la conseguenza logica, ovvero abbia le proprietà di correttezza e completezza, le quali, senza una dimostrazione formale, non possiamo dare per scontate.

  Per quanto riguarda il calcolo di deduzione naturale però, come dimostreremo in seguito, è possibile dimostrare che esso è corretto e completo, e quindi possiamo usare questo tipo di ragionamento per guidarci nella scelta delle regole da applicare, avendo la certezza che se una formula è valida, allora sarà dimostrabile, e viceversa.
\end{note}

\subsubsection*{Dimostrazione di \(P \lor \neg P\).}\label{subsubsec:excluded_middle_dem}

In questo caso, il connettivo principale è \(\lor\), e quindi un idea ragionevole è quella di iniziare \(\lor I\) per ottenere \(P \lor \neg P\) a partire da \(P\) o \(\neg P\).
Per ora iniziamo a ottenerlo assumendo \(P\) e usando \(\lor I^l\).

\begin{prooftree}
  \AxiomC{\(P\)}
  \RuleLabel{\(\lor I^l\)}
  \UnaryInfC{\(P \lor \neg P\)}
\end{prooftree}

Non avendo alcuna assunzione scaricabile, la strada più promettente è quella di usare \(\RAA\) per ottenere \(P\) a partire da \(\bot\) scaricando l'assunzione di \(\neg P\).

\begin{prooftree}
  \AxiomC{\(\bot\)}
  \RuleLabel{\(\RAA_1\)}[\({[\neg P]}_1\)]
  \UnaryInfC{\(P\)}
  \RuleLabel{\(\lor I^l\)}
  \UnaryInfC{\(P \lor \neg P\)}
\end{prooftree}

Ora però dobbiamo dimostrare \(\bot\) avendo a disposizione \(\neg P\) come assunzione scaricabile. Come già visto, l'unica regola che ci permette di dedurre \(\bot\) è \(\neg E\), ma per poterla applicare in questo caso è necessario assumere \(P\) come premessa andando a creare un circolo vizioso.

\begin{prooftree}
  \AxiomC{\(P\)}
  \AxiomC{\({[\neg P]}_1\)}
  \RuleLabel{\(\neg E\)}
  \BinaryInfC{\(\bot\)}
  \RuleLabel{\(\RAA_1\)}
  \UnaryInfC{\(P\)}
  \RuleLabel{\(\lor I^l\)}
  \UnaryInfC{\(P \lor \neg P\)}
\end{prooftree}

Dobbiamo quindi pensare a un altro modo per dedurre \(\bot\) a partire da \(\neg P\). Come detto prima, \(\lor I\) può essere usato per dedurre \(P \lor \neg P\) a partire sia da \(P\) che da \(\neg P\). Possiamo quindi dedurre \(P \lor \neg P\) a partire da \(\neg P\) usando \(\lor I^r\), in modo da poter dedurre \(\bot\) assumendo \(\neg(P\lor\neg P)\).

\begin{prooftree}
  \AxiomC{\({[\neg P]}_1\)}
  \RuleLabel{\(\lor I^r\)}
  \UnaryInfC{\(P \lor \neg P\)}
  \AxiomC{\(\neg(P\lor\neg P)\)}
  \RuleLabel{\(\neg E\)}
  \BinaryInfC{\(\bot\)}
  \RuleLabel{\(\RAA_1\)}
  \UnaryInfC{\(P\)}
  \RuleLabel{\(\lor I^l\)}
  \UnaryInfC{\(P \lor \neg P\)}
\end{prooftree}

A questo punto abbiamo dimostrato \(\neg (P\lor \neg P)\vdash P \lor \neg P\). Per poter arrivare al teorema, ovvero \(\vdash P \lor \neg P\), dobbiamo scaricare l'assunzione di \(\neg (P\lor \neg P)\). Per farlo possiamo pensare di riutilizzarla, in modo da riottenere \(\bot\) usando \(\neg E\), e a quel punto scaricarla con \(\RAA\) per ottenere \(P \lor \neg P\).

\begin{prooftree}
  \AxiomC{\({[\neg P]}_1\)}
  \RuleLabel{\(\lor I^r\)}
  \UnaryInfC{\(P \lor \neg P\)}
  \AxiomC{\({[\neg(P\lor\neg P)]}_2\)}
  \RuleLabel{\(\neg E\)}
  \BinaryInfC{\(\bot\)}
  \RuleLabel{\(\RAA_1\)}
  \UnaryInfC{\(P\)}
  \RuleLabel{\(\lor I^l\)}
  \UnaryInfC{\(P \lor \neg P\)}
  \AxiomC{\({[\neg(P\lor\neg P)]}_2\)}
  \RuleLabel{\(\neg E\)}
  \BinaryInfC{\(\bot\)}
  \RuleLabel{\(\RAA_2\)}
  \UnaryInfC{\(P \lor \neg P\)}
\end{prooftree}

Questo ci mostra che nonostante la regola di deduzione finale sia \(\RAA\) anche partendo da un altra regola di deduzione (\(\lor I\)) siamo riusciti a dimostrare lo stesso \(P \lor \neg P\).

\subsubsection*{Dimostrazione di \((P \to Q) \leftrightarrow (\neg P \lor Q)\).}\label{subsubsec:material_implication_dem}

Come abbiamo già visto nell'\Namedcref{ex:notable_tautologies}, \((P \to Q) \leftrightarrow (\neg P \lor Q)\) è una tautologia, quindi ci aspettiamo di poterla dimostrare. Come abbiamo già visto per dimostrare un equivalenza basta dimostrare entrambe le implicazioni che la compongono per poi concludere usando \(\land I\).

Andiamo quindi a dimostrare separatamente \((P \to Q) \to (\neg P \lor Q)\) e \((\neg P \lor Q) \to (P \to Q)\). Banalmente però, sfruttando \(\to I\) possiamo passare al dover dimostrare i sequenti \(P \to Q \vdash \neg P \lor Q\) e \(\neg P \lor Q \vdash P \to Q\), il che ci permette di non doverci preoccupare di scaricare l'assunzione di \(P \to Q\) e \(\neg P \lor Q\) rispettivamente, e quindi di poterle usare liberamente come assunzioni.

\paragraph{Dimostrazione di \(P \to Q \vdash \neg P \lor Q\).}

Dalla dimostrazione precedente, abbiamo visto che un modo per dimostrare formule con \(\lor\) come connettivo principale è quello di usare \(\RAA\).

\begin{prooftree}
  \AxiomC{\(\bot\)}
  \RuleLabel{\(\RAA_1\)}[\({[\neg(\neg P \lor Q)]}_1\)]
  \UnaryInfC{\(\neg P \lor Q\)}
\end{prooftree}

Come già detto numerose volte, per dimostrare \(\bot\) l'idea più promettente è quella di usare \(\neg E\), per la quale ci servono due formule, una la negazione dell'altra. Dall'assunzione di \(P \to Q\) possiamo dedurre \(Q\) assumendo \(P\). Quindi speriamo di poter usare \(\neg E\) con \(Q\) e \(\neg Q\).

\begin{prooftree}
  \AxiomC{\(P\)}
  \AxiomC{\(P \to Q\)}
  \RuleLabel{\(\to E\)}
  \BinaryInfC{\(Q\)}
  \AxiomC{\(\neg Q\)}
  \RuleLabel{\(\neg E\)}
  \BinaryInfC{\(\bot\)}
  \RuleLabel{\(\RAA_1\)}[\({[\neg(\neg P \lor Q)]}_1\)]
  \UnaryInfC{\(\neg P \lor Q\)}
\end{prooftree}

A questo punto dobbiamo scaricare o dimostrare sia \(P\) che \(\neg Q\), poiché \(P \to Q\) non è necessario che venga scaricata.
Dall'assunzione scaricabile di \(\neg(\neg P \lor Q)\) però, dall'equivalenza semantica con \(P \land \neg Q\), ci aspettiamo di poter dedurre sia \(P\) che \(\neg Q\) da \(\land E\).

Iniziamo con \(\neg Q\). Per poter introdurre una negazione, possiamo utilizzare \(\neg I\), che ci permette di dedurre \(\neg Q\) a partire da \(\bot\) scaricando l'assunzione di \(Q\).

\begin{prooftree}
  \AxiomC{\(P\)}
  \AxiomC{\(P \to Q\)}
  \RuleLabel{\(\to E\)}
  \BinaryInfC{\(Q\)}
  \AxiomC{\(\bot\)}
  \RuleLabel{\(\neg I_2\)}[\({[Q]}_2\)]
  \UnaryInfC{\(\neg Q\)}
  \RuleLabel{\(\neg E\)}
  \BinaryInfC{\(\bot\)}
  \RuleLabel{\(\RAA_1\)}[\({[\neg(\neg P \lor Q)]}_1\)]
  \UnaryInfC{\(\neg P \lor Q\)}
\end{prooftree}

A questo punto per dimostrare questo \(\bot\) dobbiamo introdurre una contraddizione. Avendo a disposizione \(Q\) e \(\neg(\neg P \lor Q)\) possiamo pensare di questa seconda insieme a \(\neg P \lor Q\) per dedurre una contraddizione usando \(\neg E\).

\begin{prooftree}
  \AxiomC{\(P\)}
  \AxiomC{\(P \to Q\)}
  \RuleLabel{\(\to E\)}
  \BinaryInfC{\(Q\)}
  \AxiomC{\({[\neg(\neg P \lor Q)]}_1\)}
  \AxiomC{\(\neg P \lor Q\)}
  \RuleLabel{\(\neg E\)}
  \BinaryInfC{\(\bot\)}
  \RuleLabel{\(\neg I_2\)}[\({[Q]}_2\)]
  \UnaryInfC{\(\neg Q\)}
  \RuleLabel{\(\neg E\)}
  \BinaryInfC{\(\bot\)}
  \RuleLabel{\(\RAA_1\)}
  \UnaryInfC{\(\neg P \lor Q\)}
\end{prooftree}

Ma \(\neg P \lor Q\) è deducibile da \(Q\) tramite \(\lor I^r\), permettendoci di chiudere questo ramo di deduzione.

\begin{prooftree}
  \AxiomC{\(P\)}
  \AxiomC{\(P \to Q\)}
  \RuleLabel{\(\to E\)}
  \BinaryInfC{\(Q\)}
  \AxiomC{\({[\neg(\neg P \lor Q)]}_1\)}
  \AxiomC{\({[Q]}_2\)}
  \UnaryInfC{\(\neg P \lor Q\)}
  \RuleLabel{\(\neg E\)}
  \BinaryInfC{\(\bot\)}
  \RuleLabel{\(\neg I_2\)}
  \UnaryInfC{\(\neg Q\)}
  \RuleLabel{\(\neg E\)}
  \BinaryInfC{\(\bot\)}
  \RuleLabel{\(\RAA_1\)}
  \UnaryInfC{\(\neg P \lor Q\)}
\end{prooftree}

Ci rimane quindi solo da dimostrare \(P\). Per farlo però possiamo applicare un ragionamento simile a quello fatto per \(\neg Q\), ma utilizzando \(\RAA\) invece di \(\neg I\) per dedurre \(P\) a partire da \(\bot\) scaricando l'assunzione di \(\neg P\).

\begin{prooftree}
  \AxiomC{\({[\neg(\neg P \lor Q)]}_1\)}
  \AxiomC{\({[\neg P]}_3\)}
  \RuleLabel{\(\lor I^r\)}
  \UnaryInfC{\(\neg P \lor Q\)}
  \RuleLabel{\(\neg E\)}
  \BinaryInfC{\(\bot\)}
  \RuleLabel{\(\RAA_3\)}
  \UnaryInfC{\(P\)}
  \AxiomC{\(P \to Q\)}
  \RuleLabel{\(\to E\)}
  \BinaryInfC{\(Q\)}
  \AxiomC{\({[\neg(\neg P \lor Q)]}_1\)}
  \AxiomC{\({[Q]}_2\)}
  \UnaryInfC{\(\neg P \lor Q\)}
  \RuleLabel{\(\neg E\)}
  \BinaryInfC{\(\bot\)}
  \RuleLabel{\(\neg I_2\)}
  \UnaryInfC{\(\neg Q\)}
  \RuleLabel{\(\neg E\)}
  \BinaryInfC{\(\bot\)}
  \RuleLabel{\(\RAA_1\)}
  \UnaryInfC{\(\neg P \lor Q\)}
\end{prooftree}

L'unica assunzione non scaricata che ci rimane è \(P \to Q\), e quindi abbiamo dimostrato \(P \to Q \vdash \neg P \lor Q\).

\paragraph{Dimostrazione di \(\neg P \lor Q \vdash P \to Q\).}

Ancora una volta, per dimostrare \(P \to Q\) è ragionevole pensare di usare \(\to I\) per ottenere \(P \to Q\) a partire da \(Q\) scaricando l'assunzione di \(P\).

\begin{prooftree}
  \AxiomC{\(Q\)}
  \RuleLabel{\(\to I_1\)}[\({[P]}_1\)]
  \UnaryInfC{\(P \to Q\)}
\end{prooftree}

Ora però dobbiamo dimostrare \(Q\) avendo a disposizione \(\neg P \lor Q\) e \(P\) come assunzioni. Essendoci un \(\lor\) che contiene \(Q\) come operando, è possibile pensare di utilizzare \(\lor E\) per dedurre \(Q\) a partire da \(\neg P \lor Q\), ma per poterlo applicare è necessario dimostrare \(Q\) sia assumendo \(\neg P\) che assumendo \(Q\).

\begin{prooftree}
  \AxiomC{\(\neg P \lor Q\)}
  \AxiomC{\(Q\)}
  \AxiomC{\(Q\)}
  \RuleLabel{\(\lor E_2\)}[\({[\neg P]}_2,\ {[Q]}_2\)]
  \TrinaryInfC{\(Q\)}
  \RuleLabel{\(\to I_1\)}[\({[P]}_1\)]
  \UnaryInfC{\(P \to Q\)}
\end{prooftree}

Per quanto riguarda dimostrare \(Q\) assumendo \(Q\), questo è immediato.
\begin{prooftree}
  \AxiomC{\(\neg P \lor Q\)}
  \AxiomC{\({[Q]}_2\)}
  \UnaryInfC{\(Q\)}
  \AxiomC{\(Q\)}
  \RuleLabel{\(\lor E_2\)}[\({[\neg P]}_2\)]
  \TrinaryInfC{\(Q\)}
  \RuleLabel{\(\to I_1\)}[\({[P]}_1\)]
  \UnaryInfC{\(P \to Q\)}
\end{prooftree}

Ora rimane da dimostrare \(Q\) assumendo \(\neg P\). Per farlo, possiamo utilizzare \(\bot E\), poiché tra le assunzioni scaricabili abbiamo sia \(P\) che \(\neg P\), e quindi possiamo dedurre \(\bot\) usando \(\neg E\), per poi dedurre \(Q\) da \(\bot\).
\begin{prooftree}
  \AxiomC{\(\neg P \lor Q\)}
  \AxiomC{\({[Q]}_2\)}
  \UnaryInfC{\(Q\)}
  \AxiomC{\({[\neg P]}_2\)}
  \AxiomC{\({[P]}_1\)}
  \RuleLabel{\(\neg E\)}
  \BinaryInfC{\(\bot\)}
  \RuleLabel{\(\bot E\)}
  \UnaryInfC{\(Q\)}
  \RuleLabel{\(\lor E_2\)}
  \TrinaryInfC{\(Q\)}
  \RuleLabel{\(\to I_1\)}
  \UnaryInfC{\(P \to Q\)}
\end{prooftree}

In questo esempio, abbiamo potuto osservare come anche le regole di inferenza più controintuitive e complesse, come \(\bot E\) e \(\lor E\) possano trovare applicazione.

\subsection{Proprietà del calcolo deduttivo}

Abbiamo visto come il calcolo di deduzione naturale sia un sistema formale che ci permette sintatticamente di ottenere nuove formule a partire da un insieme di premesse, ma per poterlo usare in modo efficace è necessario assicurarsi che esso sia in grado di catturare la conseguenza logica.

In questo senso, andremo a dimostrare varie proprietà del calcolo di deduzione naturale, che lo mettono in relazione con la conseguenza logica, per fare in modo da garantire che questa tecnica di manipolazione sintattica riesca a mantenere valore semantico, e quindi che se una formula è dimostrabile, allora è anche valida, e viceversa.

Iniziamo con una proprietà strettamente relativa al calcolo che ci permette di ampliare arbitrariamente l'insieme di premesse da cui dedurre una formula, senza perdere la possibilità di dedurla.

\begin{theorem}[namedlabel={thm:monotonicity}{Monotonicità}]{Monotonicità di \(\vdash\)}
  Dato un insieme di formule \(\Gamma\) e due formule \(\phi\) e \(\psi\), si ha che
  \[\Gamma \vdash \phi \implies \Gamma \cup \{\psi\} \vdash \phi\]
\end{theorem}

È importante notare che questa proprietà risulta solo rafforzata nel caso vengano aggiunte formule che rendono \(\Gamma\) inconsistente, poiché in questo caso sarà possibile dedurre qualsiasi formula a partire da \(\Gamma\) mediante \(\bot E\). Questa è una caratteristica tipica della logica classica, e non è condivisa da tutte le logiche, come ad esempio le logiche paraconsistenti, in cui l'inconsistenza non porta alla trivialità, ma viene in qualche modo contenuta permettendo di parlare del resto dell'insieme di formule anche in presenza di contraddizioni.

\begin{note}{}
  Questa proprietà di generalizza in \(\Gamma \vdash \phi \implies \Gamma \cup \Gamma' \vdash \phi\), applicandola più volte. La dimostrazione di questa versione più generale però segue direttamente da quella della versione più semplice, e quindi non è necessario farla esplicitamente.
\end{note}

\begin{proof}
  Per definizione di deduzione, \(\Gamma \vdash \phi\) se e solo se esiste un albero di derivazione ben formato radice etichettata con \(\phi\) e le cui foglie o sono etichettate con formule di \(\Gamma\) o sono assunzioni scaricabili. 

  In tale definizione non c'è alcun obbligo di usare ogni formula di \(\Gamma\) come etichetta di una foglia. Pertanto, aggiungendo all'insieme delle premesse \(\Gamma\) una nuova formula \(\psi\) non usata nell'albero di derivazione, lo stesso identico albero di derivazione rimane valido per dimostrare \(\Gamma \cup \{\psi\} \vdash \phi\).
\end{proof}

Andiamo quindi a vedere una proprietà analoga a quella vista per la conseguenza logica nel \Namedcref{thm:semantic_deduction}, il vero e proprio \Namedonlyhref{thm:deduction}.

\begin{theorem}[namedlabel={thm:deduction}{Teorema di Deduzione}]{Teorema di Deduzione}
  Dato un insieme di formule \(\Gamma\) e due formule \(\phi\) e \(\psi\), si ha che:
  \[\Gamma, \psi \vdash \phi \iff \Gamma \vdash \psi \to \phi\]
\end{theorem}

\begin{proof}
  Andiamo a dimostrare entrambi i versi.
  
  Iniziamo con \(\Gamma, \psi \vdash \phi \implies \Gamma \vdash \psi \to \phi\).
  Per definizione di deduzione, esiste un albero che dimostra \(\Gamma, \psi \vdash \phi\). Allora è sufficiente applicare la regola \(\to I\) in radice di tale albero, scaricando l'assunzione \(\psi\) e ottenendo così un albero che dimostra \(\Gamma \vdash \psi \to \phi\).

  Nell'altro verso, \(\Gamma \vdash \psi \to \phi \implies \Gamma, \psi \vdash \phi\), l'idea è duale.
  Per definizione di deduzione, esiste un albero che dimostra \(\Gamma \vdash \psi \to \phi\). A questo punto, aggiungendo l'assunzione \(\psi\) all'albero, possiamo applicare la regola \(\to E\) per ottenere un albero che dimostra \(\Gamma, \psi \vdash \phi\).
\end{proof}

Già da questi primi due teoremi è possibile notare un forte parallelismo tra la conseguenza logica e la deduzione, la quale verrà ulteriormente rafforzata dai prossimi teoremi, per poi essere dimostrata in modo definitivo con i teoremi di correttezza e completezza.

\begin{theorem}{Regola del taglio}
  Dato un insieme di formule \(\Gamma\) e due formule \(\phi\) e \(\psi\), si ha che:
  
  Se \(\Gamma, \psi \vdash \phi\) e \(\Gamma' \vdash \psi\) allora \(\Gamma \cup \Gamma' \vdash \phi\)
\end{theorem}

Intuitivamente, questa regola ci dice se abbiamo un albero di derivazione che dimostra un'assunzione, possiamo sostituire a ogni foglia etichettata con quell'assunzione l'albero di derivazione che dimostra quell'assunzione, mantenendo una derivazione valida.


\begin{proof}
  Possiamo formalmente dimostrare questa regola tramite il teorema di deduzione e la monotònicità di \(\vdash\).
  Infatti per monotònicità valgono sia \(\Gamma, \Gamma', \psi \vdash \phi\), che \(\Gamma \cup \Gamma'\vdash \psi\).
  A questo punto, usando il teorema di deduzione, otteniamo \(\Gamma\cup \Gamma' \vdash \psi \to \phi\).

  A questo punto, possiamo costruire un albero di derivazione con \(\phi\) in radice dedotta tramite \(\to E\), i cui due sottoalberi sono l'albero relativo a \(\Gamma \cup \Gamma' \vdash \psi\) e l'albero relativo a \(\Gamma\cup \Gamma' \vdash \psi \to \phi\).
\end{proof}

Prima di spostarci a dimostrare le proprietà di correttezza e completezza, è utile introdurre formalmente l'equivalente sintattico del concetto di insieme insoddisfacibile, ovvero quello di insieme inconsistente.

\begin{definition}{Inconsistenza}
  Un insieme di formule \(\Gamma\) è \keyword{inconsistente} se \(\Gamma \vdash \bot\).

  Equivalentemente, \(\Gamma\) è inconsistente, per definizione, se \(\Gamma \vdash \phi\), \(\forall \phi \in \PL\).
\end{definition}

Vedere che le due definizioni sono equivalenti è immediato.
Infatti, se \(\Gamma \vdash \bot\), possiamo applicare \(\bot E\) alla derivazione di \(\bot\), ottenendo così \(\Gamma \vdash \phi\) per ogni formula \(\phi\). Viceversa, se \(\Gamma \vdash \phi\) per ogni formula \(\phi\), allora in particolare si ha che \(\Gamma \vdash \bot\), poiché \(\bot\) è una formula di queste \(\phi\).

Veniamo dunque alle proprietà di correttezza e completezza, che intuitivamente esprimono relazioni di inclusione tra la conseguenza logica e la deduzione. 

\begin{definition}{Completezza}
  Un sistema di deduzione è \keyword{completo} se, per ogni insieme di formule \(\Gamma\) e ogni formula \(\phi\), si ha che:
  \[\Gamma \models \phi \implies \Gamma \vdash \phi\]
\end{definition}

\begin{definition}{Correttezza}
  Un sistema di deduzione è \keyword{corretto} se, per ogni insieme di formule \(\Gamma\) e ogni formula \(\phi\), si ha che:
  \[\Gamma \vdash \phi \implies \Gamma \models \phi\]
\end{definition}

Essendo che sia \(\vdash \subseteq 2^{\PL}\times\PL\) che \(\models \subseteq 2^{\PL}\times\PL\), è possibile pensare a queste due proprietà come a delle relazioni di inclusione tra i due insiemi. In particolare, diremo che il calcolo è corretto se \(\vdash \subseteq \models\), e completo se \(\models \subseteq \vdash\).

Intuitivamente, la completezza di un calcolo ci dice che tutto ciò che è conseguenza logica è anche deducibile, e quindi che il calcolo è in grado di catturare tutta la conseguenza logica. La correttezza invece ci dice che tutto ciò che è deducibile è anche conseguenza logica, e quindi che il calcolo non permette di dedurre formule che non sono conseguenza logica.

Queste due proprietà prese in isolamento però non ci dicono molto, poiché è possibile soddisfarle in modo triviale. Per la completezza, ad esempio, è sufficiente prendere come relazione di deduzione l'intero insieme \(2^{\PL}\times\PL\), in modo da poter dedurre qualsiasi formula a partire da qualsiasi insieme di formule, e quindi soddisfare la condizione di completezza, ma non quella di correttezza. Per la correttezza invece è sufficiente prendere come relazione di deduzione l'insieme vuoto, in modo da non poter dedurre nessuna formula a partire da nessun insieme di formule, e quindi soddisfare la condizione di correttezza, ma non quella di completezza.

Prima di poter dimostrare però che il calcolo di deduzione naturale è corretto e completo, sarà necessario introdurre alcune definizioni e risultati intermedi, che ci permetteranno di dimostrare più facilmente queste due proprietà.

Iniziamo col definire un concetto che ci aiuterà nella dimostrazione di correttezza.

\begin{definition}{Denotazione}
  Data una formula \(\phi\), la \keyword{denotazione} di \(\phi\) è l'insieme di tutti i modelli che soddisfano \(\phi\), ovvero:
  \[\den{\phi} = \set{m \in 2^{\AP}}[ m \models \phi]\]

  Dato un insieme di formule \(\Gamma\), la \keyword{denotazione} di \(\Gamma\) è l'insieme di tutti i modelli che soddisfano tutte le formule di \(\Gamma\), ovvero:
  \[\den{\Gamma} = \set{m \in 2^{\AP}}[ m \models \Gamma] = \bigcap_{\phi \in \Gamma} \den{\phi}\]
\end{definition}

Possiamo dunque procedere alla dimostrazione della correttezza del calcolo di deduzione naturale, che intuitivamente ci dice che tutto ciò che è provabile è semanticamente fondato, e che quindi il calcolo di deduzione naturale ci permette di dedurre solo formule che sono conseguenza logica delle premesse.

\begin{theorem}{Correttezza del calcolo di deduzione naturale}
  Per ogni insieme di formule \(\Gamma\) e ogni formula \(\phi\) si ha che:
  \[\Gamma \vdash \phi \implies \Gamma \models \phi\]
\end{theorem}

\begin{proof}

  Per definizione di \(\Gamma\vdash\phi\), esiste un albero di deduzione con radice \(\phi\) e con foglie che sono o assunzioni scaricate o formule di \(\Gamma\).

  Dunque dimostrare la tesi equivale a dimostrare che, per ogni albero di deduzione in deduzione naturale per \(\phi\) da \(\Gamma\), \(\Gamma\models\phi\).

  Possiamo facilmente vederlo per induzione sulla struttura (o altezza) dell'albero di deduzione.
  
  Il caso base è una deduzione di altezza \(0\), ovvero \(\phi \in \Gamma\). In questo caso, \(\Gamma\models\phi\) è ovviamente verificato.

  A questo punto sia l'altezza maggiore di \(0\). In questo caso avremo che l'albero avrà in radice \(\phi\) dedotta da una certa regola \(\Rule \in \Rules\) e con figli \(\psi_1,\ldots,\psi_k \in \PL\), ognuno di questi a sua volta radice di un sottoalbero di altezza strettamente minore di quella dell'albero di deduzione di \(\phi\). Per ipotesi induttiva, \(\Gamma\models\psi_i\) per ogni \(i=1,\ldots,k\).

  A questo punto quindi ci basta mostrare con un analisi per casi che, per ogni regola di deduzione naturale, se \(\set{\psi_1,\ldots,\psi_k} \vdash_{\Rule} \phi\) allora \(\set{\psi_1,\ldots,\psi_k} \models_{\Rule} \phi\).

  \begin{description}
    \item[\(\to E\).] Abbiamo da ipotesi induttiva che \(\Gamma\models\psi\) e \(\Gamma\models\psi\to\phi\).
      Ovvero, ogni modello che soddisfa \(\Gamma\) soddisfa anche \(\psi\) e \(\psi\to\phi\).
      
      Dalla semantica dell'implicazione, abbiamo che, se un modello \(m\) soddisfa \(\psi\to\phi\) allora o non soddisfa \(\psi\) o, se la soddisfa, allora soddisfa \(\phi\). Ma essendo che \(m\) soddisfa \(\psi\), allora \(m\) deve soddisfare anche \(\phi\). Dunque ogni modello che soddisfa \(\Gamma\) soddisfa anche \(\phi\), e quindi \(\Gamma\models\phi\).
    \item[\(\neg E\).] Abbiamo da ipotesi induttiva che \(\Gamma\models\neg\psi\) e \(\Gamma\models\psi\). Dunque ogni modello che soddisfa \(\Gamma\) soddisfa anche \(\neg\psi\) e \(\psi\). Per semantica del \(\neg\), allora \(m\) non soddisfa \(\psi\), ma abbiamo detto che \(m\) soddisfa \(\psi\), assurdo. Dunque non esiste alcun modello che soddisfa \(\Gamma\), e dunque \(\Gamma\) è insoddisfacibile.

    Una definizione alternativa di conseguenza logica che avevamo dato era che \(\Gamma\models\phi\) se \(\Gamma\cup\set{\neg\phi}\) è insoddisfacibile. Ma allora \(\Gamma\models\bot\) se \(\Gamma\cup\set{\neg \bot}\) è insoddisfacibile. Ma \(\neg\bot\) è tautologicamente vero, dunque \(\Gamma\cup\set{\neg \bot}\) è insoddisfacibile se e solo se \(\Gamma\) è insoddisfacibile. Cosa che abbiamo appena dimostrato nel caso in cui \(\Gamma\models\neg\psi\) e \(\Gamma\models\psi\). Dunque \(\Gamma\models\bot\).

    \item[\(\RAA\).] Abbiamo da ipotesi induttiva che \(\Gamma,\neg\phi\models\bot\). Ovvero \(\Gamma\cup\set{\neg\phi}\) è insoddisfacibile. Dunque \(\Gamma\models\phi\).
    \item[\(\lor E\).] Abbiamo da ipotesi induttiva che \(\Gamma,\psi_1\models\phi\), \(\Gamma,\psi_2\models\phi\) e \(\Gamma\models\psi_1\lor\psi_2\).
    
    Possiamo uniformare le assunzioni usando il teorema di deduzione semantica, ottenendo \(\Gamma\models\psi_1\to\phi\) e \(\Gamma\models\psi_2\to\phi\). Dunque ogni modello che soddisfa \(\Gamma\) soddisfa anche \(\psi_1\to\phi\), \(\psi_2\to\phi\) e \(\psi_1\lor\psi_2\). Ma essendo che \(m\) soddisfa \(\psi_1\lor\psi_2\), allora \(m\) soddisfa almeno una tra \(\psi_1\) e \(\psi_2\). Se \(m\) soddisfa \(\psi_1\), soddisfando \(\psi_1\to\phi\) soddisfa anche \(\phi\). Se \(m\) soddisfa \(\psi_2\), soddisfacendo \(\psi_2\to\phi\) soddisfa anche \(\phi\). Dunque \(m\) soddisfa \(\phi\) in ogni caso.
    
    Alternativamente, possiamo dimostrare questo punto usando \(\den{\phi}\) e \(\den{\Gamma}\).

    Essendo che \(\Gamma\models \psi_1 \lor \psi_2\), ogni modello che soddisfa \(\Gamma\) soddisfa anche \(\psi_1\lor\psi_2\). Dunque \(\forall m \in \den{\Gamma}, m \in \den{\psi_1\lor\psi_2}\). Alternativamente, \(\den{\Gamma}\subseteq \den{\psi_1\lor\psi_2}\). Con lo stesso ragionamento, dalla semantica della disgiunzione, \(\forall m \in \den{\psi_1\lor\psi_2}, m \in \den{\psi_1}\cup \den{\psi_2}\). Inoltre, abbiamo che \(\den{\Gamma}\subseteq\den{\psi_1\to\phi}\) e \(\den{\Gamma}\subseteq\den{\psi_2\to\phi}\), e dalla semantica dell'implicazione, si ha che \(\den{\Gamma} \cap \den{\psi_1}\subseteq\den{\phi}\) e \(\den{\Gamma}\cap\den{\psi_2}\subseteq\den{\phi}\), poiché ogni modello che soddisfa \(\Gamma\) e soddisfa \(\psi_1\) o \(\psi_2\) deve soddisfare anche \(\phi\). In particolare quindi
    \[(\den{\Gamma}\cap\den{\psi_1})\cup(\den{\Gamma}\cap\den{\psi_2})\subseteq \den{\phi}\]
    dalla distributività dell'intersezione rispetto all'unione, si ha che
    \[\den{\Gamma} \cap (\den{\psi_1}\cup\den{\psi_2})\subseteq\den{\phi}\]
    Avendo già mostrato però che \(\den{\Gamma} \subseteq \den{\psi_1}\cup\den{\psi_2}\) la loro intersezione è proprio \(\den{\Gamma}\),
    dunque \(\den{\Gamma}\subseteq \den{\phi}\), ovvero ogni modello che soddisfa \(\Gamma\) soddisfa anche \(\phi\), e dunque \(\Gamma\models\phi\).
    \item[] Per questioni di brevità non verranno trattate tutte le regole, ma solo quelle più interessanti. Gli altri casi sono analoghi a quelli già trattati, e possono essere dimostrati con un ragionamento simile a quello fatto per le regole già trattate.
  \end{description}
\end{proof}

Conclusa la dimostrazione di correttezza, possiamo passare al verso più interessante, ovvero quello di completezza, che ci dice che tutto ciò che è conseguenza logica è anche deducibile, e quindi che il calcolo di deduzione naturale è in grado di catturare tutta la conseguenza logica.

Questo risulta un po' più complesso da dimostrare, poiché, a differenza della correttezza che chiedeva semplicemente di verificare che le prove generate dal calcolo esibissero una certa proprietà, ovvero l'essere conseguenze logiche delle premesse, la completezza ci chiede di verificare che l'intero spazio delle conseguenze logiche sia coperto dal calcolo, e quindi non sarà più sufficiente verificare in maniera strutturale le prove.

Per dimostrare il teorema di correttezza infatti faremo uso di un paio di risultati intermedi, che ci permetteranno di dimostrare la completezza in modo più semplice.

\begin{lemma}[namedlabel={lem:lindenbaum}{Lindenbaum}]{Lindenbaum}
  Sia \(\Sigma\subseteq\PL\) un insieme consistente di formule. Allora è possibile costruire un insieme di formule \(\Sigma^*\) tale che \(\Sigma\subseteq\Sigma^*\) e \(\Sigma^*\) è massimalmente consistente, ovvero per ogni formula \(\psi\), o \(\psi\in\Sigma^*\) o \(\neg\psi\in\Sigma^*\).
\end{lemma}

\begin{proof}
  Consideriamo un'arbitraria enumerazione delle formule di \(\PL\), ovvero una sequenza \(\phi_1,\phi_2,\ldots\) che contiene tutte le formule di \(\PL\).
  
  Si costruisca dunque una sequenza di insiemi di formule \(\Sigma_0,\Sigma_1,\ldots\) tale che \(\Sigma_0 = \Sigma\) e che
  
  \[ \Sigma_{i+1} = \begin{cases}
                      \Sigma_i\cup\set{\phi_{i+1}} & \text{se } \Sigma_i\cup\set{\phi_{i+1}} \text{ è consistente} \\
                      \Sigma_i & \text{altrimenti}
                    \end{cases}
                    \]

  In questo modo è garantito che, se \(\Sigma\) è consistente, per induzione ogni \(\Sigma_i\) è consistente.

  Definiamo allora \(\Sigma^* = \bigcup_{i\geq 0}\Sigma_i\). Dobbiamo allora mostrare che \(\Sigma^*\) è massimalmente consistente.

  È importante notare che, nonostante \(\Sigma^*\) sia il limite di una successione di cui ogni elemento presenta una proprietà, non è detto che \(\Sigma^*\) presenti quella proprietà, e quindi è necessario dimostrarlo esplicitamente. Si pensi ad esempio a una qualsiasi successione numerica divergente. Ogni elemento di questa successione ha la proprietà di essere finito, ma il limite di questa successione è infinito, e quindi non presenta la proprietà di essere finito.

  Per dimostrare dunque che \(\Sigma^*\) è massimalmente consistente dovremo prima verificare che \(\Sigma^*\) è consistente, e poi verificare che \(\Sigma^*\) è massimale.

  Per assurdo, assumiamo che non sia consistente, ovvero \(\Sigma^*\vdash\bot\). Per definizione di deduzione, deve esistere una deduzione naturale di \(\bot\) da \(\Sigma^*\) che chiamiamo \(\pi\). Una deduzione però è un oggetto finito, dunque \(\pi\) contiene solo un numero finito di foglie, e dunque solo un numero finito di assunzioni non scaricate. Dunque dev'esistere un \(\Gamma\) finito tale che \(\Gamma\subset\Sigma^*\) e \(\Gamma\vdash\bot\). Ma essendo che ogni \(\Sigma_i\) contiene tutti i \(\Sigma_j\) con \(j<i\), posso prendere un \(i\) pari al massimo indice di una formula di \(\Gamma\) rispetto all'enumerazione scelta. Dunque \(\Gamma\subseteq\Sigma_i\), e dunque \(\Sigma_i\vdash\bot\), ovvero \(\Sigma_i\) non è consistente, assurdo.\footnote{Ciò che in questo caso ci permette di dimostrare che il limite mantiene la proprietà degli elementi della successione è il fatto che la deduzione è un oggetto finito, e quindi un'eventuale violazione della proprietà di consistenza da parte del limite deve essersi verificata già in un punto della successione, e quindi in un elemento della successione.}

  Passando dunque alla massimalità, dobbiamo dimostrare che, per ogni formula \(\psi\), o \(\psi\in\Sigma^*\) o \(\neg\psi\in\Sigma^*\) e non entrambe.

  Il fatto che non possano essere entrambe è immediato, poiché se \(\psi\) e \(\neg\psi\) fossero entrambe in \(\Sigma^*\), allora \(\Sigma^*\vdash\bot\) usando \(\neg E\), e quindi \(\Sigma^*\) sarebbe inconsistente, assurdo.

  Per vedere che almeno una delle due deve essere in \(\Sigma^*\), possiamo supporre senza perdita di generalità che \(\phi\) occorra prima di \(\neg\phi\) nell'enumerazione scelta e 
  per assurdo nessuna delle due appartenga a \(\Sigma^*\).

  Sia quindi \(\phi=\psi_j\) e \(\neg\phi=\psi_k\) con \(j<k\).
  Ma allora si ha che
  \[\phi\notin\Sigma^* \implies \psi_j \notin \Sigma_j \iff \Sigma_{j-1}\cup\set{\phi} \vdash \bot\]
  Essendo però che \(j<k\) abbiamo che \(\Sigma_{j-1}\subseteq\Sigma_{k-1}\) e quindi \(\Sigma_{k-1}\cup\set{\phi} \vdash \bot\).
  Si ha anche però che \(\psi_k\notin\Sigma^*\). Dunque
  \[\psi_k \notin \Sigma_k \iff \Sigma_{k-1}\cup\set{\neg\phi}\vdash \bot\]

  A questo punto abbiamo che \(\Sigma_{k-1}\cup\set{\phi}\vdash\bot\) e \(\Sigma_{k-1}\cup\set{\neg\phi}\vdash\bot\). Utilizzando quindi una volta \(\RAA\) e una volta \(\neg I\), si può dedurre \(\Sigma_{k-1}\vdash\neg\phi\) e \(\Sigma_{k-1}\vdash\phi\), e quindi \(\Sigma_{k-1}\vdash\bot\) usando \(\neg E\), assurdo.\footnote{È possibile dare una dimostrazione alternativa della massimalità \(\Sigma^*\) usando la formulazione di massimalità equivalente per cui \(\Sigma^*\) è massimalmente consistente se, \(\forall \phi \in \PL\setminus\Sigma^*, \Sigma^*\cup\set{\phi}\vdash\bot\).
  Per tale dimostrazione è sufficiente mostrare per assurdo che se c'è un \(\phi=\psi_i\) che non appartiene a \(\Sigma^*\) e che non rompe la consistenza, questo porta ad avere \(\Sigma_i\cup\set{\phi}\vdash\bot\), che per monotònicità implica \(\Sigma^*\cup\set{\phi}\vdash\bot\), assurdo.}
\end{proof}

A questo punto, possiamo a dimostrare una proprietà più generale, che ci permetterà di dimostrare la completezza del calcolo di deduzione naturale in modo più semplice, ovvero che un insieme di formule consistente è soddisfacibile.

Per fare ciò pero, dobbiamo andare a introdurre un concetto simile a quello di denotazione, ma in un certo senso duale, ovvero la \keyword{teoria} di un modello.

\begin{definition}{Teoria}
  Dato un modello \(m\), la \keyword{teoria} di \(m\) è l'insieme di tutte le formule che sono soddisfatte da \(m\), ovvero:
  \[\Th{m} = \set{\phi \in \PL }[ m \models \phi]\]
\end{definition}

In un certo senso, sia la denotazione che la teoria collegano il mondo sintattico a quello semantico tramite una sorta di proprietà di chiusura. Infatti, la denotazione di una formula è l'insieme di tutti i modelli che la soddisfano, e quindi è un insieme di modelli, mentre la teoria di un modello è l'insieme di tutte le formule che sono soddisfatte da quel modello, e quindi è un insieme di formule.

Andiamo dunque a vedere qualche proprietà delle teorie di modelli, che ci permetteranno di costruire un modello a partire da un insieme di formule consistente.

\begin{proposition}{Massima consistenza della teoria di un modello}
  Sia \(m\) un modello e \(\Th{m}\) la sua teoria. Allora \(\Th{m}\) è massimalmente consistente.
\end{proposition}

\begin{proof}
  Per definizione di teoria, \(\Th{m} = \set{\phi\in\PL}[m\models\phi]\). Dunque, se \(\Th{m}\) fosse inconsistente, allora \(\Th{m}\vdash\bot\), e quindi dalla correttezza del calcolo di deduzione naturale, \(\Th{m}\models\bot\), e quindi \(m\models\bot\), assurdo. Dunque \(\Th{m}\) è consistente.

  Per dimostrare che \(\Th{m}\) è massimale, dobbiamo dimostrare che, per ogni formula \(\psi\), o \(\psi\in\Th{m}\) o \(\neg\psi\in\Th{m}\) ma non entrambe. Questo però è immediato per semantica della negazione, poiché se \(m\) soddisfa \(\psi\) allora non può soddisfare \(\neg\psi\), e viceversa. Inoltre se \(m\) non soddisfa \(\psi\) allora soddisfa \(\neg\psi\).
\end{proof}

\begin{proposition}[label=prop:max-consistent-set-are-theory]{Insiemi massimalmente consistenti sono teorie}
  Sia \(\Sigma^* \subseteq \PL\) massimalmente consistente. Allora esiste un modello \(m\) tale che \(\Sigma^* = \Th{m}\).
\end{proposition}

\begin{proof}
  Dato \(\Sigma^*\) massimalmente consistente, dobbiamo costruire un modello \(m\) tale che \(\Sigma^* = \Th{m}\).
  Per fare ciò, basta definire \(m = \Sigma^*\cap\AP\), poiché essendo \(\Sigma^*\) massimalmente consistente, contiene tutte e sole le proposizioni atomiche che non portano a inconsistenza.
  Infatti le uniche proposizioni atomiche che non possono essere contenute in \(\Sigma^*\) sono quelle che, se aggiunte a \(\Sigma^*\), rendono \(\Sigma^*\) inconsistente, ovvero appaiono in qualche formula che è già presente in \(\Sigma^*\) e che le richiede con assegnamento falso per essere soddisfatte.

  A questo punto, dobbiamo mostrare che \(\Sigma^* = \Th{m}\), ovvero che per ogni formula \(\phi\), \(\phi \in \Sigma^* \iff m \models \phi\).

  Per fare ciò possiamo mostrare che \(\Sigma^*\) è chiuso per deduzione, ovvero che se \(\Sigma^*\vdash\phi\) allora \(\phi\in\Sigma^*\).
  Sia infatti per assurdo \(\phi\not\in\Sigma^*\) tale che \(\Sigma^*\vdash\phi\). Dalla massimalità di \(\Sigma^*\) si ha che \(\neg\phi\in\Sigma^*\). Dunque \(\Sigma^*\vdash\neg\phi\) banalmente, e quindi \(\Sigma^*\vdash\bot\) usando \(\neg E\), assurdo.

  A questo punto, possiamo procedere per induzione sulla struttura di \(\phi\) per mostrare che \(\phi \in \Sigma^*\) se e solo se \(m\models\phi\).

  Il caso base è una formula atomica \(p\). Se \(p\in\Sigma^*\), allora \(p\in m\) per costruzione, e quindi \(m\models p\) e viceversa.
  A questo punto, sia \(\phi\) una formula non atomica. In questo caso, \(\phi\) è costruita a partire da formule di lunghezza minore tramite l'applicazione di una regola di formazione. L'ipotesi induttiva è dunque che per ogni \(\psi\) di lunghezza minore di \(\phi\),\(\psi \in \Sigma^*\) se e solo se \(m\models\psi\).
  \begin{description}
    \item[\(\phi = \neg\psi_1\)] Dalla semantica della negazione, \(m\models \phi \iff m \not\models \psi\), e per ipotesi induttiva \(m\not\models \psi\iff \psi \notin\Sigma^*\). Dalla definizione di massimalità di \(\Sigma^*\) si ha che \(\psi\notin\Sigma^*\) se e solo se \(\neg\psi\in\Sigma^*\), ovvero \(\phi\in\Sigma^*\).
    \item[\(\phi = \psi_1 Op \psi_2\)] Tutte le dimostrazioni per operatori binari sono analoghe e si basano sull'uso della chiusura per deduzione e sulle regole di deduzione naturale specifiche dell'operatore. Vediamo come esempio il caso di \(Op=\land\).
        Dalla semantica della congiunzione, \(m\models \phi\) se e solo se \(m\models \psi_1\) e \(m\models \psi_2\). Per ipotesi induttiva, \(m\models \psi_1\) se e solo se \(\psi_1\in\Sigma^*\), e \(m\models \psi_2\) se e solo se \(\psi_2\in\Sigma^*\). Ci basta quindi mostrare che \(\psi_1,\psi_2 \in \Sigma^*\iff \psi_1\land\psi_2\in\Sigma^*\).
        Ma essendo che \(\psi_1,\psi_2 \vdash_{\land I} \psi_1\land\psi_2\) e \(\psi_1\land\psi_2\vdash_{\land E} \psi_1\) e \(\psi_1\land\psi_2\vdash_{\land E} \psi_2\), dalla chiusura di \(\Sigma^*\) per deduzione si ha la tesi.
  \end{description}
\end{proof}

\begin{lemma}[namedlabel={lem:model-existence}{Esistenza del modello}]{Esistenza del modello}
  Dato \(\Sigma\subseteq\PL\), si ha che:
  
  \(\Sigma\) è consistente \(\iff\Sigma\) è soddisfacibile.
\end{lemma}

\begin{proof}
  Per definizione di consistenza \(\Sigma \not\vdash \bot\), ovvero non esiste alcun albero di deduzione con radice \(\bot\) con foglie che sono o assunzioni scaricate o appartengono a \(\Sigma\).
  Essere soddisfacibile significa che esiste un modello \(m\) tale che \(m\models\Sigma\).

  Dunque dobbiamo andare a costruire un modello \(m\) che soddisfa \(\Sigma\) sapendo che le formule di \(\Sigma\) sono consistenti.

  Ma a questo punto, se \(\Sigma\) è consistente, allora per il \Namedcref{lem:lindenbaum} è possibile costruire un insieme di formule \(\Sigma^*\) tale che \(\Sigma\subseteq\Sigma^*\) e \(\Sigma^*\) è massimalmente consistente.
  Dunque, per la \Cref{prop:max-consistent-set-are-theory}, esiste un modello \(m\) tale che \(\Sigma^* = \Th{m}\). Dunque, essendo \(\Sigma\subseteq\Sigma^*\), allora \(\Sigma\subseteq\Th{m}\), e quindi \(m\models\Sigma\), ovvero \(m\) soddisfa tutte le formule di \(\Sigma\), e quindi \(\Sigma\) è soddisfacibile.

  Viceversa, se \(\Sigma\) è soddisfacibile, allora esiste un modello \(m\) tale che \(m\models\Sigma\). Per definizione di teoria, \(\Sigma\subseteq\Th{m}\).
  Essendo che \(\Th{m}\) è consistente, allora anche \(\Sigma\) è consistente, poiché se \(\Sigma\vdash\bot\) allora \(\Th{m}\vdash\bot\) per monotònicità, assurdo.
\end{proof}

A questo punto, siamo finalmente in grado di dimostrare la completezza del calcolo di deduzione naturale, che ci dice che tutto ciò che è conseguenza logica è anche deducibile, e quindi che il calcolo di deduzione naturale è in grado di catturare tutta la conseguenza logica.

\begin{theorem}{Completezza del calcolo di deduzione naturale}
  Per ogni insieme di formule \(\Gamma\) e ogni formula \(\phi\) si ha che:
  \[\Gamma \models \phi \implies \Gamma \vdash \phi\]
\end{theorem}

\begin{proof}
  Come già visto precedentemente, abbiamo che \(\Gamma\models\phi \iff \Gamma \cup \set{\neg\phi}\) è insoddisfacibile.

  Inoltre, possiamo facilmente vedere che \(\Gamma\vdash \phi \iff \Gamma \cup \set{\neg\phi}\) è inconsistente.

  Infatti, se \(\Gamma \cup \set{\neg\phi}\) è inconsistente allora \(\Gamma \cup \set{\neg\phi}\vdash \bot\), ovvero esiste un albero di deduzione per \(\bot\) da \(\Gamma \cup \set{\neg\phi}\). Dunque, se scarichiamo l'assunzione \(\neg\phi\) con \(\RAA\), otteniamo un albero di deduzione per \(\phi\) da \(\Gamma\), ovvero \(\Gamma\vdash\phi\).

  Viceversa, se \(\Gamma\vdash\phi\) allora esiste un albero di deduzione per \(\phi\) da \(\Gamma\). Se aggiungiamo a questo albero di deduzione un'assunzione \(\neg\phi\) e applichiamo \(\neg E\), allora otteniamo un albero di deduzione per \(\bot\) da \(\Gamma \cup \set{\neg\phi}\), ovvero \(\Gamma \cup \set{\neg\phi}\vdash\bot\), e quindi \(\Gamma \cup \set{\neg\phi}\) è inconsistente.

  Di conseguenza, \(\Gamma\models\phi \implies \Gamma\vdash\phi\) è equivalente a \(\Gamma \cup \set{\neg\phi}\) è insoddisfacibile \(\implies\) \(\Gamma \cup \set{\neg\phi}\) è inconsistente.

  A questo punto, possiamo dimostrare la proprietà di completezza per contrapposizione, ovvero che, preso un qualsiasi insieme \(\Gamma\), \(\Gamma\) consistente \(\implies\) \(\Gamma\)soddisfacibile.

  Ma per fare ciò è sufficiente il \Namedonlyhref{lem:model-existence}, che ci dice che un insieme di formule consistente è soddisfacibile. Di conseguenza, se \(\Gamma\cup\set{\neg\phi}\) è insoddisfacibile, allora è inconsistente, ovvero se \(\Gamma\models\phi\) allora \(\Gamma\vdash\phi\).
\end{proof}

Esistono infine un ultimo paio di risultati importanti riguardo il calcolo di deduzione naturale della logica proposizionale, entrambi che in un certo qual senso semplificano la trattazione di questo calcolo, uno dal fronte semantico, riducendo la grandezza dello spazio dei modelli da considerare, e uno dal fronte sintattico, permettendo di \qi{modularizzare} le dimostrazioni, e quindi di semplificarle.

\begin{theorem}{Compattezza}
  Sia \(\Gamma \subseteq \PL\) e \(\phi \in \PL\). Allora \(\Gamma \models \phi\) se e solo se esiste un sottoinsieme finito \(\Gamma_0 \subseteq \Gamma\) tale che \(\Gamma_0 \models \phi\).
\end{theorem}

Una formulazione alternativa della proprietà di correttezza in termini di soddisfacibilità è
\[\Unsat[\Gamma] \iff \exists \Gamma_0 \subseteq \Gamma, \Gamma_0 \text{ finito }: \Unsat[\Gamma_0]\]
per dualità
\[\Sat[\Gamma] \iff \forall \Gamma_0 \subseteq \Gamma, \Gamma_0 \text{ finito }: \Sat[\Gamma_0]\]
Tale teorema è di particolare importanza, poiché ci garantisce che nella verifica dell'insoddisfacibilità di un insieme di formule, è sufficiente considerare solo modelli finiti.

\begin{proof}
  Abbiamo che 
  \[\Gamma \models \phi \iff \Gamma \vdash \phi\]
  Ciò è equivalente all'esistenza di una deduzione \(\pi\) di \(\phi\) da \(\Gamma\).
  Sia allora \(\Gamma_0 = leaf(\pi)\cap \Gamma\). Ovviamente \(\Gamma_0\subseteq\Gamma\) e \(\abs{\Gamma_0}<\infty\).
  Inoltre, per costruzione di \(\Gamma_0\), \(\Gamma_0\vdash\phi\), e quindi \(\Gamma_0\models\phi\) per correttezza.
  Sia allora \(\Gamma' = \set{\psi \in \pi}[\psi \text{ foglia non scaricata di }\pi]\).
  Ma essendo \(\pi\) finito, \(\Gamma'\) è finito, e banalmente \(\Gamma' \subseteq \Gamma\) per definizione di deduzione.
  Si ha ovviamente che \(\Gamma' \vdash \phi\) per costruzione di \(\Gamma'\), e quindi \(\Gamma' \models \phi\) per correttezza.

  Nell'altro verso è ovvio per monotònicità.
  \[\Gamma'\models\phi\iff\Gamma'\vdash\phi\implies\Gamma\vdash\phi\iff\Gamma\models\phi\]
\end{proof}

Arriviamo infine a un teorema che generalizza il \Namedcref{thm:substitution}, che riguarda la sostituzione nella conseguenza logica invece che nella validità

\begin{theorem}{Invarianza della Conseguenza Logica rispetto alla Sostituzione}
  Sia \(\Gamma \subseteq \PL\), \(\phi \in \PL\) e \(\sigma:\AP \to \PL\) sostituzione. Allora,
  \[\Gamma\models\phi\implies\Gamma\sigma\models\phi\sigma\]
\end{theorem}

\begin{proof}
  Abbiamo già visto nel teorema di sostituzione che \(\models\phi \implies \models\phi\sigma\) per ogni \(\sigma\).

  Possiamo usare in questo caso il teorema di compattezza. Infatti da \(\Gamma\models\phi\) segue che esiste un sottoinsieme finito \(\Gamma_0 \subseteq \Gamma\) tale che \(\Gamma_0\models\phi\).
  Essendo \(\Gamma_0\) finito sappiamo che \(\Gamma_0\models\phi\iff\models(\bigwedge\Gamma_0)\to\phi\), e quindi \(\models((\bigwedge_{\phi\in\Gamma_0}\phi)\to\phi)\sigma\) per il teorema di sostituzione.

  Per la definizione di applicazione di una sostituzione \(\sigma\) si ha che \(((\bigwedge\Gamma_0)\to\phi)\sigma\) è equivalente a \((\bigwedge_{\phi\in\Gamma_0}\phi\sigma)\to\phi\sigma\).
  Quindi \(\Gamma_0\sigma\models\phi\sigma\) per il teorema di deduzione semantica e quindi \(\Gamma\sigma\models\phi\sigma\) per monotònicità.
\end{proof}

\begin{corollary}{Sostituzione per la Deducibilità}
  \(\Gamma\vdash\phi\implies\Gamma\sigma\vdash\phi\sigma\), con \(\sigma:\AP \to \PL\) sostituzione.
\end{corollary}
\begin{proof}
  Segue banalmente dal teorema precedente e da completezza e correttezza.
\end{proof}

Grazie a questo teorema, è possibile generare nuove deduzioni a partire da deduzioni già esistenti, semplicemente applicando una sostituzione uniformemente a tutte le formule che compaiono nella deduzione. In questo modo, è possibile \qi{modularizzare} le dimostrazioni, ovvero dimostrare prima un risultato per una formula generica, e poi ottenere il risultato per formule specifiche semplicemente applicando una sostituzione alla deduzione della formula generica.

\section{Calcolo di Risoluzione proposizionale}

Il calcolo di deduzione naturale è un calcolo di deduzione molto potente, che come da nome suggerisce, risulta naturale per la dimostrazione di teoremi, poiché le regole di deduzione si allineano molto bene con le tecniche di dimostrazione che siamo abituati a utilizzare.

Tuttavia, nonostante sia molto naturale per il ragionamento umano, il calcolo di deduzione naturale non è particolarmente adatto per l'automatizzazione, poiché la ricerca di una deduzione per una formula con un numero elevato di regole da poter applicare risulta difficile da gestire algoritmicamente.

Andiamo quindi a introdurre un calcolo di deduzione alternativo, che risulta più adatto per l'automatizzazione, poiché composto da un'unica regola di deduzione, il \keyword{calcolo di risoluzione}. La presenza di un unica regola di deduzione elimina la necessità di dover scegliere quale regola applicare, lasciando come unica scelta da fare quella di quali formule utilizzare per applicare la regola, e quindi semplificando notevolmente la ricerca di una deduzione.

Una prima versione fu proposta da J. A. Robinson nel 1965, e da allora è stata oggetto di numerose estensioni e miglioramenti.

Tale calcolo non si concentra sulla conseguenza logica, ma sull'inconsistenza. 
Nonostante ciò, per decidere \(\Gamma\models\phi\) è possibile farlo decidendo se \(\Gamma\cup\set{\neg\phi}\) è insoddisfacibile.
In questo senso la risoluzione non è un calcolo \qi{generativo}, come lo era la deduzione naturale, in cui dato un insieme di formule veniva generata la conseguenza logica.

Inoltre, tale calcolo non si applica direttamente a formule proposizionali, ma solo a formule in \keyword{forma clausale}.

\begin{definition}{Forma clausale}
  Una formula \(\phi\) è in forma clausale se \(\phi \in 2^{2^\mathcal{L}}\), dove \(\mathcal{L}\) è un insieme di letterali.

  In altre parole, una formula è in forma clausale se è un insieme di insiemi di letterali.
\end{definition}

Ovviamente, essendo insiemi, l'ordine degli elementi non conta, ogni insieme \qi{interno} non può contenere due volte lo stesso letterale, e l'insieme esterno non può contenere due volte lo stesso insieme di letterali.

Esiste inoltre una corrispondenza biunivoca tra formule in forma clausale e formule proposizionali in forma normale congiuntiva.
È possibile infatti interpretare ogni insieme interno come una disgiunzione di letterali, e l'insieme esterno come una congiunzione di tali disgiunzioni.

Inoltre possiamo notare che dato un insieme di formule in \(\CNF\) la congiunzione di tutte le formule è una formula in \(\CNF\) equivalente a tutte le formule dell'insieme, e quindi possiamo sempre rappresentare un insieme di formule \(\Gamma\subseteq\CNF\) con la formula \(\bigwedge_{\phi\in\Gamma}\phi\). E per la forma clausale chiaramente questo equivale a fare l'unione tra formule in forma clausale.

Ovviamente quindi per ogni \(\Sigma\subseteq\PL\) esiste un \(F_\Sigma\) in forma clausale equivalente a \(\Sigma\).
Questo si riduce a dire che \(\forall\phi\in\PL\) esiste un \(F_\phi\) in forma clausale equivalente a \(\phi\), e questo è vero poiché ogni formula è equivalente a una formula in \(\CNF\) e ogni formula in \(\CNF\) è equivalente a una formula in forma clausale.

Quando \(F = \emptyset\) è detta forma clausale banale ed è valida.

Questo poiché, intuitivamente, \(m\models F\) se e soltanto se \(\forall C \in F\) insieme di letterali, \(m\models C\) cioè \(m\models l\) per qualche \(l \in C\) per ogni \(C\in F\).
Quindi \(m\models F \iff \forall C \in F ,\exists l \in C\st m\models l\). Ovviamente se \(F =\emptyset\) qualsiasi \(m\) soddisfa \(F\) poiché non c'è nessuna clausola da poter non soddisfare.


È importante notare che la forma clausale vuota \(F=\emptyset\) non coincide con la clausola vuota \(C=\emptyset\), tipicamente indicata con \(\square\). Per la definizione di soddisfacibilità, \(m\not\models \square\) per ogni \(m\), e quindi \(\square\) è insoddisfacibile, mentre \(F=\emptyset\) è soddisfatta da ogni \(m\).

Essendo \(\square\) insoddisfacibile, per valutare l'insoddisfacibilità di un insieme di formule in forma clausale, è sufficiente verificare se è possibile dedurre \(\square\) da quell'insieme di formule.

\begin{example}{}
  Sia 
  \[\phi=(P\to(Q\land R))\land(Q\lor(Q\to P)) \land((\neg P \to Q)\lor(\neg Q \to R))\]
  che in \(\CNF\) è equivalente a
  \[\phi_{\CNF} = (\neg P\lor Q)\land(\neg P\lor R)\land(Q\lor \neg Q \lor P)\land(P\lor Q \lor Q \lor R)\]
  ottenibile applicando De Morgan, distributività, implicazione materiale e semplificazione di doppie negazioni.
  La forma clausale equivalente è quindi
  \[F_\phi = \set{\set{\neg P, Q},\set{\neg P, R},\set{Q, \neg Q, P},\set{P, Q, R}}\]
  Ogni clausola che contiene \(P\) e \(\neg P\) per un qualche \(P \in AP\) è tautologica, ed essendo \(\phi\land\top\equiv\phi\) è possibile rimuovere la clausola \(\set{Q, \neg Q, P}\) da \(F_\phi\) senza alterare la soddisfacibilità di \(F_\phi\).

  Da ora in poi quindi assumeremo che le forme clausali non contengano clausole banali.
\end{example}

Per andare a definire formalmente il calcolo di risoluzione, ovvero definire la sua unica regola di deduzione, anch'essa chiamata risoluzione, è necessario prima definire una proprietà di clausole che sarà la condizione di applicabilità della regola di risoluzione, ovvero la proprietà di essere in \keyword{contrasto} o in \keyword{clash}.

\begin{definition}{Clausole in contrasto}
  Siano \(C_1\) e \(C_2\) due clausole. Diremo che \(C_1\) e \(C_2\) sono in \keyword{contrasto}, o in \keyword{clash}, se esiste un \(l \in \mathcal{L}\) tale che \(l \in C_1\) e \(\tilde{l} \in C_2\) con
  
  \[\tilde{l} = \begin{cases}
                \neg P & \text{se } l = P\\
                P & \text{se } l = \neg P
              \end{cases}
  \]
\end{definition}

In altre parole, due clausole sono in clash se contengono un letterale su cui non sono concordi. Intuitivamente, per soddisfarle entrambe, è necessario che almeno una delle due \qi{rinunci} ad avere ragione su quel letterale, dovendo aver soddisfatto un letterale diverso da quello di clash. Non potendo però sapere a priori quale delle due clausole rinuncerà ad avere quel letterale soddisfatto, l'unica cosa che possiamo intuitivamente dedurre è che almeno uno dei letterali della prima clausola o almeno uno della seconda dovrà essere soddisfatto. Tale ragionamento è proprio ciò che caratterizza la regola di risoluzione

\begin{definition}{Regola di risoluzione}
  Siano \(C_1\) e \(C_2\) due clausole in contrasto su un letterale \(l\), con \(l \in C_1\) e \(\tilde{l} \in C_2\). Allora, è possibile dedurre la clausola \(R\) detta \keyword{risolvente} di \(C_1\) e \(C_2\) su \(l\), definita come:

  \[\infer[\text{Risoluzione}]{R = (C_1 \cup C_2)\setminus\set{l, \hat{l}}}{C_1 \quad C_2}\]

\end{definition}

Tale regola di deduzione segue esattamente il ragionamento intuitivo fatto in precedenza applicato alle forme clausali, in cui le clausole sono considerate in congiunzione, ovvero da soddisfare entrambe, e i letterali al loro interno in disgiunzione. L'insieme \((C_1 \cup C_2)\setminus\set{l, \hat{l}}\) è esattamente l'insieme contenente tutti i letterali delle due clausole tranne quelli in clash. Interpretandoli in disgiunzione, tale insieme corrisponde semanticamente esattamente al richiedere che ne sia soddisfatto almeno uno, che intuitivamente dovrà appartenere alla clausola che non avrà il letterale di clash soddisfatto.

Andiamo dunque a formalizzare la fondatezza semantica di tale regola.

\begin{proposition}{}
  La regola di risoluzione preserva la verità
\end{proposition}

\begin{proof}
  Formalmente, una regola di inferenza preserva la verità se ogni modello che soddisfa le premesse soddisfa anche la conclusione della regola.
  Chiaramente, il preservare la verità implica il preservare anche la conseguenza logica.\footnote{Questa cosa risulta banale nella logica proposizionale, me non sarà scontata nel prim'ordine.}

  Sia \(m\) un modello che soddisfa \(C_1\) e \(C_2\) che sono in clash su \(l\).
  Per definizione di \(l\) e \(\tilde{l}\), \(m\) soddisfa esattamente uno tra \(l\) e \(\tilde{l}\). Senza perdita di generalità, supponiamo che \(l\in C_1\) e \(\tilde{l}\in C_2\).

  Consideriamo il caso \(m\models l\). Essendo che \(m\models C_2\), allora \(\exists l' \in C_2\setminus{\tilde{l}}\) tale che \(m\models l'\). Dunque \(l' \in R\), ovvero \(m\models R\).
  Nel caso in cui \(m \models \tilde{l}\), essendo che \(m\models C_1\), allora \(\exists l' \in C_1\setminus{l}\) tale che \(m\models l'\). Dunque \(l' \in R\), ovvero \(m\models R\).
\end{proof}

Per quanto riguarda il verso opposto invece, l'unica proprietà che è preservata è la soddisfacibilità, ovvero se \(R\) risolvente di \(C_1\) e \(C_2\) su \(l\) è soddisfacibile, allora anche \(C_1\) e \(C_2\) sono soddisfacibili. Tale proprietà è ovviamente più debole della preservazione della verità, poiché l'unica garanzia che fornisce è che se esiste un modello in cui \(R\) è soddisfatto ne esiste uno, potenzialmente diverso, in cui \(C_1\) e \(C_2\) sono congiuntamente soddisfatte.

Non è possibile richiedere di più poiché \(R\) contiene un numero di proposizioni atomiche strettamente inferiore sia di \(C_1\) che di \(C_2\), e di conseguenza gli assegnamenti che soddisfano \(R\) non sono assegnamenti validi per soddisfare \(C_1\) e \(C_2\). Andiamo però a mostrare che dato \(\alpha: \AP[R] \to \set{0,1}\) tale che \(\Val{\alpha}[R] = 1\), esiste un assegnamento \(\alpha':\AP[C_1]\cup\AP[C_2] \to \set{0,1}\) tale che \(\Val{\alpha'}[C_1] = 1\) e \(\Val{\alpha'}[C_2] = 1\).

\begin{proposition}{}
  Sia \(R\) risolvente di \(C_1\) e \(C_2\) in clash su \(l\). Allora, se \(R\) è soddisfacibile, allora anche \(C_1\) e \(C_2\) sono soddisfacibili.
\end{proposition}

\begin{proof}
  Sia \(\alpha: \AP[R] \to \set{0,1}\) tale che \(\Val{\alpha}[R] = 1\).

  Dato che \(\Val{\alpha}[R] = 1\), per la semantica della forma clausale, \(\exists l' \in R\) tale che \(\Val{\alpha}[l'] = 1\).
  Per costruzione di \(R\), \(l'\) è un letterale che appartiene a \(C_1\) o a \(C_2\) o a entrambe.

  Consideriamo per primo il caso in cui \(l'\) appartiene a entrambe. In tal caso l'assegnamento \(\alpha': \AP[C_1\cup C_2]\to\set{0,1}\) definita come
  \[\alpha'(P) = \begin{cases}
                \alpha(P) & \text{se } P \in AP(R)\\
                0 & \text{se } P \not\in AP(R) 
              \end{cases}\]
              soddisfa sia \(C_1\) che \(C_2\), poiché \(l'\) appartiene a \(R\) e quindi tale letterale soddisfa sia \(C_1\) che \(C_2\).

  Consideriamo ora il caso in cui \(l'\) appartiene a \(C_1\) ma non a \(C_2\). In tal caso, l'assegnamento \(\alpha': \AP[C_1\cup C_2]\to\set{0,1}\) definito come
  \[\alpha'(P) = \begin{cases}
                \alpha(P) & \text{se } P \in AP(R)\\
                0 & \text{se } \tilde{l} = \neg P\\
                1 & \text{se } \tilde{l} = P
              \end{cases}
  \]
  soddisfa ancora entrambe, poiché \(C_1\) sarà soddisfatta da \(l'\), mentre \(C_2\) sarà soddisfatta da \(\tilde{l}\), poiché il suo assegnamento è costruito per essere soddisfatto qualunque sia la sua polarità.

  Analogamente, se \(l'\) appartiene a \(C_2\) ma non a \(C_1\), allora l'assegnamento \(\alpha': \AP[C_1\cup C_2]\to\set{0,1}\) definito come
  \[\alpha'(P) = \begin{cases}
                \alpha(P) & \text{se } P \in AP(R)\\
                0 & \text{se } l = \neg P\\
                1 & \text{se } l = P
              \end{cases}
  \]
  soddisfa ancora entrambe, per un ragionamento duale rispetto al caso precedente.
\end{proof}

Avendo ora un insieme di regole di inferenza, in questo caso composto da una regola sola, e un insieme di assunzioni, in questo caso la forma clausale, possiamo andare a definire il calcolo di deduzione di risoluzione attraverso la definizione di una deduzione.

\begin{definition}{Deduzione di risoluzione}
  Sia \(F\) una forma clausale. Una deduzione di risoluzione \(F\vdash_{\Rule }C\) da è un albero binario finito, la cui radice è etichettata con \(C\), tutte le foglie sono etichettate con clausole di \(F\) e ogni nodo interno è etichettato con il risolvente dei suoi figli.
\end{definition}

\begin{note}{}
  Essendo che in risoluzione \(\Rules = \set{\Rule}\), con \(\Rule\) la regola di risoluzione, per distinguere tale calcolo da quello di deduzione naturale indicheremo le deduzioni di risoluzione con \(\vdash_{\Rule}\) invece che con \(\vdash\). Quando sarà necessario esplicitare l'utilizzo del calcolo di deduzione naturale invece, per dissipare ogni ambiguità, verrò utilizzato \(\vdash_{ND}\).
\end{note}

Essendo che la regola di risoluzione preserva la verità, e quindi la conseguenza logica, è facile vedere per induzione strutturale la correttezza del calcolo, ovvero che date \(S,C\) forme clausali, \(S\vdash_{\Rule} C \implies S\models C\).

Sarà possibile dimostrare anche una forma più debole della completezza, legata alla soddisfacibilità, ovvero che \(S\not\vdash_R \square \implies \Sat[S]\), o equivalentemente \(\Unsat[S] \implies S\vdash_R \square\). Come già anticipato, non è possibile dimostrare la totale completezza, perché esistono conseguenze logiche di clausole che non sono in alcun modo generabili dal calcolo.

Consideriamo infatti la forma clausale \(F = \set{\set{A},\set{B}}\). Dall'interpretazione semantica data, è ovvio che \(R=\set{A,B}\) sia conseguenza logica di \(F\). Essendo però che non esistono letterali in clash in \(F\), non è possibile applicare la regola di risoluzione, impedendo di derivare alcunché.

È possibile però in parte aggirare questo problema utilizzando la definizione alternativa di conseguenza logica, notando che \(S\models C \iff S\cup\set{\neg C} \in \Unsat\).
Applicando tale riformulazione all'esempio di prima, potremo dimostrare che \(F\models R\) cercando di mostrare che \(F' = \set{\set{A}}, \set{B}, \set{\neg A, \neg B}\) è insoddisfacibile. In questo caso, essendo che \(C_1 = \set{A}\) e \(C_2 = \set{\neg A, \neg B}\) sono in clash su \(A\), è possibile applicare la regola di risoluzione per dedurre \(C_3 = \set{\neg B}\), da cui è possibile dedurre \(C_4 = \square\) con \(C_2\), che è in clash con \(C_3\) su \(B\). Dunque, \(F'\) è insoddisfacibile per correttezza, e quindi \(F\models R\) per la definizione di conseguenza logica.

Utilizzando questo tipo di riformulazione, e questa completezza \qi{debole}, è possibile risolvere il problema di conseguenza logica usando risoluzione, passando per il problema di soddisfacibilità.

Prima della dimostrazione formale di questa forma di completezza debole però, andiamo a vedere come automatizzare il calcolo di risoluzione tramite un algoritmo.

\begin{algorithmdesc}{Risoluzione}
    \begin{algorithmic}[1]
      \Statex{} \textbf{signature} \(\textsc{Resolution:} \quad [Clause] \to \Unsat \mid \Sat\)
      \Function{\(\textsc{Resolution}\)}{$S$} % chktex 46
          \State{} \(S_0\gets S\)
          \State{} \(E \gets \emptyset\)
          \Comment{\(E \subseteq S_0 \times S_0\)}
          \Repeat{}
            \State{} Pick \((C_1, C_2) \in (S_0 \times S_0)\setminus E \st \operatorname{clash}(C_1, C_2)\)
            \State{} \(R \gets \operatorname{resolve}(C_1, C_2)\)
            \If{\(R\) non banale}
              \State{} \(S_0 = S_0 \cup \set{R}\)
            \EndIf{}
            \State{} \(E \gets E \cup \set{(C_1, C_2)}\)
          \Until{\(R = \square \lor E = S_0 \times S_0 \cap \set{(C_1,C_2) \in S_0\times S_0}[\operatorname{clash}(C_1, C_2)] \)}
          \If{\(R = \square\)}
            \State{} \Return{\(\Unsat\)}
          \Else{}
            \State{} \Return{\(\Sat\)}
          \EndIf{}
      \EndFunction{}
    \end{algorithmic}
\end{algorithmdesc}

Per quanto riguarda la correttezza dell'algoritmo, questa risulta quasi scontata dal fatto che è l'esatta trasposizione di una deduzione di risoluzione. Riguardo alla terminazione invece, la cosa è meno scontata. Questo poiché la condizione di terminazione è e dev'essere valutata sull'insieme \(S_0\), il quale non è statico, ma cresce dinamicamente durante l'esecuzione. In linea di principio quindi, senza mostrare che esista un limite superiore al numero di clausole che è possibile generare mediante risoluzione, non è possibile garantire la terminazione di questo algoritmo.

Fortunatamente, tale limite superiore esiste, e ne vedremo dimostrazione formale dopo un esempio di applicazione di tale algoritmo.

\begin{example}{}
  Consideriamo l'\Cref{ex:logical_consequence} visto nel primo capitolo, formalizzato usando proposizioni atomiche per rappresentare le affermazioni legate al nome della persona coinvolta in ognuna di esse. In particolare sia:

  \begin{itemize}
    \item \(C\) la proposizione atomica che rappresenta l'affermazione \qi{Carlo è italiano}
    \item \(G\) la proposizione atomica che rappresenta l'affermazione \qi{Giovanni è spagnolo}
    \item \(R\) la proposizione atomica che rappresenta l'affermazione \qi{Rachele è tedesca}
    \item \(L\) la proposizione atomica che rappresenta l'affermazione \qi{Lucia è francese}
  \end{itemize}

  Possiamo formalizzare il problema creando l'insieme di formule \(\Gamma\):
  \begin{itemize}
    \item \((C\land \neg G)\to R\)
    \item \(R \to (L\lor G)\)
    \item \(\neg L \to C\)
    \item \(\neg G\)
  \end{itemize}

  Volevamo dimostrare che tale insieme avesse come conseguenza logica \(L\).
  Con risoluzione, prendiamo \(\Gamma \cup \set{\neg L}\) e verificare che sia insoddisfacibile.
  Dovremo trasformare quindi le formule in forma clausale passando per la \(\CNF\), e otteniamo:
  \begin{enumerate}
    \item \(\neg(C\land\neg G) \lor R \equiv \neg C \lor G \lor R\), ovvero \(\set{\neg C,G,R}\)
    \item \(\neg R \lor L \lor G\), ovvero \(\set{\neg R, L, G}\)
    \item \(L \lor C\), ovvero \(\set{L, C}\)
    \item \(\neg G\), ovvero \(\set{\neg G}\)
    \item \(\neg L\), ovvero \(\set{\neg L}\)
  \end{enumerate}

  Dobbiamo allora vedere se da queste \(5\) clausole è possibile derivare la clausola vuota \(\square\) tramite risoluzione.

  In primo luogo ho diverse scelte, prendiamo ad esempio le clausole \(2\) e \(4\) che fanno clash su \(G\). Il risolvente è \(\set{\neg R, L}\), chiamiamo questo risolvente \(6\)

  Da \(1\) e \(4\) che fanno clash su \(G\) ottengo \(\set{\neg C, R}\), chiamiamo questo risolvente \(7\).
  Da \(3\) e \(7\) otteniamo \(\set{L, R}\), chiamiamo questo risolvente \(8\).
  Da \(6\) e \(8\) otteniamo \(\set{L}\), chiamiamo questo risolvente \(9\).
  Da \(5\) e \(9\) otteniamo \(\square\).

  Ovviamente tale esecuzione non è unica. Ad esempio, si può iniziare da \(1\) e \(2\) in clash su \(R\) per ottenere \(\set{\neg C, G, L}\), che chiamiamo \(6\).

  Possiamo quindi risolvere \(6\) con \(3\) ottenendo \(\set{G, L}\), che chiamiamo \(7\).
  A questo punto risolvendo prima con \(4\), e poi il risolvente con \(5\) otteniamo \(\square\).
\end{example}

È fondamentale notare che l'applicazione della regola di risoluzione richiede la selezione di due clausole in clash su uno specifico letterale, e che quindi una singola applicazione della regola porti all'eliminazione di un solo letterale alla volta.

Non è possibile estendere banalmente la regola per permettere la risoluzione di più letterali contemporaneamente, poiché in tal caso è facile costruire controesempi in cui non è preservata la verità. Prendiamo infatti le clausole \(C_1 = \set{A, B}\) e \(C_2 = \set{\neg A, \neg B}\). Se fosse possibile risolvere entrambe le coppie di letterali in clash contemporaneamente, allora sarebbe possibile dedurre la clausola vuota \(\square\) da \(C_1\) e \(C_2\), che è insoddisfacibile, mentre invece \(C_1\) e \(C_2\) sono congiuntamente soddisfacibili da un assegnamento che assegna \(1\) ad \(A\) e \(0\) a \(B\), o viceversa. Andando ad applicare la regola di risoluzione in modo corretto, è possibile dedurre \(\set{A,\neg A}\) da \(C_1\) e \(C_2\) risolvendo su \(B\), ottenendo la clausola banale \(\set{A, \neg A}\), che è tautologica.

Fatte queste precisazioni, possiamo ora dimostrare la terminazione dell'algoritmo di risoluzione, mostrando che esiste un limite superiore al numero di clausole che tale algoritmo può generare, limitando quindi la grandezza dell'insieme \(S_0\) su cui viene valutata la condizione di terminazione.

\begin{theorem}{}
  \(\textsc{Resolution}(S)\) termina sempre e l'insieme di saturazione \(S_0\) avrà al più \(3^{\abs{\AP[S]}}\) clausole.
\end{theorem}

\begin{proof}
  Sappiamo che l'insieme di clausole in ingresso \(S\) è finito. Sia dunque \(\AP[S] = \set{P \in \AP[\phi]}[\phi \in S]\) l'insieme di proposizioni atomiche che compaiono in \(S\). Chiaramente \(\AP[S]\) è finito, poiché \(S\) è finito e ogni formula contiene un numero finito di proposizioni atomiche.

  Inoltre, avendo escluso le clausole banali, si ha che \(\forall C \in S, \abs{C}\leq\AP[S]\). Questo poiché ogni clausola è un insieme di letterali, che può essere presente o in forma positiva o in forma negativa, ma non entrambe, e quindi per ogni proposizione atomica in \(\AP[S]\) al più una delle due forme può essere presente in una clausola.

  Quindi, in ogni clausola, una data proposizione atomica \(P\) può essere presente in forma positiva, in forma negativa, o non essere presente affatto. Di conseguenza ogni clausola può essere vista come una funzione da \(\AP[S]\) a \(\set{-1,0,1}\), dove \(0\) rappresenta l'assenza di \(P\) nella clausola, \(1\) rappresenta la presenza di \(P\) in forma positiva, e \(-1\) rappresenta la presenza di \(P\) in forma negativa.

  Dunque, il numero di clausole che è possibile generare con le proposizioni atomiche in \(\AP[S]\) è al più il numero di funzioni da \(\AP[S]\) a \(\set{-1,0,1}\), che è \(3^{\abs{\AP[S]}}\).
\end{proof}

A questo punto, avendo dimostrato la correttezza e la terminazione dell'algoritmo di risoluzione, rimane da dimostrarne solo la completezza. Essendo che l'algoritmo per com'è definito genera ogni formula che il calcolo di risoluzione è in grado di generare, è sufficiente dimostrare la completezza del calcolo di risoluzione, ovvero che se \(S\) è insoddisfacibile allora \(S\vdash_R \square\).

\begin{theorem}{Completezza debole del calcolo di risoluzione}
  Se \(S\) è insoddisfacibile, allora \(S\vdash_R \square\).
\end{theorem}

\begin{proof}
  Per fare ciò, ragioniamo induttivamente sulla cardinalità di \(\AP[S]\).

  Il caso base è rappresentato da \(\abs{\AP[S]} = 0\), ovvero \(\AP[S] = \emptyset\). In questo caso è chiaro che l'unica formula che è possibile formare è la clausola vuota \(\square\), che è insoddisfacibile, e quindi \(S\vdash_R \square\).\footnote{È importante notare che \(\AP[S]=\emptyset\centernot\implies S = \emptyset\) poiché \(S\) è un insieme di insiemi di letterali. \(\AP[S]=\emptyset \implies  S = \emptyset \lor S = \set{\emptyset} = \set{\square}\). Avendo però che \(S\) è insoddisfacibile per ipotesi, siamo nella seconda situazione, e quindi \(S = \set{\square}\).}

  Per \(\abs{\AP[S]}>0\), si ha l'ipotesi induttiva che per ogni \(S'\) con \(\abs{\AP[S']} < \abs{\AP[S]}\), se \(S'\) è insoddisfacibile allora \(S'\vdash_R \square\).

  Per poter applicare tale ipotesi induttiva è quindi necessario in qualche modo rimuovere una proposizione atomica da \(S\), mostrare che il nuovo insieme di clausole è ancora insoddisfacibile, e da questa costruzione usando l'ipotesi induttiva mostrare che è possibile dedurre \(\square\) anche da \(S\).

  Iniziamo col fissare la proposizione atomica \(P\in\AP[S]\) che andremo a escludere. Costruiamo dunque gli insiemi \(S_P\) e \(S_{\neg P}\) come segue:
  \[S_P = \set{C\setminus P}[C\in S \text{ e } \neg P \in C]\]
  \[S_{\neg P} =  \set{C\setminus \neg P}[C\in S \text{ e }  P \in C]\]

  In altre parole, \(S_P\) è l'insieme di clausole che non contengono il letterale negativo di \(P\), a cui se presente è stato rimosso il letterale positivo di \(P\) e viceversa per \(S_{\neg P}\). In questo modo è chiaro che \(P\) non è presente in \(S_P\) e \(\neg P\) non è presente in \(S_{\neg P}\), e che inoltre \(\AP[S_P] \subset \AP[S]\) e \(\AP[S_{\neg P}]\subset \AP[S]\) poiché non sono state aggiunte nuove proposizioni atomiche, ma solo eventualmente rimosse.

  Di conseguenza \(\AP[S_P]<\AP[S]\) e \(\AP[S_{\neg P}]<\AP[S]\).
  Inoltre questi due insiemi non sono necessariamente disgiunti. In particolare infatti, ogni clausola in \(S\) che non conteneva ne \(P\) ne \(\neg P\) è presente sia in \(S_P\) che in \(S_{\neg P}\). In ogni caso, tutte le clausole dell'intersezione tra \(S_P\) e \(S_{\neg P}\) erano già presenti in \(S\), poiché se una clausola è stata ottenuta da \(S\) rimuovendo \(P\) o \(\neg P\), può essere ottenuta in tal modo solo in uno dei due insiemi, e quindi se è presente in entrambi, allora era già presente in \(S\) senza bisogno di rimuovere \(P\) o \(\neg P\).

  In aggiunta, tali insiemi sono anche certamente non vuoti. Supponiamo infatti per assurdo che \(S_P = \emptyset\). Allora, necessariamente, ogni clausola di \(S\) contiene \(\neg P\). Essendo però che per soddisfare una clausola è sufficiente soddisfare un suo letterale, allora ogni assegnamento che assegna \(0\) a \(P\) soddisfa tutte le clausole di \(S\), rendendo \(S\) soddisfacibile, in contraddizione con l'ipotesi che \(S\) sia insoddisfacibile. Vale ovviamente un ragionamento duale per \(S_{\neg P}\).

  Per poter applicare l'ipotesi induttiva a questi due insiemi ci rimane da mostrare che sono entrambi insoddisfacibili.
  Iniziamo ragionando sull'insoddisfacibilità di \(S_P\) supponendo per assurdo che \(S_P\) sia soddisfacibile. Allora, esiste un assegnamento \(\alpha: \AP[S_P] \to \set{0,1}\) tale che \(\Val{\alpha}[C] = 1\) per ogni \(C \in S_P\). Estendendo banalmente tale assegnamento a \(S\) con:
  \[\alpha^P(X) = \begin{cases}
    0 &, X=P\\
    \alpha(X) &, X\neq P
  \end{cases}\]
  possiamo facilmente vedere che \(\Val{\alpha^P}[S] = 1\) portando all'assurdo. Questo si ha sempre perché per soddisfare una clausola è sufficiente soddisfare un suo letterale, e ogni clausola di \(S\) o è presente in \(S_P\) e quindi soddisfatta da \(\alpha\), o contiene \(\neg P\) e quindi è soddisfatta da \(\alpha^P\) assegnando \(0\) a \(P\).
  Formalmente, \(\forall C \in S\), \(\neg P \in C \implies \Val{\alpha^P}[C] = 1\) perché \(\Val{\alpha^P}[\neg P] = 1\) e 
  \[\neg P \notin C \implies C\setminus\set{P} \in S_P \implies \exists l \in C\setminus\set{P} : \Val{\alpha^P}[l] = \Val{\alpha^P}[l] = 1\]

  Per quanto riguarda \(S_{\neg P}\) il ragionamento è analogo costruendo un assegnamento che assegni \(1\) a \(P\).
  A questo punto possiamo finalmente applicare l'ipotesi induttiva a \(S_P\) e \(S_{\neg P}\), avendo che entrambi sono insoddisfacibili, e quindi \(S_P\vdash_R \square\) e \(S_{\neg P}\vdash_R \square\).

  In particolare quindi esistono due alberi di derivazione del calcolo di risoluzione che hanno uno dei due insiemi come foglie e \(\square\) in radice. Chiamiamoli \(\pi_P\) e \(\pi_{\neg P}\) e valutiamo i diversi casi:
  \begin{itemize}
    \item Se vale \(S_P\subseteq S\) o \(S_{\neg P}\subseteq S\), ovvero una tra \(P\) e \(\neg P\) non è presente in nessuna clausola di \(S\), allora \(S\vdash_R \square\) banalmente per monotònicità.
    \item Alternativamente, esiste almeno una foglia in \(\pi_P\) etichettata con una clausola \(C_P\) e una in \(\pi_{\neg P}\) etichettata con una clausola \(C_{\neg P}\) tali che \(C_P\) e \(C_{\neg P}\) non sono in \(S\).
    In altre parole c'è almeno una foglia per albero deduzione che è stata ottenuta dalla sottrazione del letterale.
    Consideriamo quindi un'applicazione di risoluzione tra \(C_P\) e una qualsiasi clausola \(D \in S_P\) in clash con \(C_P\).
    Chiaramente \(P\notin C_P \land P \notin D \implies P \notin R\) risolvente di \(C_P\) e \(D\).
    A questo punto però aggiungendo \(P\) sia a \(C_P\), che ricordiamo per costruzione originariamente conteneva \(P\) in \(S\), e a \(R\), otterrei ancora un applicazione corretta di risoluzione, poiché applicata a un altro letterale in clash.
    \[\infer[\text{Risoluzione}]{R}{C_P \quad D}\implies\infer[\text{Risoluzione}]{R\cup\set{P}}{C_P\cup\set{P} \quad D}\]
    Con questa idea possiamo aggiungere \(P\) a tutte le foglie di \(\pi_P\) che non appartengono a \(S\), propagando \(P\) ai risolventi delle deduzioni in cui sono coinvolte.
    \begin{center}
      \begin{minipage}{0.42\textwidth}
        \centering
        \begin{prooftree}
          \tiny
          \AxiomC{\(C_P\)}
          \AxiomC{\(D\)}
          \BinaryInfC{\(R\)}
          \AxiomC{\(\vdots\)}
          \BinaryInfC{\(\vdots\)}
          \AxiomC{\(\vdots\)}
          \AxiomC{\(C_P'\)}
          \AxiomC{\(D'\)}
          \BinaryInfC{\(R'\)}
          \BinaryInfC{\(\vdots\)}
          \BinaryInfC{\(\square\)}
        \end{prooftree}
      \end{minipage}
      \hspace{0.3cm}
      \(\xrightarrow{P}\)
      \hspace{0.3cm}
      \begin{minipage}{0.42\textwidth}
        \centering
        \begin{prooftree}
          \tiny
          \AxiomC{\(C_P\cup\set{P}\)}
          \AxiomC{\(D\)}
          \BinaryInfC{\(R\cup\set{P}\)}
          \AxiomC{\(\vdots\)}
          \BinaryInfC{\(\vdots\)}
          \AxiomC{\(\vdots\)}
          \AxiomC{\(C_P'\cup\set{P}\)}
          \AxiomC{\(D'\)}
          \BinaryInfC{\(R'\cup\set{P}\)}
          \BinaryInfC{\(\vdots\)}
          \BinaryInfC{\(\set{P}\)}
        \end{prooftree}
      \end{minipage}
    \end{center}
    In questo modo, abbiamo costruito una nuova deduzione di risoluzione che, a partire da \(S\) essendo tute le foglie etichettate con clausole in \(S\), deduce \(\set{P}\).

    Applicando il ragionamento duale a \(S_{\neg P}\) otteniamo da \(S\) anche una deduzione per \(\set{\neg P}\).
        \begin{center}
      \begin{minipage}{0.42\textwidth}
        \centering
        \begin{prooftree}
          \tiny
          \AxiomC{\(C_{\neg P}\)}
          \AxiomC{\(D\)}
          \BinaryInfC{\(R\)}
          \AxiomC{\(\vdots\)}
          \BinaryInfC{\(\vdots\)}
          \AxiomC{\(\vdots\)}
          \AxiomC{\(C_{\neg P}'\)}
          \AxiomC{\(D'\)}
          \BinaryInfC{\(R'\)}
          \BinaryInfC{\(\vdots\)}
          \BinaryInfC{\(\square\)}
        \end{prooftree}
      \end{minipage}
      \hspace{0.2cm}
      \(\xrightarrow{\neg P}\)
      \hspace{0.2cm}
      \begin{minipage}{0.42\textwidth}
        \centering
        \begin{prooftree}
          \tiny
          \AxiomC{\(C_{\neg P}\cup\set{\neg P}\)}
          \AxiomC{\(D\)}
          \BinaryInfC{\(R\cup\set{\neg P}\)}
          \AxiomC{\(\vdots\)}
          \BinaryInfC{\(\vdots\)}
          \AxiomC{\(\vdots\)}
          \AxiomC{\(C_{\neg P}'\cup\set{\neg P}\)}
          \AxiomC{\(D'\)}
          \BinaryInfC{\(R'\cup\set{\neg P}\)}
          \BinaryInfC{\(\vdots\)}
          \BinaryInfC{\(\set{\neg P}\)}
        \end{prooftree}
      \end{minipage}
    \end{center}
    A questo punto, applicando un ultima volta la regola di risoluzione, è possibile ottenere una deduzione da \(S\) a \(\square\).

    \begin{prooftree}
      \tiny
      \AxiomC{\(C_P\cup\set{P}\)}
      \AxiomC{\(D\)}
      \BinaryInfC{\(R\cup\set{P}\)}
      \AxiomC{\(\vdots\)}
      \BinaryInfC{\(\vdots\)}
      \AxiomC{\(\vdots\)}
      \AxiomC{\(C_P'\cup\set{P}\)}
      \AxiomC{\(D'\)}
      \BinaryInfC{\(R'\cup\set{P}\)}
      \BinaryInfC{\(\vdots\)}
      \BinaryInfC{\(\set{P}\)}
    
      \AxiomC{\(C_{\neg P}\cup\set{\neg P}\)}
      \AxiomC{\(D\)}
      \BinaryInfC{\(R\cup\set{\neg P}\)}
      \AxiomC{\(\vdots\)}
      \BinaryInfC{\(\vdots\)}
      \AxiomC{\(\vdots\)}
      \AxiomC{\(C_{\neg P}'\cup\set{\neg P}\)}
      \AxiomC{\(D'\)}
      \BinaryInfC{\(R'\cup\set{\neg P}\)}
      \BinaryInfC{\(\vdots\)}
      \BinaryInfC{\(\set{\neg P}\)}
      \BinaryInfC{\(\square\)}
    \end{prooftree}
  \end{itemize}
\end{proof}

Per concludere, andiamo a mostrare un esempio della costruzione presente in questa dimostrazione.

\begin{example}{}
  Supponiamo di avere
  \[S = \set{\set{P,Q}, \set{P,\neg Q,\neg R}, \set{R}, \set{\neg P,\neg R}, \set{\neg P,Q,R} }\]
  Fissando \(P\) possiamo creare gli insiemi
  \[S_P=\set{\set{Q},\set{\neg Q, \neg R}, \set{R}}\]
  \[S_{\neg P} = \set{\set{R}, \set{\neg R}, \set{Q,R}}\]

  Ora è facile vedere che \(S_P \models_R\square\) e \(S_{\neg P} \models_R \square\):
  \begin{prooftree}
    \tiny
    \AxiomC{\(\set{Q}\)}
    \AxiomC{\(\set{\neg Q, \neg R}\)}
    \BinaryInfC{\(\set{\neg R}\)}
    \AxiomC{\(\set{R}\)}
    \BinaryInfC{\(\square\)}
  \end{prooftree}

  \begin{prooftree}
    \tiny
    \AxiomC{\(\set{R}\)}
    \AxiomC{\(\set{\neg R}\)}
    \BinaryInfC{\(\square\)}
  \end{prooftree}

  Per queste deduzioni abbiamo che \(\set{\neg R}\) in \(S_{\neg P}\) non apparteneva a \(S\), così come \(\set{\neg Q, \neg R}\).
  Aggiungiamo quindi alle deduzioni \(P\) e \(\neg P\), e propaghiamole.

  \begin{prooftree}
    \tiny
    \AxiomC{\(\set{Q}\)}
    \AxiomC{\(\set{\neg Q, \neg R,P}\)}
    \BinaryInfC{\(\set{\neg R,P}\)}
    \AxiomC{\(\set{R}\)}
    \BinaryInfC{\(\set{P}\)}
  \end{prooftree}

  \begin{prooftree}
  \tiny
    \AxiomC{\(\set{R}\)}
    \AxiomC{\(\set{\neg R,\neg P}\)}
    \BinaryInfC{\(\set{\neg P}\)}
  \end{prooftree}

  A questo punto possiamo dedurre \(\square\) facilmente col clash di \(P\) e \(\neg P\).

  \begin{prooftree}
    \tiny
    \AxiomC{\(\set{Q}\)}
    \AxiomC{\(\set{\neg Q, \neg R,P}\)}
    \BinaryInfC{\(\set{\neg R,P}\)}
    \AxiomC{\(\set{R}\)}
    \BinaryInfC{\(\set{P}\)}
    \AxiomC{\(\set{R}\)}
    \AxiomC{\(\set{\neg R,\neg P}\)}
    \BinaryInfC{\(\set{\neg P}\)}
    \BinaryInfC{\(\square\)}
  \end{prooftree}
\end{example}

\chapter{Logica del prim'ordine}

La logica proposizionale si concentra essenzialmente sulle proprietà degli operatori logici ed è detta proposizionale perché i suoi \q{atomi} (gli inglesi direbbero \emph{building blocks}) sono le proposizioni. Questo si riflette nella traduzione dal linguaggio naturale, in cui la separazione tra proposizioni avveniva tramite congiunzioni, disgiunzioni, negazioni, implicazioni, ecc.

Tale granularità però non è sufficiente per esprimere relazioni tra elementi del discorso interni alle proposizioni. Difatti, consideriamo nuovamente l'esempio del sillogismo aristotelico \q{Ogni uomo è mortale, Socrate è un uomo, dunque Socrate è mortale}. Una traduzione in logica proposizionale potrebbe essere \((p\land q)\to r\) poiché le tre frasi internamente non hanno a loro volta una struttura proposizionale. Chiaramente, la validità di questa espressione non può dipendere dalla struttura proposizionale presentata, poiché essa non è una tautologia. Ciò che invece collega le frasi e rende vero il sillogismo sono elementi interni alle proposizioni, come il fatto che \q{Socrate} è un \q{uomo} e che ogni \q{uomo} è \q{mortale}. Di conseguenza, per esprimere la validità del sillogismo è necessario analizzare gli atomi in modo più granulare, per scoprire che essi presentano a loro interno una struttura.
Ciò inizia a farci capire che gli atomi non sono veramente indivisibili come nella logica proposizionale. 

In questo caso, la relazione che ci viene nascosta è il fatto che \q{uomo} è una proprietà, vera o falsa, degli oggetti, così come \q{mortale}. Inoltre, \q{Socrate} è un oggetto specifico del nostro universo. Quello che quindi ci dicono queste frasi è che un oggetto del nostro dominio che è \q{uomo} è anche \q{mortale}, e che, in particolare, l'oggetto \q{Socrate} è un \q{uomo} e dunque se ne deduce che \q{Socrate} è \q{mortale}. Quindi, il salto concettuale da logica proposizionale a logica del prim'ordine è l'arricchimento del concetto intuitivo di situazione da un associazione tra frasi e valori di verità all' associazione tra gli oggetti del discorso e le proprietà di questi.

Senza entrare subito nel dettaglio della semantica, per esprimere tale granularità avremo bisogno dei nomi degli oggetti di una situazione, i quali potranno essere, come vedremo, finiti anche in un universo di oggetti infiniti. I nomi non sono, come nella logica proposizionale, oggetti sintattici che vengono interpretati in valori di verità, bensì oggetti sintattici che vengono interpretati in oggetti del dominio. Dunque, anche il concetto di interpretazione cambia, diventando più complesso.

Sarà inoltre necessario parlare delle proprietà di tali oggetti, o di relazioni tra essi nel caso siano coinvolti più oggetti nel discorso. Per proprietà e relazioni, il concetto che esse esprimono è molto vicino al concetto matematico di proprietà, appunto. Ad esempio, una proprietà dei numeri naturali è quella di essere pari. Ovviamente, tale proprietà, da sola, non ha molto senso, ma, se applicata ad uno specifico oggetto del dominio, come ad esempio \(5\), ovvero un numero naturale, allora è interpretabile con un valore di verità: vero se l'oggetto ha quella proprietà, falso altrimenti. Ovviamente, tale assegnamento dipende dall'interpretazione che diamo alla stringa \q{pari}, che, dal punto di vista logico, non è altro che un simbolo sintattico (il nome di una proprietà) a cui possiamo dare una semantica arbitraria. Senza un modello, \q{pari} è quindi solo un simbolo e non ha alcun significato.

Infine, nel linguaggio, capita di riferirsi ad un oggetto senza introdurre direttamente il suo nome, tramite una dipendenza con altri oggetti. Considerando ad esempio l'espressione \q{Il cane di Socrate aveva 3 gambe}, la locuzione \q{il cane di Socrate} identifica un oggetto del dominio come farebbe un nome, ma senza nominarlo esplicitamente, bensì identificandolo in funzione di un altro oggetto. Dunque possiamo rappresentare, sintatticamente, questo modo di rifersi ad un oggetto con una funzione \(f\), che prende in ingresso un oggetto del dominio e restituisce un altro oggetto del dominio. Per quanto questo possa sembrare molto simile ad una proprietà, è importante notare che una proprietà è interpretata in un valore di verità, mentre una funzione in un oggetto del dominio. 

La logica del prim'ordine permette dunque di parlare di oggetti di un dominio direttamente (nomi, detti costanti) o indirettamente a partire da altri oggetti (mediante funzioni), e di parlare di proprietà di tali oggetti o relazioni tra essi (predicati). In questo modo, è possibile creare un insieme infinito di nomi a partire da un insieme finito di costanti e funzioni.\footnote{Per avere un insieme infinito di termini sono sufficienti una costante e una funzione unaria. Interpretando la costante con il numero \(0\) e la funzione come la funzione successore, è possibile infatti rappresentare tutto \(\N\).} L'insieme di questi oggetti è detto insieme dei \keyword{termini}. Tali termini possono essere usati come argomenti di proprietà, relazioni o altre funzioni creando nuovi termini.

\section{Sintassi della Logica del Prim'Ordine}

Come già accennato nell'introduzione del capitolo, l'obiettivo è arricchire la logica proposizionale per ottenere una granularità che ci permette di parlare, oltre che di fatti a sé stanti, anche di oggetti, direttamente o indirettamente, e delle loro proprietà e relazioni. Per fare ciò dobbiamo iniziare dalla sintassi, la quale avrà bisogno in primo luogo di nuovi simboli \textit{non-logici} che colgano le differenze tra i vari tipi di componenti del discorso.

\begin{definition}{Signature di un linguaggio del prim'ordine}
  Sia \(L\) un linguaggio del prim'ordine. La sua \keyword{signature} (segnatura) è una tripla \(\langle \ConstSyms, \PredSyms, \FuncSyms \rangle\) con 
  \begin{itemize}
    \item \(\ConstSyms\) insieme dei \keyword{simboli di costante}
    \item \(\PredSyms\) insieme dei \keyword{simboli di predicato}
    \item \(\FuncSyms\) insieme dei \keyword{simboli di funzione}
  \end{itemize}
\end{definition}

La signature del linguaggio funge quindi da alfabeto di base per la costruzione di stringhe dedite alla rappresentazione di oggetti, proprietà e relazioni. È importante notare che a questo livello si sta parlando ancora di semplici simboli sintattici, e non di oggetti, proprietà o relazioni veri e propri. Un simbolo \(f \in \FuncSyms\) non è una funzione, bensì un simbolo che può essere usato per costruire stringhe che rappresentano funzioni.

Per definire propriamente tali simboli, in particolare quelli di funzione e di predicato, abbiamo bisogno di descrivere formalmente quanti oggetti del discorso coinvolgono, ovvero, da un punto di vista informatico, di quanti argomenti prendono in ingresso. Tale quantità è detta \keyword{arietà} del simbolo, e viene formalmente definita come segue.

\begin{definition}{Arietà di funzioni e predicati}
  Sia \(L\) un linguaggio del prim'ordine con signature \(\langle \ConstSyms, \PredSyms, \FuncSyms \rangle\). 
  
  Definiamo l'\keyword{arietà} come una funzione \(\arty: \FuncSyms\cup\PredSyms\to\N\). 
  
  L'arietà è un numero naturale che indica il numero di argomenti che la funzione o il predicato richiede per essere applicato.
\end{definition}

In alcune formulazioni della logica del prim'ordine i simboli di costante sono trattati congiuntamente a quelli di funzione, in quanto possono essere visti come funzioni di arietà zero. In questo modo, si ha un'unica categoria di simboli che rappresentano nomi di oggetti, siano essi costanti o funzioni. Nel nostro caso, invece, abbiamo deciso di trattarli separatamente, poiché è più intuitivo pensare a simboli che rappresentano oggetti specifici del dominio (costanti) e simboli che rappresentano modi di costruire nuovi oggetti a partire da altri (funzioni).

Con i simboli a disposizione, è possibile costruire stringhe che rappresentano oggetti del dominio, proprietà e relazioni tra essi. Tuttavia, finora possiamo descrivere soltanto oggetti specifici del dominio, quando invece non è così raro riferirsi ad oggetti generici di un dominio o ad un gruppo di essi. Non saremmo infatti in grado di esprimere frasi come \q{Esiste un uomo che è mortale} poiché le costanti permettono solo di parlare di uomini specifici, e le funzioni permettono solo di costruire nuovi oggetti a partire da altri.

Di conseguenza, per concludere l'insieme dei simboli non logici, è necessario introdurre un ultimo insieme di simboli, ovvero quello delle variabili. Le variabili sono simboli che non rappresentano né oggetti specifici del dominio, né modi di costruire nuovi oggetti, ma rappresentano oggetti generici del dominio, di cui si dovrà specificare l'identità successivamente. Intuitivamente, rappresentano quelle parti del discorso in cui non è specificato a quale oggetto ci si riferisce, ma si parla in generale di un oggetto, o di un insieme di oggetti.

\begin{definition}{Variabili}
  Sia \(L\) un linguaggio del prim'ordine con signature \(\langle \ConstSyms, \PredSyms, \FuncSyms \rangle\). 
  
  Definiamo l'insieme delle \keyword{variabili} come un insieme infinito e numerabile di simboli, disgiunto dagli altri simboli della signature, indicato con \(\VarSyms\).
\end{definition}

A questo punto possiamo finalmente introdurre le due categorie sintattiche di cui la logica del prim'ordine si compone, ovvero i \keyword{termini} e le \keyword{formule ben formate}. I termini sono stringhe che rappresentano oggetti del dominio, costruiti a partire da costanti, variabili e funzioni, mentre le formule ben formate sono stringhe che rappresentano fatti, e che quindi saranno interpretate in valori di verità.

\begin{definition}{Termini di un linguaggio del prim'ordine}
  Sia \(L\) un linguaggio del prim'ordine con signature \(\Sigma = \langle \ConstSyms, \PredSyms, \FuncSyms \rangle\). Definiamo l'insieme \(\Terms\) dei termini su \(\Sigma\) e \(\VarSyms\) insieme delle variabili tale che
  \begin{itemize}
    \item \(\VarSyms\subseteq \Terms\)
    \item \(\ConstSyms\subseteq \Terms\)
    \item Se \(f\in\FuncSyms\), \(\arty(f)=n\) e \(t_1,\dots,t_n\in \Terms\), allora \(f(t_1,\dots,t_n)\in \Terms\)
  \end{itemize}

  I termini che contengono almeno una variabile sono detti \keyword{termini aperti}, perché non è possibile associarli ad un oggetto concreto del dominio, se non con un assegnamento\footnote{Nel contesto della logica del prim'ordine, un assegnamento è un qualcosa che assegna a ciascuna variabile un oggetto del dominio.}, di cui parleremo nella prossima sezione.

  Dato un termine \(t\), indicheremo con \(\Var(t)\) l'insieme delle variabili che compaiono in \(t\). Se \(\Var(t)=\emptyset\), allora \(t\) è detto \keyword{termine chiuso} o \keyword{ground}, perché non dipende da un'assegnamento, ma è associato direttamente a un oggetto del dominio tramite l'interpretazione dei simboli di costante e di funzione che compaiono in \(t\).
\end{definition}

A questo punto, con i termini a disposizione, è possibile costruire le formule ben formate, che rappresentano fatti. Le formule sono costruite a partire da predicati e termini, che rappresentano le formule atomiche, che possono essere combinate tra loro tramite i simboli logici.

\begin{definition}{Formule atomiche di un linguaggio del prim'ordine}
  Sia \(L\) un linguaggio del prim'ordine con signature \(\Sigma = \langle \ConstSyms, \PredSyms, \FuncSyms \rangle\).
  
  Sia quindi \(\VarSyms\) l'insieme delle variabili e \(\Terms\) l'insieme dei termini su \(\Sigma\) e \(\VarSyms\).

  Definiamo l'insieme di \keyword{formule atomiche} su \(\Sigma\) e \(\VarSyms\) come l'insieme\footnote{Quando la signature e l'insieme delle variabili sono chiari dal contesto o non rilevanti, indicheremo semplicemente con \(\AF\) l'insieme delle formule atomiche.}
  \[\AF[\Sigma,\VarSyms] = \set{P(t_1,\dots,t_n)}[ P\in\PredSyms\land \arty(P)=n\land t_1,\dots,t_n\in \Terms]\]

\end{definition}

Per quanto riguarda i simboli logici, rispetto alla logica proposizionale, l'unica differenza è l'introduzione dei quantificatori, connessi al concetto di variabile. I quantificatori sono infatti simboli logici che permettono di caratterizzare le variabili riferendosi a tutti gli oggetti del dominio (quantificatore universale \(\forall\)) o ad almeno un oggetto del dominio (quantificatore esistenziale \(\exists\)) e permettendo dunque di dare un significato alle variabili senza dover passare per un'interpretazione che assegna un oggetto del dominio a ciascuna di esse.

Possiamo dunque definire le formule ben formate come segue.
\begin{definition}{Formule ben formate di un linguaggio del prim'ordine}
  Sia \(L\) un linguaggio del prim'ordine con signature \(\Sigma = \langle \ConstSyms, \PredSyms, \FuncSyms \rangle\). Diremo che \(\phi\) è una formula ben formata, ovvero \(\phi \in FOL\), dove \(\phi \in \FOL[\Sigma,\VarSyms]\)\footnote{Quando la signature e l'insieme delle variabili sono chiari dal contesto o non rilevanti, indicheremo semplicemente con \(\FOL\) l'insieme delle formule ben formate.} è l'insieme delle formule ben formate su \(\Sigma\) e \(\VarSyms\), se vale una delle seguenti condizioni:
  \begin{itemize}
    \item \(\phi\in\AF[\Sigma,\VarSyms]\);
    \item \(\phi = \neg\psi\) con \(\psi\in\FOL[\Sigma,\VarSyms]\);
    \item \(\phi = \psi\circ\theta\) con \(\psi,\theta\in\FOL[\Sigma,\VarSyms]\) e \(\circ\in\{\land,\lor,\to\}\);
    \item \(\phi = Qx.\psi\) con \(Q\in\{\forall,\exists\}\), \(x\in\VarSyms\) e \(\psi\in\FOL[\Sigma,\VarSyms]\).
  \end{itemize}

  Data \(\phi\in\FOL\), denoteremo con \(\Var(\phi)\) l'insieme delle variabili che compaiono in \(\phi\). Se \(\Var(\phi)=\emptyset\), allora \(\phi\) è una formula \keyword{ground}.
\end{definition}

A differenza della logica proposizionale, in cui era possibile esprimere le formule con un unico albero sintattico, in questo caso sarà necessario distinguere tra alberi sintattici dei termini e delle formule, nonostante si comportino in modo molto simile.

L'albero sintattico di un termine ha le foglie etichettate con costanti o variabili, e i nodi interni etichettati con simboli di funzione. L'albero sintattico di una formula, invece, ha le foglie etichettate con formule atomiche e i nodi interni con operatori logici o quantificatori.

Possiamo inoltre notare come, poiché le formule ground non contengono variabili, esse non aggiungano molto rispetto alla logica proposizionale, poiché, sostituendo ogni formula atomica con un simbolo proposizionale, è possibile ottenere una formula proposizionale equivalente. La vera potenza della logica del prim'ordine è dunque data dalla possibilità di esprimere formule con variabili, che possono essere interpretate in oggetti del dominio.

Consideriamo però ora due formule con variabili \(\forall x, \forall y, P(x,y)\) e \(\forall y, P(x,y)\).
La prima formula esprime una relazione che mette in relazione tutti gli oggetti del dominio, ed è facile da interpretare.
Per la seconda invece l'interpretazione non è possibile immediatamente. La formula è infatti ambigua, perché l'unica cosa che è possibile dire su di essa è che ogni oggetto \(y\) è in relazione \(P\) con qualcosa individuato dalla variabile \(x\), che non si sa a priori cosa sia.

Intuitivamente, la differenza sta nel fatto che, nella prima formula, tutte le variabili sono legate da un quantificatore, ovvero si trovano nell'ambito (o \emph{scope}) di un quantificatore che parla di quelle variabili, dove l'ambito è la più piccola sottoformula che contiene il quantificatore, o, alternativamente, facendo riferimento all'albero sintattico, il sottoalbero radicato nel quantificatore. Nella seconda formula, invece, la variabile \(x\) è libera, ovvero non è nell'ambito di nessun quantificatore che parla di \(x\).

\begin{definition}{Variabili libere e legate di una formula}
  Sia \(\phi\) una formula. Definiamo l'insieme \(FV(\phi)\) come:
  \[FV(\phi) = \begin{cases}
    \bigcup_{i=1}^{n} \Var(t_i), & \phi = P(t_1,\ldots,t_n) \text{ e } \arty(P) = n\\
    FV(\psi), & \phi = \neg \psi\\
    FV(\psi)\cup FV(\theta), & \phi = \psi\circ\theta, \circ \in \set{\land,\lor,\to}\\
    FV(\psi)\setminus\set{x}, & \phi = Q x. \psi, Q\in\set{\forall,\exists}
  \end{cases}\]

  Definiamo invece \(BV(\phi)\) come l'insieme delle variabili legate di \(\phi\):
  \[BV(\phi) = \begin{cases}
    \emptyset, & \phi = P(t_1,\ldots,t_n)\\
    BV(\psi), & \phi = \neg \psi\\
    BV(\psi)\cup BV(\theta), & \phi = \psi\circ\theta, \circ \in \set{\land,\lor,\to}\\
    BV(\psi)\cup\set{x}, & \phi = Q x. \psi, Q\in\set{\forall,\exists}
  \end{cases}\]
\end{definition}

Intuitivamente, l'insieme delle variabili libere è costruibile bottom-up con una visita in post-order dell'albero sintattico, aggiungendo tutte le variabili che si incontrano e rimuovendo quelle che vengono legate da un quantificatore.
Analogamente, l'insieme delle variabili legate è costruibile con una visita in pre-order dell'albero sintattico, aggiungendo tutte le variabili che vengono legate da un quantificatore.
Possiamo notare, inoltre, che, in generale, \(FV(\phi)\cap BV(\phi)\neq\emptyset\). Questo perché, riprendendo l'analogia della vista dell'albero sintattico, una stessa variabile può apparire libera in un ramo dell'albero e legata in un altro ramo. Infatti sarebbe più corretto parlare di occorrenze libere ed occorrenze legate di variabili piuttosto che di variabili libere e legate, poiché la singola occorrenza di una variabile non può essere sia libera che legata, ma, poiché una variabile può avere più occorrenze, potrebbero esservi sia occorrenze libere che occorrenze legate per la stessa variabile. Per esempio nella formula \(\forall x,(P(x,y)\land \forall y, Q(y,x))\) abbiamo che \(x\) è legata, mentre \(y\) è sia libera che legata.

\begin{definition}{Formule chiuse}
  Data una formula \(\phi\), se \(FV(\phi)=\emptyset\), allora \(\phi\) è una formula \keyword{chiusa} (\keyword{sentence}).
  Se \(FV(\phi)\neq\emptyset\), allora \(\phi\) è una formula \keyword{aperta}.
\end{definition}

Intuitivamente, per dare semantica alle formule chiuse sarà sufficiente un dominio del discorso e un'interpretazione che associ simboli sintattici a oggetti del dominio, proprietà e relazioni tra essi. Per le formule aperte, invece, sarà necessario anche un'assegnamento che associ le variabili a oggetti del dominio, in modo da poter interpretare le variabili libere. Prima di passare alla trattazione formale della semantica, però, vediamo come tale linguaggio ci permetta di esprimere concetti che non sono esprimibili in logica proposizionale, ignorando temporaneamente il concetto di interpretazione. Consideriamo come nostro dominio la teoria elementare dei numeri naturali, \((\N,+,\cdot,<,=)\).
Il linguaggio per parlare di \(\N\) ha tipicamente la signature composta da \(\ConstSyms = \set{\overline{0}}, \FuncSyms=\set{\overline{\cdot}, \overline{+}, \overline{s}}\) e \(\PredSyms = \set{\overline{<}, \overline{=}}\). I simboli sono stati sopralineati\footnote{Nonostante le numerose ricerche in merito risulta che nessun parlante italiano si sia mai trovato nella difficile condizione di apporre una linea \textbf{sopra} una parola e non poter usare quindi il termine \textbf{sotto}\textit{lineato} per riferirsi a essa. Spero di non essere il primo a essersi trovato in quest'annosa questione e che un termine corretto per questa circostanza esista, ma finché non lo avrò trovato spero vada bene il neologismo. Infine, esiste: \url{https://www.treccani.it/vocabolario/sopralineare/}} per evidenziare la differenza tra simboli che hanno solo connotato sintattico, senza alcuna interpretazione, e l'interpretazione usuale di tali simboli. Sia le funzioni che i predicati hanno arietà \(2\), ad eccezione della funzione \(\overline{s}\) che ha arietà \(1\). In questo linguaggio abbiamo un'unica costante, ovvero \(0\), unico elemento del dominio di cui si parla direttamente. Possiamo riferirci a tutti gli altri numeri applicando la funzione \(\overline{s}\) a \(0\) per ottenere \(1\), a \(1\) per ottenere \(2\), e così via. Ad esempio l'espressione \qi{\(2+2=4\)} è esprimibile come:
\[\overline{+}(\overline{s}(\overline{s}(\overline{0})), \overline{s}(\overline{s}(\overline{0})))\overline{=} \overline{s}(\overline{s}(\overline{s}(\overline{s}(0))))\]


Tale formula è chiusa, ed in particolare è ground.
Ovviamente, per ora non abbiamo modo di garantire che il simbolo \(\overline{s}\) si comporti effettivamente come la funzione successore. Infatti, ci servirà una sua interpretazione semantica o un modo per costruire propriamente una teoria, ovvero sia specificarne le proprietà tramite altre formule.

Se volessimo esprimere l'espressione \qi{Ogni numero è maggiore o uguale a \(0\)}, potremmo usare la formula \(\forall x, \overline{0} \overline{<} x \lor \overline{0}\overline{=}x\). Ovviamente, tale concetto sarebbe esprimibile equivalentemente con \(\forall x, \neg(x\overline{<}\overline{0})\). Tali formule, contenendo variabili, non sono ground, ma sono chiuse.

Possiamo pensare di caratterizzare la funzione successore esprimendo l'espressione \qi{Ogni numero diverso da \(0\) è il successore di qualche numero} con la formula \(\forall x, (\neg(\overline{0}\overline{=}x)\to\exists y\st \overline{s}(y)\overline{=}x)\). Se il simbolo \(\forall\) cattura intuitivamente il concetto di \q{ogni}, è meno chiaro da dove arrivi il simbolo \(\exists\). In questo caso, tale simbolo esprime il concetto di \qi{qualche}. 

Andiamo ora a caratterizzare il simbolo \(\overline{<}\). Potremmo esprimere le caratteristiche di una relazione d'ordine (stretto) totale: antiriflessività, transitività, totalità. L'antiriflessività è esprimibile con \qi{Nessun numero è minore di se stesso}, ovvero con la formula \(\neg \exists x\st x \overline{<} x\). La transitività con \qi{Se \(x\) è minore di \(y\) e \(y\) è minore di \(z\), allora \(x\) è minore di \(z\)}, ovvero con la formula \(\forall x, \forall y, \forall z, \left((x \overline{<} y \land y \overline{<} z)\to x \overline{<} z\right)\), e la totalità con \qi{Per ogni coppia di numeri, o il primo è minore del secondo, o sono uguali, o il secondo è minore del primo}, ovvero con la formula \(\forall x, \forall y, ((x \overline{<} y)\lor(x\overline{=}y)\lor(y\overline{<} x))\).

Il potere espressivo della logica del prim'ordine è quindi molto maggiore rispetto alla logica proposizionale, ma ciò porta anche delle complicazioni, come ad esempio l'indecidibilità dei problemi di soddisfacibilità e validità, che vedremo più avanti.

\section{Semantica della Logica del Prim'Ordine}

Come anticipato, per caratterizzare la semantica della logica del prim'ordine sarà necessario introdurre in primo luogo il dominio del discorso e un'interpretazione che associ simboli sintattici a oggetti del dominio, proprietà e relazioni tra essi. Per fare ciò, introduciamo il concetto di struttura logica, che intuitivamente svolge un po' il ruolo del modello nella logica proposizionale, ma con una complessità maggiore, poiché la caratterizzazione semantica dei simboli sintattici da parte del modello non è più implicita nella definizione di soddisfacibilità, ma è esplicitamente definita tramite l'interpretazione.\footnote{Nella logica proposizionale, un modello si limita ad assegnare valori di verità alle proposizioni, trattando queste ultime come proposizioni già semanticamente valide. Invece, nella logica del prim'ordine, è necessario definire una struttura logica, ovvero una caratterizzazione semantica degli elementi delle formule (e.g., il dominio degli elementi, la definizione delle relazioni), prima di poter assegnare un valore di verità.} 

\begin{definition}{Struttura \(\FOL\)}
  Una \keyword{struttura} semantica per un linguaggio del prim'ordine è una coppia \(M = \langle D, \Interp \rangle\) con \(D\) un insieme non vuoto detto \keyword{dominio}, e \(\Interp\) una \keyword{funzione di interpretazione} tale che:
  \begin{itemize}
    \item Per ogni \(c\in\ConstSyms\), \(\Interp(c)\in D\)
    \item Per ogni \(f\in\FuncSyms\) con \(\arty(f)=n\), \(\Interp(f) \in D^{(D^n)}\),  ovvero \(\Interp(f): D^n \to D\)
    \item Per ogni \(P\in\PredSyms\) con \(\arty(P)=n\), \(\Interp(P)\subseteq D^n\) relazione \(n\)-aria su \(D\). Le \(n-\)uple appartenenti a \(\Interp(P)\) soddisfano \(P\) in \(M\).
  \end{itemize}
\end{definition}

Come anticipato, tale struttura è sufficiente a descrivere il significato e valutare ogni formula ground del linguaggio, ma non è sufficiente per valutare formule aperte, poiché in esse compaiono variabili libere che non sono interpretate direttamente dalla struttura, ma dipendono da un'interpretazione che associa le variabili a oggetti del dominio. Per valutare dunque formule aperte, è necessario formalizzare anche il concetto di assegnamento.

\begin{definition}{Assegnamento di variabili}
  Sia \(M = \langle D, \Interp \rangle\) una struttura semantica per un linguaggio del prim'ordine. Un'assegnamento è una funzione \(\sigma: \VarSyms\to D\) che assegna a ogni variabile un oggetto del dominio.
\end{definition}

Adesso, con una struttura e un assegnamento, siamo in grado di valutare la verità di ogni formula del linguaggio, e quindi di definire il concetto di soddisfazione \(M,\sigma\models \phi\). Per poterlo fare, però, avendo definito la struttura sintattica delle formule a partire da quella dei termini, dovremo prima di tutto definire l'interpretazione dei termini per poter definire quella delle formule. Per semplicità di notazione, fissata una struttura \(M=\langle D,\Interp\rangle\), denoteremo \(\Interp(c)\) con \(c^M\), \(\Interp(f)\) con \(f^M\) e \(\Interp(P)\) con \(P^M\).

\begin{definition}{Interpretazione di un termine}
  Fissato un linguaggio del prim'ordine con signature \(\Sigma = \langle \ConstSyms, \PredSyms, \FuncSyms \rangle\) e un insieme di variabili \(\VarSyms\),
  sia \(t\in \Terms\), \(M = \langle D, \Interp \rangle\) una struttura semantica per un linguaggio del prim'ordine e \(\sigma: \VarSyms\to D\) un assegnamento. Definiamo l'interpretazione di \(t\) in \(M\) con assegnamento \(\sigma\), denotata con \(\den{t}[M][\sigma]\), come segue:
  \[\den{t}[M][\sigma] = \begin{cases}
    \sigma(x), & t = x\in\VarSyms\\
    c^M, & t = c\in\ConstSyms\\
    f^M(\den{t_1}[M][\sigma], \ldots, \den{t_n}[M][\sigma]), & t = f(t_1,\ldots,t_n), f\in\FuncSyms, \arty(f)=n
  \end{cases}\]
\end{definition}

A questo punto, possiamo finalmente definire il concetto di soddisfazione \(M,\sigma\models \phi\) per ogni formula \(\phi\) del linguaggio, e dunque il concetto di soddisfacibilità e validità.

\begin{definition}{Relazione di soddisfazione}
  Data una formula \(\phi\), una struttura \(M = \langle D, \Interp \rangle\) e un assegnamento \(\sigma\) definiamo la relazione di soddisfazione \(M,\sigma\models \phi\) con le seguenti regole:
  \begin{itemize}
    \item Se \(\phi = P(t_1,\ldots,t_n)\), allora \(M,\sigma\models \phi\) se e solo se \((\den{t_1}[M][\sigma], \ldots, \den{t_n}[M][\sigma]) \in P^M\)
    \item Se \(\phi = \neg \psi\), allora \(M,\sigma\models \phi\) se e solo se \(M,\sigma\not\models \psi\)
    \item Se \(\phi = \psi\land\theta\), allora \(M,\sigma\models \phi\) se e solo se \(M,\sigma\models \psi\) e \(M,\sigma\models \theta\)
    \item Se \(\phi = \psi\lor\theta\), allora \(M,\sigma\models \phi\) se e solo se \(M,\sigma\models \psi\) o \(M,\sigma\models \theta\)
    \item Se \(\phi = \forall x, \psi\), allora \(M,\sigma\models \phi\) se e solo se\footnote{Con \(\sigma[x\mapsfrom d]\) si intende l'assegnamento che associa a \(x\) l'oggetto \(d\) e a ogni altra variabile \(x'\) associa \(\sigma(x')\). Tale operazione può essere vista come una \q{sovrascrittura} dell'assegnamento originale in un punto ed è un modo per imporre uno specifico assegnamento ad una variabile senza alterare l'assegnamento delle altre.} \(M,\sigma[x \mapsfrom d]\models \psi\), per ogni \(d\in D\)
    \item Se \(\phi = \exists x\st \psi\), allora \(M,\sigma\models \phi\) se e solo se \(M,\sigma[x \mapsfrom d]\models \psi\), per qualche \(d\in D\)
  \end{itemize}
\end{definition}

Come per la logica proposizionale, per la relazione di soddisfazione, oltre che a una definizione basata su teoria degli insiemi, è possibile dare una definizione più algebrica, basata su assegnamenti e una funzione di valutazione che, data una formula e un assegnamento di verità, restituisce un valore di verità. 

\begin{definition}{Valutazione}
Sia \(M = \langle D, \Interp \rangle\) una struttura semantica per un linguaggio del prim'ordine e \(\sigma\) un assegnamento. Definiamo la funzione di valutazione \(\Val{M,\sigma}: \FOL \to \set{0,1}\) come segue:
\begin{itemize}
  \item \[\Val{M,\sigma}[P(t_1,\ldots,t_n)] = \begin{cases}
    1, & (\den{t_1}[M][\sigma], \ldots, \den{t_n}[M][\sigma]) \in P^M \\
    0, & \text{altrimenti}
  \end{cases}\]
  \item \(\Val{M,\sigma}[\neg \psi] = 1 - \Val{M,\sigma}[\psi]\)
  \item \(\Val{M,\sigma}[\psi\land\theta] = \min\{\Val{M,\sigma}[\psi], \Val{M,\sigma}[\theta]\}\)
  \item \(\Val{M,\sigma}[\psi\lor\theta] = \max\{\Val{M,\sigma}[\psi], \Val{M,\sigma}[\theta]\}\)
  \item \(\Val{M,\sigma}[\forall x, \psi] = \min \set{\Val{M,\sigma[x \mapsfrom d]}[\psi] }[ d\in D]\)
  \item \(\Val{M,\sigma}[\exists x\st \psi] = \max \set{\Val{M,\sigma[x \mapsfrom d]}[\psi] }[ d\in D]\)
\end{itemize}
\end{definition}

Questa definizione mette in evidenza la correlazione tra \(\land\) e \(\forall\) e tra \(\lor\) e \(\exists\), che è alla base della dualità tra i quantificatori universale ed esistenziale.

In altre parole, il quantificatore universale \(\forall\) può essere visto come una congiunzione potenzialmente infinita di formule, una per ogni elemento del dominio, mentre il quantificatore esistenziale \(\exists\) può essere visto come una disgiunzione potenzialmente infinita di formule, una per ogni elemento del dominio.

Dualmente, \textbf{su domini finiti}, la logica del prim'ordine si riduce alla logica proposizionale, in quanto ogni formula con quantificatori può essere riscritta come una formula senza quantificatori, sostituendoli con congiunzioni o disgiunzioni di formule ground.

Data la definizione di soddisfazione, è possibile definire i concetti di soddisfacibilità, verità e validità per il prim'ordine. Rispetto alla logica proposizionale, in questo caso la soddisfazione ha \q{due gradi di libertà}, in quanto dipende sia dalla struttura \(M\) che dall'assegnamento \(\sigma\). Di conseguenza, i concetti di soddisfacibilità, verità e validità devono essere definiti tenendo conto di entrambi questi aspetti.

\begin{definition}{Soddisfacibilità e verità in un modello}
  Data \(\phi \in \FOL\) e una struttura \(M = (D, \Interp) \) diremo che:

  \begin{itemize}
    \item \(\phi\) è \keyword{soddisfacibile in \(M\)} se esiste un assegnamento \(\sigma\) tale che \(M,\sigma\models \phi\).
    \item \(\phi\) è \keyword{vera in \(M\)} se per ogni assegnamento \(\sigma\), \(M,\sigma\models \phi\). Tale relazione è denotata con \(M\models \phi\). In tali casi si dice che \(M\) è un \keyword{modello} di \(\phi\).
  \end{itemize}

  Queste due definizioni sono equivalenti se \(\phi\) è una formula chiusa, in quanto in questo caso non ci sono variabili libere, e quindi l'interpretazione di \(\phi\) non dipende dall'assegnamento. 

\end{definition}

\begin{definition}{Soddisfacibilità e validità}
  Data \(\phi \in \FOL\) diremo che:

  \begin{itemize}
    \item \(\phi\) è \keyword{soddisfacibile} se esiste una struttura \(M\) e un assegnazione \(\sigma\) tali che \(M,\sigma\models \phi\).
    \item \(\phi\) è \keyword{valida} se per ogni struttura \(M\) e ogni assegnazione \(\sigma\), \(M,\sigma\models \phi\).
  \end{itemize}
  
\end{definition}

Alternativamente, diremo che \(\phi \in \FOL\) è soddisfacibile se esiste una struttura \(M\) tale che \(\phi\) è soddisfacibile in \(M\), e diremo che \(\phi\) è valida se per ogni struttura \(M\), \(\phi\) è vera in \(M\). 

\begin{definition}{Conseguenza logica}
  Siano \(\Gamma\subseteq \FOL\) e \(\phi\in\FOL\).
  Diremo che \(\phi\) è una \keyword{conseguenza logica} di \(\Gamma\), e scriveremo \(\Gamma\models \phi\), se per ogni struttura \(M\) e ogni assegnazione \(\sigma\) che soddisfa tutte le formule in \(\Gamma\), \(M,\sigma\models \phi\).
\end{definition}

Come già vista per la logica proposizionale, per \(\Gamma = \emptyset\), la conseguenza logica è equivalente alla validità.

Alle volte tra i simboli logici si annovera anche il simbolo \(=\), con semantica fissata a prescindere dall'interpretazione del modello. La sua interpretazione è quella solita di \keyword{identità}, ovvero \(M,\sigma\models t_1 = t_2\) se e solo se \(\den{t_1}[M][\sigma] = \den{t_2}[M][\sigma]\). Ovviamente, il primo \(=\) non è l'identità tra termini, ma solo un simbolo sintattico la cui interpretazione è l'identità tra oggetti del dominio, interpretazioni di termini. Ad esempio, il termine \(3+4-2\) \textbf{non} è identico al termine \(5\), ma l'interpretazione usuale di \(\N\) ci porterebbe ad avere vera \(3+4-2 = 5\), poiché le interpretazioni di \(3+4-2\) e \(5\) sono identiche, ovvero \(5\).

Tale aggiunta è fatta poiché, per quanto sia possibile considerare \(=\) come simbolo non logico ed assiomatizzarlo\footnote{Per assiomatizzazione, si intende l'aggiunta di formule --- che un modello deve soddisfare --- che descrivono le proprietà di un qualche simbolo non logico. Ciò permette, in un certo senso, di \qi{forzare} l'interpretazione di un simbolo, in quanto, avendo delle formule fissate che descrivono le proprietà di tale simbolo, gli unici modelli che si potranno prendere in considerazione saranno quelli che soddisfano le proprietà assiomatizzate del simbolo in questione. Una forma di assiomatizzazione è stata già presentata quando abbiamo \q{forzato} il significato di \(\overline{<}\) nell'esempio sulla teoria dei numeri.}, dimodoché si forzi ogni modello ad interpretarlo come una relazione di equivalenza, non è altrettanto possibile assiomatizzarlo affinché ogni modello lo interpreti esattamente come l'identità --- la relazione di equivalenza più restrittiva, poiché richiede che un oggetto sia in relazione solo ed esclusivamente con sè stesso (vi è una classe d'equivalenza per ogni oggetto del dominio contenente solo l'oggetto stesso) --- in quanto non è possibile escludere che vi siano modelli in cui \(=\) è interpretato come una relazione di equivalenza meno restrittiva (come, ad esempio, il valore assoluto oppure il modulo da divisione euclidea). Questo è il motivo per cui, spesso, in contesti matematici, si considera \(=\) come un simbolo logico con interpretazione fissata all'identità, sebbene questa scelta non aggiunga potenza espressiva al linguaggio --- non permette di dimostrare più teoremi --- rispetto al caso in cui si scelga di assiomatizzare \(=\) come una semplice relazione di equivalenza. La scelta è dettata principalmente dalla comodità di indicare la relazione di identità sempre con il simbolo \(=\).

\section{Proprietà e Sostituzioni nel Prim'Ordine}

Passiamo dunque a dimostrare alcune proprietà della logica del prim'ordine, che ci porteranno a dimostrare risultati simili ad alcuni ottenuti per la logica proposizionale, come ad esempio la possibilità di limitarsi a considerare le sole variabili presenti nelle formule, o il concetto di sostituzione.


\begin{proposition}[label=prop:term_interpretation_independence]{}
  Sia \(t \in \Terms\) tale che \(\Var(t) = \set{x_1, \ldots, x_n}\) e siano \(\sigma_1\) e \(\sigma_2\) assegnamenti tali che \(\sigma_1(x_i) = \sigma_2(x_i)\) per ogni \(i\in\{1,\ldots,n\}\). Allora \(\den{t}[M][\sigma_1] = \den{t}[M][\sigma_2]\) per ogni struttura \(M\).
\end{proposition}

In termini intuitivi, questa proposizione ci sta dicendo che l'interpretazione di un termine dipende solo dall'interpretazione delle variabili che compaiono in quel termine, e non da altre variabili che potrebbero essere presenti in altri termini o formule. Di conseguenza, per valutare l'interpretazione di un termine, è sufficiente considerare solo le variabili che compaiono in quel termine, e non tutte le variabili del linguaggio.

\begin{proof}
  Possiamo dimostrare la proposizione per induzione sulla struttura sintattica di \(t\).
  \begin{itemize}
    \item Se \(t\in C\), \(\den{t}[M][\sigma_1] = C^M = \den{t}[M][\sigma_2]\) per ogni struttura \(M\), dalla definizione di interpretazione di un termine costante.
    \item Se \(t\in\VarSyms\), allora \(\den{t}[M][\sigma_1] = \sigma_1(t)= \sigma_2(t) = \den{t}[M][\sigma_2]\) per ogni struttura \(M\), dall'ipotesi che \(\sigma_1\) e \(\sigma_2\) assegnano lo stesso valore a tutte le variabili che compaiono in \(t\).
    \item Se \(t = f(t_1,\ldots,t_m)\), allora \(\den{t}[M][\sigma_1] = f^M(\den{t_1}[M][\sigma_1], \ldots, \den{t_m}[M][\sigma_1])\). Essendo \(t_i\) sottotermine di \(t\), e quindi \(\Var(t_i)\subseteq \Var(t)\), per ipotesi induttiva si ha che \(\den{t_i}[M][\sigma_1] = \den{t_i}[M][\sigma_2]\) per ogni \(i\in\{1,\ldots,m\}\), e quindi \(\den{t}[M][\sigma_1] = f^M(\den{t_1}[M][\sigma_2], \ldots, \den{t_m}[M][\sigma_2]) = \den{t}[M][\sigma_2]\) per ogni struttura \(M\).
  \end{itemize}
\end{proof}

Estendiamo dunque questo risultato alle formule, dimostrando che, fissata una struttura \(M\), la relazione di soddisfazione \(M,\sigma\models \phi\) dipende solo dall'assegnamento delle variabili libere che compaiono in \(\phi\), e non da altre variabili che potrebbero essere presenti in altri termini o formule.

\begin{proposition}[namedlabel={prop:satisfaction_independence}{Indipendenza della soddisfazione dalle variabili non libere}]{Indipendenza della soddisfazione dalle variabili non libere}
  Sia \(\phi\in\FOL\) con \(FV(\phi) = \set{x_1, \ldots, x_n}\) e \(\sigma_1,\sigma_2\) con \(\sigma_1(x_i) = \sigma_2(x_i)\) per ogni \(i\in\{1,\ldots,n\}\), allora:
  \[M,\sigma_1\models \phi \iff M,\sigma_2 \models \phi\]
  per ogni struttura \(M\).
\end{proposition}

\begin{proof}
  Fissata una generica struttura \(M\), si ha che, per tutti i casi induttivi che non coinvolgono quantificatori, la tesi segue dal fatto che la soddisfacibilità di \(\phi\) dipende solo dalla soddisfacibilità delle sue sottoformule, e quindi si può applicare l'ipotesi induttiva in maniera diretta.
  
  Ad esempio, per \(\phi = \psi_1 \land \psi_2\), si ha che \(M,\sigma_1\models \phi\) se e solo se \(M,\sigma_1\models \psi_1\) e \(M,\sigma_1\models \psi_2\), che per ipotesi induttiva equivale a \(M,\sigma_2\models \psi_1\) e \(M,\sigma_2\models \psi_2\), che equivale a \(M,\sigma_2\models \phi\). Gli altri casi, che considerano altri simboli logici al di fuori dei quantificatori, sono del tutto analoghi.

  Per \(\phi = \forall x, \psi\), si ha che \(M,\sigma_1\models \phi\) se e solo se \(M,\sigma_1[x\mapsfrom d]\models \psi\) per ogni \(d\in D\). Quello che cerchiamo di mostrare è che ciò è equivalente a \(M,\sigma_2\models \phi\), ovvero che \(M,\sigma_2[x\mapsfrom d]\models \psi\) per ogni \(d\in D\). Di conseguenza, ci basta mostrare che \(\sigma_1(x_i) = \sigma_2(x_i)\) per ogni \(i\in\{1,\ldots,n\}\) implica che \(\sigma_1[x\mapsfrom d](x_i) = \sigma_2[x\mapsfrom d](x_i)\) per ogni \(i\in\{1,\ldots,n\}\) e per ogni \(d\in D\).
  La variabile \(x\) non appartiene a \(FV(\phi)\) per definizione di variabili libere. A questo punto, se \(x \notin FV(\psi)\) la tesi segue banalmente poiché la soddisfazione di \(\phi\) si riduce alla soddisfazione di \(\psi\), per cui vale l'ipotesi induttiva.
  Sia invece \(x\in FV(\psi)\). Fissato un \(d\), si ha che \(\sigma_1[x\mapsfrom d](x) = d = \sigma_2[x\mapsfrom d](x)\), e per ogni \(i\in\{1,\ldots,n\}\), \(\sigma_1[x\mapsfrom d](x_i) = \sigma_2[x\mapsfrom d](x_i)\) per ipotesi induttiva, e quindi \(M,\sigma_1[x\mapsfrom d]\models \psi\) se e solo se \(M,\sigma_2[x\mapsfrom d]\models \psi\) per ogni \(d\in D\), e quindi \(M,\sigma_1\models \phi\) se e solo se \(M,\sigma_2\models \phi\).

  Per il caso \(\phi = \exists x\st \psi\) il ragionamento è del tutto analogo, considerando un qualche \(d\in D\) invece che ogni \(d\in D\).
\end{proof}


Come già accennato in precedenza, questa proprietà è in un certo senso l'analogo della proprietà degli assegnamenti proposizionali per cui due assegnamenti che assegnano lo stesso valore almeno alle variabili proposizionali che compaiono in una formula, assegnano lo stesso valore alla formula. In altre parole, la soddisfacibilità di una formula dipende solo dall'interpretazione delle variabili libere che compaiono in essa.

\begin{corollary}{}
  Sia \(\phi\in\FOL\) chiusa. Allora per ogni struttura \(M\) e per ogni coppia di assegnamenti \(\sigma, \sigma'\),  \(M,\sigma\models \phi \iff M,\sigma'\models \phi\). In altre parole, la soddisfacibilità in \(M\) equivale alla verità in \(M\) per le formule chiuse.
\end{corollary}

\begin{proof}
  Poiché \(\phi\) è chiusa, \(FV(\phi) = \emptyset\). La tesi segue dunque banalmente dalla proposizione precedente, poiché per ogni coppia di assegnamenti \(\sigma, \sigma'\) si ha che \(\sigma(x) = \sigma'(x)\) per ogni \(x \in FV(\phi)\).
\end{proof}

Passiamo dunque a definire il concetto di sostituzione, che ci permetterà di sostituire, all'interno di una formula, una variabile con un termine. Successivamente, dimostreremo che tale operazione è ben definita, ma che dovrà essere usata con attenzione. Iniziamo come sempre dai termini --- definiremo la sostituzione di una variabile con un termine, all'interno di un termine --- per poi estendere il concetto alle formule.

\begin{definition}{Sostituzione di un termine a una variabile in un termine}
  Sia \(x\in\VarSyms\) e \(t,t' \in \Terms\). Definiamo la sostituzione, in \(t\), di \(x\) con \(t'\), denotata con \(t[t'/x]\) o \(\subs{t}{x}{t'}\), come segue:
  \[\subs{t}{x}{t'} = \begin{cases}
    t', & t = x\\
    t, & t \in \ConstSyms \cup (\VarSyms\setminus\set{x})\footnote{Si intende: se \(t\) non contiene la variabile \(x\), allora la sostituzione lascia \(t\) inalterato.}\\
    f(\subs{t_1}{x}{t'},\ldots,\subs{t_n}{x}{t'}), & t = f(t_1,\ldots,t_n), f\in\FuncSyms, \arty(f)=n
  \end{cases}\]
\end{definition}

In altre parole, sostituire, in \(t\), \(x\) con \(t'\) significa: sostituire ricorsivamente tutte le occorrenze di \(x\) con \(t\), percorrendo l'albero sintattico del termine in post-order. 
Consideriamo, ad esempio, il termine \(t = f(g(x,y), h(z,x))\). Applicando la trasformazione \(\subs{t}{x}{f(0)}\), otteniamo \(f(g(f(0),y), h(z,f(0)))\).

Proprio come la sostituzione nella logica proposizionale, anche la sostituzione nella logica del prim'ordine deve essere vista come simultanea.
Pertanto, una sostituzione di più variabili non andrebbe definita come composizione di sostituzioni, bensì come \({(t)}^{x_1,\ldots,x_k}_{t_1',\ldots,t_k'}\), la cui formalizzazione è analoga al caso singolo, ma con un caso base per variabile. 
Ad esempio, sempre considerando il termine \(t = f(g(x,y), h(z,x))\), applicando la sostituzione multipla \({(t)}^{x,y}_{y,f(0)}\) --- ovvero sostituendo \(x\) con \(y\) e \(y\) con \(f(0)\), si ottiene il termine \(f(g(y,f(0)), h(z,y))\).

Definita l'applicazione di sostituzioni ai termini, possiamo facilmente definire l'estensione dell'applicazione di sostituzioni alle formule.

\begin{definition}{Sostituzione di un termine a una variabile in una formula}
  Sia \(x\in\VarSyms\), \(t\in\Terms\) e \(\phi\in\FOL\). Definiamo la sostituzione, in \(\phi\), di \(t\) con \(x\), denotata con \(\phi[t/x]\) o \(\subs{\phi}{x}{t}\), come segue:
  \[\subs{\phi}{x}{t} = \begin{cases}
    \phi, & x \notin FV(\phi)\\
    P({(t_1)}^x_{t}, \ldots, {(t_n)}^x_{t}), & \phi = P(t_1,\ldots,t_n), P\in\PredSyms, \arty(P)=n, t_1,\ldots,t_n \in \Terms\\
    \neg {\psi}^x_{t}, & \phi = \neg \psi\\
    {\psi}^x_{t}\circ{\theta}^x_{t}, & \phi = \psi\circ\theta, \circ\in\set{\land,\lor,\to}\\
    Q y. {\psi}^x_{t}, & \phi = Q y. \psi, Q\in\set{\forall,\exists}, (y\neq x)
  \end{cases}\]
\end{definition}

\textbf{La sostituzione è applicata solo alle occorrenze libere} di \(x\) in \(\phi\), e non a quelle legate da un quantificatore, poiché altrimenti si andrebbe a cambiare inevitabilmente la semantica della formula.

Consideriamo, ad esempio, la formula \(\phi = \exists y\st R(x,y) \land P(x,z)\). Applicando la sostituzione \(\subs{\phi}{x}{f(z)}\), otteniamo \(\exists y\st R(f(z),y) \land P(f(z),z)\).
Utilizzando invece la sostituzione \(\subs{\phi}{x}{y}\) otteniamo \(\exists y\st R(y,y) \land P(y,z)\).

Questa seconda sostituzione però, sebbene sia permessa dalla definizione data, è problematica, in quanto \(y\) è una variabile legata, e quindi la sostituzione ha trasformato una variabile libera in una variabile legata. Questo tipo di sostituzione è dibattuta nella comunità logica. Parte della letteratura, infatti, esclude la possibilità di tale sostituzione direttamente nella definizione di sostituzione. Nel nostro caso, invece, manterremo la possibilità di avere tali sostituzioni, ma dettaglieremo successivamente come gestirle correttamente.

Infine, applicando la sostituzione \(\subs{\phi}{y}{f(z)}\) avremmo \(\exists y\st R(x,y) \land P(x,z)\), poiché \(y\) non è una variabile libera in \(\phi\), e quindi la sostituzione non ha effetto su \(\phi\).

Passiamo dunque al descrivere alcune proprietà della sostituzione così definita.

\begin{lemma}{Sostituzione}
  Dati \(t, t' \in \Terms\), \(\sigma\) un assegnamento, e \(M\) una struttura semantica, allora
  \[\den{{\subs{t}{x}{t'}} }[M][\sigma] = \den{t}[M][\sigma[x \mapsfrom \den{t'}[M][\sigma]]]\]
  Dove la prima sostituzione è quella definita sopra, ovvero la sostituzione sintattica, mentre la seconda è semantica.
\end{lemma}

Tale lemma intuitivamente ci dice che applicare una sostituzione ai termini a livello sintattico equivale ad applicare la sostituzione a livello semantico, ovvero a modificare l'assegnamento in modo che assegni a \(x\) l'interpretazione di \(t'\).

\begin{proof}{}
  Possiamo dimostrare la tesi per induzione sulla struttura sintattica di \(t\).
  Il caso induttivo, una volta risolti i casi base, risulterà banale.
  Il primo caso base si verifica quando il termine \(t\) da sostituire è proprio \(x\), mentre il secondo si verifica quando \(t\) è una costante oppure una variabile diversa da \(x\).
  Si noti che entrambi tali casi prevedono \(t\in \VarSyms\cup\ConstSyms\), e quello che li differenzia è se \(t\) sia uguale a \(x\) o meno.

  Nel caso \(t\neq x\) avremo \(\den{\subs{t}{x}{t'} }[M][\sigma] = \den{t}[M][\sigma]\) per la definizione del caso base. Inoltre abbiamo già visto che l'unica informazione che ha impatto sull'interpretazione di \(t\) è l'interpretazione delle variabili che compaiono in \(t\), e quindi essendo \(x\notin \Var(t)\) abbiamo che \(\den{t}[M][\sigma] = \den{t}[M][\sigma[x\mapsfrom d]]\) per ogni \(d\in D\), e quindi in particolare per \(d = \den{t'}[M][\sigma]\), e quindi \(\den{t}[M][\sigma] = \den{t}[M][\sigma[x\mapsfrom \den{t'}[M][\sigma]]]\).

  Altrimenti, se \(t=x\) avremo \(\den{\subs{t}{x}{t'} }[M][\sigma] = \den{\subs{x}{x}{t'} }[M][\sigma] = \den{t'}[M][\sigma]\) per la definizione del caso base. D'altro canto, però, 
  \[\den{t}[M][\sigma[x\mapsfrom \den{t'}[M][\sigma]]] = \den{x}[M][\sigma[x\mapsfrom \den{t'}[M][\sigma]]] = \sigma[x\mapsfrom \den{t'}[M][\sigma]](x) = \den{t'}[M][\sigma]\]
  che abbiamo visto essere uguale a \(\den{\subs{t}{x}{t'} }[M][\sigma]\), e quindi anche in questo caso la tesi è verificata.

  Per il caso induttivo, si avrà che \(t = f(t_1,\ldots,t_n)\) con \(f\in\FuncSyms\), \(\arty(f)=n\) e \(t_1,\ldots,t_n\in \Terms\).
  Dunque, applicando la definizione di sostituzione nei termini al caso induttivo e successivamente la definizione di interpretazione per i simboli funzionali, si ha
  \[\den{\subs{t}{x}{t'} }[M][\sigma] = \den{f(\subs{t_1}{x}{t'},\ldots,\subs{t_n}{x}{t'})}[M][\sigma] = f^M(\den{\subs{t_1}{x}{t'} }[M][\sigma],\ldots,\den{\subs{t_n}{x}{t'} }[M][\sigma])\]
  Inoltre, si ha che 
  \[\den{t}[M][\sigma[x\mapsfrom \den{t'}[M][\sigma]]] = \den{f(t_1,\ldots,t_n)}[M][\sigma[x\mapsfrom \den{t'}[M][\sigma]]] = f^M(\den{t_1}[M][\sigma[x\mapsfrom \den{t'}[M][\sigma]]],\ldots,\den{t_n}[M][\sigma[x\mapsfrom \den{t'}[M][\sigma]]])\]
  Per ipotesi induttiva, si ha che \(\den{\subs{t_i}{x}{t'} }[M][\sigma] = \den{t_i}[M][\sigma[x\mapsfrom \den{t'}[M][\sigma]]]\) per ogni \(i\in\{1,\ldots,n\}\), allora si ha che 
  \[\den{t}[M][\sigma[x\mapsfrom \den{t'}[M][\sigma]]]  =  f^M(\den{t_1}[M][\sigma[x\mapsfrom \den{t'}[M][\sigma]]],\ldots,\den{t_n}[M][\sigma[x\mapsfrom \den{t'}[M][\sigma]]]) = f^M(\den{\subs{t_1}{x}{t'} }[M][\sigma],\ldots,\den{\subs{t_n}{x}{t'} }[M][\sigma]) = \den{\subs{t}{x}{t'} }[M][\sigma]\]
\end{proof}

Per passare alla sostituzione nelle formule, si potrebbe pensare di poter estendere banalmente il lemma al teorema seguente:

\begin{theorem}{Sostituzione Naïve}
  Dati \(\phi\in\FOL\), \(t\in \Terms\), \(\sigma\) un assegnamento, e \(M\) una struttura semantica e \(x\in\VarSyms\) allora
  \[M,\sigma \models \subs{\phi}{x}{t} \iff M,\sigma[x\mapsfrom \den{t}[M][\sigma]] \models \phi\]
\end{theorem}

Sfortunatamente, però, tale teorema non è vero in generale, e ciò è facilmente mostrabile con un semplice controesempio, usando come struttura \(\N\) e la sua interpretazione usuale.
Presa infatti \(\phi = \exists y \st x<y\), avremo una formula che è vera in \(\N\) per ogni assegnamento, poiché, per ogni \(d\in\N\), esiste un \(d'\in\N\) tale che \(d<d'\). Applicando la sostituzione \(\subs{\phi}{x}{y}\), otteniamo \(\exists y \st y<y\), che è invece falsa in \(\N\) per ogni assegnamento, poiché, nell'interpretazione standard di \(\N\), il simbolo \(<\) rappresenta una disuguaglianza stretta. Chiaramente, il problema di questo esempio è lo stesso riscontrato negli esempi di sostituzione in formule visti in precedenza, e si verifica quando una variabile libera viene sostituita con una variabile legata, cambiando la semantica della formula.

Per evitare questo problema, è necessario restringere le ipotesi del teorema, richiedendo che, in caso di sostituzione di una variabile libera, le variabili introdotte da \(t\) siano ancora libere.

\begin{definition}{Termine libero per una variabile}
  Sia \(t\in \Terms\), \(x\in\VarSyms\) e \(\phi \in \FOL\). Diremo che \(t\) è \keyword{libero per \(x\) in \(\phi\)} se nessuna variabile di \(t\) è legata in \(\phi\) da un quantificatore che ha \(x\) libera nel suo ambito.
\end{definition}

In altre parole diremo che \(t\) è libero per \(x\) in \(\phi\) se tutte le variabili presenti in \(t\) sono libere in \(\phi\) nei contesti in cui è libera \(x\).

\begin{example}{}
  Sia \(\phi = \exists y \, (x < y)\) e \(t_1 = y, t_2 = z, t_3 = x\).
  \begin{itemize}
    \item \q{\(t_1\) è libero per \(x\) in \(\phi\)?} No, perché tra le variabili di \(t_1\) vi è \(y\), che è legata da \(\exists\) in \(\phi\).
    \item \q{\(t_2\) è libero per \(x\) in \(\phi\)?} Sì, perché la sua unica variabile è \(z\), che non è legata ad alcun quantificatore in \(\phi\).
    \item \q{\(t_3\) è libero per \(y\) in \(\phi\)?} No, perché \(y\) non è libera in \(\phi\) e, da definizione di sostituzione, non sarebbe nemmeno possibile sostituirla.
  \end{itemize}

  In definitiva, chiedersi \q{\(t\) è libero per \(x\) in \(\phi\)?} è l'esatto contrario di chiedersi \q{Se sostituisco, in \(\phi\), la variabile \(x\) con il termine \(t\), sto introducendo qualche variabile che è legata nel contesto di \(x\)?}.
  Se la risposta all'ultima domanda è affermativa, allora \(t\) \emph{non} è libero per \(x\) in \(\phi\), ovvero non siamo \q{liberi} di sostituire \(x\) con \(t\) se vogliamo mantenere invariata la semantica di \(\phi\).
\end{example}

A questo punto, possiamo riformulare il teorema nella sua versione corretta.

\begin{theorem}[namedlabel={thm:fol_substitution}{Sostituzione}]{Sostituzione}
  Dati \(\phi\in\FOL\), \(t\in \Terms\), \(\sigma\) un assegnamento, e \(M\) una struttura semantica e \(x\in\VarSyms\) allora
  \[M,\sigma \models \subs{\phi}{x}{t} \iff M,\sigma[x\mapsfrom \den{t}[M][\sigma]] \models \phi\]
  se \(t\) è libero per \(x\) in \(\phi\).
\end{theorem}

\begin{proof}
  Come al solito procediamo per induzione sulla struttura sintattica di \(\phi\) partendo dal caso base.

  Sia \(\phi = P(t_1,\ldots,t_n)\) con \(P\in\PredSyms\), \(\arty(P)=n\) e \(t_1,\ldots,t_n\in \Terms\). In questo caso,
  \[M,\sigma\models \subs{\phi}{x}{t} \iff (\den{\subs{t_1}{x}{t}}[M][\sigma], \ldots, \den{\subs{t_n}{x}{t}}[M][\sigma]) \in P^M\]
  Dal lemma precedente, si ha che \(\den{\subs{t_i}{x}{t} }[M][\sigma] = \den{t_i}[M][\sigma[x\mapsfrom \den{t}[M][\sigma]]]\) per ogni \(i\in\{1,\ldots,n\}\), e quindi 
  \[M,\sigma\models \subs{\phi}{x}{t} \iff (\den{t_1}[M][\sigma[x\mapsfrom \den{t}[M][\sigma]]], \ldots, \den{t_n}[M][\sigma[x\mapsfrom \den{t}[M][\sigma]]]) \in P^M \iff M,\sigma[x\mapsfrom \den{t}[M][\sigma]]\models \phi\]

  Chiaramente, se la formula \(\phi\) non contiene quantificatori, allora ogni termine è libero per ogni variabile, e quindi la tesi segue direttamente dal lemma precedente ragionando induttivamente sulla struttura sintattica di \(\phi\).

  Consideriamo il caso \(\phi = \forall y, \psi\).
  Si ha, per ipotesi, che \(M,\sigma\models \subs{\forall y, \psi}{x}{t}\). Nel caso \(x=y\), si ha che \(x\not\in FV(\psi)\) e quindi tentativo di sostituire \(x\) sarà inefficace, per definizione di sostituzione di un termine in una formula. Quindi, si osserverà \(\subs{\forall y, \psi}{x}{t} = \forall y, \psi\), e quindi \(M,\sigma\models \subs{\forall y, \psi}{x}{t}\) se e solo se \(M,\sigma\models \forall y, \psi\) --- ed in particolare --- se e solo se \(M,\sigma[x\mapsfrom \den{t}[M][\sigma]]\models \forall y, \psi\) poiché la soddisfacibilità di una formula dipende solo dall'assegnamento sulle variabili libere.

  Veniamo dunque al caso \(x\neq y\). In questo caso, per definizione di sostituzione, si ha che
  \[M,\sigma\models \subs{\forall y, \psi}{x}{t} \iff M,\sigma\models \forall y, \subs{\psi}{x}{t}\]
  e quest'ultima espressione, per semantica di \(\forall\), è equivalente a \(M,\sigma[y\mapsfrom d]\models \subs{\psi}{x}{t}\) per ogni \(d\in D\). Essendo \(\psi\) sottoformula di \(\phi\) è possibile applicare l'ipotesi induttiva, e quindi, per ogni \(d\in D\), avere\footnote{La notazione \(\sigma[x\mapsfrom d][y\mapsfrom d']\) indica la concatenazione di due cambiamenti di assegnamento. Essendo L'operazione di cambiamento di assegnamento intuitivamente una \q{sovrascrittura} dell'assegnamento in un punto, se più cambiamenti di assegnamento sovrascrivono lo stesso punto, allora l'ultimo cambiamento di assegnamento è quello che ha effetto. Di conseguenza, il cambiamento di assegnamento non è commutativo.}
  \[M,\sigma[y\mapsfrom d]\models \subs{\psi}{x}{t} \iff M,\sigma[y\mapsfrom d]\left[x\mapsfrom \den{t}[M][\sigma[y\mapsfrom d]]\right]\models \psi\]

  Vorremmo arrivare a mostrare che tale oggetto è equivalente a \(M,\sigma[x\mapsfrom \den{t}[M][\sigma]]\models \forall y,\psi\) che per semantica equivale a \(M,\sigma[x\mapsfrom \den{t}[M][\sigma]][y\mapsfrom d]\models \psi\) per ogni \(d \in D\).
  
  Questa nuova scrittura è molto simile alla precedente, ma varia nell'ordine di applicazione dei cambiamenti di assegnamento, oltre che nell'assegnamento applicato nell'interpretazione di \(t\). Di conseguenza, non è banale mostrare che siano uguali. Tuttavia, se \(y\) non è in \(t\), allora l'ordine diventa irrilevante\footnote{Ciò accade poiché sostituendo prima \(x\) non verranno aggiunte istanze di \(y\) da sostituire, e sostituire prima \(y\) introduce solo termini ground, su cui non è possibile applicare alcuna sostituzione}, e risulta che \(\llbracket t\rrbracket^M_{\sigma[y\mapsfrom d]} = \llbracket t\rrbracket^M_{\sigma}\) grazie alla \Cref{prop:term_interpretation_independence}.

  Per ipotesi, sappiamo che \(t\) è libero per \(x\) in \(\phi\). Ma poiché \(y\) è legata in \(\phi\) da \(\forall\), allora \(t\) non può contenere \(y\) come variabile, altrimenti \(t\) conterrebbe una variabile legata in \(\phi\) in uno scope in cui \(x\) non lo è. Di conseguenza, come appena notato, l'ordine dei cambi di assegnamento diventa irrilevante e gli assegnamenti usati per l'interpretazione di \(t\) sono equivalenti, e quindi si ha la tesi.

  Il caso \(\phi = \exists y\st \psi\) è analogo al caso \(\phi = \forall y, \psi\) a meno della locuzione \q{per ogni \(d\in D\)} che diventa \q{per qualche \(d\in D\)}, avendo la tesi anche in questo caso.
\end{proof}


\subsection{Equivalenze notevoli}

Consideriamo adesso delle equivalenze notevoli tra formule di \(\FOL\), ovvero formule soddisfatte esattamente dalle stesse coppie \(M, \sigma\) di strutture e assegnamenti.
Alcune risultano abbastanza ovvie, come la dualità dei quantificatori.

\subsubsection{Dualità tra i quantificatori}
Dimostriamo che valgono le equivalenze \(\forall x,\phi \equiv \neg \exists x\st \neg \phi\) e \(\exists x\st \phi \equiv \neg\forall x, \neg \phi\).

Infatti, \(M,\sigma\models \forall x,\phi\) se e solo se per ogni \(d\in D\), \(M,\sigma[x\mapsfrom d]\models \phi\), il che è equivalente a dire che non esiste alcun \(d\in D\) tale che \(M,\sigma[x\mapsfrom d]\not\models \phi\), ovvero non esiste alcun \(d\in D\) tale che \(M,\sigma[x\mapsfrom d]\models \neg \phi\), il che significa precisamente \(M,\sigma\not\models \exists x\st \neg \phi\), da cui \(M,\sigma\models \neg \exists x\st \neg \phi\). L'equivalenza duale segue per ragionamento analogo.

\subsubsection{Distributività di \(\forall\) su \(\land\)}
\(\forall x, (\phi_1 \land \phi_2) \equiv \forall x, \phi_1 \land \forall x, \phi_2\), ovvero che il quantificatore universale distribuisce sulla congiunzione. Tale equivalenza è ovvia se si considera la semantica tramite funzione di valutazione, poiché l'interpretazione del quantificatore universale era un minimo, in questo caso di una congiunzione, che è anch'essa un minimo. Ci troviamo quindi davanti a due minimi di minimi, di cui chiaramente l'ordine non ha importanza. 

\subsubsection{Distributività di \(\exists\) su \(\lor\)}
Si ha un'equivalenza analoga per il quantificatore esistenziale con la disgiunzione, ovvero \(\exists x\st (\phi_1 \lor \phi_2) \equiv \exists x\st \phi_1 \lor \exists x\st \phi_2\).

\subsubsection{Implicazione per mancata distributività di \(\forall\) su \(\lor\)}
Chiaramente \(\forall x, (\phi_1 \lor \phi_2) \not\equiv \forall x, \phi_1 \lor \forall x, \phi_2\) perché la prima è più restringente della seconda, ma vale una relazione di implicazione, ovvero \(\forall x,\phi_1 \lor \forall x, \phi_2 \implies \forall x, (\phi_1 \lor \phi_2)\). Ciò si verifica in quanto, presi qualsiasi \(M\) e \(\sigma\), se \(M,\sigma\models\forall x,\phi_1 \lor \forall x, \phi_2\) allora per semantica \(M,\sigma\models \forall x,\phi_1\) o \(M,\sigma\models\forall x, \phi_2\). 

Senza perdita di generalità, supponiamo che \(M,\sigma\models\forall x,\phi_1\). Allora, per semantica, per ogni \(d\in D\), \(M,\sigma[x\mapsfrom d]\models \phi_1\), e quindi \(M,\sigma[x\mapsfrom d]\models \phi_1 \lor \phi_2\) per ogni \(d\in D\), e quindi \(M,\sigma\models \forall x, (\phi_1 \lor \phi_2)\). 

Per mostrare che l'altra implicazione non vale, dobbiamo passare per la semantica di \(\implies\), che in particolare è \(\phi \implies\psi\) se e solo se, per ogni coppia \(M, \sigma\), si ha che \(M,\sigma\models \phi \implies M,\sigma\models \psi\). 

Ovviamente la sua negazione sarà equivalente, per qualche coppia \(M, \sigma\), a \(M,\sigma\models \phi\) e \(M,\sigma\not\models \psi\). Di conseguenza, ci basta trovare il cosiddetto \keyword{contromodello}, ovvero un modello che soddisfa un controesempio alla tesi. 

In particolar modo, basterà trovare un \(M\) che contenga un \(d \in D\) che soddisfi \(\phi_1\) e non \(\phi_2\) e un \(d' \in D\) che soddisfi \(\phi_2\) e non \(\phi_1\). Chiaramente, \(M,\sigma\models \forall x, (\phi_1 \lor \phi_2)\) ma \(M,\sigma\not\models \forall x, \phi_1 \lor \forall x, \phi_2\).


\subsubsection{Implicazione per mancata distributività di \(\exists\) su \(\land\)}
È possibile fare un ragionamento duale per mostrare che 
\[\exists x\st( \phi_1 \land \phi_2) \implies \exists x\st \phi_1 \land \exists x\st \phi_2\]
e
\[\exists x\st (\phi_1 \land \phi_2 )\centernot\impliedby \exists x\st \phi_1 \land \exists x\st \phi_2\]

\subsubsection{Distribuire \(\forall\) e \(\exists\) su simboli logici non affini}
In alcuni casi, però, è possibile distribuire i quantificatori sugli operatori logici non a loro affini. Consideriamo, ad esempio,
\[\exists x\st(\phi_1 \land \phi_2) \equiv \exists x\st( \phi_1) \land \phi_2\]
con \(x \notin \Var(\phi_2)\)\footnote{In particolare, \(x \notin FV(\phi_2)\)}.

Ciò è vero, perché quantificare su una variabile non presente in una formula non influenza la verità di quest'ultima.

Per definizione semantica, si ha \(M,\sigma \models \exists x\st\phi_2 \iff M,\sigma[x\mapsfrom d] \models \phi_2 \) per qualche \(d\in D\). Se \(x \notin FV(\phi_2)\), allora \(\sigma(y)=\sigma[x\mapsfrom d](y)\) per ogni \(y \in FV(\phi_2)\), dunque dalla \Cref{prop:satisfaction_independence} si ha \(M,\sigma \models \phi_2\) se e solo se \(M,\sigma[x\mapsfrom d] \models \phi_2\). Di conseguenza, \(M,\sigma \models \exists x\st\phi_2\) se e solo se \(M,\sigma \models \phi_2\).

La stessa cosa accade col quantificatore universale e \(\lor\), ovvero
\[\forall x, (\phi_1 \lor \phi_2) \equiv \forall x, \phi_1 \lor \phi_2\]
con \(x \notin FV(\phi_2)\).

Inoltre si ha che, con \(x \notin FV(\phi_1)\),
\[\forall x, (\phi_1 \to \phi_2) \equiv \phi_1 \to \forall x, \phi_2\]
e analogamente
\[\exists x\st (\phi_1 \to \phi_2) \equiv \phi_1 \to \exists x\st \phi_2\]

Se invece \(x \notin FV(\phi_2)\), allora valgono
\[\forall x, (\phi_1 \to \phi_2) \equiv (\exists x\st \phi_1 )\to \phi_2\]
e
\[\exists x\st (\phi_1 \to \phi_2) \equiv (\forall x, \phi_1 )\to \phi_2\]

Questo perché 
\[\forall x, (\phi_1 \to \phi_2) \equiv \forall x, (\neg \phi_1 \lor \phi_2) \equiv \forall x, (\neg \phi_1 )\lor \phi_2 \equiv \neg (\forall x, \neg \phi_1) \to \phi_2 \equiv (\exists x\st \phi_1)\to \phi_2\]

Il ragionamento è ovviamente analogo per il caso esistenziale.

\subsection{Dimostrare validità e soddisfacibilità di una formula}
In generale, dimostrare la validità di una formula è indecidibile, e lo vedremo più avanti codificando il problema della fermata. Tuttavia, in alcuni casi, la regolarità dei modelli che possono soddisfare una tale formula consente di esplorarli in tempo finito.

Prendiamo, ad esempio, la formula

\[\phi \eqdef (\exists x\st A(x)\to \exists x\st B(x))\to \forall x, (A(x)\to B(x))\]

Essendo la validità una proprietà universale, il modo più facile per dimostrarla è verificare che non esistano contromodelli.

Avendo \(\phi\) l'implicazione come operatore principale, un contromodello \(M\) deve soddisfare \(\exists x\st A(x)\to \exists x\st B(x)\) e non soddisfare \(\forall x, (A(x)\to B(x))\).

Inoltre, per soddisfare \(\exists x\st A(x)\to \exists x\st B(x)\), dobbiamo far sì che valga \(M,\sigma\not\models \exists x\st A(x)\) oppure \(M,\sigma \models \exists x\st B(x)\). 

Ora, per non soddisfare \(\exists x\st A(x)\) la struttura deve avere \(A^M =\emptyset\) ma con dominio non vuoto. Ora, però, in tali strutture chiaramente nessun punto soddisfa \(A\), quindi è sicuramente soddisfatta \(\forall x, \neg A(x)\) e quindi è soddisfatta la conclusione \(\forall x, (A(x)\to B(x))\). 
Di conseguenza, se un contromodello esiste, non è grazie al primo caso.

Passando al secondo, l'unica cosa che sappiamo da \(M,\sigma\models \exists x\st B(x)\) è che \(M\) dev'essere tale per cui \(B^M \neq \emptyset\). Ora, possiamo trovare una struttura con queste proprietà che non soddisfa la conclusione \(\forall x, (A(x)\to B(x))\), in particolare con \(M\) che contiene un elemento \(d\) che soddisfa \(A\) e non \(B\). A questo punto, abbiamo trovato un contromodello, da cui si ha che la formula non è valida. Poiché il primo caso ci ha detto che vi sono delle strutture che la soddisfano, \(\phi\) è soddisfacibile.

Consideriamo adesso la formula 
\[\phi \eqdef (\exists x\st A(x)\to \forall x\st B(x))\to \forall x, (A(x)\to B(x))\]
e analizziamola in modo analogo.

Il primo caso è rimasto identico, ma nel secondo i modelli che soddisfano \(\forall x\st B(x)\) sono tali che \(B^M = D\). Ma questo chiaramente soddisfa \(\forall x, (\neg A(x)\lor B(x))\) perché \(\forall x, (\neg A(x)) \lor \forall x, B(x) \implies \forall x, (\neg A(x)\lor B(x))\). Quindi non è possibile trovare un contromodello, e, di conseguenza, la formula è valida. Ovviamente, è anche soddisfacibile, poiché ogni formula valida è anche soddisfacibile.

Consideriamo ora due formule, di cui la prima ha un albero sintattico che ha la radice etichettata con un simbolo logico unario.
Con \(Hat\) predicato vero per individui che indossano un cappello, consideriamo le due formule
\[\exists x, (Hat(x) \to \forall y, Hat(y))\]
\[(\exists x, Hat(x))\to(\forall y, Hat(y))\]

La prima formula potrebbe essere la traduzione di \qi{C'è un individuo tale che se egli ha il cappello allora tutti hanno il cappello}. L'altra invece potrebbe essere la traduzione di \qi{Se c'è un individuo che ha il cappello, allora tutti hanno il cappello}. In linguaggio naturale, possono essere percepite come con lo stesso significato, ma non è così. 

Verifichiamo la validità della prima formula, valutando se è possibile l'esistenza di un contromodello. Chiamiamo \(\phi_1 \eqdef Hat(x)\) e \(\phi_2 \eqdef \forall y, Hat(y)\). Notiamo da subito che, poiché \(x \notin FV(\phi_2)\), vale l'equivalenza \(\exists x, (\phi_1 \to \phi_2) \equiv \forall x, \phi_1 \to \phi_2\), ovvero la formula equivale a
\[\forall x, Hat(x) \to \forall y, Hat(y)\]
che è evidentemente valida. Infatti, se esistesse un contromodello \(M\), allora si dovrebbe avere, contemporaneamente, \(Hat^M = D\) e \(Hat^M \neq D\), il che è impossibile.

Adesso, verifichiamo se la seconda formula è valida.
Un contromodello \(M\) dovrebbe essere t.c. \(M,\sigma \models \phi_1\) e \(M,\sigma \not\models \phi_2\).
Per osservare \(M,\sigma \models \phi_1\), si deve avere \(Hat^M \neq \emptyset\), mentre, per osservare \(M,\sigma \not\models \phi_2\), si deve avere \(Hat^M \neq D\).
Pertanto, basta prendere, ad esempio, \(D=\set{a, b}\) e \(Hat^M=\set{a}\), per non soddisfare la formula. Pertanto, un contromodello esiste e quindi la formula non è valida.
Inoltre, la formula è soddisfacibile in quanto soddisfatta da ogni struttura con \(Hat^M=D\) --- la conseguenza dell'implicazione è verificata --- o \(Hat^M=\emptyset\) --- la premessa dell'implicazione è falsa, dunque l'implicazione è banalmente verificata.

Infine, prima di passare alla trattazione del calcolo dei predicati, analizziamo, con brevi esempi, come la logica del prim'ordine ci permetta, una volta fissata una segnatura, di poter formalizzare e tradurre il linguaggio naturale. Quest'analisi ci permetterà di evidenziare alcune caratteristiche e correlazioni tipiche tra simboli logici, che ci aiuteranno ad avere una maggiore dimestichezza con la logica del prim'ordine, e a comprendere meglio le sue potenzialità e limitazioni.

\section{Esempi di formalizzazione del linguaggio naturale}

Cominciamo con un primo semplice esempio di traduzione. Consideriamo la frase:

\begin{quote}
  \q{I cani che abbaiano non hanno un proprietario}
\end{quote}

In primo luogo, va stabilita la segnatura del linguaggio, ovvero i simboli utilizzabili. In questo saranno necessari un simbolo di predicato unario \(C\) interpretato vero per gli oggetti del dominio che sono cani, un simbolo di predicato unario \(A\) interpretato vero per gli oggetti che abbaiano e infine il concetto di proprietario, che può essere rappresentato sia come proprietà del singolo oggetto, ma in maniera più completa anche come una relazione binaria \(P\) che lega un cane al suo proprietario.
A questo punto possiamo facilmente esprimere la frase in logica del prim'ordine con la formula
\[\forall x, ((C(x) \land A(x)) \to \neg\exists y\st P(y,x))\]

È importante notare che in questa frase è presente un artificio tipico del linguaggio naturale, ovvero la presenza di una quantificazione universale implicita. Quando si parla di una proprietà di una categoria, come in questo caso i cani che abbaiano, si sottintende che ogni individuo della categoria soddisfa quella proprietà, e quindi si ha una quantificazione universale implicita, che tendenzialmente quantifica un'implicazione, la cui premessa denota la categoria, e la conseguenza la proprietà vorremmo essere vera. In questo caso, la frase \q{I cani che abbaiano non hanno un proprietario} è equivalente a \q{Tutti i cani che abbaiano non hanno un proprietario}. I questo caso inoltre, anche la quantificazione esistenziale è implicita, poiché introdotta dalla locuzione \q{un}, che comunica il riferimento a un individuo.

In generale, le frasi con la struttura \q{Gli \(A\) hanno la proprietà \(P\)} configurano una situazione in cui l'insieme degli elementi con la proprietà \(A\) è un sottoinsieme dell'insieme degli elementi con la proprietà \(P\). Di conseguenza, la traduzione in logica del prim'ordine di una frase con questa struttura è sempre della forma \(\forall x, (A(x) \to P(x))\).
Ovviamente non è detto che \(A\) e \(P\) debbano essere atomiche, ma possono essere anche formule più complesse come nell'esempio precedente, in cui \(A(x)\) è \(C(x)\land A(x)\) e \(P(x)\) è \(\not\exists y\st P(y,x)\).

Consideriamo ora la frase 
\begin{quote}
  \q{C'è un cane che ha due proprietari}
\end{quote}

La locuzione \q{C'è} comunica la presenza di una quantificazione esistenziale. In questo caso per indicare che ci sono due individui distinti è necessario utilizzare la negazione del simbolo di identità, ottenendo la formula
\[\exists x\st\left( C(x) \land \exists y \exists z\st P(x,y) \land P(x,z) \land \neg(y = z)\right)\]

Tecnicamente, la frase è ambigua e potrebbe essere interpretata anche con l'avere esattamente due proprietari, e quindi andrebbe specificata differentemente la formula.

Inoltre, possiamo notare che in questo caso la presenza di una quantificazione esistenziale porta all'utilizzo di una congiunzione invece che di un implicazione. Questo è dato dal fatto che la quantificazione esistenziale sceglie un elemento del dominio, e quindi è necessario che tale elemento soddisfi tutte le proprietà richieste, mentre la quantificazione universale non sceglie un elemento specifico ma li seleziona tutti, e quindi è necessario filtrare via quelli di cui non mi interessa parlare con una implicazione.

In particolare, considerando il predicato unario \(HP\) vero per un cane che ha padrone, una formula del tipo \(\forall x, (C(x) \land HP(x))\) è una proprietà fortissima, che è falsa in ogni dominio che non ha solo cani non randagi. D'altro canto \(\exists x\st C(x) \to HP(x)\) è una proprietà quasi nulla, poiché banalmente vera per ogni dominio che ha almeno un oggetto non cane.

In generale quindi, nella formalizzazione del linguaggio naturale in logica del prim'ordine, i quantificatori universali tendono a essere associati a implicazioni, mentre i quantificatori esistenziali tendono a essere associati a congiunzioni. Ovviamente, non è detto che questa sia una regola fissa.

Consideriamo ora la frase 
\begin{quote}
  \q{I cani che hanno un proprietario non abbaiano}.
\end{quote}

Questa frase è molto simile alla frase già vista \q{I cani che abbaiano non hanno un proprietario}, che abbiamo tradotto come
\[\forall x, ((C(x) \land A(x)) \to \neg\exists y\st P(y,x))\]

Un modo naturale per tradurre la nuova frase potrebbe essere
\[\forall x, ((C(x) \land \exists y\st( P(y,x) )) \to \not A(x))\]

La cosa che le rende simile è che in linguaggio naturale risultano essere modi diversi per esprimere la stessa cosa, e quindi ci si aspetta che le loro traduzioni siano equivalenti. Un modo molto semplice per mostrare che sono equivalenti è sfruttare le proprietà della logica proposizionale che abbiamo visto in precedenza. Ignorando infatti la quantificazione e concentrandosi sulla mera struttura proposizionale, possiamo trasformare la seconda formula in
\[(A \land B )\to C\]
e di conseguenza la prima formula sarà equivalente a
\[(A \land \neg C) \to \neg B\]

Ma ora usando l'equivalenza \((P \to Q) \equiv (\neg P \lor Q)\) e De Morgan abbiamo
\[(A \land B )\to C\equiv \neg (A \land B) \lor C\equiv \neg A \lor \neg B \lor C \equiv (\neg A \lor C) \lor \neg B \equiv \neg (\neg A \lor C) \to \neg B \]
che è proprio la forma della prima formula.

Tornando alla formalizzazione, nella frase vista in precedenza, \q{C'è un cane che ha due proprietari} abbiamo notato come il linguaggio naturale contenga ambiguità riguardo all'espressione di quantità. Precedentemente, infatti, abbiamo interpretato tale frase come \q{C'è un cane che ha almeno due proprietari}, anche se non è la semantica usuale intesa nel parlato.

Consideriamo infatti la frase \q{Ogni genitore ha due figli}. Ancor di più questa frase è naturalmente percepita come intendere \q{Ogni genitore ha esattamente due figli}. In questo caso per poter effettuare una traduzione ci sarà necessaria una nuova segnatura. Consideriamo quindi come predicati un predicato unario \(G\) vero per i genitori e un predicato binario \(F\) vero per la relazione tra un genitore e un figlio.

In questo caso, il modo più naturale per tradurre la frase è
\[\forall x,( G(x) \to \exists y \st(\exists z\st (F(x,y) \land F(x,z) \land \neg(y = z) \land \forall w, (w = y \lor w = z \lor \neg F(x,w)))))\]

Chiaramente, una formulazione alternativa, ma equivalente, è
\[\forall x,( G(x) \to \exists y \st(\exists z\st (F(x,y) \land F(x,z) \land \neg(y = z) \land \neg \exists w\st (F(x,w) \land \neg(w = y) \land \neg(w = z)))))\]

poiché vale l'equivalenza \(\forall w, P(w) \equiv \neg \exists w, \neg P(w)\).
Un occhio di riguardo va dato all'uso esaustivo delle parentesi, che, per quanto rendano tutto molto verboso, rimuovono eventuali ambiguità sintattiche, che possono indurre a errori di interpretazione. 

Consideriamo ora una frase simile, ovvero \q{Ci sono precisamente tre numeri primi minori di 7}. In questo caso, come predicati useremo \(N\) vero per i numeri, \(P\) vero per i numeri primi e \(<\) e \(=\) con l'interpretazione usuale. Inoltre, sarà necessario un simbolo di costante \(7\) da interpretare come il numero \(7\). Considerando quindi la formula \(\phi_{PN}(x,y,z) \eqdef P(x) \land P(y) \land P(z) \land N(x) \land N(y) \land N(z) \) possiamo tradurre la frase con la formula
\[\exists x :(\exists y : (\exists z :(\neg (x=y) \land \neg (x=z) \land \neg (y=z) \land \phi_{PN}(x,y,z) \land x<7 \land y<7 \land z<7)))\]

Questa formula però non coglie il \q{precisamente} della frase, poiché è soddisfatta anche da un dominio in cui ci sono almeno tre numeri primi minori di \(7\). Alla catena di congiunzioni va quindi aggiunta la formula
\[\forall k, ((N(k)\land P(k) \land k<7) \to (k=x\lor k=y \lor k=z))\]

Questa frase ci mostra come generalizzare il concetto di \q{esattamente \(n\)} a qualsiasi \(n\). In particolare, per esprimere che ci sono esattamente \(n\) elementi che soddisfano una certa proprietà \(P\), è necessario introdurre \(n\) variabili distinte \(x_1, \dots, x_n\) e congiungere la formula che afferma che tutte queste variabili soddisfano \(P\) e sono distinte tra loro, con la formula che afferma che ogni elemento che soddisfa \(P\) è uguale a uno di questi \(n\) elementi.

Consideriamo ora la frase \q{Ci sono infiniti numeri primi}. Normalmente, il concetto di infinitezza non è esprimibile in logica del prim'ordine, ma, avendo a disposizione un predicato binario \(<\) con l'interpretazione usuale, è possibile esprimere questa frase con la formula
\[\forall x,( (P(x) \land N(x)) \to \exists y, (P(y) \land N(y) \land x<y))\]
Per poter esprimere questa infinitezza, ci siamo dovuti affidare alla presenza di un predicato che fosse antiriflessivo, transitivo e totale. Potendo, come visto, assiomatizzare il comportamento di \(<\), è possibile esprimere in questo caso l'infinità. In generale, possiamo parlare di infinitezza solo su strutture ordinabili, ovvero per le quali esiste un predicato binario che soddisfa le proprietà di antiriflessività, transitività e totalità.

Passiamo quindi alla frase \q{I numeri naturali divisori di un numero sono finiti}. Un modo alternativo di formulare lo stesso concetto è \q{Ogni numero naturale ha un numero finito di divisori}. Sempre nella stessa segnatura di prima, con l'aggiunta del predicato binario \(D\) per la divisibilità, questa frase può essere tradotta con la formula
\[\forall x, (N(x) \to \exists y: (D(y,x) \land N(y) \land \forall z, ((D(z,x) \land N(z)) \to (z < y \lor z = y))))\]

Vediamo infine qualche schema di traduzione generale per le locuzioni \q{almeno}, \q{al più} ed \q{esattamente}.

Iniziamo con \q{Ci sono almeno \(n\) elementi}. Ovviamente \(n\) è un parametro della formula da cui dipenderà la lunghezza della formula. Sarà necessario solo \(=\) come predicato.
\[\phi_{\geq n} \eqdef \exists x_1 \dots \exists x_n: (\bigwedge_{i=1}^n \bigwedge_{j=i+1}^n \neg (x_i = x_j))\]

Per poter dire invece \q{Ci sono al più \(n\) elementi} usando solo \(=\) come predicato, possiamo sia negare una formula esistenziale, oppure usare la formula
\[\phi_{\leq n} \eqdef \forall x_1 \dots \forall x_{n+1}: (\bigvee_{i=1}^{n+1} \bigvee_{j=i+1}^{n+1} (x_i = x_j))\equiv \neg \phi_{\geq n+1}\]

Infine, per esprimere \q{Ci sono esattamente \(n\) elementi} è sufficiente congiungere le due formule precedenti, ottenendo
\[\phi_{= n} \eqdef \phi_{\geq n} \land \phi_{\leq n}\]

Questo però porta a una formula molto lunga. Lo stesso concetto può essere espresso in maniera più compatta come segue
\[\phi_{= n} \eqdef \exists x_1 \dots \exists x_n ((\bigwedge_{i=1}^n \bigwedge_{j=i+1}^n \neg (x_i = x_j)) \land \forall x_{n+1}, \bigvee_{i=1}^{n} (x_{n+1} = x_i) )\]

\chapter{Calcolo Deduttivo del Prim'Ordine}

Come osservato nella logica proposizionale, oltre a sintassi e semantica, parte fondamentale di una logica è il suo calcolo deduttivo, ovvero il sistema di regole che ci permette, partendo da assiomi e assunzioni, di dimostrare teoremi. In questo capitolo analizzeremo il calcolo per la logica del prim'ordine, definendone le regole e dimostrando la sua correttezza e completezza.

\section{Calcolo di deduzione naturale per il prim'ordine}

Essendo il calcolo di deduzione naturale definito in modo che ogni regola catturi la semantica attesa di ogni operatore logico, è naturale pensare che il calcolo di deduzione naturale per la logica del prim'ordine sia definito come estensione di quello proposizionale, con l'aggiunta di regole per i quantificatori e, eventualmente, per l'uguaglianza se considerata simbolo logico.

Di conseguenza, le regole per i simboli proposizionali rimangono invariate, come rimangono invariate le loro proprietà.

Essendo i quantificatori operatori unari, avremo per ognuno due regole, una di introduzione e una di eliminazione. Abbiamo inoltre già osservato come i quantificatori siano generalizzazioni di congiunzione e disgiunzione, e quindi è naturale aspettarsi che le regole di introduzione ed eliminazione per i quantificatori siano generalizzazioni di quelle per la congiunzione e la disgiunzione. In particolare, ci aspettiamo che le regole per il quantificatore universale siano generalizzazioni di quelle per la congiunzione, e che le regole per il quantificatore esistenziale siano generalizzazioni di quelle per la disgiunzione.

Inoltre, abbiamo già osservato nel capitolo sul calcolo proposizionale di come il calcolo funga in un certo senso tra \textit{ponte} tra sintassi e semantica, nel senso che, tramite operazioni di pura manipolazione sintattica, si tenta di catturare la semantica della conseguenza logica. Di conseguenza, non sorprenderà che, essendo che la semantica dei quantificatori è definita istanziando le variabili quantificate tramite sostituzioni semantiche, le regole a loro associate faranno uso delle sostituzioni sintattiche e sfrutteranno le proprietà che le collegano a quelle semantiche viste nel capitolo precedente. 

Di fatto, ciò è quello che accade in una possibile dimostrazione di uno dei primi esempi di ragionamento deduttivo che abbiamo visto, ovvero \q{Tutti gli uomini sono mortali, Socrate è un uomo, quindi Socrate è mortale}. In questo caso, la formalizzazione di queste affermazioni è la seguente:
\begin{itemize}
  \item \q{Tutti gli uomini sono mortali}:\(\forall x, (U(x)\to M(x))\)
  \item \q{Socrate è un uomo}:\( U(S)\)
  \item \q{Socrate è mortale}:\(M(S)\)
\end{itemize}

Chiaramente, per dimostrare \(M(S)\) a partire da \(\forall x, (U(x)\to M(x))\) e \(U(S)\), è necessario in qualche modo istanziare la variabile \(x\) con il termine \(S\) per poter applicare \(\to I\).

In particolare, tali proprietà valevano sotto opportune condizioni, che per le regole rappresenteranno dei vincoli di applicabilità. A differenza di quanto accadeva per i connettivi proposizionali quindi, le regole che andremo a vedere non saranno sempre applicabili quando si presenta un certo pattern di struttura dell'albero di dimostrazione, ma andranno verificate tali condizioni.
Nella definizione di tali regole saranno ovviamente fondamentali le sostituzioni, che siano sintattiche o semantiche, e le loro proprietà viste nei capitoli precedenti. 

Vediamo dunque tali regole.

\paragraph{Regole per \(\forall\):}

Iniziamo con le regole per il quantificatore universale. I vincoli di applicabilità per queste regole sono riportati di fianco a ogni regole. Verificheremo in seguito che tali vincoli sono necessari e sufficienti per garantire la correttezza del calcolo.

\begin{equation*}
  \infer[(\forall E), \text{ se } t \text{ è libero per } x \text{ in } \phi]{\subs{\phi}{x}{t}}{\forall x, \phi}
\end{equation*}
\begin{equation*}
  \infer[(\forall I), \text{ se } \substack{ y \text{ non occorre in } \phi, \\ y \text{ non occorre libera in assunzioni non scaricate per } {\subs{\phi}{x}{y}}}]{\forall x, \phi}{\subs{\phi}{x}{y}}
\end{equation*}


Seppure i vincoli per \(\forall I\) possano risultare contorti, questi catturano esattamente il tipico ragionamento per dimostrare una proprietà universale, ovvero si prende un elemento generico \(y\) e si mostra che soddisfa la proprietà, da cui è possibile concludere che tutti gli elementi soddisfano la proprietà, non avendo fatto alcuna assunzione specifica su \(y\). Se \(y\) non fosse generico, non potremmo concludere niente su tutti gli elementi.

L'analogia con le regole per la congiunzione è abbastanza chiara. Infatti, considerando sempre \(\forall\) come congiunzione infinita su tutte le istanziazioni della variabile, la regola di eliminazione seleziona un congiunto specifico, proprio come \(\land E\), mentre quella di introduzione, mediante i vincoli, pretende che sia dimostrato un congiunto generico, ovvero tutti i congiunti, proprio come \(\land I\) pretende che siano dimostrati due congiunti specifici.

\paragraph{Regole per \(\exists\):}

Le regole per il quantificatore esistenziale sono le seguenti:

\begin{equation*}
  \infer[(\exists I), \text{ se } t \text{ è libero per } x \text{ in } \phi]{\exists x\st \phi}{\subs{\phi}{x}{t}}
\end{equation*}
\begin{equation*}
  \infer[(\exists E ), \text{ se }
          \substack{  y \text{ non occorre in }\phi, \\
            y \text{ non occorre libera in assunzioni non scaricate per } \psi\\
            y \text{ non occorre libera in } \psi            
            }
        ]{\psi}{
          \begin{array}[b]{c}
            \exists x\st \phi
            \end{array} &
            \begin{array}[b]{c}
            {[\subs{\phi}{x}{y}]}_1 \\
            \vdots \\
            \psi
          \end{array} 
  }
\end{equation*}

Qui l'analogia con le regole per la disgiunzione è ancora più chiara. Infatti, considerando sempre \(\exists\) come disgiunzione infinita su tutte le istanziazioni della variabile, la regola di introduzione ci permette di dedurre un esistenziale a partire da un qualsiasi disgiunto, proprio come \(\lor I\), mentre quella di eliminazione, mediante i vincoli e la struttura dell'albero di dimostrazione che ricorda \(\lor E\), richiede di dimostrare la conclusione assumendo uno alla volta ogni disgiunto, facendo di fatto un ragionamento per casi infinitario.

Introdotte le regole, andiamo ad analizzare perché i vincoli di applicabilità associati a esse sono necessari. Per la sufficienza invece, delegheremo la dimostrazione alla dimostrazione di correttezza del calcolo in toto, che vedremo in seguito.

Per verificare la necessità dei vincoli, sarà sufficiente trovare dei controesempi, in questo caso applicazioni di regole con un singolo vincolo violato, che portano a dedurre conclusioni che non sono conseguenze logiche delle premesse. In questi casi, l'unico motivo per cui la deduzione è errata è proprio la violazione del vincolo, e quindi si avrà che il vincolo è necessario.

Iniziamo con applicazioni errate di \(\forall E\). Consideriamo \(\forall x, \exists y\st x < y\). Ora, per \(t=y\) si ha \(\subs{\exists y\st x < y}{x}{y} = \exists y\st y < y\).

Quindi applicando \(\forall E\) avremmo

\begin{prooftree}
  \AxiomC{\(\forall x, \exists y\st x < y\)}
  \RuleLabel{\(\forall E\) errata}
  \UnaryInfC{\(\exists y\st y < y\)}
\end{prooftree}

Tale conclusione è chiaramente non valida, perché già solo considerando come struttura \(\N\) con interpretazione usuale la premessa è soddisfatta, mentre la conclusione è falsificata, essendo \(y\) un elemento di \(\N\) e quindi non è vero che \(y < y\).

Questo si ha perché \(y\) non è libero per \(x\) in \(\exists y\st x < y\), poiché \(y\) è legato da un quantificatore esistenziale, e quindi non essendo soddisfatte le condizioni per l'applicazione di \(\forall E\), ovvero le premesse del \Namedcref{thm:fol_substitution}, non è garantito che la sostituzione sintattica di \(y\) per \(x\) in \(\exists y\st x < y\) sia semanticamente equivalente alla formula originale, e quindi non è garantito che la deduzione sia corretta.

Passiamo ora a \(\forall I\). Consideriamo \(\forall x, (x=0\to x+1=1)\), che è vera in \(\N\) con interpretazione usuale. Ora, applicando \(\forall E\) correttamente otteniamo:

\begin{prooftree}
  \AxiomC{\(\forall x, (x=0\to x+1=1)\)}
  \RuleLabel{\(\forall E\) corretta}
  \UnaryInfC{\(y=0\to y+1=1\)}
\end{prooftree}

Ora assumendo \(y=0\) e applicando \(\to E\) correttamente otteniamo:

\begin{prooftree}
  \AxiomC{\(\forall x, (x=0\to x+1=1)\)}
  \RuleLabel{\(\forall E\) corretta}
  \UnaryInfC{\(y=0\to y+1=1\)}
  \AxiomC{\(y=0\)}
  \RuleLabel{\(\to E\) corretta}
  \BinaryInfC{\(y+1=1\)}
\end{prooftree}

Se ora applicassimo \(\forall I\) erratamente, violando il secondo vincolo, ovvero avendo \(y\) che occorre libera in un'assunzione non scaricata, ovvero \(y=0\), avremmo:

\begin{prooftree}
  \AxiomC{\(\forall x, (x=0\to x+1=1)\)}
  \RuleLabel{\(\forall E\) corretta}
  \UnaryInfC{\(y=0\to y+1=1\)}
  \AxiomC{\(y=0\)}
  \RuleLabel{\(\to E\) corretta}
  \BinaryInfC{\(y+1=1\)}
  \RuleLabel{\(\forall I\) errata}
  \UnaryInfC{\(\forall x, (x+1=1)\)}
\end{prooftree}
che chiaramente, sempre considerando come struttura \(\N\) con interpretazione usuale, è falsificata.

Questo esempio mostra chiaramente l'analogia tra questa regola e il processo usuale di dimostrazione di una proprietà universale. Una versione informale di questa dimostrazione infatti sarebbe.
\begin{quote}
  \q{
    Sappiamo che per ogni \(x\in\N\), se \(x=0\) allora \(x+1=1\). Dunque, qualsiasi sia \(y\in\N\), se \(y=0\) allora \(y+1=1\). Consideriamo dunque \(y=0\). Dato che \(y=0\), allora \(y+1=1\).
  }
\end{quote}

Fino a qui sarebbe tutto corretto, ma se volessimo trarne conclusioni universali, dovremmo aggiungere la locuzione \q{Essendo \(y\) un elemento generico di \(\N\), possiamo concludere}, che ovviamente non potremmo fare essendo \(y=0\), quindi tutt'altro che generico.

Per il secondo vincolo, consideriamo \(\forall x, x=x\), anch'essa vera in \(\N\) con interpretazione usuale\footnote{Tale formula è vera in ogni struttura con \(=\) interpretata come equivalenza}. Ancora, per \(\forall E\) applicata correttamente, otteniamo:

\begin{prooftree}
  \AxiomC{\(\forall x, x=x\)}
  \RuleLabel{\(\forall E\) corretta}
  \UnaryInfC{\(z=z\)}
\end{prooftree}

Ora, considerando \(\phi \eqdef x=z\), possiamo notare che \(z=z\) è proprio \(\subs{\phi}{x}{z}\), e quindi, se applicassimo \(\forall I\) erratamente, violando il primo vincolo, ovvero avendo \(z\) che occorre in \(\phi\), avremmo:
\begin{prooftree}
  \AxiomC{\(\forall x, x=x\)}
  \RuleLabel{\(\forall E\) corretta}
  \UnaryInfC{\(z=z\)}
  \RuleLabel{\(\forall I\) errata}
  \UnaryInfC{\(\forall x, x=z\)}
\end{prooftree}

A questo punto applicando nuovamente \(\forall I\), questa volta correttamente\footnote{In questo caso \(\phi \eqdef \forall x, x=y\), dunque \(z\) non occorre in \(\phi\), ed è possibile di conseguenza dedurre \(\forall y, \phi\) da \(\subs{\phi}{y}{z}\), ovvero \(\forall x, x=y\)}, avremmo:
\begin{prooftree}
  \AxiomC{\(\forall x, x=x\)}
  \RuleLabel{\(\forall E\) corretta}
  \UnaryInfC{\(z=z\)}
  \RuleLabel{\(\forall I\) errata}
  \UnaryInfC{\(\forall x, x=z\)}
  \RuleLabel{\(\forall I\) corretta}
  \UnaryInfC{\(\forall y, \forall x, x=y\)}
\end{prooftree}
Ovvero, dal fatto che ogni oggetto è uguale a se stesso, possiamo dedurre che ogni oggetto è uguale a ogni altro oggetto, il che è chiaramente falso in \(\N\) con interpretazione usuale, e quindi la deduzione è errata.

Per passare alle regole per \(\exists\), possiamo iniziare a notare che l'ultimo esempio in un caso esistenziale sarebbe assolutamente legittimo.
Infatti da \(z=z\) è corretto dedurre \(\exists x\st x=z\), da cui è corretto dedurre \(\exists x\st \exists y\st x=y\).

Questo perché l'unica cosa da verificare è che \(z\) sia libera per \(x\) in \(x=z\), e questo è vero poiché \(z\) è libera in \(x=z\), e quindi a maggior ragione è libera per \(x\).

Consideriamo però un caso in cui si viola il vincolo di applicabilità di \(\exists I\). Chiaramente, vista la somiglianza dei vincoli di applicabilità, questo sarà molto simile a quello di \(\forall E\). Consideriamo infatti \(\forall y, y = y\), che abbiamo già osservato essere vera in \(\N\) con interpretazione usuale. Ora, applicando \(\exists I\) erratamente, si potrebbe pensare di dedurre \(\exists x\st\forall y, x = y\), che è chiaramente falsa in ogni struttura con più di un elemento. Possiamo notare infatti che per \(\phi \eqdef \forall y, x = y\), si ha che \(\subs{\phi}{x}{y} = \forall y, y = y\), ed è quindi rispettata la struttura di \(\exists I\). A non essere soddisfatto però è il vincolo di applicabilità, poiché \(y\) non è libera per \(x\) in \(\forall y, x = y\), poiché \(y\) è legato da un quantificatore universale.

Passiamo ora alla regola più articolata, ovvero quella di eliminazione di \(\exists\). 
Consideriamo \(\exists x\st\exists y \st x<y\), anch'essa chiaramente vera in \(\N\) con interpretazione usuale. A questo punto, assumiamo \(\exists y\st y<y\), che è chiaramente falsa in \(\N\) con interpretazione usuale. Possiamo notare però che per \(\phi \eqdef \exists y \st x<y\), si ha che \(\subs{\phi}{x}{y} = \exists y \st y<y\). Di conseguenza, essendo rispettata la struttura della regola si potrebbe considerare valida la seguente deduzione:
\begin{prooftree}
  \AxiomC{\({[\exists y\st y<y]}_1\)}
  \UnaryInfC{\(\exists y\st y<y\)}
  \AxiomC{\(\exists x\st\exists y \st x<y\)}
  \RuleLabel{\(\exists E_1\) errata}
  \BinaryInfC{\(\exists y\st y<y\)}
\end{prooftree}
che è chiaramente errata, essendo la premessa vera e la conclusione falsa in \(\N\) con interpretazione usuale. Il vincolo violato in questo caso è che \(y\) occorre nella formula quantificata.

Consideriamo ora \(\exists x \st x \leq 5\). Chiaramente con \(\phi \eqdef x \leq 5\) si ha che \(\subs{\phi}{x}{y} = y \leq 5\), e quindi assumendo \(y \leq 5\) e applicando \(\exists E\) erratamente otterremmo:
\begin{prooftree}
  \AxiomC{\({[y \leq 5]}_1\)}
  \UnaryInfC{\(y \leq 5\)}
  \AxiomC{\(\exists x \st x \leq 5\)}
  \RuleLabel{\(\exists E_1\) errata}
  \BinaryInfC{\(y \leq 5\)}
\end{prooftree}

A questo punto, applicando correttamente \(\forall I\) otterremmo:

\begin{prooftree}
  \AxiomC{\({[y \leq 5]}_1\)}
  \UnaryInfC{\(y \leq 5\)}
  \AxiomC{\(\exists x \st x \leq 5\)}
  \RuleLabel{\(\exists E_1\) errata}
  \BinaryInfC{\(y \leq 5\)}
  \RuleLabel{\(\forall I\) corretta}
  \UnaryInfC{\(\forall z, z \leq 5\)}
\end{prooftree}
che è chiaramente falso in \(\N\) con interpretazione usuale, e quindi la deduzione è errata. Il vincolo che si è violato in questo caso è che \(y\) occorre libera nella conclusione \(\psi \eqdef y < 5\).

Infine, assumiamo \(z=5\) e \(\neq (z = 5)\). Chiaramente si ha:

\begin{prooftree}
  \AxiomC{\(z=5\)}
  \AxiomC{\(\neq (z=5)\)}
  \RuleLabel{\(\neg E\)}
  \BinaryInfC{\(\bot\)}
\end{prooftree}

A questo punto, considerando \(\phi \eqdef x = 5\), si ha che \(\subs{\phi}{x}{z}\) è \( z = 5\), dunque, assumendo \(\exists x\st x = 5\) per \(\exists E\) possiamo dedurre \(\bot\) scaricando l'assunzione \(z=5\).

\begin{prooftree}\label{prooft:wrong_fol_exists_e}
  \AxiomC{\({[z=5]}_1\)}
  \AxiomC{\(\neq (z=5)\)}
  \RuleLabel{\(\neg E\)}
  \BinaryInfC{\(\bot\)}
  \AxiomC{\(\exists x\st x = 5\)}
  \RuleLabel{\(\exists E_1\) errata}
  \BinaryInfC{\(\bot\)}
\end{prooftree}

Ora, tramite \(\RAA\), potremmo scaricare \(\neq (z=5)\) deducendo \(z=5\) e, come nel precedente esempio, mediante \(\forall I\) ottenere \(\forall x, x=5\).

\begin{prooftree}
  \AxiomC{\({[z=5]}_1\)}
  \AxiomC{\({[\neq (z=5)]}_2\)}
  \RuleLabel{\(\neg E\)}
  \BinaryInfC{\(\bot\)}
  \AxiomC{\(\exists x\st x = 5\)}
  \RuleLabel{\(\exists E_1\) errata}
  \BinaryInfC{\(\bot\)}
  \RuleLabel{\(\RAA_2\)}
  \UnaryInfC{\(z=5\)}
  \RuleLabel{\(\forall I\) corretta}
  \UnaryInfC{\(\forall x, x=5\)}
\end{prooftree}

Si è dunque dimostrato che, assumendo \(\exists x\st x=5\), si può dedurre che \(\forall x, x = 5\). Tale deduzione è chiaramente errata, essendo la premessa vera e la conclusione falsa in \(\N\) con interpretazione usuale. Qui notare la violazione del vincolo è più sottile, poiché a deduzione conclusa a una prima occhiata tutte le assunzioni sembrano essere state scaricate. È importante notare però che per l'applicabilità di una regola, è necessario che i vincoli siano soddisfatti già al momento dell'applicazione della regola, ovvero nel sottoalbero di deduzione radicato nella regola. Come si può vedere chiaramente nel secondo albero di deduzione di questo caso, al momento dell'applicazione di \(\exists E\), \(\neg (z = 5)\) non è ancora stata scaricata da \(\RAA\) ed è quindi ancora presente come assunzione non scaricata, contenendo \(z\) libera. Dunque è violato l'ultimo vincolo di applicabilità di \(\exists E\), e rendendo la deduzione è errata.

Mostrati dunque esempi errati di deduzione naturale in prim'ordine, per evidenziare la necessità dei vincoli di applicabilità delle regole per i quantificatori, passiamo ora a mostrare esempi di deduzioni corrette.


\begin{example}{Deduzione naturale per \(\FOL\) 1}
  Vogliamo dimostrare \(\vdash (\forall x, (P(x) \land R(x)))\leftrightarrow (\forall x, P(x) \land \forall x, R(x))\).

  Per farlo dimostriamo entrambe le direzioni, che per semplicità mostreremo come \((\forall x, (P(x) \land R(x)))\vdash (\forall x, P(x) \land \forall x, R(x))\) e \((\forall x, P(x) \land \forall x, R(x))\vdash (\forall x, (P(x) \land R(x)))\).

  Iniziamo con la prima.

  L'operatore principale è \(\land\), quindi è ragionevole iniziare con \(\land I\).

  \begin{prooftree}
    \AxiomC{\(\forall x, P(x)\)}
    \AxiomC{\(\forall x, R(x)\)}
    \RuleLabel{\(\land I\)}
    \BinaryInfC{\(\forall x, P(x) \land \forall x, R(x)\)}
  \end{prooftree}

  Abbiamo dunque da dimostrare due formule con \(\forall\) come operatore principale. In linea di principio, per mostrare tali formule si può procedere o per contraddizione, o usando \(\forall I\). In questo caso la contraddizione risulta poco utile, non avendo a disposizione negazioni per generarla, quindi è ragionevole procedere con \(\forall I\).
  
  \begin{prooftree}
    \AxiomC{\(P(y)\)}
    \RuleLabel{\(\forall I\)}
    \UnaryInfC{\(\forall x, P(x)\)}
    \AxiomC{\(R(y)\)}
    \RuleLabel{\(\forall I\)}
    \UnaryInfC{\(\forall x, R(x)\)}
    \RuleLabel{\(\land I\)}
    \BinaryInfC{\(\forall x, P(x) \land \forall x, R(x)\)}
  \end{prooftree}

  A questo punto con l'assunzione \(\forall x, (P(x) \land R(x))\) possiamo dedurre \(P(y) \land R(y)\) per \(\forall E\) da cui, con \(\land E\) otteniamo i congiunti che ci servivano a chiudere la dimostrazione.

   \begin{prooftree}
        \AxiomC{\(\forall x, (P(x) \land R(x))\)}
        \RuleLabel{\(\forall E\)}
        \UnaryInfC{\(P(y) \land R(y)\)}
        \RuleLabel{\(\land E^l\)}
        \UnaryInfC{\(P(y)\)}
        \RuleLabel{\(\forall I\)}
        \UnaryInfC{\(\forall x, P(x)\)}

        \AxiomC{\(\forall x, (P(x) \land R(x))\)}
        \RuleLabel{\(\forall E\)}
        \UnaryInfC{\(P(y) \land R(y)\)}
        \RuleLabel{\(\land E^r\)}
        \UnaryInfC{\(R(y)\)}
        \RuleLabel{\(\forall I\)}
        \UnaryInfC{\(\forall x, R(x)\)}
      \RuleLabel{\(\land I\)}
      \BinaryInfC{\(\forall x, P(x) \land \forall x, R(x)\)}
  \end{prooftree}

  Finita la prova è importante controllare che tutti i vincoli di applicabilità siano rispettati. In questo caso, avendo usato \(\forall I\) per dimostrare \(\forall x, P(x)\) e \(\forall x, R(x)\), è necessario che \(y\) non occorra in \(\phi\), ovvero in \(P(x)\) e \(R(x)\), e che \(y\) non occorra libera in assunzioni non scaricate per \(\forall x, P(x)\) e \(\forall x, R(x)\). Per quanto riguarda \(\forall E\) invece è necessario che \(y\) sia libero per \(x\) in \(P(x) \land R(x)\), e questo è vero poiché \(y\) non occorre in \(P(x) \land R(x)\).
  \end{example}
  \begin{example}{Deduzione naturale per \(\FOL\) 2}
  Passiamo all'altro verso, ovvero \((\forall x, P(x) \land \forall x, R(x))\vdash (\forall x, (P(x) \land R(x)))\).

  Questa volta partiamo con una formula quantificata. Ancora una volta, è ragionevole iniziare con \(\forall I\), non avendo negazioni con cui produrre una contraddizione.

  \begin{prooftree}
    \AxiomC{\(P(y) \land R(y)\)}
    \RuleLabel{\(\forall I\)}
    \UnaryInfC{\(\forall x, (P(x) \land R(x))\)}
  \end{prooftree}

  A questo punto possiamo ottenere la congiunzione tramite \(\land I\) dovendo dimostrare \(P(y)\) e \(R(y)\).
  
  \begin{prooftree}
    \AxiomC{\(P(y)\)}
    \AxiomC{\(R(y)\)}
    \RuleLabel{\(\land I\)}
    \BinaryInfC{\(P(y) \land R(y)\)}
    \RuleLabel{\(\forall I\)}
    \UnaryInfC{\(\forall x, (P(x) \land R(x))\)}
  \end{prooftree}

  Possiamo notare che avendo usato \(\forall I\) non possiamo più usare assunzioni non scaricate con \(y\) libero. Fortunatamente, l'assunzione che abbiamo non contiene \(y\). Proprio da tale assunzione infatti possiamo dedurre i due congiunti \(\forall x, P(x)\) e \(\forall x, R(x)\) mediante \(\land E\), dai quali, con \(\forall E\) potremmo dedurre \(P(y)\) e \(R(y)\) che ci servono per concludere la dimostrazione.

  \begin{prooftree}
      \AxiomC{\(\forall x, P(x) \land \forall x, R(x)\)}
      \RuleLabel{\(\land E^l\)}
      \UnaryInfC{\(\forall x, P(x)\)}
      \RuleLabel{\(\forall E\)}
      \UnaryInfC{\(P(y)\)}
      
      \AxiomC{\(\forall x, P(x) \land \forall x, R(x)\)}
      \RuleLabel{\(\land E^r\)}
      \UnaryInfC{\(\forall x, R(x)\)}
      \RuleLabel{\(\forall E\)}
      \UnaryInfC{\(R(y)\)}
    \RuleLabel{\(\land I\)}
    \BinaryInfC{\(P(y) \land R(y)\)}
    \RuleLabel{\(\forall I\)}
    \UnaryInfC{\(\forall x, (P(x) \land R(x))\)}
  \end{prooftree}

  Anche qui possiamo verificare che i vincoli di \(\forall E\) sono rispettati, poiché \(y\) è libero per \(x\) in \(P(x)\) e \(R(x)\) non apparendo.
\end{example}

\subsection{Applicazione del calcolo naturale con i quantificatori}

Come fatto per la logica proposizionale, andiamo ad analizzare nel dettaglio alcuni esempi di deduzioni naturali che esploreranno tutte le nuove regole di inferenza che abbiamo introdotto, in modo da avere un'idea più chiara di come funziona questo calcolo, e di come scegliere le regole da applicare per dedurre formule più complesse.

\subsubsection*{Deduzione di \(\vdash (\exists x\st (P(x)\lor R(x)) )\to (\exists x\st P(x) \lor \exists x\st R(x))\)}

Come fatto in precedenza, per semplicità ci limiteremo a dimostrare \(\exists x\st (P(x)\lor R(x)) \vdash \exists x\st P(x) \lor \exists x\st R(x)\), dalla quale, tramite \(\to I\) si potrà dedurre la formula che ci interessa.

Una prima osservazione che è possibile fare è che, avendo come premessa una formula con un quantificatore esistenziale come operatore principale, è possibile contare tra le assunzioni disponibili anche la formula che quel quantificatore esistenziale quantifica, ovvero \(P(y)\lor R(y)\), con \(y\) variabile arbitraria, poiché sarà scaricabile mediante \(\exists E\) a partire da \(\exists x\st (P(x)\lor R(x))\).

Dovendo dimostrare una disgiunzione, un primo approccio possibile sembrerebbe quello di utilizzare \(\lor I\), e per fare ciò sarebbe necessario dimostrare \(\exists x\st P(x)\) o \(\exists x\st R(x)\), per i quali servirebbe dimostrare \(P(t)\) o \(R(t)\) singolarmente.


\begin{center}
  \begin{minipage}{0.42\textwidth}
    \centering
    \begin{prooftree}
      \AxiomC{\(P(y)\)}
      \RuleLabel{\(\exists I\)}
      \UnaryInfC{\(\exists x\st P(x)\)}

      \RuleLabel{\(\lor I\)}
      \UnaryInfC{\(\exists x\st P(x) \lor \exists x\st R(x)\)}
    \end{prooftree}
  \end{minipage}
  \hspace{0.3cm}
  \begin{minipage}{0.42\textwidth}
    \centering
    \begin{prooftree}
      \AxiomC{\(R(y)\)}
      \RuleLabel{\(\exists I\)}
      \UnaryInfC{\(\exists x\st R(x)\)}

      \RuleLabel{\(\lor I\)}
      \UnaryInfC{\(\exists x\st P(x) \lor \exists x\st R(x)\)}
    \end{prooftree}
  \end{minipage}
\end{center}

Possiamo ora notare che abbiamo dimostrato la stessa formula, ovvero \(\exists x\st P(x) \lor \exists x\st R(x)\), assumendo una prima volta \(P(y)\) e una seconda volta \(R(y)\). Proprio questa situazione è quella che ci permette di applicare \(\lor E\) per dedurre \(\exists x\st P(x) \lor \exists x\st R(x)\) assumendo solo \(P(y) \lor R(y)\)


\begin{prooftree}
  \AxiomC{\({[P(y)]}_1\)}
  \RuleLabel{\(\exists I\)}
  \UnaryInfC{\(\exists x\st P(x)\)}

  \RuleLabel{\(\lor I\)}
  \UnaryInfC{\(\exists x\st P(x) \lor \exists x\st R(x)\)}

  \AxiomC{\({[R(y)]}_2\)}
  \RuleLabel{\(\exists I\)}
  \UnaryInfC{\(\exists x\st R(x)\)}

  \RuleLabel{\(\lor I\)}
  \UnaryInfC{\(\exists x\st P(x) \lor \exists x\st R(x)\)}
  \AxiomC{\(P(y) \lor R(y)\)}
  \RuleLabel{\(\lor E_{1,2}\)}
  \TrinaryInfC{\(\exists x\st P(x) \lor \exists x\st R(x)\)}
\end{prooftree}

A questo punto, come anticipato, possiamo scaricare \(P(y) \lor R(y)\) mediante \(\exists E\) con \(\exists x\st (P(x)\lor R(x))\), senza dover dunque assumere \(P(y)\lor R(y)\) come assunzione, e concludere la dimostrazione.

\begin{prooftree}
  \small
  \AxiomC{\({[P(y)]}_1\)}
  \RuleLabel{\(\exists I\)}
  \UnaryInfC{\(\exists x\st P(x)\)}

  \RuleLabel{\(\lor I\)}
  \UnaryInfC{\(\exists x\st P(x) \lor \exists x\st R(x)\)}

  \AxiomC{\({[R(y)]}_2\)}
  \RuleLabel{\(\exists I\)}
  \UnaryInfC{\(\exists x\st R(x)\)}

  \RuleLabel{\(\lor I\)}
  \UnaryInfC{\(\exists x\st P(x) \lor \exists x\st R(x)\)}

  \AxiomC{\({[P(y) \lor R(y)]}_3\)}
  
  \RuleLabel{\(\lor E_{1,2}\)}
  \TrinaryInfC{\(\exists x\st P(x) \lor \exists x\st R(x)\)}
  \AxiomC{\(\exists x\st (P(x)\lor R(x))\)}
  \RuleLabel{\(\exists E_3\)}
  \BinaryInfC{\(\exists x\st P(x) \lor \exists x\st R(x)\)}
\end{prooftree}

Andando dunque a controllare i vincoli di applicabilità, per entrambe le \(\exists I\) si ha che \(y\) non è in \(P(x)\) e \(R(x)\), e quindi è banalmente libero per \(x\) in \(P(x)\) e \(R(x)\).
Per \(\exists E\) dobbiamo verificare che \(y\) non occorre in \(\phi\), ovvero in \(P(x)\lor R(x)\), che è verificato, che \(y\) non appartenga a nessuna assunzione non scaricata, e questo è verificato poiché l'unica assunzione fatta è scaricata da \(\exists E\), e che \(y\) non occorra libera in \(\psi\).

\begin{warning}{}
  L'ordine di applicazione delle regole è molto importante in questo esempio. Si potrebbe infatti pensare, al posto di scaricare \(P(y)\lor R(y)\) con \(\exists E\) di dimostrarla, rendendola così non più un assunzione con una deduzione del tipo:
  
  \begin{prooftree}
    \tiny
    \AxiomC{\({[P(y)]}_1\)}
    \RuleLabel{\(\exists I\)}
    \UnaryInfC{\(\exists x\st P(x)\)}

    \RuleLabel{\(\lor I\)}
    \UnaryInfC{\(\exists x\st P(x) \lor \exists x\st R(x)\)}

    \AxiomC{\({[R(y)]}_2\)}
    \RuleLabel{\(\exists I\)}
    \UnaryInfC{\(\exists x\st R(x)\)}

    \RuleLabel{\(\lor I\)}
    \UnaryInfC{\(\exists x\st P(x) \lor \exists x\st R(x)\)}

    \AxiomC{\({[P(y) \lor R(y)]}_3\)}
    \UnaryInfC{\(P(y) \lor R(y)\)}
    \AxiomC{\(\exists x\st (P(x)\lor R(x))\)}
    \RuleLabel{\(\exists E_3\)}
    \BinaryInfC{\(P(y) \lor R(y)\)}

    \RuleLabel{\(\lor E_{1,2}\)}
    \TrinaryInfC{\(\exists x\st P(x) \lor \exists x\st R(x)\)}
  \end{prooftree}

  Se analizziamo bene il sottoalbero di \(\exists E\) però notiamo che si viola il vincolo di applicabilità.
  \begin{prooftree}
    \AxiomC{\({[P(y) \lor R(y)]}_3\)}
    \UnaryInfC{\(P(y) \lor R(y)\)}
    \AxiomC{\(\exists x\st (P(x)\lor R(x))\)}
    \RuleLabel{\(\exists E_3\)}
    \BinaryInfC{\(P(y) \lor R(y)\)}
  \end{prooftree}

  Infatti, \(y\) occorre libera in \(\psi\), ovvero in \(P(y) \lor R(y)\), e quindi non è possibile applicare \(\exists E\) stesso. Nonostante le deduzioni siano estremamente simili, e in particolare usino le stesse regole, il solo scambiare l'ordine di applicazione di due di queste ci porta a una deduzione errata.
\end{warning}

\subsubsection*{Dimostrazione di \(\vdash (\exists x, P(x) \lor \exists x, R(x))\to (\exists x, (P(x)\lor R(x)) )\)}

Passiamo a dimostrare il verso opposto, ovvero \(\exists x\st P(x) \lor \exists x\st R(x) \vdash \exists x\st (P(x)\lor R(x))\).

Essendo che la premessa, e quindi unica assunzione che è possibile non scaricare, è una disgiunzione, è ragionevole pensare di provare a dedurre la conclusione per casi assumendo un congiunto della disgiunzione alla volta, per poi chiudere la dimostrazione con \(\lor E\).

Assumendo \(\exists x\st P(x)\) potremo scaricare \(P(y)\) per \(\exists E\), di conseguenza possiamo assumere \(P(y)\), da cui, con \(\lor I\) otteniamo \(P(y)\lor R(y)\). A questo punto, con \(\exists I\) otteniamo \(\exists x\st (P(x)\lor R(x))\).

\begin{prooftree}
  \AxiomC{\({[P(y)]}_1\)}
  \RuleLabel{\(\lor I\)}
  \UnaryInfC{\(P(y)\lor R(y)\)}
  \RuleLabel{\(\exists I\)}
  \UnaryInfC{\(\exists x\st (P(x)\lor R(x))\)}
  \AxiomC{\(\exists x\st P(x)\)}
  \RuleLabel{\(\exists E_1\)}
  \BinaryInfC{\(\exists x\st (P(x)\lor R(x))\)}
\end{prooftree}

È importante notare che è possibile applicare \(\exists E\) solo alla fine di questo ramo di dimostrazione e non prima, poiché in ogni altro punto del sottoalbero radicato in \(\exists E\) si avrebbe \(y\) libera nella conclusione raggiunta, e quindi si violerebbe il vincolo di applicabilità di \(\exists E\).

A questo punto possiamo ripetere la stessa dimostrazione assumendo \(\exists x\st R(x)\) e scaricando \(R(y)\) per \(\exists E\). Infine, per \(\lor E\), possiamo scaricare i due disgiunti, lasciando come unica assunzione non scaricata l'unica premessa, ovvero \(\exists x\st P(x) \lor \exists x\st R(x)\), e concludere la dimostrazione.

\begin{prooftree}
  \tiny
  \AxiomC{\({[P(y)]}_1\)}
  \RuleLabel{\(\lor I\)}
  \UnaryInfC{\(P(y)\lor R(y)\)}
  \RuleLabel{\(\exists I\)}
  \UnaryInfC{\(\exists x\st (P(x)\lor R(x))\)}
  \AxiomC{\({[\exists x\st P(x)]}_3\)}
  \RuleLabel{\(\exists E_1\)}
  \BinaryInfC{\(\exists x\st (P(x)\lor R(x))\)}

  \AxiomC{\({[\exists x\st R(x)]}_4\)}
  \AxiomC{\({[R(y)]}_2\)}
  \RuleLabel{\(\lor I\)}
  \UnaryInfC{\(P(y)\lor R(y)\)}
  \RuleLabel{\(\exists I\)}
  \UnaryInfC{\(\exists x\st (P(x)\lor R(x))\)}
  \RuleLabel{\(\exists E_2\)}
  \BinaryInfC{\(\exists x\st (P(x)\lor R(x))\)}

  \AxiomC{\(\exists x\st P(x) \lor \exists x\st R(x)\)}
  \RuleLabel{\(\lor E_{3,4}\)}
  \TrinaryInfC{\(\exists x\st (P(x)\lor R(x))\)}
\end{prooftree}

\subsubsection*{Dimostrazione di \(\forall x, (A(y) \to B(x)) \dashv\vdash A(y) \to \forall x, B(x)\)}

Continuiamo con la dimostrazione di equivalenze semantiche che coinvolgono quantificatori, dimostrando \(\forall x, (A(y) \to B(x)) \dashv\vdash A(y) \to \forall x, B(x)\).

\paragraph{Dimostrazione di \(\forall x, (A(y) \to B(x)) \vdash A(y) \to \forall x, B(x)\)}

Iniziamo col primo verso, ovvero \(\forall x, (A(y) \to B(x)) \vdash A(y) \to \forall x, B(x)\).

Qui l'operatore principale è l'implicazione, quindi come solito è ragionevole iniziare con \(\to I\).

\begin{prooftree}
  \AxiomC{\(\forall x, B(x)\)}
  \RuleLabel{\(\to I\)}[\(A(y)\)]
  \UnaryInfC{\(A(y) \to \forall x, B(x)\)}
\end{prooftree}

Ora per dimostrare una formula quantificata universalmente potremmo procedere con \(\forall I\), ovvero cercando di dimostrare \(B(y)\) per una variabile arbitraria \(y\).

\begin{prooftree}
  \AxiomC{\(B(z)\)}
  \RuleLabel{\(\forall I\)}
  \UnaryInfC{\(\forall x, B(x)\)}
  \RuleLabel{\(\to I\)}[\(A(y)\)]
  \UnaryInfC{\(A(y) \to \forall x, B(x)\)}
\end{prooftree}

Chiaramente, se ci fermassimo qui, l'applicazione di \(\forall I\) non sarebbe legittima, poiché \(z\) occorre libera in un'assunzione non scaricata, ovvero \(B(z)\). In ogni caso questo non è il punto d'arrivo sperato, poiché la deduzione esposta fin'ora è quella del sequente \(B(z) \vdash A(y) \to \forall x, B(x)\), che è molto lontano da quello che vogliamo dimostrare.

Per rendere legittima l'applicazione di \(\forall I\) sarò necessario o scaricare \(B(z)\) o dimostrarlo senza assumerlo. Avendo a disposizione le assunzioni \(\forall x, A(y)\to B(x)\) e \(A(y)\), per dimostrare \(B(z)\) è sufficiente applicare \(\forall E\) a \(\forall x, A(y)\to B(x)\) per ottenere \(A(y)\to B(z)\), da cui, con \(A(y)\) e \(\to E\) otteniamo \(B(z)\), chiudendo la dimostrazione.

\begin{prooftree}
  \AxiomC{\({[A(y)]}_1\)}
  \AxiomC{\(\forall x, A(y)\to B(x)\)}
  \RuleLabel{\(\forall E\)}
  \UnaryInfC{\(A(y)\to B(z)\)}
  \RuleLabel{\(\to E\)}
  \BinaryInfC{\(B(z)\)}
  \RuleLabel{\(\forall I\)}
  \UnaryInfC{\(\forall x, B(x)\)}
  \RuleLabel{\(\to I_1\)}
  \UnaryInfC{\(A(y) \to \forall x, B(x)\)}
\end{prooftree}

Un modo per facilitare la verifica dei vincoli di applicabilità è, come in questo caso, di usare variabili nuove dove possibile, in modo da evitare le occorrenze di tali variabili nelle formule coinvolte nelle regole. Da verificare in questo caso c'è, per \(\forall E\), che \(z\) non sia libera per \(x\) in \(A(y)\to B(x)\), e per \(\forall I\) che \(z\) non occorra in \(B(x)\) e che \(z\) non occorra libera in assunzioni non scaricate per \(B(z)\). Per quanto riguarda l'occorrere di \(z\) in \(B(x)\) e in \(A(y) \to B(x)\) la verifica è immediata avendo usato una variabile nuova, mentre per la verifica che \(z\) non occorra libera in assunzioni non scaricate per \(B(z)\) è sufficiente notare che \(A(y)\) non contiene \(z\) e che \(\forall x, A(y)\to B(x)\) non contiene variabili libere, e quindi \(z\) non occorre libera in nessuna assunzione.

È importante notare che i vincoli di applicabilità considerando non scaricata \(A(y)\), poiché, essendo che verrà scaricata da \(\to I\) in radice, nel momento in cui si applica \(\forall I\) l'assunzione non è ancora scaricata.

\paragraph{Dimostrazione di \(A(y) \to \forall x, B(x) \vdash \forall x, (A(y) \to B(x))\)}

Continuando con la dimostrazione dell'equivalenza, passiamo a dimostrare \(A(y) \to \forall x, B(x) \vdash \forall x, (A(y) \to B(x))\).

Ancora una volta ci troviamo a dimostrare una formula quantificata universalmente, e quindi è ragionevole pensare di procedere con \(\forall I\) dimostrando \(A(y) \to B(z)\) per una variabile arbitraria \(z\), ovvero, per \(\to I\), dimostrare \(B(z)\) potendo scaricare \(A(y)\).

\begin{prooftree}
  \AxiomC{\(B(z)\)}
  \RuleLabel{\(\to I\)}[\(A(y)\)]
  \UnaryInfC{\(A(y) \to B(z)\)}
  \RuleLabel{\(\forall I\)}
  \UnaryInfC{\(\forall x, (A(y) \to B(x))\)}
\end{prooftree}

Come prima, dobbiamo fare in modo di scaricare o dimostrare \(B(z)\) senza assumerlo, altrimenti l'applicazione di \(\forall I\) risulterà illegittima. Possiamo notare però che, avendo a disposizione \(A(y) \to \forall x, B(x)\) e \(A(y)\), è sufficiente applicare \(\to E\) per ottenere \(\forall x, B(x)\), da cui, con \(\forall E\) otteniamo \(B(z)\), chiudendo la dimostrazione.

\begin{prooftree}
  \AxiomC{\({[A(y)]}_1\)}
  \AxiomC{\(A(y) \to \forall x, B(x)\)}
  \RuleLabel{\(\to E\)}
  \BinaryInfC{\(\forall x, B(x)\)}
  \RuleLabel{\(\forall E\)}
  \UnaryInfC{\(B(z)\)}
  \RuleLabel{\(\to I_1\)}
  \UnaryInfC{\(A(y) \to B(z)\)}
  \RuleLabel{\(\forall I\)}
  \UnaryInfC{\(\forall x, (A(y) \to B(x))\)}
\end{prooftree}

Venendo quindi alla verifica dei vincoli di applicabilità, per \(\forall E\) dobbiamo verificare che \(z\) sia libero per \(x\) in \(B(x)\), e questo è verificato poiché \(z\) non occorre in \(B(x)\). Per \(\forall I\) invece dobbiamo verificare che \(z\) non occorra in \(A(y)\to B(x)\), che è verificato poiché \(z\) non occorre in \(A(y)\to B(x)\), e che \(z\) non occorra libera in assunzioni non scaricate per \(B(z)\). Nel momento dell'applicazione di \(\forall I\) le assunzioni non scaricate sono \(A(y)\) e \(A(y) \to \forall x, B(x)\), e nessuna delle due contiene \(z\).

\subsubsection*{Dimostrazione di \(\forall x, P(x) \dashv \vdash \neg \exists x\st \neg P(x)\)}

Passiamo ora a un esempio più complesso, ovvero alla dimostrazione della dualità dei quantificatori. Essendo questa la generalizzazione di una legge di De Morgan, è ragionevole aspettarsi che la sua dimostrazione abbia una struttura simile a quella di una legge di De Morgan vista in precedenza.

\paragraph{Dimostrazione di \(\forall x, P(x) \dashv \neg \exists x\st \neg P(x)\)}

Iniziamo con la prima direzione, ovvero \(\forall x, P(x) \dashv \neg \exists x\st \neg P(x)\).

Qui l'unica assunzione a disposizione è \(\neg \exists x\st \neg P(x)\), che essendo una negazione non è molto usabile se non avendo a disposizione la versione positiva.

Per dimostrare \(\forall x, P(x)\) possiamo iniziare con \(\forall I\), in modo da rimuovere il quantificatore e dover dimostrare solo \(P(y)\) per una variabile arbitraria \(y\).
Chiaramente a questo punto non c'è altra strada che quella di procedere per contraddizione, ovvero applicando \(\RAA\), che ci permette di scaricare \(\neg P(y)\).

\begin{prooftree}
  \AxiomC{\(\bot\)}
  \RuleLabel{\(\RAA\)}[\(\neg P(y)\)]
  \UnaryInfC{\(P(y)\)}
  \RuleLabel{\(\forall I\)}
  \UnaryInfC{\(\forall x, P(x)\)}
\end{prooftree}

A questo punto dobbiamo ottenere una contraddizione avendo a disposizione \(\neg \exists x\st \neg P(x)\) e \(\neg P(y)\). Possiamo subito notare però che \(\neg P(y)\) è proprio la formula che \(\exists x\st \neg P(x)\) quantifica, e che quindi possiamo ottenere la versione positiva di \(\neg \exists x\st \neg P(x)\) applicando \(\exists I\) a \(\neg P(y)\). A questo punto, con \(\neg E\) otteniamo la contraddizione che ci serve per chiudere la dimostrazione.

\begin{prooftree}
  \AxiomC{\({[\neg P(y)]}_1\)}
  \RuleLabel{\(\exists I\)}
  \UnaryInfC{\(\exists x\st \neg P(x)\)}
  \AxiomC{\(\neg \exists x\st \neg P(x)\)}
  \RuleLabel{\(\neg E\)}
  \BinaryInfC{\(\bot\)}
  \RuleLabel{\(\RAA_1\)}
  \UnaryInfC{\(P(y)\)}
  \RuleLabel{\(\forall I\)}
  \UnaryInfC{\(\forall x, P(x)\)}
\end{prooftree}

Veniamo dunque alla verifica dei vincoli di applicabilità. Per \(\exists I\) l'unico vincolo da verificare è che \(y\) sia libero per \(x\) in \(\neg P(x)\), e questo è verificato poiché \(y\) non occorre in \(\neg P(x)\). Per \(\forall I\) invece dobbiamo verificare che \(y\) non occorra in \(P(x)\), il che è banale, e che \(y\) non occorra libera in assunzioni non scaricata per \(P(y)\). Nel momento in cui si applica \(\forall I\) l'unica assunzione non scaricata è \(\neg \exists x\st \neg P(x)\), che non contiene \(y\), e quindi il vincolo è verificato.

L'analogo proposizionale di questa dimostrazione è la dimostrazione del sequente \(\neg(\neg A \lor \neg B) \vdash A \land B\). La dimostrazione è come segue:

\begin{prooftree}
  \AxiomC{\({[\neg A]}_1\)}
  \RuleLabel{\(\lor I\)}
  \UnaryInfC{\(\neg A \lor \neg B\)}
  \AxiomC{\(\neg(\neg A \lor \neg B) \)}
  \RuleLabel{\(\neg E\)}
  \BinaryInfC{\(\bot\)}
  \RuleLabel{\(\RAA_1\)}
  \UnaryInfC{\(A\)}
  
  \AxiomC{\({[\neg B]}_2\)}
  \RuleLabel{\(\lor I\)}
  \UnaryInfC{\(\neg A \lor \neg B\)}
  \AxiomC{\(\neg(\neg A \lor \neg B) \)}
  \RuleLabel{\(\neg E\)}
  \BinaryInfC{\(\bot\)}
  \RuleLabel{\(\RAA_2\)}
  \UnaryInfC{\(B\)}

  \RuleLabel{\(\land I\)}
  \BinaryInfC{\(A \land B\)}
\end{prooftree}

Possiamo notare che la dimostrazione è del tutta analoga a quella appena vista con i quantificatori, in cui le regole per \(\forall\) lasciano spazio a quelle per \(\land\) e quelle per \(\exists\) lasciano spazio a quelle per \(\lor\). L'unica differenza un po' più sostanziale è che, essendo \(\land\) e \(\lor\) operatori binari, è necessario lavorare su due rami di dimostrazione, mentre con i quantificatori è sufficiente lavorare su un solo ramo, essendo unari.

\paragraph{Dimostrazione di \(\forall x, P(x) \vdash \neg \exists x\st \neg P(x)\)}

Vediamo dunque l'altro verso dell'equivalenza, ovvero \(\forall x, P(x) \vdash \neg \exists x\st \neg P(x)\).

In questo caso dobbiamo dimostrare una negazione, e quindi è immediato pensare di procedere con \(\neg I\), la quale ci fornisce come assunzione scaricabile \(\exists x, \neg P(x)\), dalla quale potremmo poi, mediante \(\exists E\), scaricare anche \(\neg P(y)\).

\begin{prooftree}
  \AxiomC{\(\bot\)}
  \RuleLabel{\(\neg I\)}
  \UnaryInfC{\(\neg \exists x\st \neg P(x)\)}
\end{prooftree}

A questo punto, da \(\forall x, P(x)\) possiamo dedurre \(P(y)\) per \(\forall E\) e, mediante \(\neg E\) assumendo \(\neg P(y)\), ottenere la contraddizione cercata.

\begin{prooftree}
  \AxiomC{\(\neg P(y)\)}
  \AxiomC{\(\forall x, P(x)\)}
  \RuleLabel{\(\forall E\)}
  \UnaryInfC{\(P(y)\)}
  \RuleLabel{\(\neg E\)}
  \BinaryInfC{\(\bot\)}
  \RuleLabel{\(\neg I\)}[\(\exists x\st \neg P(x)\)]
  \UnaryInfC{\(\neg \exists x\st \neg P(x)\)}
\end{prooftree}

Come anticipato, per scaricare \(\neg P(y)\) andremo a usare \(\exists E\) con \(\exists x\st \neg P(x)\). È importante notare come l'unico punto adatto in cui è possibile applicare \(\exists E\) è subito prima della \(\neg I\), poiché applicarlo più in alto non ci permetterebbe di scaricare \(\exists x\st \neg P(x)\), e applicarlo più in bassi ci chiederebbe di dedurre tramite \(\exists E\) una formula che contiene \(y\) libera, e quindi violerebbe il vincolo di applicabilità di \(\exists E\). Di conseguenza, la deduzione completa è la seguente:

\begin{prooftree}
  \AxiomC{\({[\neg P(y)]}_2\)}
  \AxiomC{\(\forall x, P(x)\)}
  \RuleLabel{\(\forall E\)}
  \UnaryInfC{\(P(y)\)}
  \RuleLabel{\(\neg E\)}
  \BinaryInfC{\(\bot\)}
  \AxiomC{\({[\exists x\st \neg P(x)]}_1\)}
  \RuleLabel{\(\exists E_2\)}
  \BinaryInfC{\(\bot\)}
  \RuleLabel{\(\neg I_1\)}
  \UnaryInfC{\(\neg \exists x\st \neg P(x)\)}
\end{prooftree}

Analogamente a come visto prima, sostituendo le regole per i quantificatori con quelle per le congiunzioni e disgiunzioni, è possibile ottenere una dimostrazione del sequente \(A \land B \vdash \neg(\neg A \lor \neg B)\) del tutto analoga a quella appena vista.

\subsubsection*{Dimostrazione di \(\neg \exists x \st P(x)\vdash \exists x \st \neg P(x)\)} 

Passiamo ora a dimostrare \(\neg \exists x \st P(x)\vdash \exists x \st \neg P(x)\).\footnote{È importante notare che questa formula è vera vista la semantica dei modelli che abbiamo definito, ovvero di insiemi non vuoti. In questo caso infatti, se \(\neg \exists x \st P(x)\) è vera, essendo che almeno un elemento del esiste, l'unico motivo per cui \(\neg \exists x \st P(x)\) è vera è che \(P(x)\) è falsa per ogni elemento, e quindi ce n'è almeno uno per cui \(\neg P(x)\) è vera, e quindi \(\exists x \st \neg P(x)\) è vera. Se nella semantica ammettessimo anche insiemi vuoti, invece, questa formula non sarebbe vera, poiché se l'insieme è vuoto, \(\neg \exists x \st P(x)\) è vera, ma \(\exists x \st \neg P(x)\) è falsa.}

Partendo da una formula quantificata esistenzialmente, è ragionevole pensare di procedere con \(\exists I\).

\begin{prooftree}
  \AxiomC{\(\neg P(y)\)}
  \RuleLabel{\(\exists I\)}
  \UnaryInfC{\(\exists x \st \neg P(x)\)}
\end{prooftree}

Avendo ora da dimostrare una negazione che non è presente neanche indirettamente nelle assunzioni che abbiamo a disposizione, è ragionevole pensare di procedere con \(\neg I\), che ci fornisce come assunzione scaricabile \(P(y)\).

\begin{prooftree}
  \AxiomC{\(\bot\)}
  \RuleLabel{\(\neg I\)}[\(P(y)\)]
  \UnaryInfC{\(\neg P(y)\)}
  \RuleLabel{\(\exists I\)}
  \UnaryInfC{\(\exists x \st \neg P(x)\)}
\end{prooftree}

A questo punto è semplice ottenere la contraddizione, poiché ancora per \(\exists I\) possiamo ottenere \(\exists x \st P(x)\) da \(P(y)\), ovvero la forma positiva di \(\neg \exists x \st P(x)\), e quindi, con \(\neg E\) otteniamo la contraddizione che ci serve per chiudere la dimostrazione.

\begin{prooftree}
  \AxiomC{\({[P(y)]}_1\)}
  \RuleLabel{\(\exists I\)}
  \UnaryInfC{\(\exists x \st P(x)\)}
  \AxiomC{\(\neg \exists x \st P(x)\)}
  \RuleLabel{\(\neg E\)}
  \BinaryInfC{\(\bot\)}
  \RuleLabel{\(\neg I_1\)}
  \UnaryInfC{\(\neg P(y)\)}
  \RuleLabel{\(\exists I\)}
  \UnaryInfC{\(\exists x \st \neg P(x)\)}
\end{prooftree}

Qui verificare i vincoli di applicabilità è immediato avendo usato una variabile nuova.

Chiaramente la direzione opposta, ovvero \(\exists x \st \neg P(x) \vdash \neg \exists x \st P(x)\), è falsa, poiché un qualsiasi modello con un oggetto che soddisfa \(P(x)\) e un oggetto che soddisfa \(\neg P(x)\) è un controesempio per questa direzione. Se il calcolo fosse inconsistente, però, sarebbe possibile dimostrare anche questa direzione. Vediamo dunque se è possibile dimostrare \(\exists x \st \neg P(x) \vdash \neg \exists x \st P(x)\).

Chiaramente come primo passo è ragionevole pensare di procedere con \(\neg I\) per dimostrare \(\neg \exists x \st P(x)\), e quindi assumere \(\exists x \st P(x)\).

\begin{prooftree}
  \AxiomC{\(\bot\)}
  \RuleLabel{\(\neg I\)}[\(\exists x \st P(x)\)]
  \UnaryInfC{\(\neg \exists x \st P(x)\)}
\end{prooftree}

A questo punto, l'unico modo per ottenere una contraddizione è mediante \(\neg E\) con \(P(y)\) e \(\neg P(y)\). Fortunatamente però tali assunzioni non sono scaricabili mediante \(\exists E\) dalle assunzioni che abbiamo, poiché nell'applicare \(\exists E\) dovremmo dedurre \(\psi = P(y)\) e \(\psi = \neg P(y)\), e in entrambi i casi \(y\) occorre libera in \(\psi\), violando così il vincolo di applicabilità di \(\exists E\). Di conseguenza, non è possibile ottenere la contraddizione che ci serve per chiudere la dimostrazione, e quindi non è possibile dimostrare \(\exists x \st \neg P(x) \vdash \neg \exists x \st P(x)\).

\subsubsection*{Dimostrazione di \(\vdash \exists x \st (C(x)\to \forall y, C(y))\)}

Vediamo ora una dimostrazione di una formula che a primo impatto può non sembrare vera, ovvero \(\vdash \exists x \st (C(x)\to \forall y, C(y))\). Possiamo notare però che questa formula è equivalente a \(\exists x\st (\neg C(x) \lor \forall y, C(y))\). A questo punto è chiaro che, preso un qualsiasi modello, o questo contiene un oggetto che non soddisfa \(C(x)\), e quindi \(\neg C(x)\) è vera per almeno un oggetto, e quindi \(\exists x\st \neg C(x)\) è vera, oppure \(C(x)\) è vera per ogni oggetto, e quindi \(\forall y, C(y)\) è vera, e quindi \(\exists x\st \forall y, C(y)\) è vera. In entrambi i casi, dunque, \(\exists x\st (\neg C(x) \lor \forall y, C(y))\) è vera.

Cerchiamo dunque di dimostrare questo sequente.
Partendo da una formula quantificata esistenzialmente, è ragionevole pensare di procedere con \(\exists I\), seguendo con \(\to I\) per dimostrare \(C(z)\to \forall y, C(y)\).

\begin{prooftree}
  \AxiomC{\(\forall y, C(y)\)}
  \RuleLabel{\(\to I\)}[\(C(z)\)]
  \UnaryInfC{\(C(z)\to \forall y, C(y)\)}
  \RuleLabel{\(\exists I\)}
  \UnaryInfC{\(\exists x \st (C(x)\to \forall y, C(y))\)}
\end{prooftree}

A questo punto, per dimostrare \(\forall y, C(y)\) si potrebbe procedere sia per contraddizione che per \(\forall I\). La seconda opzione può sembrare più diretta, ma purtroppo è errata. Per farla infatti si potrebbe pensare di dimostrare \(\forall y, C(y)\) con la seguente deduzione:

\begin{prooftree}
  \AxiomC{\({[C(z)]}_1\)}
  \RuleLabel{\(\forall I\)}
  \UnaryInfC{\(\forall y, C(y)\)}
  \RuleLabel{\(\to I_1\)}
  \UnaryInfC{\(C(z)\to \forall y, C(y)\)}
  \RuleLabel{\(\exists I\)}
  \UnaryInfC{\(\exists x \st (C(x)\to \forall y, C(y))\)}
\end{prooftree}

Purtroppo però, nel momento in cui si applica \(\forall I\) l'assunzione \(C(z)\) non è ancora scaricata, e quindi \(z\) occorre libera in un'assunzione non scaricata, violando così il vincolo di applicabilità di \(\forall I\). Di conseguenza, questa deduzione è errata, e dobbiamo necessariamente procedere per contraddizione per \(\bot E\).

\begin{prooftree}
  \AxiomC{\(\bot\)}
  \RuleLabel{\(\bot E\)}
  \UnaryInfC{\(\forall y, C(y)\)}
  \RuleLabel{\(\to I\)}[\(C(z)\)]
  \UnaryInfC{\(C(z)\to \forall y, C(y)\)}
  \RuleLabel{\(\exists I\)}
  \UnaryInfC{\(\exists x \st (C(x)\to \forall y, C(y))\)}
\end{prooftree}

Avendo a disposizione \(C(z)\) scaricabile possiamo pensare di ottenere la contraddizione che ci serve tramite \(\neg E\) con \(\neg C(z)\).

\begin{prooftree}
  \AxiomC{\({[C(z)]}_1\)}
  \AxiomC{\(\neg C(z)\)}
  \RuleLabel{\(\neg E\)}
  \BinaryInfC{\(\bot\)}
  \RuleLabel{\(\bot E\)}
  \UnaryInfC{\(\forall y, C(y)\)}
  \RuleLabel{\(\to I_1\)}
  \UnaryInfC{\(C(z)\to \forall y, C(y)\)}
  \RuleLabel{\(\exists I\)}
  \UnaryInfC{\(\exists x \st (C(x)\to \forall y, C(y))\)}
\end{prooftree}

Dobbiamo quindi capire come scaricare \(\neg C(z)\). Si può pensare di scaricarlo mediante \(\RAA\) ottenendo \(C(z)\). Per fare ciò però necessaria una nuova contraddizione, questa volta con \(\exists x \st (C(x) \to \forall y, C(y))\), e di conseguenza assumendo \(\neg \exists x \st (C(x) \to \forall y, C(y))\).

\begin{prooftree}
  \AxiomC{\({[C(z)]}_1\)}
  \AxiomC{\({[\neg C(z)]}_2\)}
  \RuleLabel{\(\neg E\)}
  \BinaryInfC{\(\bot\)}
  \RuleLabel{\(\bot E\)}
  \UnaryInfC{\(\forall y, C(y)\)}
  \RuleLabel{\(\to I_1\)}
  \UnaryInfC{\(C(z)\to \forall y, C(y)\)}
  \RuleLabel{\(\exists I\)}
  \UnaryInfC{\(\exists x \st (C(x)\to \forall y, C(y))\)}
  \AxiomC{\(\neg \exists x \st (C(x) \to \forall y, C(y))\)}
  \RuleLabel{\(\neg E\)}
  \BinaryInfC{\(\bot\)}
  \RuleLabel{\(\RAA_2\)}
  \UnaryInfC{\(C(z)\)}
\end{prooftree}

Possiamo notare che ora, essendo che \(C(z)\) non è più assunzione, bensì è stata dimostrata, possiamo riprender l'idea di dimostrare \(\forall y, C(y)\) con \(\forall I\) senza violare il vincolo di applicabilità.

\begin{prooftree}
  \AxiomC{\({[C(z)]}_1\)}
  \AxiomC{\({[\neg C(z)]}_2\)}
  \RuleLabel{\(\neg E\)}
  \BinaryInfC{\(\bot\)}
  \RuleLabel{\(\bot E\)}
  \UnaryInfC{\(\forall y, C(y)\)}
  \RuleLabel{\(\to I_1\)}
  \UnaryInfC{\(C(z)\to \forall y, C(y)\)}
  \RuleLabel{\(\exists I\)}
  \UnaryInfC{\(\exists x \st (C(x)\to \forall y, C(y))\)}
  \AxiomC{\(\neg \exists x \st (C(x) \to \forall y, C(y))\)}
  \RuleLabel{\(\neg E\)}
  \BinaryInfC{\(\bot\)}
  \RuleLabel{\(\RAA_2\)}
  \UnaryInfC{\(C(z)\)}
  \RuleLabel{\(\forall I\)}
  \UnaryInfC{\(\forall y, C(y)\)}
  \RuleLabel{\(\to I\)}[\(C(z)\)]
  \UnaryInfC{\(C(z)\to \forall y, C(y)\)}
  \RuleLabel{\(\exists I\)}
  \UnaryInfC{\(\exists x \st (C(x)\to \forall y, C(y))\)}
\end{prooftree}

Ci manca solo da scaricare \(\neg \exists x \st (C(x) \to \forall y, C(y))\). Per fare ciò però è sufficiente riutilizzarla con un ultima applicazione di \(\RAA\).

\begin{prooftree}
  \small
  \AxiomC{\({[C(z)]}_1\)}
  \AxiomC{\({[\neg C(z)]}_2\)}
  \RuleLabel{\(\neg E\)}
  \BinaryInfC{\(\bot\)}
  \RuleLabel{\(\bot E\)}
  \UnaryInfC{\(\forall y, C(y)\)}
  \RuleLabel{\(\to I_1\)}
  \UnaryInfC{\(C(z)\to \forall y, C(y)\)}
  \RuleLabel{\(\exists I\)}
  \UnaryInfC{\(\exists x \st (C(x)\to \forall y, C(y))\)}
  \AxiomC{\({[\neg \exists x \st (C(x) \to \forall y, C(y))]}_3\)}
  \RuleLabel{\(\neg E\)}
  \BinaryInfC{\(\bot\)}
  \RuleLabel{\(\RAA_2\)}
  \UnaryInfC{\(C(z)\)}
  \RuleLabel{\(\forall I\)}
  \UnaryInfC{\(\forall y, C(y)\)}
  \RuleLabel{\(\to I\)}
  \UnaryInfC{\(C(z)\to \forall y, C(y)\)}
  \RuleLabel{\(\exists I\)}
  \UnaryInfC{\(\exists x \st (C(x)\to \forall y, C(y))\)}
  \AxiomC{\({[\neg \exists x \st (C(x) \to \forall y, C(y))]}_3\)}
  \RuleLabel{\(\neg E\)}
  \BinaryInfC{\(\bot\)}
  \RuleLabel{\(\RAA_3\)}
  \UnaryInfC{\(\exists x \st (C(x)\to \forall y, C(y))\)}
\end{prooftree}

È importante ricordare, come abbiamo già visto per la logica proposizionale, che il fatto che una regola permetta di scaricare un assunzione non ci obbliga a usarla per scaricarla.

Analizziamo ora invece una differenza importante tra le deduzioni in logica proposizionale e in logica del prim'ordine, dovuta dai vincoli di applicabilità delle regole per i quantificatori. Consideriamo di avere una deduzione \(\pi_1\) per una formula \(\psi\) a partire da un assunzione \(P\), e un altra deduzione indipendente \(\pi_2\) che dimostra sempre \(\psi\) a partire da un assunzione \( \neg P\). In logica proposizionale, è sempre possibile combinare queste due deduzioni in una deduzione di \(\psi\) senza assunzioni, semplicemente applicando \(\lor E\) a \(P \lor \neg P\), che abbiamo dimostrato in precedenza (\ref{subsubsec:excluded_middle_dem}).

\begin{prooftree}
  \AxiomC{\({[P]}_1\)}
  \UnaryInfC{\(\vdots\)}
  \RuleLabel{\(\pi_1\)}
  \UnaryInfC{\(\psi\)}
  
  \AxiomC{\({[\neg P]}_2\)}
  \UnaryInfC{\(\vdots\)}
  \RuleLabel{\(\pi_2\)}
  \UnaryInfC{\(\psi\)}

  \AxiomC{\ref{subsubsec:excluded_middle_dem}}

  \UnaryInfC{\(P \lor \neg P\)}
  \RuleLabel{\(\lor E_{1,2}\)}
  \TrinaryInfC{\(\psi\)}
\end{prooftree}

Avendo sempre a disposizione \(\pi_1\) e \(\pi_2\), è possibile ottenere una dimostrazione di \(\psi\) senza assunzioni anche ragionando per contraddizioni come segue:

\begin{prooftree}
  \AxiomC{\({[P]}_1\)}
  \UnaryInfC{\(\vdots\)}
  \RuleLabel{\(\pi_1\)}
  \UnaryInfC{\(\psi\)}
  \AxiomC{\({[\neg \psi]}_2\)}
  \RuleLabel{\(\neg E\)}
  \BinaryInfC{\(\bot\)}
  \RuleLabel{\(\neg I_1\)}
  \UnaryInfC{\(\neg P\)}
  \UnaryInfC{\(\vdots\)}
  \RuleLabel{\(\pi_2\)}
  \UnaryInfC{\(\psi\)}

  \AxiomC{\({[\neg \psi]}_2\)}
  \RuleLabel{\(\neg E\)}
  \BinaryInfC{\(\bot\)}
  \RuleLabel{\(\RAA_{2}\)}
  \UnaryInfC{\(\psi\)}
\end{prooftree}

Un esempio di questo schema è la dimostrazione di \(P\to Q \vdash \neg P \lor Q\) vista in precedenza (\ref{subsubsec:material_implication_dem}), anche se in forma diversa.

In particolare, il primo schema applicato a questo caso è:

\begin{prooftree}
  \AxiomC{\({[P]}_1\)}
  \AxiomC{\(P\to Q\)}
  \RuleLabel{\(\to E\)}
  \BinaryInfC{\(Q\)}
  \RuleLabel{\(\lor I\)}
  \UnaryInfC{\(\neg P \lor Q\)}

  \AxiomC{\({[\neg P]}_2\)}
  \RuleLabel{\(\lor I\)}
  \UnaryInfC{\(\neg P \lor Q\)}

  \AxiomC{\ref{subsubsec:excluded_middle_dem}}
  \UnaryInfC{\(P \lor \neg P\)}
  \RuleLabel{\(\lor E_{1,2}\)}
  \TrinaryInfC{\(\neg P \lor Q\)}
\end{prooftree}

Mentre il secondo schema è:

\begin{prooftree}
  \AxiomC{\({[P]}_1\)}
  \AxiomC{\(P\to Q\)}
  \RuleLabel{\(\to E\)}
  \BinaryInfC{\(Q\)}
  \RuleLabel{\(\lor I\)}
  \UnaryInfC{\(\neg P \lor Q\)}
  \AxiomC{\({[\neg (\neg P \lor Q)]}_2\)}
  \RuleLabel{\(\neg E\)}
  \BinaryInfC{\(\bot\)}
  \RuleLabel{\(\neg I_1\)}
  \UnaryInfC{\(\neg P\)}
  \RuleLabel{\(\lor I\)}
  \UnaryInfC{\(\neg P \lor Q\)}
  \AxiomC{\({[\neg (\neg P \lor Q)]}_2\)}
  \RuleLabel{\(\neg E\)}
  \BinaryInfC{\(\bot\)}
  \RuleLabel{\(\RAA_2\)}
  \UnaryInfC{\(\neg P \lor Q\)}
\end{prooftree}

Passando al prim'ordine questa equivalenza di dimostrazioni purtroppo si perde. In particolare, se abbiamo una deduzione della prima forma, essendo che \(\pi_1\) e \(\pi_2\) sono deduzioni indipendenti, è sempre possibile combinare le due deduzioni in una deduzione della seconda forma mentre se abbiamo una deduzione della seconda forma, non è sempre possibile combinare le due deduzioni in una deduzione della prima forma, poiché, scomponendo l'unico ramo di deduzione in due rami indipendenti si rischia di violare dei vincoli di applicabilità.

Nell'ultima dimostrazione vista, ovvero \(\vdash \exists x \st (C(x)\to \forall y, C(y))\), ci troviamo proprio in questo secondo schema, con \(\pi_1\) che corrisponde a:

\begin{prooftree}
  \AxiomC{\({[C(z)]}_1\)}
  \AxiomC{\(\neg C(z)\)}
  \RuleLabel{\(\neg E\)}
  \BinaryInfC{\(\bot\)}
  \RuleLabel{\(\bot E\)}
  \UnaryInfC{\(\forall y, C(y)\)}
  \RuleLabel{\(\to I_1\)}
  \UnaryInfC{\(C(z)\to \forall y, C(y)\)}
  \RuleLabel{\(\exists I\)}
  \UnaryInfC{\(\exists x \st (C(x)\to \forall y, C(y))\)}
\end{prooftree}
dove \(P = \neg C(z)\) e \(\psi = \exists x \st (C(x)\to \forall y, C(y))\), e \(\pi_2\) che corrisponde a:

\begin{prooftree}
  \AxiomC{\(C(z)\)}
  \RuleLabel{\(\forall I\)}
  \UnaryInfC{\(\forall y, C(y)\)}
  \RuleLabel{\(\to I\)}
  \UnaryInfC{\(C(z)\to \forall y, C(y)\)}
  \RuleLabel{\(\exists I\)}
  \UnaryInfC{\(\exists x \st (C(x)\to \forall y, C(y))\)}
\end{prooftree}
con \(\neg P = C(z)\).

In questo caso, se proviamo ad applicare il primo schema otterremo la seguente deduzione:

\begin{prooftree}
  \AxiomC{\({[C(z)]}_1\)}
  \AxiomC{\({[\neg C(z)]}_3\)}
  \RuleLabel{\(\neg E\)}
  \BinaryInfC{\(\bot\)}
  \RuleLabel{\(\bot E\)}
  \UnaryInfC{\(\forall y, C(y)\)}
  \RuleLabel{\(\to I_1\)}
  \UnaryInfC{\(C(z)\to \forall y, C(y)\)}
  \RuleLabel{\(\exists I\)}
  \UnaryInfC{\(\exists x \st (C(x)\to \forall y, C(y))\)}

  \AxiomC{\({[C(z)]}_2\)}
  \RuleLabel{\(\forall I\)}
  \UnaryInfC{\(\forall y, C(y)\)}
  \RuleLabel{\(\to I\)}
  \UnaryInfC{\(C(z)\to \forall y, C(y)\)}
  \RuleLabel{\(\exists I\)}
  \UnaryInfC{\(\exists x \st (C(x)\to \forall y, C(y))\)}

  \AxiomC{\ref{subsubsec:excluded_middle_dem}}
  \UnaryInfC{\(C(z) \lor \neg C(z)\)}
  \RuleLabel{\(\lor E_{2,3}\)}
  \TrinaryInfC{\(\exists x \st (C(x)\to \forall y, C(y))\)}
\end{prooftree}

Come abbiamo già osservato però \(\pi_2\) è legittima nel secondo schema poiché \(C(z)\) non è un assunzione, ma una formula dimostrata, mentre non è legittima nel primo schema poiché \(C(z)\) è un assunzione, e quindi \(z\) occorre libera in un'assunzione non scaricata, violando così il vincolo di applicabilità di \(\forall I\).

\subsection{Regole di deduzione per la logica del prim'ordine con uguaglianza}

Come già accennato numerose volte, spesso \(=\) è considerato un simbolo logico, e quindi come tale sarà necessario introdurre delle regole di deduzione per poterlo usare nelle dimostrazioni. L'alternativa sarebbe quella di considerare \(=\) come un simbolo non logico, assiomatizzando il suo comportamento.

Formalmente, va arricchito il linguaggio con un concetto di \keyword{teoria}, ovvero un insieme di formule vere in ogni interpretazione. Tale teoria è rappresentabile se esiste un insieme di formule dette \keyword{assiomi} assunte come vere, da cui ogni altra formula della teoria è dimostrabile. Nonostante abbiamo sempre considerato l'insieme di assiomi \(\Axi\) vuoto, nella \Namedcref{def:deducibility} abbiamo già definito il concetto di deducibilità a partire da un insieme di assiomi, e quindi non è necessario modificare la definizione di deducibilità per poter considerare teorie con assiomi.

In particolare, per la deduzione naturale, la differenza principale sarà nel poter lasciare come assunzioni non scaricate oltre alle dell'insieme delle premesse \(\Gamma\) anche quelle presenti in \(\Axi\).

Tratteremo più nel dettaglio successivamente il concetto di teoria, ma possiamo iniziare subito con un esempio di teoria semplice, ovvero la teoria dell'uguaglianza.

In particolare, l'insieme degli assiomi necessari per la teoria dell'uguaglianza è il seguente:

\begin{enumerate}
  \item \(\forall x, x=x\)
  \item \(\forall x_1,\ldots, \forall x_n, \forall y_1, \ldots, \forall y_n, (\bigwedge_{i=1}^n x_i = y_i) \to (f(x_1, \ldots, x_n) = f(y_1, \ldots, y_n))\) per ogni \(f \in \FuncSyms\) con \(\arty(f) = n\).
  \item \(\forall x_1,\ldots, \forall x_n, \forall y_1, \ldots, \forall y_n, (\bigwedge_{i=1}^n x_i = y_i) \to (P(x_1, \ldots, x_n) \to P(y_1, \ldots, y_n))\) per ogni \(P \in \PredSyms\) con \(\arty(P) = n\).
\end{enumerate}

Tale insieme è sufficiente per garantire che \(=\) si comporti come una relazione di equivalenza, e che sia rispettata la sostituzione di termini uguali in ogni formula.

In particolare ci si aspetterebbe che valgano simmetria, transitività e che l'ultimo assioma presenti una doppia implicazione. Andiamo a vedere uno alla volta come queste tre proprietà siano dimostrabili.

\paragraph{Simmetria}

Per dimostrare la simmetria vorremmo mostrare che vale \(\vdash \forall x, \forall y, x=y \to y=x\). È importante notare che, avendo assiomatizzato l'uguaglianza, si ha che \(=\in\PredSyms\) e di conseguenza il terzo schema di assiomi è applicabile anche a \(P\eqdef =\) e quindi in particolare \(P(x_1,x_2)\eqdef x_1=x_2\).

Allora dal terzo assioma istanziato con \(x_1\eqdef x\), \(y_1\eqdef y\), \(x_2 \eqdef x\) \(y_2 \eqdef x\) abbiamo:
\todo{Fai controllare al prof}
\begin{prooftree}
  \tiny
  \AxiomC{\(\forall x_1, x_2, y_1, y_2, (x_1 = y_1 \land x_2 = y_2) \to (P(x_1, x_2) \to P(y_1, y_2))\)}
  \RuleLabel{\(\forall E^*\)}
  \UnaryInfC{\(z=w \land z = z \to (P(z,z) \to P(w,z))\)}
  \RuleLabel{\(P\eqdef =\)}
  \UnaryInfC{\(z=w \land z = z \to (z=z \to w=z)\)}

  \AxiomC{\(\forall x, x=x\)}
  \RuleLabel{\(\forall E\)}
  \UnaryInfC{\(z=z\)}
  \AxiomC{\({[z=w]}_1\)}
  \RuleLabel{\(\land I\)}
  \BinaryInfC{\(z=w \land z=z\)}
  \RuleLabel{\(\to E\)}
  \BinaryInfC{\(z=z \to w=z\)}
  \AxiomC{\(\forall x, x=x\)}
  \RuleLabel{\(\forall E\)}
  \UnaryInfC{\(z=z\)}
  \RuleLabel{\(\to E\)}
  \BinaryInfC{\(w=z\)}
  \RuleLabel{\(\to I_1\)}
  \UnaryInfC{\(z=w \to w=z\)}
  \RuleLabel{\(\forall I\)}
  \UnaryInfC{\(\forall y, z=y \to y=z\)}
  \RuleLabel{\(\forall I\)}
  \UnaryInfC{\(\forall x, \forall y, x=y \to y=x\)}
\end{prooftree}

Dove con \(\forall E^*\) indichiamo l'applicazione di \(\forall E\) ripetuta più volte, e con \(P\eqdef =\) indichiamo l'applicazione dell'aver definito \(P\) come \(=\).

\paragraph{Transitività}

Una volta dimostrata la simmetria, è immediato dimostrare la transitività, ovvero \(\vdash \forall x, \forall y, \forall z, (x=y \land y=z) \to x=z\).

\begin{prooftree}
  \tiny
  \AxiomC{\(\forall x_1, x_2, y_1, y_2, (x_1 = y_1 \land x_2 = y_2) \to (P(x_1, x_2) \to P(y_1, y_2))\)}
  \RuleLabel{\(\forall E^*\)}
  \UnaryInfC{\(t=s \land t=r \to (P(t,t) \to P(s,r))\)}
  \RuleLabel{\(P\eqdef =\)}
  \UnaryInfC{\(t=s \land t=r \to (t=t \to s=r)\)}

  \AxiomC{\({[s=t \land t=r]}_1\)}
  \RuleLabel{\(\land E^r\)}
  \UnaryInfC{\(t=r\)}

  \AxiomC{\({[s=t \land t=r]}_1\)}
  \RuleLabel{\(\land E^l\)}
  \UnaryInfC{\(s=t\)}
  \RuleLabel{\(Symm\)}
  \UnaryInfC{\(t=s\)}
  \RuleLabel{\(\land I\)}
  \BinaryInfC{\(t=s \land t=r\)}
  \RuleLabel{\(\to E\)}
  \BinaryInfC{\(t=t \to s=r\)}

  \AxiomC{\(\forall x, x=x\)}
  \RuleLabel{\(\forall E\)}
  \UnaryInfC{\(t=t\)}
  \RuleLabel{\(\to E\)}
  \BinaryInfC{\(s=r\)}

  \RuleLabel{\(\to I_1\)}
  \UnaryInfC{\((s=t \land t=r) \to (s=r)\)}
  \RuleLabel{\(\forall I\)}
  \UnaryInfC{\(\forall z, (s=t \land t=z) \to (s=z)\)}
  \RuleLabel{\(\forall I\)}
  \UnaryInfC{\(\forall y, \forall z, (s=y \land y=z) \to (s=z)\)}
  \RuleLabel{\(\forall I\)}
  \UnaryInfC{\(\forall x, \forall y, \forall z, (x=y \land y=z) \to x=z\)}
\end{prooftree}

Infine, possiamo mostrare l'altra direzione dell'implicazione nel terzo schema di assioma, ovvero \(\forall x_1,\ldots, \forall x_n, \forall y_1, \ldots, \forall y_n, (P(x_1, \ldots, x_n) \to P(y_1, \ldots, y_n))\to (\bigwedge_{i=1}^n x_i = y_i)\).

Ancora una volta mostreremo il caso più semplice, ovvero con \(n=1\), che risulta essere
\[\forall x,\forall y, (x = y) \to (P(y) \to P(x))\]

\begin{prooftree}
  \AxiomC{\(\forall x_1, y_1, (x_1 = y_1) \to (P(x_1) \to P(y_1))\)}
  \RuleLabel{\(\forall E^*\)}
  \UnaryInfC{\(s=t \to (P(s) \to P(t))\)}

  \AxiomC{\({[t=s]}_1\)}
  \RuleLabel{\(Symm\)}
  \UnaryInfC{\(s=t\)}
  \RuleLabel{\(\to E\)}
  \BinaryInfC{\(P(s) \to P(t)\)}
  \RuleLabel{\(\to I_1\)}
  \UnaryInfC{\(t=s \to (P(s) \to P(t))\)}
  \RuleLabel{\(\forall I\)}
  \UnaryInfC{\(\forall y, t=y \to (P(s) \to P(y))\)}
  \RuleLabel{\(\forall I\)}
  \UnaryInfC{\(\forall x, \forall y, x=y \to (P(x) \to P(y))\)}
\end{prooftree}

Passando quindi all'utilizzo di \(=\) come simbolo logico, è chiaro per la natura del calcolo di deduzione naturale vorremmo avere una regola di deduzione per ogni assioma descritto precedentemente, dovendo le regole catturare la semantica dei simboli a cui si applicano. Di conseguenza, con \(t_i, s_i \in \Terms\) le regole di inferenza per \(=\) sono le seguenti:

\[\infer[Refl]{t=t}{}\]
\[\infer[Sub_f]{\subs{f(x_1, \ldots, x_n)}{x_1,\ldots,x_n}{t_1,\ldots, t_n} = \subs{f(x_1, \ldots, x_n)}{x_1,\ldots,x_n}{s_1,\ldots,s_n}}{t_1 = s_1 & \ldots & t_n = s_n}\]
\[\infer[Sub_P]{\subs{P(x_1, \ldots, x_n)}{x_1,\ldots,x_n}{t_1,\ldots,t_n}\to\subs{P(x_1, \ldots, x_n)}{x_1,\ldots,x_n}{s_1,\ldots, s_n}}{t_1 = s_1 & \ldots & t_n = s_n}\]

È facile vedere come, applicando queste regole di deduzione dove sono stati usati gli assiomi nelle dimostrazioni precedente, è possibile dimostrare simmetria, transitività e direzione inversa della terza regola. Chiaramente, tali dimostrazioni sono decisamente più succinte non dovendo passare dall'eliminazione dei quantificatori dagli assiomi. Avremo infatti per \(x=y\vdash y = x\):

\begin{prooftree}
  \AxiomC{\(x=y\)}
  \AxiomC{}
  \RuleLabel{\(Refl\)}
  \UnaryInfC{\(x=x\)}
  \RuleLabel{\(Sub_P\)}
  \BinaryInfC{\(P(x,x) \to P(y,x)\)}
  \RuleLabel{\(P\eqdef =\)}
  \UnaryInfC{\(x=x \to y=x\)}
  \AxiomC{}
  \RuleLabel{\(Refl\)}
  \UnaryInfC{\(x=x\)}
  \RuleLabel{\(\to E\)}
  \BinaryInfC{\(y=x\)}
\end{prooftree}

Per \(x=y, y=z \vdash x=z\):

\begin{prooftree}
  \AxiomC{\(x=y\)}
  \RuleLabel{\(Symm\)}
  \UnaryInfC{\(y=x\)}
  \AxiomC{\(y=z\)}
  \RuleLabel{\(Sub_P\)}
  \BinaryInfC{\(P(y,y) \to P(x,z)\)}
  \RuleLabel{\(P\eqdef =\)}
  \UnaryInfC{\(y=y \to x=z\)}
  \AxiomC{}
  \RuleLabel{\(Refl\)}
  \UnaryInfC{\(y=y\)}
  \RuleLabel{\(\to E\)}
  \BinaryInfC{\(x=z\)}
\end{prooftree}

Infine per \(\forall x, \forall y, (x=y) \to (P(y) \to P(x))\):

\begin{prooftree}
  \AxiomC{\({[w=z]}_1\)}
  \RuleLabel{\(Symm\)}
  \UnaryInfC{\(w=z\)}
  \RuleLabel{\(Sub_P\)}
  \UnaryInfC{\(P(w) \to P(z)\)}
  \RuleLabel{\(\to I_1\)}
  \UnaryInfC{\(w=z \to (P(w) \to P(z))\)}
  \RuleLabel{\(\forall I\)}
  \UnaryInfC{\(\forall y, w=y \to (P(y) \to P(w))\)}
  \RuleLabel{\(\forall I\)}
  \UnaryInfC{\(\forall x, \forall y, x=y \to (P(x) \to P(y))\)}
\end{prooftree}

Inoltre, le regole di sostituzione \(Sub_f\) e \(Sub_P\) possono essere generalizzate rispettivamente a termini e formule arbitrari, ottenendo così le seguenti regole:

\[\infer[Sub_t]{\subs{t}{x_1,\ldots,x_n}{t_1,\ldots,t_n} = \subs{t}{x_1,\ldots,x_n}{s_1,\ldots,s_n}}{t_1 = s_1 & \ldots & t_n = s_n}\]
\[\infer[Sub_\phi]{\subs{\phi}{x_1,\ldots,x_n}{t_1,\ldots,t_n}\to\subs{\phi}{x_1,\ldots,x_n}{s_1,\ldots, s_n}}{t_1 = s_1 & \ldots & t_n = s_n}\]

È possibile dimostrare che queste regole equivalgono all'applicazione iterata di \(Sub_f\) e \(Sub_P\) rispettivamente, e che quindi la loro validità deriva direttamente da quella di \(Sub_f\) e \(Sub_P\). In particolare possiamo dimostrare la validità di tali regole mediante induzione strutturale.

\begin{proposition}{Validità di \(Sub_t\)}
  Sia \(\Gamma = \set{r_1 = s_1, \ldots, r_n = s_n}\), con \(r_1, \ldots, r_n, s_1,\ldots s_n \in \Terms\). Dato \(t \in \Terms\) si ha che
  \[\Gamma \vdash \subs{t}{x_1,\ldots,x_n}{s_1,\ldots,s_n} = \subs{t}{x_1,\ldots,x_n}{r_1,\ldots,r_n}\]
\end{proposition}

\begin{proof}
  Procediamo per induzione strutturale su \(t\).

  I casi base sono \(t = x\in \VarSyms\) e \(t = c \in \ConstSyms\)
  Per il primo caso, iniziamo considerando \(t = x\) con \(x \in \set{x_1,\ldots,x_n}\). Senza perdita di generalità sia \(x = x_1\). In questo caso, abbiamo che \(\subs{x}{x_1,\ldots,x_n}{r_1,\ldots,r_n} = \subs{x}{x_1,\ldots,x_n}{s_1,\ldots,s_n}\equiv r_1 = s_1\), che è una delle assunzioni.

  Se invece \(t = y\) con \(y\notin\set{x_1,\ldots, x_n}\)  o \(t = c \in \ConstSyms\) si ha che
  \[\subs{y}{x_1,\ldots,x_n}{r_1,\ldots,r_n} = \subs{y}{x_1,\ldots,x_n}{s_1,\ldots,s_n}\equiv y=y\]
  e che
  \[\subs{c}{x_1,\ldots,x_n}{r_1,\ldots,r_n} = \subs{c}{x_1,\ldots,x_n}{s_1,\ldots,s_n}\equiv c=c\]
  che sono entrambe dimostrabili con \(Refl\).

  Venendo infine al caso induttivo, ovvero \(t=f(t_1,\ldots, t_k)\) con \(f \in \FuncSyms\) di \(\arty(f)=k\) e \(t_1, \ldots, t_k \in \Terms\).
  Per ipotesi induttiva sappiamo che per ogni sottotermine \(t_i\) con \(i \in \set{1,\ldots,k}\) vale che \(\Gamma \vdash \subs{t_i}{x_1,\ldots,x_n}{s_1,\ldots,s_n} = \subs{t_i}{x_1,\ldots,x_n}{r_1,\ldots,r_n}\).

  Per definizione di applicazione di sostituzione a un termine funzionale, si ha che
  \[\subs{f(t_1,\ldots,t_k)}{x_1,\ldots,x_n}{s_1,\ldots,s_n} \eqdef f(\subs{t_1}{x_1,\ldots,x_n}{s_1,\ldots,s_n},\ldots,\subs{t_k}{x_1,\ldots,x_n}{s_1,\ldots,s_n}) \eqdef \subs{f(y_1,\ldots,y_k)}{y_1,\ldots,y_k}{\subs{t_1}{x_1,\ldots,x_n}{s_1,\ldots,s_n},\ldots,\subs{t_k}{x_1,\ldots,x_n}{s_1,\ldots,s_n}}\]

  A questo punto è possibile applicare \(Sub_f\) come segue:

  \begin{prooftree}
    \AxiomC{\(\Gamma\)}
    \UnaryInfC{\(\vdots\)}
    \RuleLabel{IH}
    \UnaryInfC{\(\subs{t_1}{x_1,\ldots,x_n}{s_1,\ldots,s_n} = \subs{t_1}{x_1,\ldots,x_n}{r_1,\ldots,r_n}\)}

    \AxiomC{\(\ldots\)}

    \AxiomC{\(\Gamma\)}
    \UnaryInfC{\(\vdots\)}
    \RuleLabel{IH}
    \UnaryInfC{\(\subs{t_k}{x_1,\ldots,x_n}{s_1,\ldots,s_n} = \subs{t_k}{x_1,\ldots,x_n}{r_1,\ldots,r_n}\)}

    \RuleLabel{\(Sub_f\)}
    \TrinaryInfC{\(\subs{f(y_1,\ldots,y_k)}{y_1,\ldots,y_k}{\subs{t_1}{x_1,\ldots,x_n}{s_1,\ldots,s_n},\ldots,\subs{t_k}{x_1,\ldots,x_n}{s_1,\ldots,s_n}} = \subs{f(y_1,\ldots,y_k)}{y_1,\ldots,y_k}{\subs{t_1}{x_1,\ldots,x_n}{r_1,\ldots,r_n},\ldots,\subs{t_k}{x_1,\ldots,x_n}{r_1,\ldots,r_n}}\)}
  \end{prooftree}
\end{proof}

\begin{proposition}{Validità di \(Sub_\phi\)}
  Sia \(\Gamma = \set{r_1 = s_1, \ldots, r_n = s_n}\), con \(r_1, \ldots, r_n, s_1,\ldots s_n \in \Terms\). Dato \(\phi \in \FOL\) si ha che
  \[\Gamma \vdash \subs{\phi}{x_1,\ldots,x_n}{s_1,\ldots,s_n} \to \subs{\phi}{x_1,\ldots,x_n}{r_1,\ldots,r_n}\]
\end{proposition}

\begin{proof}
  Procediamo per induzione strutturale su \(\phi\).

  Per il caso base \(\phi = P(t_1,\ldots,t_k)\) con \(P \in \PredSyms\) di \(\arty(P) = k\) e \(t_1, \ldots, t_k \in \Terms\). In questo caso, analogamente al caso induttivo di \(Sub_P\) è sufficiente usare \(Sub_P\) per ogni \(t_i\) con \(i \in \set{1,\ldots,k}\) e infine applicare \(Sub_P\).

  I casi induttivi saranno uno per ogni simbolo logico. La struttura di dimostrazione resta sempre bene o male la stessa, di conseguenza ci concentreremo solo su i più interessanti.

  Sia \(\phi \eqdef \neg \psi\) con \(\psi \in \FOL\). Per ipotesi induttiva, si ha che \(\Gamma \vdash \subs{\psi}{x_1,\ldots,x_n}{s_1,\ldots,s_n} \to \subs{\psi}{x_1,\ldots,x_n}{r_1,\ldots,r_n}\). A questo punto, tramite \(\to E\) assumendo \(\subs{\psi}{x_1,\ldots,x_n}{s_1,\ldots,s_n}\) è possibile ottenere \(\subs{\psi}{x_1,\ldots,x_n}{r_1,\ldots,r_n}\).


  \begin{prooftree}
    \AxiomC{\(\Gamma\)}
    \UnaryInfC{\(\vdots\)}
    \RuleLabel{IH}
    \UnaryInfC{\(\subs{\psi}{x_1,\ldots,x_n}{s_1,\ldots,s_n} \to \subs{\psi}{x_1,\ldots,x_n}{r_1,\ldots,r_n}\)}

    \AxiomC{\(\subs{\psi}{x_1,\ldots,x_n}{s_1,\ldots,s_n}\)}
    \RuleLabel{\(\to E\)}
    \BinaryInfC{\(\subs{\psi}{x_1,\ldots,x_n}{r_1,\ldots,r_n}\)}
  \end{prooftree}
  Ora possiamo notare, per definizione di sostituzione per formule negate che
  \[\subs{\neg \psi}{x_1,\ldots,x_n}{s_1,\ldots,s_n} = \neg \subs{\psi}{x_1,\ldots,x_n}{s_1,\ldots,s_n}\]

  Dunque assumendo \(\subs{\neg \psi}{x_1,\ldots,x_n}{s_1,\ldots,s_n}\) riusciamo a ottenere una contraddizione, e quindi tramite \(\neg I\) otteniamo \(\neg \subs{\psi}{x_1,\ldots,x_n}{r_1,\ldots,r_n}\). Infine, scarichiamo l'ultima assunzione tramite \(\to I\), ottenendo il risultato sperato.

  \begin{prooftree}
    \AxiomC{\(\Gamma\)}
    \UnaryInfC{\(\vdots\)}
    \RuleLabel{IH}
    \UnaryInfC{\(\subs{\psi}{x_1,\ldots,x_n}{s_1,\ldots,s_n} \to \subs{\psi}{x_1,\ldots,x_n}{r_1,\ldots,r_n}\)}

    \AxiomC{\({[\subs{\psi}{x_1,\ldots,x_n}{s_1,\ldots,s_n}]}_1\)}
    \RuleLabel{\(\to E\)}
    \BinaryInfC{\(\subs{\psi}{x_1,\ldots,x_n}{r_1,\ldots,r_n}\)}
    \AxiomC{\({[\neg\subs{\psi}{x_1,\ldots,x_n}{s_1,\ldots,s_n}]}_2\)}
    \RuleLabel{\(\neg E\)}
    \BinaryInfC{\(\bot\)}
    \RuleLabel{\(\neg I_1\)}
    \UnaryInfC{\(\neg \subs{\psi}{x_1,\ldots,x_n}{r_1,\ldots,r_n}\)}
    \RuleLabel{\(\to I_2\)}
    \UnaryInfC{\(\neg\subs{ \psi}{x_1,\ldots,x_n}{s_1,\ldots,s_n} \to \neg\subs{\psi}{x_1,\ldots,x_n}{r_1,\ldots,r_n}\)}
  \end{prooftree}

  Ora possiamo notare, per definizione di sostituzione per formule negate, che
  \[\subs{\phi}{x_1,\ldots,x_n}{s_1,\ldots,s_n} \eqdef \subs{\neg \psi}{x_1,\ldots,x_n}{s_1,\ldots,s_n} \eqdef \neg \subs{\psi}{x_1,\ldots,x_n}{s_1,\ldots,s_n}\]

  Dunque l'ultima deduzione ottenuta è proprio \(\Gamma \vdash \subs{\phi}{x_1,\ldots,x_n}{s_1,\ldots,s_n} \to \subs{\phi}{x_1,\ldots,x_n}{r_1,\ldots,r_n}\).

  Tutti gli altri casi induttivi, ovvero quelli per \(\land\), \(\lor\), \(\to\), \(\forall\) e \(\exists\) sono analoghi, ovvero partono dall'ipotesi induttiva, applicando \(\to E\) ottengono la sottoformula sostituita con \(r_1,\ldots,r_n\) e poi, tramite regole di introduzione ed eliminazione specifiche per il simbolo logico in questione otteniamo la formula sostituita con \(r_1,\ldots,r_n\), scaricando infine l'assunzione tramite \(\to I\) per ottenere il risultato sperato.

  Un facile esempio si ha con \(\phi = \psi_1 \land \psi_2\). Per questione di brevità andremo a indicare le sequenze di simboli in notazione vettoriale, ovvero \(\overline{x} = x_1,\ldots,x_n\), \(\overline{s} = s_1,\ldots,s_n\) e \(\overline{r} = r_1,\ldots,r_n\)

  \begin{prooftree}
    \small
    \AxiomC{\(\Gamma\)}
    \UnaryInfC{\(\vdots\)}
    \RuleLabel{IH}
    \UnaryInfC{\(\subs{\psi_1}{\overline{x}}{\overline{s}} \to \subs{\psi_1}{\overline{x}}{\overline{r}}\)}

    \AxiomC{\({[\subs{\psi_1}{\overline{x}}{\overline{s}} \land \subs{\psi_2}{\overline{x}}{\overline{s}}]}_1\)}
    \RuleLabel{\(\land E^l\)}
    \UnaryInfC{\(\subs{\psi_1}{\overline{x}}{\overline{s}}\)}
    \RuleLabel{\(\to E\)}
    \BinaryInfC{\(\subs{\psi_1}{\overline{x}}{\overline{r}}\)}

    \AxiomC{\({[\subs{\psi_1}{\overline{x}}{\overline{s}} \land \subs{\psi_2}{\overline{x}}{\overline{s}}]}_1\)}
    \RuleLabel{\(\land E^r\)}
    \UnaryInfC{\(\subs{\psi_2}{\overline{x}}{\overline{s}}\)}
    \RuleLabel{\(\to E\)}

    \AxiomC{\(\Gamma\)}
    \UnaryInfC{\(\vdots\)}
    \RuleLabel{IH}
    \UnaryInfC{\(\subs{\psi_2}{\overline{x}}{\overline{s}} \to \subs{\psi_2}{\overline{x}}{\overline{r}}\)}

    \BinaryInfC{\(\subs{\psi_2}{\overline{x}}{\overline{r}}\)}
    \RuleLabel{\(\land I\)}
    \BinaryInfC{\(\subs{\psi_1}{\overline{x}}{\overline{r}} \land \subs{\psi_2}{\overline{x}}{\overline{r}}\)}
    \RuleLabel{\(\to I_1\)}
    \UnaryInfC{\((\subs{\psi_1}{\overline{x}}{\overline{s}} \land \subs{\psi_2}{\overline{x}}{\overline{s}}) \to (\subs{\psi_1}{\overline{x}}{\overline{r}} \land \subs{\psi_2}{\overline{x}}{\overline{r}})\)}
  \end{prooftree}
  Notando che
  \[\subs{\phi}{x_1,\ldots,x_n}{s_1,\ldots,s_n} \eqdef \subs{\psi_1\land\psi_2}{x_1,\ldots,x_n}{s_1,\ldots,s_n} \eqdef \subs{\psi_1}{x_1,\ldots,x_n}{s_1,\ldots,s_n} \land \subs{\psi_2}{x_1,\ldots,x_n}{s_1,\ldots,s_n}\]
  \todo{Ma riguardo i quantificatori? Non bisogna stare attenti ai vincoli?}
\end{proof}

Infine, andiamo come anticipato a mostrare come i vincoli di applicabilità per le regole dei quantificatori oltre a essere necessari sono sufficienti. Come già accennato precedentemente, sarà sufficiente mostrare la correttezza del calcolo di deduzione naturale ovvero, come fatto nella logica proposizionale, mostrare che dati \(\phi \in \FOL\) e \(\Gamma \subseteq \FOL\) si ha \(\Gamma \vdash \phi \implies \Gamma \models \phi\).

\begin{theorem}{Correttezza del calcolo di deduzione naturale nel prim'ordine}
  Dati \(\phi \in \FOL\) e \(\Gamma \subseteq \FOL\) si ha che
  \[\Gamma \vdash \phi \implies \Gamma \models \phi\]
\end{theorem}

\begin{proof}
  Come visto per la logica proposizionale, per dimostrare la correttezza andremo a fare induzione strutturale sull'albero di deduzione.

  Per alberi di altezza \(0\) si ha chiaramente che l'unico nodo è sia foglia che radice, ovvero è etichettato da \(\phi\) e appartiene a \(\Gamma\). Chiaramente se \(\phi \in \Gamma\) si ha che \(\Gamma \models \phi\) per definizione di soddisfacibilità.

  Passando al caso induttivo, supponiamo di avere una deduzione con \(\phi\) in radice, foglie non scaricate tutte presenti in \(\Gamma\). Chiaramente, ogni sottalbero della radice è una deduzione di altezza minore per un certo \(\psi\) con assunzioni non scaricate in \(\Gamma\) per cui vale l'ipotesi induttiva che \(\Gamma \models \psi\) per ogni sottoalbero.

  Per dimostrare la tesi sarà necessario quindi analizzare tutte le regole che possono aver prodotto \(\phi\) a partire dalle \(\psi\) dei sottoalberi, e mostrare che l'applicazione di tale regole mantiene correttezza semantica. Ovviamente, per le regole comuni alla logica proposizionale non c'è alcuna differenza con la dimostrazione già vista.

  Iniziamo con \(\forall E\). In tal caso, l'albero di deduzione avrà la forma

  \begin{prooftree}
    \AxiomC{\(\Gamma\)}
    \UnaryInfC{\(\vdots\)}
    \UnaryInfC{\(\forall x, \psi\)}
    \RuleLabel{\(\forall E\)}
    \UnaryInfC{\(\subs{\psi}{x}{t}\)}
  \end{prooftree}
  
  Ovviamente, per applicabilità di \(\forall E\), \(t\) è libero per \(x\) in \(\psi\). Da ipotesi induttiva si ha che \(\Gamma\models\forall x,\psi\) e quindi per ogni coppia \(M,\sigma\) si ha, per ogni \(d \in D\):
  \[M,\sigma\models \Gamma \implies M,\sigma \models \forall x, \psi \models M,\sigma[x\mapsfrom d]\models \psi\]

  Ovviamente \(\den{t}[M][\sigma] \in D\) e quindi, per \(d = \den{t}[M][\sigma]\), si ha che \(M,\sigma \models \Gamma \implies M,\sigma[x\mapsfrom \den{t}[M][\sigma]]\models \psi\). Ora, dal \Namedcref{thm:fol_substitution}, si ha che \(M,\sigma[x\mapsfrom \den{t}[M][\sigma]]\models \psi \iff M,\sigma\models \subs{\psi}{x}{t}\) se \(t\) è libero per \(x\) in \(\psi\). Essendo ciò vero per ipotesi, abbiamo che per ogni \(M,\sigma\):
  \[M,\sigma\models \Gamma \implies M,\sigma\models \subs{\psi}{x}{t}\]
  e quindi \(\Gamma \models \subs{\psi}{x}{t} \eqdef \phi\).
\end{proof}


\chapter{Applicazioni ed Espressività della logica}

Infine, andiamo ad analizzare come i linguaggi logici e i sistemi di deduzione possano essere applicati nell'informatica nella risoluzione di problemi. Una logica col suo calcolo di deduzione infatti fornisce un sistema formale di rappresentazione dei problemi oltre che una sistema di inferenza per risolverli. 
In tal senso, non si è molto distanti da un modello di calcolo computazionale.
Infatti, esistono linguaggi di programmazione logica, come Prolog, che si basano su questi principi per esprimere programmi e risolvere problemi.

Tra le varie applicazioni, quella su cui ci concentreremo è la specifica di problemi decisionali, ovvero problemi che possono essere risolti con una risposta sì o no.

In particolare, andremo a vedere come le sintassi e le semantiche dei linguaggi logici possano essere utilizzate per rappresentare gli input, gli output e una relazione tra di essi, in modo tale che la congiunzione delle formule risultanti sia soddisfacibile se e solo se esiste il problema da risolvere ha risposta affermativa.

In tal senso questo tipo di computazione si discosta molto dai paradigmi più noti, poiché puramente dichiarativa. Infatti, ciò che ci permette di risolvere un problema non è una serie di operazioni o manipolazioni da eseguire, ma la verifica di alcune condizioni e vincoli che legano input e output.

\section{Applicazione della logica proposizionale}

Avendo trovato un algoritmo corretto, completo e terminante per risolvere la soddisfacibilità di formule proposizionali, se è possibile riuscire a trovare una \keyword{riduzione} di un qualsiasi problema a \(\Sat\) possiamo avere un algoritmo per risolvere quel problema.

Per riduzione si intende una codifica di un'istanza qualsiasi del problema in un'istanza di \(\Sat\), ovvero una formula proposizionale, tale che la soluzione dell'istanza di \(\Sat\) ci fornisca la soluzione dell'istanza del problema originale.

Tale codifica per essere utile deve essere efficiente, ovvero deve essere eseguibile in tempo polinomiale rispetto alla dimensione dell'istanza del problema originale. Questo poiché essendo \(\Sat\) un problema NP-completo, riduzioni inefficienti porterebbero, oltre che a poter rappresentare problemi che non sono in NP, a non avere un algoritmo efficiente per risolvere il problema originale, la cui complessità è nascosta proprio nella riduzione.

Consideriamo qualche esempio semplice di riduzione a \(\Sat\), concentrandoci principalmente su problemi su grafi, che sono tra i più comuni.

\subsection{Ciclicità}
Il problema della ciclicità è il seguente: dato un grafo \(G=(V,E)\), con \(V\) insieme dei vertici e \(E\subseteq V\times V\) insieme degli archi, esiste un ciclo in \(G\)?

Un ciclo in un grafo è ovviamente definito come un percorso che contiene più di una volta lo stesso vertice, o riducendoci al caso dei cicli semplici, un percorso che inizia e finisce nello stesso vertice. Per percorso chiaramente si intende una sequenza di vertici \(v_1,v_2,\dots,v_k\) tale che \((v_i,v_{i+1})\in E\) per ogni \(i=1,\dots,k-1\).

Come già detto, per operare la riduzione, è necessario trovare una rappresentazione dell'istanza dell'input, ovvero del grafo \(G\), e dell'output, ovvero un potenziale ciclo, e dei vincoli che legano i due, in modo tale che la congiunzione di tutte le formule risultanti sia soddisfacibile se e solo se esiste un ciclo in \(G\).

Per la rappresentazione del grafo è necessario andare a rappresentare i vertici e gli archi. Essendo che \(E\subseteq V\times V\) possiamo rappresentare ogni elemento di \((u,v)\in V\times V\) con una variabile \(X_{u,v}\) e richiedere che \(X_{u,v}\) sia vera se e solo se \((u,v)\in E\).

\[\phi_E \eqdef \bigwedge_{(u,v)\in E} X_{u,v} \land \bigwedge_{(u,v)\notin E} \neg X_{u,v}\]

Tale formula intuitivamente rappresenta un assegnamento di verità alle variabili \(X_{u,v}\) che è vero se e solo se rappresenta correttamente il grafo \(G\). Di conseguenza, andando a congiungere questa formula con altre che rappresentano i vincoli che legano l'input all'output, possiamo essere certi che qualunque assegnamento soddisfi la formula risultante questo corrisponda a proprietà del grafo \(G\) in input.

Per quanto riguarda l'output, un modo alternativo alla sequenza di vertici per descrivere un ciclo è quello di rappresentarlo come sequenza di archi. Vorremo dunque mostrare che esista un \(S\subseteq E\) tale che \(S\) contiene un ciclo. Partiamo da definire le variabili per l'output come \(Y_{u,v}\) per ogni \(u,v \in V\). Dunque possiamo definire 
\[\phi_S \eqdef \bigwedge_{u,v \in V} Y_{u,v}\to X_{u,v}\]
Tale formula ci dice che ogni elemento di \(S\) deve stare in \(E\). Inoltre, per far si che \(S\) contenga un ciclo non può essere vuoto, e questo è esprimibile con
\[\phi_{\neq \emptyset} \eqdef \bigvee_{u,v\in V} Y_{u,v} \]

Ci manca solo da definire i vincoli che legano \(S\) al fatto che contenga un ciclo. Seppur in modo non estremamente intuitivo, la formula che ci permette di esprimere questo vincolo è la seguente:
\[\phi_{cyc} \eqdef \bigwedge_{u,v} (Y_{u,v \in V}\to \bigvee_{z \in V} Y_{v,z})\]

Riflettendoci, essendo che \(\phi_{\neq \emptyset}\) ci dice che esiste almeno un \(Y_{u,v}\) vero, \(\phi_{cycle}\) porta ad avere una catena di \(Y_{u,v}\) veri, che può chiudersi solo se esiste un ciclo. Infatti, l'unico modo per far si che la catena di \(Y_{u,v}\) veri non si chiuda è poter scegliere ogni volta un \(Y_{v,z}\) diverso da soddisfare, ma essendo che \(V\) è finito, prima o poi si sarà costretti a chiudere la catena, e dunque ad avere un ciclo. Equivalentemente, il percorso indotto dalle variabili di output vere per non contenere un ciclo dovrebbe essere infinito, ma essendo che \(V\) è finito, questo non è possibile. Viceversa, se esiste un ciclo, allora esiste un assegnamento che soddisfa \(\phi_{cycle}\) e \(\phi_{\not \emptyset}\) che è quello che rende vere le variabili \(Y_{u,v}\) corrispondenti agli archi del ciclo.

Dunque, concludendo, la formula che rappresenta l'istanza di \(\Sat\) a cui ridurre il problema della ciclicità è la seguente:
\[\phi \eqdef \phi_E \land \phi_S \land \phi_{\neq \emptyset} \land \phi_{cyc}\]

Per tale problema sono noti anche algoritmi polinomiali, dunque questa riduzione non è particolarmente utile, ma è un esempio semplice di come si possa rappresentare un problema con una formula proposizionale.

Passiamo ora a un altro problema, più complesso, per cui non sono noti algoritmi polinomiali, ovvero il problema della colorazione di grafi.

\subsection{K-colorazione}
Il problema della \(k\)-colorazione è il seguente: dato un grafo \(G=(V,E)\) e un intero \(k\leq \abs{V}\) numero di colori, è possibile colorare i vertici di \(G\) con al più \(k\) colori in modo tale che nessuna coppia di vertici adiacenti abbia lo stesso colore?

Più formalmente, \(\exists f:V\to C = \set{1,\dots,k} \st \forall (u,v)\in E, f(u)\neq f(v)\)?
Questo problema è molto interessante perché rappresenta una generalizzazione di un problema di allocazione di risorse. Inoltre, questo problema è molto più complicato del precedente perché si dimostra essere NP-completo.

Vediamo dunque una sua riduzione a \(\Sat\). Per rappresentare il grafo \(G\), possiamo rifarci alla rappresentazione già vista per il problema della ciclicità. Questo porta ad avere un insieme di proposizioni atomiche \(\AP[\phi_E]\) con \(\abs{\AP} = \abs{V}^2\) per rappresentare il grafo \(G\).
Per la funzione di colorazione \(f\) invece, possiamo vederla come una relazione \(f\subseteq V\times C\) con un vincolo funzionale. Dunque possiamo rappresentarla estensionalmente tramite le coppie del prodotto cartesiano.

In particolare, possiamo avere \(Y_v^c\) con \(v\in V\) e \(c\in C\), in cui ogni assegnamento di tali variabili va interpretato come una colorazione del vertice \(v\) del colore \(c\) se \(Y_v^c\) vale \(1\). Dunque \(\AP' = \set{Y_v^c}[c\in C \land v \in V]\) e quindi \(\abs{\AP'} = \abs{V}k\).

Ci rimane da imporre i vincoli per rendere tale relazione una funzione, ovvero la funzionalità e la totalità, e il vincolo che impone che vertici adiacenti non abbiano lo stesso colore.

Per la funzionalità e la totalità, ovvero che un vertice sia colorato con uno e un solo colore, possiamo imporre il seguente vincolo:
\[\phi_{f} \eqdef \bigwedge_{v\in V}\left(\bigvee_{c\in C} \left(Y_v^c \land \bigwedge_{c'\in C\setminus\set{c}} \neg Y_v^{c'}\right)\right)\]

Tale vincolo infatti impone che per ogni vertice \(v\) esista almeno un colore \(c\) tale che \(Y_v^c\) sia vero, e che per ogni colore \(c'\) diverso da \(c\), \(Y_v^{c'}\) sia falso. Dunque, ogni vertice è colorato con uno e un solo colore.

Per il vincolo che impone che vertici adiacenti non abbiano lo stesso colore, possiamo imporre il seguente vincolo:
\[\phi_{adj} \eqdef \bigwedge_{(u,v)\in E} \bigwedge_{c\in C} \neg (Y_u^c \land Y_v^c)\]

Possiamo osservare che \(\abs{\phi_{f}} = \BigO{\abs{V}k^2}\) e \(\abs{\phi_{adj}} = \BigO{\abs{V}^2k}\), dunque la formula risultante dalla congiunzione di queste due formule sarà \(\BigO{\abs{V}^3}\) essendo \(k\leq \abs{V}\). Questa codifica del problema è però problematica se l'input dei colori è rappresentato usando solo \(k\) e non l'intero insieme \(C\), poiché per rappresentare \(C\) è necessario rappresentare ogni elemento di \(C\), il che richiede spazio lineare in \(k\), mentre per rappresentare \(k\) è sufficiente rappresentare il numero \(k\) in notazione posizionale, che richiede spazio logaritmico in \(k\). In tale scenario, una lunghezza lineare in \(k\) è esponenzialmente più grande di quella dell'input e dunque la codifica non è efficiente. 

Esiste fortunatamente modo di ovviare a tale problema sfruttando proprio la rappresentazione in notazione posizionale di \(k\). In particolare, essendo che \(k\) si può rappresentare con \(\lceil \log_2(k) \rceil\) bit, è possibile rappresentare ogni colore \(c\) con una stringa di \(\lceil \log_2(k) \rceil\) bit, e dunque utilizzando \(\lceil \log_2(k) \rceil\) variabili. Dunque, invece di avere \(Y_v^c\) con \(v\in V\) e \(c\in C\), possiamo avere \(Y_v^i\) con \(v\in V\) e \(i=1,\dots,\lceil \log_2(k) \rceil\), in cui ogni assegnamento di tali variabili va interpretato come una colorazione del vertice \(v\) del colore rappresentato dalla stringa di bit formata da tali variabili. Dunque \(\AP' = \set{Y_v^i}[i=1,\dots,\lceil \log_2(k) \rceil \land v \in V]\) e quindi \(\abs{\AP'} = \BigO{\abs{V}\log k}\).

È interessante notare che questa nuova rappresentazione non necessita l'espressione del vincolo di funzionalità e totalità, poiché per ogni vertice esiste un solo insieme di variabili \(Y_v^i\) con \(i=1,\dots,\lceil \log_2(k) \rceil\) che può essere vero, e dunque ogni vertice è colorato con uno e un solo colore. Dunque, l'unico vincolo da esprimere è quello che impone che vertici adiacenti non abbiano lo stesso colore, che è il seguente:
\[\phi_2' = \bigwedge_{(u,v)\in E}\left( \bigvee_{i=0}^{\log_2(k)-1} (Y_u^i \land \neg Y_v^i) \lor (\neg Y_u^i \land Y_v^i)\right)\] 

In altre parole, quello che richiediamo adesso è che ogni coppia di vertici adiacenti abbia almeno un bit diverso nella rappresentazione del colore. Questo può essere espresso in maniera leggermente più compatta utilizzando l'operatore XOR, che è vero se e solo se i disgiunti sono diversi:
\[\phi_2' = \bigwedge_{(u,v)\in E}\left( \bigvee_{i=0}^{\log_2(k)-1} (Y_u^i \oplus Y_v^i)\right)\]

Dunque, concludendo, la formula che rappresenta l'istanza di \(\Sat\) a cui ridurre il problema della \(k\)-colorazione è la seguente:

\[\phi' \eqdef \phi_E \land \phi_2'\]

Tale codifica chiaramente occupa spazio \(\abs{V}^2\log_2(k)\), che è lineare nell'input.
Possiamo facilmente notare in oltre che in questa codifica, per \(k=2\) otteniamo tutte formule che contengono al più 2 variabili.
Tale frammento ristretto della logica proposizionale è noto come \(2-\Sat\) e per esso il problema della soddisfacibilità è risolvibile in tempo polinomiale. Dunque, questa codifica ci dice anche che \(2-\)colorazione è risolvibile in tempo polinomiale.

\subsection{Ciclo Hamiltoniano}

Un altro problema molto interessante è il problema del ciclo Hamiltoniano, che è il seguente: dato un grafo \(G=(V,E)\), esiste un ciclo semplice che contiene ogni vertice di \(G\) esattamente una volta?

Anche qui, per la rappresentazione del grafo \(G\) utilizzeremo \(\phi_E\) vista in precedenza.

Per avere un ciclo hamiltoniano, ovviamente, ci serve una sequenza di vertici lunga quanto l'insieme dei vertici, tale che ogni vertice sia presente esattamente una volta e che ci sia un arco tra ogni coppia di vertici adiacenti nella sequenza e tra l'ultimo e il primo vertice della sequenza.

Quindi per \(\abs{V}=n\) bisogna avere una sequenza di vertici \(v_0, \ldots, v_{n-1}\) tale che \(v_i \neq v_j\) per \(i\neq j\) e \((v_i,v_{i+1 \mod n})\in E\) per \(i=0,\ldots, n-1\).

È possibile vedere tale sequenza come una permutazione dell'insieme dei vertici \(V\), ovvero una funzione biettiva \(\pi:V\to \set{0,\ldots,n-1}\) che associa a ogni vertice la sua posizione nella sequenza. Dunque, per rappresentare l'output, ci basta rappresentare tale permutazione e poi imporre il vincolo che vertici adiacenti nella sequenza siano adiacenti nel grafo.

Per semplicità ci restringeremo al caso in cui \(n\) è una potenza di \(2\). Chiaramente, è possibile estendere la codifica al caso generale, ma richiede diversi accorgimenti tecnici che non sono particolarmente interessanti ai fini di questa trattazione.

Chiaramente per la rappresentazione di \(\pi\) possiamo prendere spunto da quanto fatto per la \(k\)-colorazione, e utilizzare per ogni vertice \(v\) un insieme di variabili \(Y_v^i\) con \(i=1,\dots,\log_2(n)\), in cui ogni assegnamento di tali variabili va interpretato come la rappresentazione binaria della posizione del vertice \(v\) nella sequenza. Qui l'ipotesi che \(n\) sia una potenza di \(2\) è fondamentale, poiché in questo modo è garantito che non ci possano essere vertici assegnati a posizioni che non esistono. Questo infatti garantisce che l'insieme dei vertici e l'insieme dei numeri rappresentati dalle variabili \(Y_v^i\) siano della stessa cardinalità, e dunque che sia possibile avere una corrispondenza biunivoca tra i due insiemi, che è ciò che ci serve per rappresentare una permutazione. Inoltre, questo ci permette di poter esprimere la biettività della funzione \(\pi\) utilizzando solo l'iniettività, poiché l'iniettività garantisce anche la suriettività quando i due insiemi hanno la stessa cardinalità. Dunque, per esprimere la biettività di \(\pi\) è sufficiente esprimere il vincolo di iniettività, che è il seguente:

\[\phi_{bij} \eqdef \bigwedge_{u,v\in V\st u\neq v} \bigvee_{i=0}^{\log_2(n)-1} (Y_u^i \oplus Y_v^i)\]

In altre parole, quello che tale vincolo richiede è che \(\forall u,v \in V, u\neq v \implies \pi(u)\neq\pi(v)\).
Quello che ci rimane da fare è imporre il vincolo che vertici adiacenti nella sequenza siano adiacenti nel grafo.

Iniziamo assicurandoci che tale permutazione sia effettivamente un percorso. Per fare ciò è sufficiente che ogni coppia di vertici consecutivi nella sequenza sia effettivamente un arco. Per fare ciò è sufficiente imporre che:
\[\phi_{path} \eqdef \bigwedge_{u,v\in V} \phi_{succ(u,v)} \to X_{u,v}\]
Dove \(\phi_{succ(u,v)}\) è una formula che è vera se e solo se \(\pi(v) = \pi(u)+1 \mod n\). Per definire tale formula, è necessario avere modo di codificare l'incremento di \(1\) modulo \(n\) nella rappresentazione binaria. Per fare ciò possiamo pensare a un circuito logico che prende in input la rappresentazione binaria di un numero \(x\) e restituisce in output la rappresentazione binaria di \(x+1 \mod n\). Tale circuito è composto, per ogni bit, da una porta XOR che calcola il nuovo valore del bit, e da una porta AND che calcola se è necessario portare \(1\) al bit successivo come riporto. Volendo generalizzare, è possibile vedere l'incremento di \(1\) come riporto iniziale portando a una struttura di questo tipo:

\begin{figure}[h!]
	\centering
	\begin{tikzpicture}[
			x=1cm,
			y=1cm,
			>=Latex,
			snodo/.style={draw, rectangle, minimum width=1.4cm, minimum height=0.9cm, font=\small},
			stage/.style={font=\small, align=center}
		]
		% Nodi di stadio
		\node[snodo] (b1) at (0,0) {};
		\node[snodo] (b2) at (4.2,0) {};

		% Input dall'alto
		\draw[->] (0,1.6) -- (b1.north) node[midway, right, stage] {\(x_i\)};
		\draw[->] (4.2,1.6) -- (b2.north) node[midway, right, stage] {\(x_{i+1}\)};

		% Carry in da sinistra e propagazione a destra
    \node[stage] at (-2.5,0) {\(\cdots\)};
		\draw[->] (-2.0,0) -- (b1.west) node[midway, above, stage] {\(c_i\)};
		\draw[->] (b1.east) -- (b2.west) node[midway, above, stage] {\(c_{i+1} = x_i \land c_i\)};
		\draw[->] (b2.east) -- (8,0) node[midway, above, stage] {\(c_{i+2} = x_{i+1} \land c_{i+1}\)};
		\node[stage] at (8.5,0) {\(\cdots\)};

		% Output dal basso
		\draw[->] (b1.south) -- (0,-2.0) node[midway, right, stage] {\(y_i = x_i \oplus c_i\)};
		\draw[->] (b2.south) -- (4.2,-2.0) node[midway, right, stage] {\(y_{i+1} = x_{i+1} \oplus c_{i+1}\)};

		% Carry iniziale
		\node[stage] at (-1.6,-2.4) {Incremento di 1: \(c_0 = 1\)};
	\end{tikzpicture}
	\caption{Schema a blocchi dell'incrementatore binario con propagazione del riporto.}\label{fig:binary-incrementer}
\end{figure}

Per implementare ciò sarà sufficiente introdurre le variabili per i riporto \(C_i^{u,v}\), richiedere che \(C_0^{u,v}\) sia vero, e poi imporre i vincoli che definiscono il funzionamento del circuito, ovvero \(C_{i+1}^{u,v} \leftrightarrow (Y_i^v \land C_i^{u,v})\) e \(Y_v^i \leftrightarrow (Y_u^i \oplus C_i^{u,v})\). Dunque, la formula che impone che vertici adiacenti nella sequenza siano adiacenti nel grafo è la seguente:
\[\phi_{succ(u,v)} \eqdef C_0^{u,v}\land \bigwedge_{i=0}^{\log_2(n)-1} (Y_i^u \leftrightarrow (Y_i^v \oplus C_i^{u,v}))\land (C_{i+1}^{u,v} \leftrightarrow (Y_{i}^v \land C_{i}^{u,v}))\]

A questo punto sostituendo opportunamente le varie \(\phi_{succ(u,v)}\) in \(\phi_{path}\) al variare di \(u,v\) possiamo ottenere la riduzione completa con la formula:
\[\phi = \phi_E \land \phi_{perm} \land \phi_{path}\]

\subsection{Copertura dei vertici}

Un ultimo problema che andiamo ad analizzare, anche se meno nel dettaglio, è quello della copertura dei vertici. Dato un grafo \(G=(V,E)\) e un intero \(k\leq \abs{V}\), esiste un insieme \(S\subseteq V\) di cardinalità al più \(k\) tale che ogni arco di \(E\) sia incidente ad almeno un vertice di \(S\)?

Per semplicità, analizzeremo una variante del problema che richiede il vincolo di uguaglianza sulla cardinalità di \(S\) più che la disuguaglianza. Tale semplificazione non intacca particolarmente il problema poiché se è possibile trovare una copertura di esattamente \(k\) vertici, tale copertura rispetta anche il vincolo di averne almeno \(k\).

Per la rappresentazione dell'input, come solito, possiamo rappresentare il grafo con \(\phi_E\) e \(k\) in binario.

Dobbiamo dunque definire come rappresentare l'output. Per ogni \(v \in S\) definiamo una variabile \(S_v\) la cui interpretazione sarà la relazione di appartenenza, ovvero \(S_v \iff v \in S\).

A questo punto per verificare la copertura è sufficiente imporre 

\[\phi_{cover} = \bigwedge_{(u,v) \in E} (S_v \lor S_u)\]

Tale parte è banale poiché è banale individuare una copertura qualsiasi. Infatti \(V\) è banalmente una copertura di \(G\) per definizione. Non a caso tale formula è soddisfatta assegnando a tutte le variabili \(S_v\) il valore \(1\). Il vero problema è invece quello di imporre il vincolo che \(S\) abbia esattamente \(k\) elementi.

Per fare ciò dovremmo innanzitutto usare le variabili predisposte per \(k\) per codificare una formula che sia vera se e solo se l'assegnazione di tali variabili corrisponda alla rappresentazione binaria di \(k\). Per fare ciò possiamo banalmente costruire l'insieme \(P_k\) delle posizioni della rappresentazione binaria di \(k\) che sono pari a \(1\) ed esprime il vincolo come segue:

\[\phi_k = \bigwedge_{i \in P_k} K_i \land \bigwedge_{i \notin P_k} \neg K_i\]

A questo punto per rappresentare \(\abs{S}\) abbiamo bisogno di un contatore che conti quanti elementi sono in \(S\). Per fare ciò però è sufficiente replicare la codifica dell'incrementatore visto per il ciclo hamiltoniano una volta per ogni vertice usando come riporto iniziale la variabile che rappresenta la presenza del vertice nella soluzione \(S_v\), in modo da incrementare il contatore solo se il vertice è presente. A questo punto avendo lo stato del contatore fissato a zero, ottenibile richiedendo che le variabili del primo input siano tutte false usando la negazione, confrontando l'output di ogni incrementatore con l'input del successivo in modo che siano uguali e di confrontare l'output dell'ultimo incrementatore con la rappresentazione di \(k\). Per questioni di brevità verrà omessa la rappresentazione in formule di questa costruzione, che è però facile da implementare.

\section{Espressività e complessità del primo ordine}

Passando al prim'ordine, per quanto sia possibile rappresentare problemi decisionali come nella logica proposizionale, andremo a vedere che i problemi decisionali della logica, ovvero VALIDITY e SATISFIABILITY, sono tutti indecidibili in principio.

In particolare VALIDITY è RE (ricorsivamente enumerabile) e SATISFIABILITY è coRE (il suo complemento è ricorsivamente enumerabile), ma nessuno dei due è ricorsivo. Dunque, come vedremo successivamente, esistono algoritmi che terminano in tempo finito per l'accettazione di formule valide ma possono divergere per formule non valide, e algoritmi che terminano in tempo finito per l'accettazione di formule insoddisfacibili ma possono divergere per formule soddisfacibili.

Piuttosto che approfondire quindi come ridurre problemi decisionali a VALIDITY o SATISFIABILITY, prima di mostrare come questi siano indecidibili riducendo il problema della fermata di una macchina di Turing a VALIDITY, ci concentreremo su un altro aspetto pratico interessante, ovvero la rappresentazione della conoscenza e di conseguenza l'espressività della logica del primo ordine come linguaggio.

Abbiamo già visto per la logica proposizionale alcuni problemi che potevano essere espressi usando la logica proposizionale su un input finito. In particolare tutte le proprietà finitarie erano in qualche modo esprimibili. Al prim'ordine invece possiamo esprimere proprietà con una granularità molto fine avendo a che fare con domini potenzialmente infiniti. 

Andiamo a vedere quindi alcuni esempi e tecniche molto rudimentali per dimostrare caratteristiche di espressività della logica del primo ordine. In particolare, saremo interessati a descrivere alcune proprietà di oggetti del discorso e capire se tali proprietà sono esprimibili o meno in logica del primo ordine, ovvero se è possibile trovare una formula del primo ordine che sia soddisfatta esclusivamente dai domini che soddisfano quella proprietà.

Ovviamente quest'analisi dipenderà dal linguaggio che stiamo utilizzando, ovvero dalle costanti, funzioni e predicati di cui disponiamo. 

Iniziamo definendo propriamente il concetto di proprietà e di espressibilità.

\begin{definition}{Proprietà}
  Sia \(L\) un linguaggio di \(\FOL\) e sia \(\mathcal{M}\) la classe delle strutture con \(M \in \mathcal{M}\) se \(M = \langle D,\Interp\rangle\) per \(L\). Una \textbf{proprietà} \(P\) è un qualsiasi un sott'insieme di \(\mathcal{M}\).
\end{definition}

Intuitivamente, in ottica descrittivista del linguaggio, una proprietà non è altro che l'insieme degli oggetti del discorso che la rispettano. Ad esempio, se stiamo parlando di persone, la proprietà di essere un uomo è l'insieme di tutte le persone che sono uomini.

\begin{definition}{Espressibilità}
  Sia \(L\) un linguaggio di \(\FOL\) e sia \(\mathcal{M}\) la classe delle strutture con \(M \in \mathcal{M}\) se \(M = \langle D,\Interp\rangle\) per \(L\). Una proprietà \(P\) è \textbf{esprimibile} in \(L\) se esiste una formula \(\phi \in \FOL\) tale che \(P=\den{\phi}=\set{M \in \mathcal{M} \mid M \models \phi}\)
\end{definition}

Questo problema è particolarmente interessante per l'informatica poiché ci permette di capire quanto potere espressivo sia necessario per esprimere determinate proprietà nella modellazione di basi di conoscenza. Intuitivamente l'espressività di un linguaggio è inversamente proporzionale alla complessità computazionale dei problemi di decisione di tale linguaggio, quindi è importante capire quali sono le proprietà che possiamo esprimere con un linguaggio e quali no per trovare un giusto equilibrio tra potere espressivo e complessità computazionale.

Un modo molto semplice per capire se una proprietà è esprimibile o meno è quello di cercare di costruire una formula che la esprima. 

Un esempio l'abbiamo già visto, ovvero la proprietà di essere un ordinamento lineare.
\[\phi_{lin}\eqdef\forall x, \neg (x=x) \land \forall x,y (x=y\lor x<y \lor y<x) \land \forall x,y,z ((x<y \land y<z) \to x<z)\]
con chiaramente un linguaggio che abbia come minimo un simbolo di relazione binaria \(<\) e un simbolo di uguaglianza \(=\).

Chiaramente quindi \(\den{\phi_{lin}}\) è proprio la proprietà di essere linearmente ordinato.

Sempre sullo stesso linguaggio si può pensare anche di esprimere la proprietà \(DSC = \set{M\in\mathcal{M}}[ M \text{ è ordinamento lineare discreto}]\). Sicuramente \(\phi_{lin}\) sarà parte della formula che esprime \(DSC\) ma non è sufficiente. Dovremo congiungere a \(\phi_{lin}\) anche una proprietà che esprima la discretezza. Intuitivamente quello che differenzia insiemi discreti e non discreti è il concetto di successore. Bisogna però stare attenti al fatto che un insieme discreto può ammettere massimo, e che dunque non è richiesto che ogni elemento abbia un successore. Si dovranno dunque trattare come casi particolari anche gli elementi che non hanno successore. Intuitivamente quindi quello che dovremo esprimere è che ogni elemento che ha degli elementi successivi ha un minimo elemento successivo, che è proprio il successore. La formula che esprime tutto ciò è la seguente:
\[\phi' \eqdef \forall x, ((\exists y \st x< y) \to \exists z \st(x<z \land \forall w, (x<w \to \neg(w<z))))\]
E quindi avremo che \(\den{\phi_{lin} \land \phi'} = DSC\).

Questa \(\phi_{lin} \land \phi'\) però chiaramente non caratterizza \((\N,<)\) poiché ha altri modelli che la soddisfano. Un esempio è l'unione disgiunta di \(\N\) con se stesso, ovvero \(\N + \N\), formato da coppie \((n,i)\) con \(n \in \N\) e \(i \in \set{0,1}\) e con l'ordinamento definito come segue:
\[(n,i) < (m,i) \text{ se } n < m \in \N\]
\[(n,0) < (m,1) \text{ per ogni } n,m \in \N\]

Passiamo invece alla proprietà \(DNS = \set{M\in\mathcal{M}}[ M \text{ è ordinamento lineare denso}]\). Anche in questo caso \(\phi_{lin}\) è parte della formula che esprime \(DNS\) ma non è sufficiente. Dovremo congiungere a \(\phi_{lin}\) anche
\[\phi''\eqdef \forall x,y (x<y \to \exists z (x<z \land z<y))\]

Si ha quindi \(\den{\phi_{DNS} \eqdef\phi_{lin} \land \phi'} = DNS\). Chiaramente \((\Q,<)\models \phi_{DNS}\) e \((\R,<)\models \phi_{DNS}\).

È interessante notare che \(\phi'\not\equiv\neg \phi''\). Possiamo considerare ad esempio \(([0,1]\cup \set{2},<)\). Questo chiaramente non è denso poiché per \(x = 1\) e \(y = 2\) non esiste un elemento \(z\) tale che \(1 < z < 2\). D'altra parte però non è discreto poiché per qualsiasi elemento in \([0,1]\) non esiste un successore. 

Passando a qualche risultato negativo, abbiamo che \((\Q,<)\) è indistinguibile da \((\R,<)\), ovvero non esiste una formula \(\phi\) tale che \((\Q,<)\models \phi\) e \((\R,<)\not\models \phi\). Questa è conseguenza diretta del teorema di Löwenheim-Skolem che recita quanto segue:
\begin{theorem}{Löwenheim-Skolem}
  Se \(\Gamma\) è un insieme di formule chiuse al più numerabile e soddisfacibile, allora \(\Gamma\) ha un modello al più numerabile.
\end{theorem}

Sia infatti \(\phi\) la potenziale formula che distingue \((\Q,<)\) da \((\R,<)\). Dunque \((\Q,<)\models \phi\) e \((\R,<)\not\models \phi\). 
Equivalentemente, \((\R,<)\models \neg \phi\) e \((\Q,<)\not\models \neg \phi\). Dunque \(\set{\neg \phi, \phi_{DNS}}\) è un insieme di formule chiuse al più numerabile e soddisfatto da \((\R,<)\). Dal teorema di Löwenhaim-Skolem esiste però un modello al più numerabile che soddisfa \(\set{\neg \phi, \phi_{DNS}}\). Tale modello non può essere \((\R,<)\) poiché non è al più numerabile, quindi deve essere \((\Q,<)\). Dunque \((\Q,<)\models \neg \phi\) e \((\Q,<)\not\models \neg \phi\), il che è assurdo. 

Intuitivamente il teorema di Löwenhaim-Skolem vale per la costruzione vista per il teorema di completezza, poiché la costruzione del modello dell'insieme soddisfacibile \(\Gamma\) portava ad avere un dominio al più numerabile, poiché abbiamo aggiunto un numero enumerabile di insiemi enumerabile di costanti che è ancora un insieme enumerabile se lo era quello di partenza.

Questo quindi ci dice che ci sono limiti in ciò che il prim'ordine può esprimere.

\begin{theorem}{Compattezza}
  Per ogni \(\Gamma\subseteq\FOL\) (formule chiuse) e per ogni formula \(\phi\) (chiusa), vale che
  
  \[\Gamma \models \phi \iff \text{ esiste } \Gamma_0 \subseteq \Gamma \text{ finito tale che } \Gamma_0 \models \phi\]
\end{theorem}

\begin{proof}
  Si ha per completezza che \(\Gamma \models \phi \implies \Gamma \vdash \phi\). Se \(\Gamma \vdash \phi\) allora esiste una dimostrazione di \(\phi\) a partire da \(\Gamma\) che è finita, quindi coinvolge solo un numero finito di formule di \(\Gamma\), ovvero esiste un sottoinsieme finito \(\Gamma_0\) di \(\Gamma\) tale che \(\Gamma_0 \vdash \phi\). Per correttezza, se \(\Gamma_0 \vdash \phi\) allora \(\Gamma_0 \models \phi\).

  L'altra direzione è banale per monotònicità
\end{proof}

Una formulazione alternativa della compattezza è la seguente:
\begin{theorem}{Compattezza Alternativa}
  Per ogni \(\Gamma\subseteq\FOL\) (formule chiuse),

  \[\Gamma \in\Unsat \iff \text{ esiste } \Gamma_0 \subseteq \Gamma \text{ finito tale che } \Gamma_0 \in\Unsat\]
  o equivalentemente se
  \[\Gamma \in\Sat \iff \text{ per ogni } \Gamma_0 \subseteq \Gamma \text{ finito, } \Gamma_0 \in\Sat\]
\end{theorem}

Il teorema di compattezza è uno strumento particolarmente usato, soprattutto in logica matematica, per dimostrare risultati di espressività. Ad esempio \(FIN = \set{M \in \mathcal{M}}[\abs{M} < \infty]\) non è esprimibile in \(\FOL\) semplice, ovvero su linguaggio \(L = (\emptyset, \emptyset, \set{=})\). Supponiamo per assurdo che esista una formula \(\phi_{FIN}\) tale che \(\den{\phi_{FIN}} = FIN\). Consideriamo ora l'insieme parametrico di formule al variare di \(n\) definite come
\[\exists x_1, \ldots, x_n \bigwedge_{\substack{1 \leq i \leq n \\ 1 \leq j \leq n \\ i \neq j}} \neg(x_i = x_j)\]
Intuitivamente ognuna di queste formule esprime che esistono almeno \(n\) elementi distinti.

Consideriamo quindi \(\Gamma = \set{\phi_n}[n\in\N^*]\). Chiaramente questo insieme è numerabile. Sia ora \(\Sigma = \Gamma \cup \set{\phi_{FIN}}\). Ora se esiste un \(m\) che soddisfa \(\Gamma\) necessariamente deve essere infinito, poiché deve avere almeno \(n\) elementi distinti per ogni \(n\in \N\). Dunque \(\Sigma\) è insoddisfacibile perché qualsiasi modello di \(\Sigma\) dovrebbe essere sia finito che infinito. Ma ora considerando qualsiasi sott'insieme finito \(\Sigma_0\) di \(\Sigma\).

Chiaramente \(\Sigma_0\) è soddisfacibile, poiché basta un dominio finito di cardinalità maggiore o uguale al massimo indice \(n\) che compare in \(\Sigma_0\) per soddisfare tutte le formule di \(\Sigma_0\). Quindi abbiamo un insieme insoddisfacibile che però non ha nessun sottoinsieme finito insoddisfacibile, il che è assurdo per il teorema di compattezza.

Un altro esempio di particolare interesse per gli informatici è quello della non esistenza di una formula che caratterizza esattamente la classe dei grafi connessi. Chiaramente ciò non vuol dire che dato uno specifico grafo connesso non esista una formula che lo caratterizza o che esprima se quello specifico sia connesso, ma semplicemente che non esiste una formula che caratterizza tutti i grafi connessi e solo quelli.

Sia \(CONN = \set{M\in\mathcal{M}}[ M \text{ è un grafo connesso}]\). Supponiamo per assurdo che esista una formula \(\phi_{CONN}\) tale che \(\den{\phi_{CONN}} = CONN\). Per essere un grafo mi è sufficiente avere un linguaggio con un simbolo di relazione binaria \(E\) che rappresenta l'adiacenza tra i nodi con \(D\) che rappresenta i nodi. Per comodità aggiungiamo al linguaggio due costanti \(s\) e \(t\). Costruiamo quindi l'insieme di formule \(\Gamma = \set{\phi_n}[n\in\N \land n > 1]\) con
\[\phi_n \eqdef \neg\exists x_1, \ldots, x_n ( x_1 = s \land x_n = t \land \bigwedge_{1 \leq i < n} E(x_i, x_{i+1}))\]
Costruiamo quindi \(\Gamma' \eqdef \Gamma \cup \set{\forall x, \neg E(x,x)}\). Ogni formula \(\phi_n\) ci dice che non esiste un cammino di lunghezza \(n-1\) tra \(s\) e \(t\). Chiaramente preso \(M\) tale che \(M\models\Gamma'\) si ha che \(t^M\) non sarà raggiungibile da \(s^M\) e quindi \(M\) non sarà connesso. Dunque \(\Sigma \eqdef \Gamma' \cup \set{\phi_{CONN}}\) è insoddisfacibile. Ma ora se prendo un qualsiasi sottoinsieme finito \(\Sigma_0\) di \(\Sigma\) se questo non contiene \(\phi_{CONN}\) allora è banalmente soddisfacibile da qualsiasi grafo non connesso. Se invece \(\Sigma_0\) contiene \(\phi_{CONN}\) allora dalla finitezza esiste un indice massimo di formule di \(\Gamma\) che compare in \(\Sigma_0\), ovvero esiste \(n\) tale che \(\phi_n \in \Sigma_0\) e per ogni \(m > n\) si ha che \(\phi_m \not\in \Sigma_0\). A questo punto basta prendere un qualsiasi grafo in cui \(s\) e \(t\) sono connessi da un cammino di lunghezza strettamente maggiore di \(n\) per soddisfare tutte le formule di \(\Sigma_0\). Quindi anche in questo caso abbiamo un insieme insoddisfacibile che però non ha nessun sottoinsieme finito insoddisfacibile, il che è assurdo per il teorema di compattezza.

\subsection{Indecidibilità del prim'ordine}

Come anticipato precedentemente la logica del prim'ordine è indecidibile, e per mostrare ciò è necessario prima formalizzare precisamente tale idea. Riprendendo nozioni di teoria della complessità, nel dire che la logica del prim'ordine è indecidibile si intende che i problemi di decisione a essa associati, ovvero VALIDITY, SATISFIABILITY e LOGICAL CONSEQUENCE, sono indecidibili ovvero, riformulando tali problemi come linguaggi, e quindi la decisione come appartenenza a tali linguaggi, non esistono algoritmi che, in tempo finito, possano decidere l'appartenenza di un'arbitraria formula a tali linguaggi. Dunque, in altre parole, esistono formule per cui qualsiasi algoritmo che tenti di rispondere \qi{yes} o \qi{no} a seconda che la formula sia valida\footnote{Ovviamente riformulando il problema in un problema di appartenenza questo equivale a \q{la formula appartiene o meno al linguaggio delle formule valide}}, soddisfacibile o che sia conseguenza logica di un insieme di formule, non termina in tempo finito.

Come precedentemente accennato però è dimostrabile che questi problemi non sono \q{simmetrici} nel divergere. In particolare, esistono algoritmi che terminano in tempo finito per l'accettazione\footnote{Ovvero che rispondono \qi{yes} in tempo finito} di formule valide ma possono divergere per formule non valide, e algoritmi che terminano in tempo finito per il rifiuto\footnote{Ovvero che rispondono \qi{no} in tempo finito} di formule insoddisfacibili ma possono divergere per formule soddisfacibili. 

Per caratterizzare tali problemi è necessario introdurre alcune classi di complessità più fini, che non richiedono la totale decidibilità ma che si accontentino di uno dei due lati. Tali problemi sono detti semi-decidibili e appartengono alle classi RE (ricorsivamente enumerabili) e coRE (il complemento di RE).

Un insieme è ricorsivamente enumerabile se esiste una procedura che può enumerare tutti gli elementi dell'insieme. Invece coRE se il complemento dell'insieme può essere enumerato. 

Chiaramente avendo queste due procedure è semplice capire se il linguaggio è decidibile, poiché basta fare un passo alla volta alle enumerazioni di entrambe le procedure e vedere quale mi genera l'istanza.

In termini algoritmici, un linguaggio \(L\) è in RE se esiste un algoritmo che, dato un input \(x\), risponde \qi{yes} se \(x \in L\) e diverge altrimenti. Un linguaggio \(L\) è in coRE se esiste un algoritmo che, dato un input \(x\), risponde \qi{no} se \(x \not\in L\) e diverge altrimenti. Chiaramente, per essere decidibile, o ricorsivo, un linguaggio deve essere sia in RE che in coRE, ovvero deve esistere un algoritmo che risponde \qi{yes} se \(x \in L\) e \qi{no} se \(x \not\in L\) in tempo finito. Quello che andremo a mostrare in questa ultima sezione è che VALIDITY è in RE ma non è ricorsivo, e dunque per dualità e analogia si potrà concludere che SATISFIABILITY è in coRE ma non è ricorsivo, e che LOGICAL CONSEQUENCE è in RE ma non è ricorsivo.

Per mostrare ciò dovremo innanzitutto mostrare che VALIDITY è in RE, ovvero mostrare un algoritmo che almeno per le istanze positive del problema termini sempre con risposta \qi{yes}, e poi mostrare che ogni algoritmo che tenta di decidere completamente VALIDITY, e dunque anche di rispondere \qi{no} per le istanze negative porti a una contraddizione, ovvero che esistono formule per cui tale algoritmo diverge.

Iniziando con la prima parte, ovvero mostrare che VALIDITY è in RE, è sufficiente sfruttare i teoremi di completezza e correttezza. In particolare, vale \(\models \phi \iff \vdash \phi\). Riformulando quindi il problema in termini di appartenenza a linguaggi, si ha che \(\phi \in \Valid \iff \phi \in \set{\psi \in \FOL}[\vdash \psi]\). Dunque, per mostrare che VALIDITY è in RE, è sufficiente mostrare che l'insieme \(\set{\phi \in \FOL}[\vdash \phi]\) è in RE, ovvero che esiste una procedura per enumerarne tutti gli elementi. Per definizione di dimostrazione, una formula \(\phi\) è dimostrabile se esiste una derivazione di \(\phi\) a partire da \(\emptyset\). Essendo le derivazioni alberi finiti hanno tutti un altezza finita, e dunque è possibile iniziare a enumerare tutte le derivazioni di altezza \(0\), seguite da tutte le derivazioni di altezza \(1\), e così via. Tale enumerazione è ottenibile tramite una procedura poiché è sufficiente enumerare l'insieme finito di regole di inferenza e costruire tutte le possibili combinazioni di tali regole che portano a una derivazione di altezza \(n\) per ogni \(n\). Dunque, per ogni formula \(\phi\), se \(\phi\) è dimostrabile, allora esiste una derivazione di \(\phi\) a partire da \(\emptyset\) che sarà enumerata in un tempo finito, e dunque la procedura terminerà con risposta \qi{yes}.

Se si volesse mostrare adesso che VALIDITY è ricorsivo, si potrebbe pensare di mostrare che anche il complemento di VALIDITY, ovvero l'insieme delle formule non valide, è in RE, e dunque che esiste una procedura che termina con risposta \qi{yes} per ogni formula non valida. Un approccio intuitivo potrebbe essere quello di utilizzare lo stesso procedimento con le formule negate. Il problema sorge dal fatto che se \(\phi\) è valida, sicuramente \(\neg \phi\) è non valida, ma se \(\phi\) è non valida, non è detto che \(\neg \phi\) sia valida. In altre parole non tutte le formule non valide sono negazioni di formule valide, poiché è sufficiente che una formula sia soddisfacibile ma non valida per far si che la sua negata non sia valida. In altre parole l'insieme delle formule non valide è più grande dell'insieme delle negazioni di formule valide, e dunque non è sufficiente enumerare le negazioni di formule valide per enumerare tutte le formule non valide. 

A questo punto quindi si può provare a dimostrare che \(\overline{\text{VALIDITY}}\), ovvero il complemento di VALIDITY, non è in RE\@.

La teoria della complessità descrive diverse gerarchie di complessità per caratterizzare e confrontare problemi a seconda della loro difficoltà computazionale. Una nota gerarchia è la gerarchia polinomiale, che mette in relazione problemi di decisione risolvibili in tempo polinomiale con diversi modelli di computazione, come macchine di Turing deterministiche e non deterministiche. Una gerarchia che ci potrà essere utile è invece una delle tante che classifica i problemi indecidibili, ovvero la gerarchia aritmetica, in cui le proprietà sono classificate rispetto alle proprietà in teoria dei numeri. In particolare abbiamo \(\Sigma_1^0\) insieme delle formule per cui esiste una prova, \(\Pi_1^0\) insieme delle formule per cui non esiste una prova.

Un noto problema indecidibile è il problema della fermata di una macchina di Turing, ovvero data una qualsiasi macchina di Turing stabilire se essa termini in tempo finito o diverga. Come spesso accade in teoria della complessità, per mostrare che un problema è almeno difficile quanto un altro, e dunque in questo caso indecidibile, è sufficiente mostrare una riduzione, ovvero una procedura che, data un'istanza di un problema indecidibile, costruisca un'istanza del problema che vogliamo dimostrare essere indecidibile, in modo tale che la risposta alla nuova istanza sia \qi{yes} se e solo se la risposta alla vecchia istanza era \qi{yes}. Se ora il nuovo problema in esame fosse decidibile, allora sarebbe possibile decidere anche il problema indecidibile applicando a ogni sua istanza la riduzione seguita dalla procedura di decisione, portando a una contraddizione con l'indecidibilità nota del problema originale. 

Tali procedure di riduzione non possono essere arbitrarie, ma devono loro stesse appartenere a classi di complessità che siano più deboli della classe di complessità a cui il problema oggetto della dimostrazione dovrebbe appartenere. Se si vuole mostrare che un problema sia NP-completo, allora la procedura di riduzione non può essere a sua volta NP-completa, poiché altrimenti la complessità della procedura potrebbe risiedere nella procedura di riduzione stessa e non nel problema oggetto della dimostrazione. In questo caso, trattando di semi-decidibilità, è sufficiente che la procedura sia decidibile, ovvero che termini in tempo finito per ogni istanza, non potendo così \q{imputare} a essa un eventuale divergenza.

Andiamo dunque a mostrare tale riduzione, per poi dimostrare come essa verifichi la condizione di correttezza, ovvero che la formula costruita sia valida se e solo se la macchina di Turing termina. 

Iniziamo quindi a rivedere la definizione di macchina di Turing per poi mostrare come codificarla in una formula del prim'ordine.

\begin{definition}{Macchia di Turing}
  Una macchina di Turing è una tupla \(M = (Q, \Sigma, \delta, q_0, q_{n})\) dove:
  \begin{itemize}
    \item \(Q = \set{q_0, q_1, \ldots, q_t}\) è un insieme finito di stati;
    \item \(\Sigma = \set{\sigma_0, \sigma_1, \ldots, \sigma_m}\) insieme finito di simboli di nastro con \(\sigma_0 = \sqcup\) simbolo \q{blank};
    \item \(\delta: Q \times \Sigma \to Q \times \Sigma \times \{\leftarrow, \rightarrow\}\) funzione di transizione che, dato uno stato e un simbolo letto sul nastro, restituisce un nuovo stato, un simbolo da scrivere e una direzione in cui muovere la testina;
    \item \(q_0 \in Q\) è lo stato iniziale;
    \item \(q_{n} \in Q\) è lo stato di terminazione;
  \end{itemize}
\end{definition}

La semantica di questi modelli è data dalle \keyword{esecuzioni}, sequenze di \keyword{configurazioni} della macchina, ovvero istantanee dello stato complessivo della macchina in un dato momento. Formalmente:

\begin{definition}{Configurazione}
  Una configurazione \(C\) di una macchina di Turing \(M\) è una tripla \((k, i,t)\) con:
  \begin{itemize}
    \item \(k \in \set{0,\ldots,n}\) indice di stato, con \(q_k\in Q\);
    \item \(i \in \N\) indice di posizione della testina sul nastro;
    \item \(t: \N \to \Sigma\) funzione che associa a ogni cella del nastro il simbolo presente in quella cella.
  \end{itemize}

  L'insieme di tutte le possibili configurazioni di \(M\) è denotato con \(\mathcal{C}_M\).

  Indichiamo con \(C_0\) la configurazione iniziale di \(M\), ovvero la tripla \((0,0,t_0)\) con \(t_0 : \N \to \set{\sqcup}\).
\end{definition}

\begin{definition}{Esecuzione}
  Un esecuzione di una macchina di Turing \(M\) è una funzione \(\pi: \N \to \mathcal{C}_M\) che associa a ogni passo di esecuzione la configurazione della macchina \(M\) in quel passo.
  
  Si ha che \(\pi(0) = C_0\) e che, per ogni \(n\in\N\), se \(\pi(n) = (k, i, t_n)\) e \(\delta(q_k, t(i)) = (q_{k'}, \sigma_l, d)\) allora \(\pi(n +1) = (k', i', t_{n+1})\) con
  \begin{itemize}
    \item \(i' = i + 1\) se \(d = \rightarrow\) e \(i' = i - 1\) se \(d = \leftarrow\);
    \item \(t_{n+1}(j) = t_n(j)\) per ogni \(j \neq i\) e \(t_{n+1}(i) = \sigma_l\).
  \end{itemize}

  Un esecuzione è terminante se esiste \(j\in\N\) tale che \(\pi(j) = (n, i, t_j)\), ovvero se in qualche passo di esecuzione la macchina raggiunge lo stato di terminazione \(q_n\).
\end{definition}

Trattando macchine deterministiche, per ogni macchina di Turing \(M\) esiste una e una sola esecuzione \(\pi\) di \(M\). 
Quello che dobbiamo definire quindi è una formula \(\phi_M\) tale che\footnote{Data una tupla \(v\) e un indice \(1 \leq i \leq \abs{v}\), con \({(v)}_i\) indichiamo l'elemento di indice \(i\) della tupla \(v\). Ad esempio, se \(v = (k,i,t)\), allora \({(v)}_1 = k\), \({(v)}_2 = i\) e \({(v)}_3 = t\). Essendo nel nostro caso la terza componente una funzione, con \({(v)}_3(y)\) indichiamo l'applicazione della funzione alla variabile \(y\).}:
\[\models \phi_M \iff \exists i\st {(\pi_M(i))}_1 = n\]

Per lavorare in logica del prim'ordine, è necessario definire un linguaggio che ci permetta di rappresentare le configurazioni e le esecuzioni di \(M\). Intuitivamente, avendo descritto sia il nastro che l'esecuzione come funzioni con dominio \(\N\), possiamo costruire il linguaggio assumendo di avere tale dominio e definendo un interpretazione dei simboli coerente con la macchina di Turing. In particolare, sarà necessaria una costante \(0\) che rappresenterà l'inizio dei conteggi, sia di posizioni del nastro che di passi di esecuzione, due funzioni unarie \(s\) e \(p\) che intuitivamente verranno interpretate nelle funzioni di successore e predecessore, per permettere di esprimere lo spostamento della testina sul nastro e l'evoluzione dell'esecuzione nel tempo, e infine due insiemi di predicati binari, una per gli stati e un'altra per il nastro.

Per il resto della trattazione, per semplicità di notazione, fisseremo l'utilizzo della variabile \(x\) per descrivere i passi di esecuzione, e le variabili \(y\) e \(z\) per descrivere le posizioni sul nastro.

I predicati \(q_k\), con \(k \in \set{0,\ldots,n}\), che si occuperanno di descrivere le prime due componenti di una configurazione, saranno interpretati come segue:
\[\Interp(q_k) = \set{(x,y)\in\N^2}[{(\pi(x))}_1 = k \land {(\pi(x))}_2 = y ]\]
ovvero \(\Interp(q_k)\) è l'insieme di tutte le coppie \((x,y)\) tali che alla \(x\)-esima configurazione dell'esecuzione, la macchina si trovi nello stato \(q_k\) e la testina sia posizionata sulla cella \(y\).

I predicati \(\sigma_l\), con \(l \in \set{0,\ldots,m}\), che si occuperanno di descrivere il contenuto del nastro, saranno interpretati come segue:
\[\Interp(\sigma_l) = \set{(x,y)\in\N^2}[{(\pi(x))}_3(y) = \sigma_l] \]
ovvero \(\Interp(\sigma_l)\) è l'insieme di tutte le coppie \((x,y)\) tali che alla \(x\)-esima configurazione dell'esecuzione, il nastro contenga in posizione \(y\) il simbolo \(\sigma_l\).

A questo punto, avendo definito il linguaggio, non ci resta che definire un insieme di formule che descrivano il comportamento della macchina di Turing \(M\) in modo da poter mostrare la tesi.

Le prime formule che dobbiamo definire servono a imporre i vincoli necessari in modo che gli unici comportamenti descritti saranno coerenti con un esecuzione valida di \(M\). Iniziamo con l'assiomatizzare la proprietà di \(s\) e \(p\) di essere inverse.

\begin{equation}
  \phi_1 \eqdef \forall x, (p(s(x)) = x \land (\neg(x = 0) \to s(p(x)) = x))
\end{equation}

In aggiunta a questa formula sarebbe possibile aggiungere tutte le formule che assiomatizzano la teoria dei numeri naturali, ma non sarà necessario, poiché le formule che andremo a definire in seguito, che descrivono il comportamento di \(M\), implicheranno già tutte le proprietà necessarie per \(s\) e \(p\) per essere interpretati come funzioni di successore e predecessore su \(\N\).

Un altro vincolo che la macchina di Turing deve rispettare è che le celle di memoria sono atomiche, nel senso che, a ogni passo di esecuzione, ogni posizione del nastro può contenere solo e soltanto un simbolo. Per esprimere questo vincolo, è sufficiente la formula:

\begin{equation}
  \phi_2 \eqdef \forall x, \forall y, \left(\bigwedge_{i=0}^{m} \bigwedge_{\substack{j=0\\j \neq i}}^{m} \neg\left(\sigma_i\left(x,y\right) \land \sigma_j\left(x,y\right)\right)\right)
\end{equation}

Inoltre, a ogni passo di esecuzione, la macchina può trovarsi in un solo stato.

\begin{equation}
  \phi_3 \eqdef \forall x, \forall y, \forall z, \left(\bigwedge_{i=0}^{n} \bigwedge_{\substack{j=0\\j \neq i}}^{n} \neg\left(q_i\left(x,y\right) \land q_j\left(x,z\right)\right)\right)
\end{equation}

In questo caso è necessaria una doppia quantificazione per la posizione della testina sul nastro, poiché l'unico vincolo che si vuole imporre è l'unicità dello stato in cui si trova la macchina, e non l'unicità della posizione della testina.

Infine, per completare le proprietà di consistenza della macchina, ovvero quelle che escludono comportamenti anomali, dobbiamo garantire che, a ogni passo di esecuzione, anche essendo nello stesso stato, la macchina non possa avere la testina in due posizioni diverse, ovvero non possa leggere più celle contemporaneamente:

\begin{equation}
  \phi_4 \eqdef \forall x, \forall y, \forall z, \left(\neg(y = z)\to \bigwedge_{i=0}^{n} \neg\left(q_i\left(x,y\right) \land q_i\left(x,z\right)\right)\right)
\end{equation}

Per soddisfare queste formule è quindi necessario che, fissato un \(x\), esista un unico \(q_i\) che contiene coppie con \(x\) come primo elemento, e che tutte le coppie con \(x\) come primo elemento in tale \(q_i\) abbiano lo stesso secondo elemento, ovvero la posizione della testina. Inoltre, per ogni coppia \((x,y)\), questa può essere contenuta da un unico \(\sigma_l\), ovvero quello del simbolo presente in posizione \(y\) del nastro alla \(x\)-esimo passo di esecuzione.

In un certo senso, queste prima quattro formule impongono vincoli \q{statici} della macchina, nel senso che ne descrivono invarianti necessarie per mantenerne la coerenza, e di conseguenza escludere tutti i modelli che non rappresentano esecuzioni valide di \(M\). Le formule che andremo a definire adesso invece descrivono invece il comportamento dinamico di \(M\), ovvero come evolve l'esecuzione nel tempo, e dunque come evolvono le configurazioni della macchina da un passo di esecuzione al successivo. Intuitivamente, in queste formule sarà codificata la funzione di transizione \(\delta\). Iniziamo però dal codificare la configurazione iniziale. Come già detto la configurazione iniziale è \(C_0 = (0,0,t_0)\) con \(t_0 : \N \to \set{\sqcup}\). Tale configurazione dev'essere quella presente per prima nell'esecuzione, e quindi dev'essere quella associata a \(x = 0\). Per esprimere ciò, è sufficiente la formula:

\begin{equation}
  \phi_5 \eqdef q_0(0,0) \land \forall y, \sigma_0(0,y)
\end{equation}

Grazie ai vincoli di coerenza espressi prima, per assicurarci che all'inizio dell'esecuzione la macchina si trovi nella configurazione \(C_0\) è sufficiente descrivere come dev'essere fatta tale configurazione, poiché, grazie ai vincoli precedenti, sarà l'unica configurazione che soddisfa tali descrizioni.

A questo punto, prima di codificare propriamente la funzione di transizione, è necessario definire un ultimo vincolo di coerenza, che questa volta invece che essere locale assicura la coerenza tra due configurazioni consecutive. In particolare, coi vincoli precedenti abbiamo garantito che il nastro in ogni configurazione non abbia più di un simbolo nella stessa posizione e che la testina si trovi esattamente su una sola di esse, ma è necessario garantire anche che la posizione letta dalla testina sia l'unica che può essere modificata da un passo di esecuzione al successivo, e che tutte le altre posizioni del nastro rimangano inalterate. Per esprimere ciò, è sufficiente la formula:

\begin{equation}
  \phi_6 \eqdef \forall x, \forall y, \left(\bigwedge_{i=0}^{m} \left(\left(\sigma_i\left(x,y\right) \land \neg \sigma_i\left(s\left(x\right),y\right)\right) \to \left(\bigvee_{j=0}^{n} q_j\left(x,y\right)\right)\right)\right)
\end{equation}

Questa formula può sembrare complessa, ma in definitiva quello che richiede è che se in due istanti successivi (\(x\) e \(s(x)\)) il simbolo presente in una posizione \(y\) del nastro è diverso, allora la testina deve essere posizionata in \(y\) in uno qualsiasi degli stati. Questo vincolo, unito a quello che garantisce che in un istante di tempo la testina si trovi in una sola posizione, garantisce che l'unica posizione del nastro che può essere modificata da un passo di esecuzione al successivo è quella su cui è posizionata la testina. Se infatti da un passo al successivo più celle venissero modificate avremmo gli antecedenti di due di queste implicazioni in congiunzione verificati, ma al più uno dei conseguenti potrebbe essere verificato.

Infine, possiamo passare alla formula più complessa, che codifica la funzione di transizione \(\delta\). Per semplificarne la scrittura formeremo una formula aperta per ogni applicazione di \(\delta\), che poi andremo a congiungere e quantificare. Dunque presa un applicazione di \(\delta\) del tipo \(\delta(q_k, \sigma_l) = (q_{k'}, \sigma_{l'}, d)\), con \(d\) che può essere \(\rightarrow\) o \(\leftarrow\), definiamo le formule:
\begin{equation}
  \begin{aligned}
     \phi^{\delta^{\rightarrow}}_{k,l} (x,y) &\eqdef ((q_k(x,y) \land \sigma_l(x,y)) \to q_{k'}(s(x),s(y)) \land \sigma_{l'}(s(x),y))\\
     \phi^{\delta^{\leftarrow}}_{k,l} (x,y) &\eqdef ((q_k(x,y) \land \sigma_l(x,y)) \to q_{k'}(s(x),p(y)) \land \sigma_{l'}(s(x),y))
  \end{aligned}
\end{equation}
Dove \(k\) e \(l\) variano rispettivamente tra \(0\) e \(n\) e tra \(0\) e \(m\). Di queste chiaramente, per ogni coppia \((k,l)\) considereremo solo quella con direzione coerente a \(d\).

Possiamo dunque definire
\begin{equation}
  \phi^{\delta}_{k,l} (x,y) \eqdef \begin{cases}
    \phi^{\delta^{\rightarrow}}_{k,l} (x,y) & {(\delta(q_k, \sigma_l))}_3 = \rightarrow\\
    \phi^{\delta^{\leftarrow}}_{k,l} (x,y) & {(\delta(q_k, \sigma_l))}_3 = \leftarrow
  \end{cases}
\end{equation}

Per codificare quindi l'intera funzione di transizione sarà sufficiente utilizzare la formula:
\begin{equation}
  \phi_7 \eqdef \forall x, \forall y, \left(\bigwedge_{k=0}^{n} \bigwedge_{l=0}^{m} \phi^{\delta}_{k,l} (x,y)\right)
\end{equation}

Queste sette formule quindi descrivono precisamente la macchina di Turing \(M\) e il suo comportamento, nel senso che ogni modello di queste formule rappresenta un'esecuzione valida di \(M\), e ogni esecuzione valida di \(M\) è rappresentata da un modello di queste formule. A questo punto, per esprimere la proprietà di terminazione della macchina non ci serve altro che richiedere che se la congiunzione di queste condizioni è verificata esista un passo di esecuzione in cui lo stato di terminazione è raggiunto, ovvero che esista \(x\) e \(y\) tali che \(q_n(x,y)\) sia verificato. Dunque, la formula che codifica la terminazione macchina di Turing \(M\) è la seguente:
\begin{equation}
  \phi_M \eqdef \left(\bigwedge_{i=1}^{7} \phi_i\right) \to \exists x, \exists y, q_n(x,y)
\end{equation}

A questo punto, per completare la dimostrazione dell'indecidibilità di VALIDITY, è sufficiente mostrare che \(\phi_M\) è valida se e solo se \(M\) termina.

\begin{theorem}{Indecidibilità di VALIDITY}
  Sia \(M\) una macchina di Turing. Allora \(\models \phi_M\) se e solo se \(M\) termina, ovvero se esiste un \(i \in \N\) tale che \({(\pi_M(i))}_1 = n\).
\end{theorem}
\begin{proof}
  
  Per mostrare la doppia implicazione possiamo verificare separatamente il caso in cui \(M\) non termina e quello in cui termina.
  
  Iniziamo supponendo che \(M\) non termini, ovvero che non esista \(i \in \N\) tale che \({(\pi_M(i))}_1 = n\).
  Per mostrare che \(\not\models \phi_M\) è sufficiente mostrare che esiste un modello di \(\bigwedge_{i=1}^{7} \phi_i\) in cui \(\neg\exists x, \exists y, q_n(x,y)\) è verificato. Essendo \(M\) non terminante, esiste un'esecuzione \(\pi_M\) di \(M\) in cui non esiste \(i \in \N\) tale che \({(\pi_M(i))}_1 = n\). Dunque, prendendo come modello \(M = (\N,I)\) con \(I\) che interpreta i simboli del linguaggio come descritto precedentemente, si ha che \(M\) è un modello di \(\bigwedge_{i=1}^{7} \phi_i\) in cui \(\neg\exists x, \exists y, q_n(x,y)\) è verificato, e dunque \(M\) è un modello di \(\bigwedge_{i=1}^{7} \phi_i \land \neg\exists x, \exists y, q_n(x,y)\), e dunque \(\not\models \phi_M\).
  
  L'altro verso è invece più complesso. Si vuole mostrare che se \(M\) termina, ovvero se esiste \(i \in \N\) tale che \({(\pi_M(i))}_1 = n\), allora \(\models \phi_M\). Essendo che la proprietà di terminazione è finitaria per definizione, se \(M\) termina deve esistere un prefisso finito di \(\pi_M\) di lunghezza \(i\) tale che \({(\pi_M(i-1))}_1 = n\). Possiamo quindi analizzare per induzione questa sequenza finita, mostrando che la codifica della configurazione di ogni passo è conseguenza logica della codifica della macchina di Turing e che può inoltre essere espressa a partire dalla configurazione iniziale iterando la funzione di successore. Formalmente, se a passo \(k\) si ha \(\pi_M(k) = (j, i, t_k)\) allora vale:
  
  \begin{equation}
    \bigwedge_{h=1}^{7} \phi_h \models q_j(s^k(0), s^i(0)) 
  \end{equation}
  
  E per ogni posizione \(r\) del nastro al passo \(k\), con \(t_k(r) = \sigma_l\), vale anche
  \begin{equation}
    \bigwedge_{h=1}^{7} \phi_h \models \sigma_l(s^k(0), s^r(0))
  \end{equation}
  
  Chiaramente, dimostrato questo si ha la tesi, poiché valendo questa conseguenza logica per ogni passo di esecuzione, varrà in particolare per il primo passo in cui lo stato di terminazione è raggiunto. Di conseguenza si avrà che
  \begin{equation}
    \models \bigwedge_{h=1}^{7} \phi_h \models q_n(s^x(0), s^y(0))
  \end{equation}
  per qualche \(x\) e \(y\), e dunque \(\models \phi_M\) per introduzione di \(\exists\).
  
  Per dimostrare questa invariante verifichiamo innanzitutto il caso base, ovvero \(\pi_M(0) = C_0\). Questo caso è banalmente vero grazie a \(\phi_5\), interpretando chiaramente \(s^0\) come l'identità. Per il passo induttivo, supponiamo che valga per un passo \(k\), ovvero che \(\bigwedge_{h=1}^{7} \phi_h \models q_j(s^k(0), s^i(0))\) e che per ogni posizione \(r\) del nastro al passo \(k\), con \(t_k(r) = \sigma_l\), valga anche \(\bigwedge_{h=1}^{7} \phi_h \models \sigma_l(s^k(0), s^r(0))\). Dobbiamo mostrare che allora vale anche per il passo \(k+1\).
  
  Per definizione induttiva di \(\pi_M\) si ha che se \(\pi_M(k) = (j, i, t_k)\) e \(\delta(q_j, t_k(i)) = (q_{j'}, \sigma_{l'}, d)\) allora \(\pi_M(k +1) = (j', i', t_{k+1})\) con
  \begin{itemize}
    \item \(i' = i + 1\) se \(d = \rightarrow\) e \(i' = i - 1\) se \(d = \leftarrow\);
    \item \(t_{k+1}(r) = t_k(r)\) per ogni \(r \neq i\) e \(t_{k+1}(i) = \sigma_{l'}\).
  \end{itemize}
  
  Per ipotesi induttiva si ha che
  \[\bigwedge_{h=1}^{7} \phi_h \models q_j(s^k(0), s^i(0))\]
  e che per ogni posizione \(r\) del nastro al passo \(k\), con \(t_k(r) = \sigma_l\), vale anche
  \[\bigwedge_{h=1}^{7} \phi_h \models \sigma_l(s^k(0), s^r(0))\]
  
  Rispetto alla direzione, l'unica differenza tra le configurazioni d'arrivo è che in una la posizione della testina, codificata con l'esponente \(i\), è incrementata di \(1\) e nell'altra è decrementata di \(1\). Iniziamo assumendo che \(d = \rightarrow\).
  
  Grazie a \(\phi_7\) e alla definizione di \(\phi^{\delta}_{j,l}\), si ha che
  \[\bigwedge_{h=1}^{7} \phi_h \models (q_j(s^k(0),s^i(0)) \land \sigma_l(s^k(0),s^i(0))) \to q_{j'}((s(s^{k}(0)),s(s^{i})(0)) \land \sigma_{l'}(s(s^{k}(0)),s^{i}(0)))\]
  
  Chiaramente \(s(s^{k}(0)) = s^{k+1}(0)\) e \(s(s^{i}(0)) = s^{i+1}(0)\). Dunque essendo per ipotesi veri entrambi i congiunti di questa implicazione sarà anche vero
  \[\bigwedge_{h=1}^{7} \phi_h \models q_{j'}(s^{k+1}(0),s^{i+1}(0)) \land \sigma_{l'}(s^{k+1}(0),s^{i+1}(0))\]
  
  Inoltre, grazie a \(\phi_6\) si ha che tutte le posizioni diverse da \(i\) del nastro al passo \(k+1\) contengono lo stesso simbolo che avevano al passo \(k\), e dunque per ipotesi induttiva si ha che per ogni posizione \(r\) del nastro al passo \(k+1\), con \(t_{k+1}(r) = \sigma_l\), vale anche
  \[\bigwedge_{h=1}^{7} \phi_h \models \sigma_l(s^{k+1}(0), s^r(0))\]
  
  Venendo al caso di \(d = \leftarrow\), l'unica differenza si trova della definizione di \(\phi^{\delta}_{j,l}\), che in questo caso presenta la funzione \(p\). Possiamo notare però che grazie a \(\phi_1\) si ha che \(p(s^i(0)) = s^{i-1}(0)\) ottenendo esattamente la posizione della testina corretta al passo \(k+1\). 
\end{proof}

% 

Decidibilità della logica del primo ordine.

RECAP SU DECIDIBILITÀ, CHIEDI A FABRIZIO
(Essenzialmente problema decisionale come appartenenza di un istanza, decidibile se ricorsivo, ovvero RE e coRE)

Un insieme è ricorsivamente enumerabile se esiste una procedura che può enumerare tutti gli elementi dell'insieme. Invece coRE se il complemento dell'insieme può essere enumerato. 

Chiaramente avendo queste due procedure è semplice capire se il linguaggio è decidibile, poiché basta fare un passo alla volta alle enumerazioni di entrambe le procedure e vedere quale mi genera l'istanza.

Il linguaggio che vogliamo decidere è quello delle formule valide. Chiaramente decidere questo è sufficiente per decidere anche soddisfacibilità e conseguenza logica.

Formalmente quindi abbiamo \(\Valid = \set{\phi \in \FOL}[\models_{\FOL}\phi]\). Da come abbiamo definito la decidibilità dobbiamo verificare se \(\Valid\) è RE e coRE, ovvero \(\overline{\Valid}\) è RE\@.

Da completezza e correttezza è facile vedere che \(\Valid\) è RE\@. Infatti da \(\models \phi \iff \vdash \phi\) si ha che \(\phi \in \Valid \iff \phi \in \set{\phi \in \FOL}[\vdash \phi]\). Ovviamente essendo per definizione che \(\vdash \phi \iff \exists \pi\) derivazione di \(\phi\) da \(\emptyset\) ed essendo che le derivazioni sono alberi finiti. L'insieme delle possibili derivazioni per un linguaggio numerabile è chiaramente numerabile, perché è possibile enumerare tutte le possibili derivazioni di lunghezza \(0\), seguite da tutte le possibili derivazioni di lunghezza \(1\) e così via. Più in generale ogni derivazione è una sequenza finita di simboli in un certo \(\Sigma\), quindi tutte le derivazioni sono un sottoinsieme di \(\Sigma^*\) che è chiaramente numerabile. 

Dunque quello che manca è \(\overline{\Valid}\) RE\@. Si potrebbe pensare di usare lo stesso ragionamento con le formule negate. Il problema è che se \(\phi\) è valida è insoddisfacibile, ma le formule che sono soddisfacibili ma non valide non sarebbero generate e quindi non avremmo una procedura che genera tutte le formule non valide.

Possiamo pensare di provare a mostrare allora che \(\overline{\Valid}\) non sia RE\@. 

Similmente alla gerarchia polinomiale per la decidibilità ci sono diverse gerarchia. Consideriamo la gerarchia aritmetica, in cui le proprietà sono classificate rispetto alle proprietà in teoria dei numeri. In particolare abbiamo \(\Sigma_1^0\) insieme delle formule per cui esiste una prova, \(\Pi_1^0\) insieme delle formule per cui non esiste una prova.

Come fatto in teoria della complessità quindi per mostrare che \(\Valid\) è semi-decidibile possiamo provare a formulare una riduzione da un problema noto essere semi-decidibile ma non decidibile, ad esempio il problema della fermata. Noto infatti un problema \(P'\) che è HARD per una classe \(C\), è possibile mostrare che un altro problema \(P\) è HARD per \(C\) mostrando una riduzione da \(P'\) a \(P\), poiché chiaramente se ogni problema in \(C\) è riducibile a \(P'\) e \(P'\) è riducibile a \(P\) allora ogni problema in \(C\) è riducibile a \(P\).

Chiaramente le riduzioni devono appartenere a classi meno potenti. Per la verifica della NP-HARDNESS ad esempio servono riduzioni polinomiali, mentre per la verifica della RE-HARDNESS servono riduzioni ricorsive, ovvero riduzioni che sono calcolabili da una macchina di Turing.

Come esempio di problema semi-decidibile è sufficiente quindi ridurre ogni istanza del problema della fermata a una formula del primo ordine che è valida se e solo se la macchina di Turing si ferma su quell'istanza.

>>>>>>>>>>>>>>>>>>>>>>>>>>>>>>>>>>>>>>>>>>>>>>>>>>>>>>>>>>>>>>>>>>>>>>>>

Una macchina di Turing \(M\) deterministica è definita da una tupla \((Q, \Sigma, \delta, q_0, F)\) con:
\begin{itemize}
  \item \(Q\) insieme finito di stati di controllo
  \item \(\Sigma = \set{\sigma_0, \sigma_1, \ldots, \sigma_n}\) insieme finito di simboli di nastro con \(\sigma_0 = \sqcup\) simbolo \q{blank}.
  \item \(\delta: Q \times \Sigma \to Q \times \Sigma \times \set{\rightarrow, \leftarrow}\) funzione di transizione
  \item \(q_0\) stato iniziale
  \item \(q_h\) stato di terminazione.
\end{itemize}

La semantica di tali macchine sono chiaramente le esecuzioni, ovvero sequenze di configurazioni, ovvero istantanee dello stato della macchina, che sono definite come segue triple \((k,i,t)\) con \(k\) indice di stato tale che \(q_k\in Q\), \(i\) è la posizione del nastro che la macchina osserva e \(t\) è la configurazione del nastro, ovvero una funzione \(t:\N \to \Sigma\) che a ogni posizione del nastro associa il simbolo che è presente in quella posizione. Chiamando \(\mathcal{C}_M\) l'insieme di tutte le configurazioni di \(M\), una esecuzione non è altro che una funzione \(\pi: \N \to \mathcal{C}_M\) che associa a ogni passo di esecuzione una configurazione, con la condizione che \(\pi(0) = (0,0,t_0:\N\to\set{\sqcup})\) e che per ogni \(n\) si ha che \(\pi(n+1)\) è la configurazione che si ottiene da \(\pi(n)\) applicando la funzione di transizione \(\delta\). Più precisamente si avrà \(\pi(n) = (k,i,t_n)\) e \(\pi(n+1) = (k',i',t_{n+1}')\) se \(\delta(q_k,t_n(i)) = (q_r,\sigma_l,m)\) allora \(k = l\), \(t_{n+1}(i) = \sigma_l\) e \(t_{n+1}(j) = t_n(j)\) per ogni \(j \neq i\) e \(i' = i+1\) se \(m = \rightarrow\) e \(i' = i-1\) se \(m = \leftarrow\).

Ovviamente un esecuzione è terminante se esiste \(n\) tale che \(\pi(n) = (h,i,t)\) per qualche \(i\) e \(t\).

Dunque data una \(M\) qualsiasi (che determina univocamente \(\pi\))  dovremo costruire una formula \(\phi_M\) tale che
\[\models \phi_M \iff M \text{ è terminante}\]

Innanzitutto dobbiamo definire il linguaggio \(L\) che useremo per costruire \(\phi_M\). Intuitivamente il dominio dei modelli di questa formula saranno numeri naturali che rappresentano i passi di esecuzione e le posizioni del nastro, quindi avremo bisogno di un unica costante \(0\), due funzioni binarie per muoverci tra i numeri naturali, ovvero \(s\) e \(p\) che rappresentano rispettivamente il successore e il predecessore. Infine ci servono due insiemi finiti di predicati binari \(q_k\) con \(k \in \set{0, \ldots, h}\) tali che \(q_k(x,y)\) è vera se e solo se al passo \(x\) di esecuzione la macchina è nello stato \(q_k\) e nella posizione \(y\) del nastro, e \(\sigma_l\) con \(l \in \set{0, \ldots, n}\) tali che \(\sigma_l(x,y)\) è vera se e solo se al passo \(x\) di esecuzione nella posizione \(y\) del nastro è presente il simbolo \(\sigma_l\).

Chiaramente ora su questo linguaggio dovremo scrivere un insieme di formule che rappresentino l'esecuzione della macchina. Iniziamo con
\[\phi_1 \eqdef \forall x s(p(x)) = x \land p(s(x)) = x\]

Insieme a questa formula potremmo aggiungere tutta la teoria dei numeri di Peano, ma per quello che ci interessa non sarà necessario. Vediamo ora le formule che rappresentano le caratteristiche della macchina. Per comodità di notazione, la variabile \(x\) verrà utilizzata per rappresentare i passi di esecuzione, mentre le variabili \(y\) e \(z\) verranno utilizzate per rappresentare le posizioni del nastro.

\[\phi_2 \eqdef \forall x \forall y (\bigwedge_{i=0}^{n} \bigwedge_{\substack{j=0\\j \neq i}}^{n} \neg(\sigma_i(x,y) \land \sigma_j(x,y)))\]

Questa formula semplicemente dice che in ogni passo di esecuzione non possono essere presenti due simboli diversi nella stessa posizione del nastro.

\[\phi_3 \eqdef \forall x \forall y \forall z (\bigwedge_{i=0}^{h} \bigwedge_{\substack{j=0\\j \neq i}}^{h} \neg(q_i(x,y) \land q_j(x,z)))\]

Qui abbiamo usato una variabile in più perché è necessario che la macchina sia in un solo stato in ogni passo di esecuzione a prescindere dalla posizione del nastro, quindi è necessario quantificare universalmente sulla posizione. Per concludere le proprietà di consistenza abbiamo infine

\[\phi_4 \eqdef \forall x \forall y \forall z (\neg(y = z)\to \bigwedge_{i=0}^{n} \neg(q_i(x,y) \land q_i(x,z)))\]

Quest'ultima ci dice che in ogni passo di esecuzione anche trovandosi nello stesso stato la macchina non può osservare più posizioni del nastro contemporaneamente.

Queste formule ci parlano solo di proprietà locale alle configurazioni, che intuitivamente ci descrivono appunto una configurazione ben formata.

Passiamo ora alla riduzione vera e propria. Iniziamo col descrivere la configurazione iniziale, ovvero quella in cui la macchina si trova nello stato iniziale \(q_0\) e nella posizione \(0\) del nastro, con il nastro completamente vuoto. La formula che descrive questa configurazione è la seguente:
\[\phi_5 \eqdef q_0(0,0) \land \forall y \sigma_0(0,y) \]

Ci mancano a questo punto due formule. La prima ci dice che l'unica possibilità che una cella cambi simbolo è che la macchina si trovi in quella cella.

\[\phi_6 \eqdef \forall x \forall y (\bigwedge_{i=0}^{n} (\sigma_i(x,y) \land \neg \sigma_i(s(x),y))) \to (\bigvee_{j=0}^{h} q_j(x,y))\]

Questa ci dice che solo le celle osservate dalla macchina possono cambiare simbolo.

Veniamo quindi all'ultima formula, la più complessa, che descrive di fatto la dinamica di esecuzione della macchina. Questa formula è di fatto una congiunzione di formule più semplici, ognuna delle quali descrive una possibile transizione della macchina. Avremo in particolare una formula per ogni possibile input di \(\delta\), quindi per ogni \(\delta(q_i,\sigma_j) = (q_k, \sigma_l, m)\) avremo una formula del tipo
\[\phi^{\delta}_{i,j} (x,y)\eqdef ((q_i(x,y) \land \sigma_j(x,y)) \to q_k(s(x),m(y)) \land \sigma_l(s(x),y))\]
con \(m(y) = s(y)\) se \(m = \rightarrow\) e \(m(y) = p(y)\) se \(m = \leftarrow\). Ovviamente andrà aggiunta solo la formula che rappresenta la transizione effettiva, quindi solo una tra \(m = \rightarrow\) e \(m = \leftarrow\) per ogni possibile input di \(\delta\).

A questo punto ci basta quantificare universalmente su \(x\) e \(y\) per ottenere la formula che descrive la dinamica di esecuzione della macchina, ovvero
\[\phi_7 \eqdef \forall x \forall y \bigwedge_{i=0}^{h} \bigwedge_{j=0}^{n} \phi^{\delta}_{i,j}(x,y)\]

Dunque, per descrivere la terminazione di \(M\) sarà sufficiente la formula
\[\phi_M \eqdef \bigwedge_{i=1}^7 \phi_i \to \exists x \exists y q_h(x,y)\]

Mostriamo ora che \(\models \phi_M\) se e solo se \(M\) è terminante. Dimostriamo le due implicazioni separatamente.

Per mostrare \(\models \phi_M \implies M\) termina possiamo per contrapposizione mostrare che \(M\) non termina \(\implies \not\models \phi_M\). Per mostrare che \(\not\models \phi_M\), essendo \(\phi_M\) una formula del tipo \(\psi \to \chi\) è sufficiente mostrare che esiste un modello \(m\) tale che \(m \models \psi\) e \(m \not\models \chi\). Ora però \(\psi\) è la descrizione della macchina \(M\), quindi affinché \(m \models \psi\) è necessario che \(m\) sia coerente con un esecuzione di \(M\). Scegliamo l'esecuzione non terminante, esistente per ipotesi. Dunque avremo \(m = (\N, I)\) con \(I(0) = 0\), \(I(s) = \text{succ}\), \(I(p) = \text{pred}\) e, per ogni \(k\), \(I(q_k) = \set{(x,y)\in\N^2}[\pi(x) = (k,y,t_x)]\) e, per ogni \(l\), \(I(\sigma_l) = \set{(x,y)\in\N^2}[\pi(x) = (k,y,t_x) \land t_x(y) = \sigma_l]\). Ora questo modello sta codificando esattamente l'esecuzione non terminante \(\pi\) di \(M\), quindi è chiaro che \(m \models \psi\). D'altra parte però \(m \not\models \chi\) per ipotesi.

Il viceversa è un po' più complicato, e iniziamo a vederne un intuizione. Vogliamo mostrare che se \(M\) è terminante allora \(\models \phi_M\). Se \(M\) è terminante, allora esiste un prefisso finito di lunghezza \(k\) di \(\pi_M\) tale che \(\pi_M(k) = (h,i,t)\) per qualche \(i\) e \(t\). Allora possiamo pensare di fare induzione sui passi di esecuzione per mostrare che se al passo \(p\) si ha \(\pi_M() = (j, i, t_p)\) allora \(\bigwedge_{z = 1}^{7} \phi_{z} \models q_j(s^p(0),s^i(0))\) e inoltre per ogni \(l\), se \(t_i(l) = \sigma_l\) allora \(\bigwedge_{z = 1}^{7} \phi_{z} \models \sigma_l(s^p(0),s^l(0))\). Questo vale sicuramente per \(i = 0\) poiché \(\pi_M(0) = (0,0,t_0)\) e chiaramente le due proprietà derivano direttamente da \(\phi_5\). Nel caso induttivo la cosa è più delicata, ma l'idea è che supponendo che \(\pi(i) = (k, j, t_i)\) e che \(\pi(i+1) = (k', j', t_{i+1})\) allora per ipotesi induttiva avremo che \(\bigwedge_{z = 1}^{7} \phi_{z} \models q_k(s^i(0),s^j(0))\) e inoltre per ogni \(l\), se \(t_i(l) = \sigma_l\) allora \(\bigwedge_{z = 1}^{7} \phi_{z} \models \sigma_l(s^i(0),s^l(0))\). Ora però passare da \(\pi(i)\) a \(\pi(i+1)\) farò esattamente il passo che è descritto dalla \(\delta\), ovvero \(\delta(q_k, \sigma_l) = (q_{k'}, \sigma_{l'}, m)\) per \(m\in\set{\rightarrow, \leftarrow}\).

VEDI BENE NEI DETTAGLI, l'idea è che serve \(\phi_1\).

\appendix{}
\chapter{Esercizi Logica Proposizionale}

Quest'appendice contiene una raccolta di esercizi proposti durante il corso, con relative soluzioni.

Gli esercizi verranno organizzati in sezioni, a seconda dell'argomento trattato.

\section{Deduzione naturale per \(\PL\)}\label{app_sec:natural_deduction_exercises}

\subsection{Svolti}

\subsubsection*{Dimostrazione di \(P\land Q \equiv \neg (\neg P \lor \neg Q)\).}

Per dimostrare questa equivalenza andiamo a dimostrare i due sequenti \(P \land Q \vdash \neg (\neg P \lor \neg Q)\) e \(\neg (\neg P \lor \neg Q) \vdash P \land Q\).

\paragraph{Dimostrazione di \(P \land Q \vdash \neg (\neg P \lor \neg Q)\).}

Essendo che l'operatore principale è la una contraddizione possiamo iniziare con una \(\neg I\).

\begin{prooftree}
  \AxiomC{\(\bot\)}
  \RuleLabel{\(\neg I_1\)}[\(\neg P \lor \neg Q\)]
  \UnaryInfC{\(\neg (\neg P \lor \neg Q)\)}
\end{prooftree}

Essendo che semanticamente sappiamo che da \(P \land Q\) segue che sia \(P\) che \(Q\) sono veri, è ragionevole pensare di poter ottenere una contraddizione da \(\neg P \lor \neg Q\).
Questo perché \(\neg P \lor \neg Q\) richiede che almeno uno tra \(P\) e \(Q\) sia falso, ma noi sappiamo che entrambi sono veri.
Qui ci tocca usare \(\lor E\) poiché per usare una \(\neg E\) sarebbe necessario usare ciò che sto cercando di dimostrare, e questo non è possibile.

\begin{prooftree}
  \AxiomC{\(\bot\)}
  \AxiomC{\(\bot\)}
  \AxiomC{\({[\neg P \lor \neg Q]}_1\)}
  \RuleLabel{\(\lor E_{2,3}\)}[\(\neg P\), \(\neg Q\)]
  \TrinaryInfC{\(\bot\)}
  \RuleLabel{\(\neg I_1\)}
  \UnaryInfC{\(\neg (\neg P \lor \neg Q)\)}
\end{prooftree}

A questo punto sarà necessario trovare una contraddizione sia da \(\neg P\) che da \(\neg Q\). Per entrambi sarà sufficiente usare una \(\land E\) per ottenere \(P\) e \(Q\) da \(P \land Q\).

\begin{prooftree}
  \AxiomC{\(P\land Q\)}
  \RuleLabel{\(\land E\)}
  \UnaryInfC{\(P\)}
  \AxiomC{\({[\neg P]}_2\)}
  \RuleLabel{\(\neg E\)}
  \BinaryInfC{\(\bot\)}
  \AxiomC{\(P\land Q\)}
  \RuleLabel{\(\land E\)}
  \UnaryInfC{\(Q\)}
  \AxiomC{\({[\neg Q]}_3\)}
  \RuleLabel{\(\neg E\)}
  \BinaryInfC{\(\bot\)}
  \AxiomC{\({[\neg P \lor \neg Q]}_1\)}
  \RuleLabel{\(\lor E_{2,3}\)}
  \TrinaryInfC{\(\bot\)}
  \RuleLabel{\(\neg I_1\)}
  \UnaryInfC{\(\neg (\neg P \lor \neg Q)\)}
\end{prooftree}

\paragraph{Dimostrazione di \(\neg (\neg P \lor \neg Q) \vdash P \land Q\).}

Qui il ragionamento è diverso. Avendo come operatore principale una congiunzione, è ragionevole pensare di usare \(\land I\) per ottenere \(P \land Q\).

\begin{prooftree}
  \AxiomC{\(P\)}
  \AxiomC{\(Q\)}
  \RuleLabel{\(\land I\)}
  \BinaryInfC{\(P \land Q\)}
\end{prooftree}

Dobbiamo quindi dimostrare separatamente \(P\) e \(Q\).

A questo punto con \(\RAA\) possiamo assumere \(\neg P\) e dimostrare una contraddizione, e poi assumere \(\neg Q\) e dimostrare una contraddizione. Da queste assunzioni possiamo introdurre \(\neg P \land \neg Q\) con \(\land I\), che ci permette di ottenere una contraddizione con \(\neg (\neg P \lor \neg Q)\).

\begin{prooftree}
  \AxiomC{\({[\neg Q]}_1\)}
  \RuleLabel{\(\lor I\)}
  \UnaryInfC{\(\neg P \lor \neg Q\)}
  \AxiomC{\(\neg(\neg P \lor \neg Q)\)}
  \BinaryInfC{\(\bot\)}
  \RuleLabel{\(\RAA_1\)}
  \UnaryInfC{\(Q\)}
  \AxiomC{\({[\neg P]}_2\)}
  \RuleLabel{\(\lor I\)}
  \UnaryInfC{\(\neg P \lor \neg Q\)}
  \AxiomC{\(\neg(\neg P \lor \neg Q)\)}
  \BinaryInfC{\(\bot\)}
  \RuleLabel{\(\RAA_2\)}
  \UnaryInfC{\(P\)}
  \RuleLabel{\(\land I\)}
  \BinaryInfC{\(P \land Q\)}
\end{prooftree}

\subsubsection*{Dimostrazione di \((P\to Q) \leftrightarrow (\neg Q \to \neg P)\).}

Anche in questo caso conviene andare per passi. Iniziamo con \((P\to Q) \to (\neg Q \to \neg P)\). Per dimostrare questa implicazione, possiamo usare la regola di introduzione dell'implicazione (\(\to I\)) per assumere \(P \to Q\) e poi dimostrare \(\neg Q \to \neg P\).

\begin{prooftree}
  \AxiomC{\(\neg Q \to \neg P\)}
  \RuleLabel{\(\to I\)}[\(P \to Q\)]
  \UnaryInfC{\((P\to Q) \to (\neg Q \to \neg P)\)}
\end{prooftree}

Abbiamo quindi un ulteriore implicazione da dimostrare, quindi possiamo usare di nuovo \(\to I\) per assumere \(\neg Q\) e poi dimostrare \(\neg P\).

\begin{prooftree}
  \AxiomC{\(\neg P\)}
  \RuleLabel{\(\to I\)}[\(\neg Q\)]
  \UnaryInfC{\(\neg Q \to \neg P\)}
  \RuleLabel{\(\to I\)}[\(P \to Q\)]
  \UnaryInfC{\((P\to Q) \to (\neg Q \to \neg P)\)}
\end{prooftree}

Ora per dimostrare una negazione, possiamo usare la regola di introduzione della negazione (\(\neg I\)) che ci permette di assumere \(P\) e dimostrare una contraddizione.
A questo punto possiamo usare la regola di eliminazione dell'implicazione (\(\to E\)) con \(P\) e \(P \to Q\) per ottenere \(Q\), e poi usare \(Q\) e \(\neg Q\) per ottenere la contraddizione che cercavamo.

\begin{prooftree}
  \AxiomC{\({[P]}_3\)}
  \AxiomC{\({[P \to Q]}_1\)}
  \RuleLabel{\(\to E\)}
  \BinaryInfC{\(Q\)}
  \AxiomC{\({[\neg Q]}_2\)}
  \RuleLabel{\(\neg E\)}
  \BinaryInfC{\(\bot\)}
  \RuleLabel{\(\neg I_3\)}
  \UnaryInfC{\(\neg P\)}
  \RuleLabel{\(\to I_2\)}
  \UnaryInfC{\(\neg Q \to \neg P\)}
  \RuleLabel{\(\to I_1\)}
  \UnaryInfC{\((P\to Q) \to (\neg Q \to \neg P)\)}
\end{prooftree}

Per l'altro verso, il ragionamento è analogo, ma utilizzando \(\RAA\) invece di \(\neg I\).

\begin{prooftree}
  \AxiomC{\({[\neg Q]}_3\)}
  \AxiomC{\({[\neg Q \to \neg P]}_1\)}
  \RuleLabel{\(\to E\)}
  \BinaryInfC{\(\neg P\)}
  \AxiomC{\({[P]}_2\)}
  \RuleLabel{\(\neg E\)}
  \BinaryInfC{\(\bot\)}
  \RuleLabel{\(\RAA_3\)}
  \UnaryInfC{\(Q\)}
  \RuleLabel{\(\to I_2\)}
  \UnaryInfC{\(P\to Q\)}
  \RuleLabel{\(\to I_1\)}
  \UnaryInfC{\((\neg Q \to \neg P)\to(P\to Q)\)}
\end{prooftree}


\subsubsection*{Dimostrazione di \((\neg \neg (A \to B)) \leftrightarrow (A \to B)\)}

Possiamo notare che questo teorema ha la stessa struttura di \(\neg \neg A \leftrightarrow A\), che abbiamo già dimostrato nell'\Namedcref{ex:natural_deduction}.

Possiamo allora costruire un albero di derivazione identico e sostituire ogni occorrenza di \(A\) con \(A \to B\).

Questo è possibile per il \Namedref{thm:substitution}, poiché stiamo applicando una sostituzione \(\sigma = \set{A \leftarrow A \to B}\) a una formula valida. Il teorema di sostituzione ci garantisce che la formula risultante dalla sostituzione è ancora valida, e quindi dimostrabile, per la completezza di \(\vdash\).

\subsubsection*{Dimostrazione di \((A \lor B) \land C \equiv (A \land C) \lor (B \land C)\).}

Iniziamo come sempre a dimostrare i due sequenti \( (A \lor B) \land C \vdash (A \land C) \lor (B \land C)\) e \((A \land C) \lor (B \land C) \vdash (A \lor B) \land C\).

\paragraph{Dimostrazione di \((A \lor B) \land C \vdash (A \land C) \lor (B \land C)\).}

Abbiamo una disgiunzione come operatore principale, quindi andiamo a utilizzare \(\lor I\).

\begin{prooftree}
  \AxiomC{\(A \land C\)}
  \RuleLabel{\(\lor I\)}
  \UnaryInfC{\((A \land C) \lor (B \land C)\)}
\end{prooftree}

A questo punto per dimostrare \(A \land C\) possiamo usare \(\land I\).

\begin{prooftree}
  \AxiomC{\(A\)}
  \AxiomC{\(C\)}
  \RuleLabel{\(\land I\)}
  \BinaryInfC{\(A \land C\)}
  \RuleLabel{\(\lor I\)}
  \UnaryInfC{\((A \land C) \lor (B \land C)\)}
\end{prooftree}

Dobbiamo quindi dimostrare separatamente \(A\) e \(C\). Per \(C\) è sufficiente usare una \(\land E\) per ottenere \(C\) da \((A \lor B) \land C\).

Per \(A\) invece dobbiamo sfruttare \(A \lor B\) ottenibile per \(\land E\) da \((A \lor B) \land C\). A questo punto, trovandoci con una disgiunzione, dobbiamo usare \(\lor E\) per dimostrare \(A\) sia assumendo \(A\) che assumendo \(B\).

\begin{prooftree}
  \AxiomC{\(A\)}
  \AxiomC{\(A\)}
  \AxiomC{\(A\lor B\)}
  \RuleLabel{\(\lor E_{2,3}\)}[\(A\), \(B\)]
  \TrinaryInfC{\(A\)}
  \AxiomC{\((A \lor B) \land C\)}
  \RuleLabel{\(\land E\)}
  \UnaryInfC{\(C\)}
  \RuleLabel{\(\land I\)}
  \BinaryInfC{\(A \land C\)}
  \RuleLabel{\(\lor I\)}
  \UnaryInfC{\((A \land C) \lor (B \land C)\)}
\end{prooftree}

Ottenere \(A\) assumendo \(A\) è ovvio, ma ottenere \(A\) assumendo \(B\) non è possibile.

Un alternativa per la dimostrazione di una disgiunzione, per quanto paradossale possa sembrare, è quella di usare \(\lor E\) invece di \(\lor I\). Questo è possibile se tra le assunzioni che possiamo usare c'è una disgiunzione. In questo caso, avendo \(A \lor B\) ottenibile da \((A \lor B) \land C\), possiamo usare \(\lor E\).

Ricominciamo quindi la dimostrazione usando \(\lor E\) invece di \(\lor I\).

\begin{prooftree}
  \AxiomC{\((A\lor B)\land C\)}
  \RuleLabel{\(\land E\)}
  \UnaryInfC{\(A\lor B\)}
  \AxiomC{\((A \land C) \lor (B \land C)\)}
  \AxiomC{\((A \land C) \lor (B \land C)\)}
  \RuleLabel{\(\lor E_{1,2}\)}[\(A, B\)]
  \TrinaryInfC{\((A \land C) \lor (B \land C)\)}
\end{prooftree}

A questo punto possiamo andare a dimostrare i due rami introducendo la disgiunzione dimostrando il disgiunto che contiene l'assunzione che possiamo scaricare. Per esempio, per il primo ramo, assumendo \(A\) possiamo dimostrare \(A \land C\) usando \(\land I\), e poi introdurre la disgiunzione con \(\lor I\).

\begin{prooftree}
  \AxiomC{\((A\lor B)\land C\)}
  \RuleLabel{\(\land E\)}
  \UnaryInfC{\(A\lor B\)}
  \AxiomC{\((A \lor B) \land C\)}
  \RuleLabel{\(\land E\)}
  \UnaryInfC{\(C\)}
  \AxiomC{\({[A]}_1\)}
  \RuleLabel{\(\land I\)}
  \BinaryInfC{\((A \land C) \)}
  \RuleLabel{\(\lor I\)}
  \UnaryInfC{\((A \land C) \lor (B \land C)\)}
  \AxiomC{\((A \land C) \lor (B \land C)\)}
  \RuleLabel{\(\lor E_{1,2}\)}[\(B\)]
  \TrinaryInfC{\((A \land C) \lor (B \land C)\)}
\end{prooftree}

Il secondo ramo è analogo.

\begin{prooftree}
  \AxiomC{\((A\lor B)\land C\)}
  \RuleLabel{\(\land E\)}
  \UnaryInfC{\(A\lor B\)}
  \AxiomC{\((A \lor B) \land C\)}
  \RuleLabel{\(\land E\)}
  \UnaryInfC{\(C\)}
  \AxiomC{\({[A]}_1\)}
  \RuleLabel{\(\land I\)}
  \BinaryInfC{\((A \land C) \)}
  \RuleLabel{\(\lor I\)}
  \UnaryInfC{\((A \land C) \lor (B \land C)\)}
  \AxiomC{\((A \lor B) \land C\)}
  \RuleLabel{\(\land E\)}
  \UnaryInfC{\(C\)}
  \AxiomC{\({[B]}_1\)}
  \RuleLabel{\(\land I\)}
  \BinaryInfC{\((B \land C) \)}
  \RuleLabel{\(\lor I\)}
  \UnaryInfC{\((A \land C) \lor (B \land C)\)}
  \RuleLabel{\(\lor E_{1,2}\)}
  \TrinaryInfC{\((A \land C) \lor (B \land C)\)}
\end{prooftree}

\paragraph{Dimostrazione di \((A \land C) \lor (B \land C) \vdash (A \lor B) \land C\).}

In questo caso, abbiamo da dimostrare una congiunzione, quindi è necessario iniziare con \(\land I\).

\begin{prooftree}
  \AxiomC{\(A \lor B\)}
  \AxiomC{\(C\)}
  \RuleLabel{\(\land I\)}
  \BinaryInfC{\((A \lor B) \land C\)}
\end{prooftree}

Nuovamente, avendo una disgiunzione tra le assunzioni, possiamo usare \(\lor E\) per ottenere assunzioni scaricabili che ci permettono di dimostrare \(A \lor B\) e \(C\).

Possiamo ottenere infatti \(C\) da entrambi i disgiunti usando \(\land E\)

\begin{prooftree}
  \AxiomC{\((A \land C) \lor (B \land C)\)}
  \AxiomC{\({[B \land C]}_1\)}
  \RuleLabel{\(\land E\)}
  \UnaryInfC{\(C\)}
  \AxiomC{\({[A \land C]}_1\)}
  \RuleLabel{\(\land E\)}
  \UnaryInfC{\(C\)}
  \RuleLabel{\(\lor E_{1,2}\)}
  \TrinaryInfC{\(C\)}
  \AxiomC{\(A \lor B\)}
  \RuleLabel{\(\land I\)}
  \BinaryInfC{\((A \lor B) \land C\)}
\end{prooftree}

Per dimostrare \(A \lor B\) invece, dovremmo usare nuovamente \(\lor E\), ma questa volta derivando una volta \(A\) e una volta \(B\) usando \(\land E\).

\begin{prooftree}
  \tiny
  \AxiomC{\((A \land C) \lor (B \land C)\)}
  \AxiomC{\({[B \land C]}_2\)}
  \RuleLabel{\(\land E\)}
  \UnaryInfC{\(C\)}
  \AxiomC{\({[A \land C]}_1\)}
  \RuleLabel{\(\land E\)}
  \UnaryInfC{\(C\)}
  \RuleLabel{\(\lor E_{1,2}\)}
  \TrinaryInfC{\(C\)}
  \AxiomC{\({[A \land C]}_3\)}
  \RuleLabel{\(\land E\)}
  \UnaryInfC{\(A\)}
  \RuleLabel{\(\lor I\)}
  \UnaryInfC{\(A \lor B\)}
  \AxiomC{\({[B \land C]}_4\)}
  \RuleLabel{\(\land E\)}
  \UnaryInfC{\(B\)}
  \RuleLabel{\(\lor I\)}
  \UnaryInfC{\(A \lor B\)}
  \AxiomC{\((A \land C) \lor (B \land C)\)}
  \RuleLabel{\(\lor E_{3,4}\)}
  \TrinaryInfC{\(A \lor B\)}
  \RuleLabel{\(\land I\)}
  \BinaryInfC{\((A \lor B) \land C\)}
\end{prooftree}

\subsubsection*{Dimostrazione di \((A \land B) \lor C \equiv (A \lor C) \land (B \lor C)\).}

Dividiamo la dimostrazione in due parti, come al solito. Iniziamo con \((A \land B) \lor C \vdash (A \lor C) \land (B \lor C)\).

Per dimostrare questa implicazione, avendo una congiunzione come operatore principale, è ragionevole pensare di usare \(\land I\) per ottenere \((A \lor C) \land (B \lor C)\).

\begin{prooftree}
  \AxiomC{\(A \lor C\)}
  \AxiomC{\(B \lor C\)}
  \RuleLabel{\(\land I\)}
  \BinaryInfC{\((A \lor C) \land (B \lor C)\)}
\end{prooftree}

Ancora una volta, la presenza di una disgiunzione tra le assunzioni ci porta a usare \(\lor E\) per ottenere assunzioni scaricabili che ci permettono di dimostrare \(A \lor C\) e \(B \lor C\). Iniziamo con \(A \lor C\).

\begin{prooftree}
  \AxiomC{\({[A \land B]}_1\)}
  \RuleLabel{\(\land E\)}
  \UnaryInfC{\(A\)}
  \RuleLabel{\(\lor I\)}
  \UnaryInfC{\(A \lor C\)}
  \AxiomC{\({[C]}_2\)}
  \RuleLabel{\(\lor I\)}
  \UnaryInfC{\(A \lor C\)}
  \AxiomC{\((A \land B) \lor C\)}
  \RuleLabel{\(\lor E_{1,2}\)}
  \TrinaryInfC{\(A \lor C\)}
  \AxiomC{\(B \lor C\)}
  \RuleLabel{\(\land I\)}
  \BinaryInfC{\((A \lor C) \land (B \lor C)\)}
\end{prooftree}

L'altro ramo è analogo, ma invece di dimostrare \(A \lor C\) dobbiamo dimostrare \(B \lor C\). Verrà omesso per brevità e spazio.

Per l'altro verso, \((A \lor C) \land (B \lor C) \vdash (A \land B) \lor C\), possiamo sia usare \(\land I\) o, come visto in precedenza, \(\lor E\). In questo caso, nonostante abbiamo una congiunzione come assunzione, possiamo vederla come invece due disgiunzioni grazie alla regola di eliminazione della congiunzione (\(\land E\)).

Quindi scegliendo una delle due disgiunzioni, ad esempio \(A \lor C\), possiamo usare \(\lor E\).

\begin{prooftree}
  \AxiomC{\((A \lor C) \land (B \lor C)\)}
  \RuleLabel{\(\land E\)}
  \UnaryInfC{\(A \lor C\)}
  \AxiomC{\((A \land B) \lor C\)}
  \AxiomC{\((A \land B) \lor C\)}
  \RuleLabel{\(\lor E_{1,2}\)}[\(A, C\)]
  \TrinaryInfC{\((A \land B) \lor C\)}
\end{prooftree}

Un ramo possiamo subito dimostrarlo, usando \(\lor I\) per introdurre la disgiunzione mediante \(C\).

\begin{prooftree}
  \AxiomC{\((A \lor C) \land (B \lor C)\)}
  \RuleLabel{\(\land E\)}
  \UnaryInfC{\(A \lor C\)}
  \AxiomC{\({[C]}_2\)}
  \RuleLabel{\(\lor I\)}
  \UnaryInfC{\((A \land B) \lor C\)}
  \AxiomC{\((A \land B) \lor C\)}
  \RuleLabel{\(\lor E_{1,2}\)}[\(A\)]
  \TrinaryInfC{\((A \land B) \lor C\)}
\end{prooftree}

Per l'altro ramo, invece, si potrebbe essere tentati di usare nuovamente \(\lor I\) tentando di dimostrare \(A \land B\). Possiamo però subito accorgerci che le assunzioni che abbiamo a disposizione non sono sufficienti per dimostrare direttamente \(A \land B\). Abbiamo infatti \(A\) e \(B \lor C\), dalle quali sappiamo solo che è necessariamente vero uno tra \(B\) e \(C\), e quindi non possiamo concludere con certezza che \(B\) sia vero, necessario a \(A \land B\).

Possiamo allora utilizzare nuovamente \(\lor E\), in modo da poter avere sui nuovi rami una volta \(C\), che da sola come visto ci permette di dimostrare \((A \land B) \lor C\), e una volta \(B\), che insieme ad \(A\) ci permette di dimostrare \(A \land B\) e quindi di introdurre la disgiunzione con \(\lor I\).

\begin{prooftree}
  \tiny
      \AxiomC{\((A \lor C) \land (B \lor C)\)}
      \RuleLabel{\(\land E\)}
    \UnaryInfC{\(A \lor C\)}

      \AxiomC{\({[C]}_2\)}
      \RuleLabel{\(\lor I\)}
    \UnaryInfC{\((A \land B) \lor C\)}
    
        \AxiomC{\((A \lor C) \land (B \lor C)\)}
        \RuleLabel{\(\land E\)}
      \UnaryInfC{\(B \lor C\)}
      
        \AxiomC{\({[C]}_4\)}
        \RuleLabel{\(\lor I\)}
      \UnaryInfC{\((A \land B) \lor C\)}
      
          \AxiomC{\({[A]}_2\)}
          \AxiomC{\({[B]}_3\)}
          \RuleLabel{\(\land I\)}
        \BinaryInfC{\(A \land B\)}
        \RuleLabel{\(\lor I\)}
      \UnaryInfC{\((A \land B) \lor C\)}
      \RuleLabel{\(\lor E_{3,4}\)}
    \TrinaryInfC{\((A \land B) \lor C\)}
    \RuleLabel{\(\lor E_{1,2}\)}
  \TrinaryInfC{\((A \land B) \lor C\)}
\end{prooftree}

\begin{note}{}
  Mai come questo esercizio mostra come la regola di \(\lor E\) sia sottile. Infatti, nonostante ci dicesse già nella prima applicazione che avremmo potuto scaricare \(C\), questo era possibile solo in uno dei due rami, e non in entrambi. Per poterla scaricare anche nell'altro siamo stati costretti a una seconda applicazione di \(\lor E\),
\end{note}

\subsection{Non svolti}
\[\neg (A \to B), (\neg B) \to C \vdash C\]
\[((A\land B)\lor (\neg C \lor B))\to (C \to B)\]
\[(A \to (B\lor C))\lor((\neg A \land \neg C)\to B)\]
\chapter{Esercizi Logica del Primo Ordine}
\listoftheorems{}
\listofexamples{}
\listofalgos{}

\newpage
\whitepage{}
\makebackcover{}

\end{document}